% !TEX program = xelatex
% =====================================================================
%  圆锥曲线的三维矢量方法
%  编译方式：xelatex main.tex  （连续编译两次以生成目录与交叉引用）
%  或直接运行：sh build.sh
% =====================================================================
\documentclass[11pt,openany]{ctexbook}

% !TEX program = xelatex
% =====================================================================
%  圆锥曲线的三维矢量方法 —— 导言区
%  编译方式：xelatex main.tex  （连续编译两次以生成目录与交叉引用）
% =====================================================================

\usepackage{amsmath,amssymb,amsthm}
\usepackage{geometry}
\usepackage{mathrsfs}
\usepackage{array}
\usepackage{booktabs}
\usepackage{enumitem}
\usepackage{tikz}
\usetikzlibrary{calc,intersections}
\usepackage[hidelinks,bookmarksnumbered]{hyperref}

\geometry{a4paper,left=2.6cm,right=2.6cm,top=2.8cm,bottom=2.8cm}

% ---------- 定理类环境 ----------
\theoremstyle{definition}
\newtheorem{definition}{定义}[chapter]
\newtheorem{example}[definition]{例}
\newtheorem{exercise}[definition]{习题}
\newtheorem{remark}[definition]{注}
\theoremstyle{plain}
\newtheorem{theorem}[definition]{定理}
\newtheorem{proposition}[definition]{命题}
\newtheorem{lemma}[definition]{引理}
\newtheorem{corollary}[definition]{推论}

\renewcommand{\proofname}{\bfseries 证明}
\renewcommand{\contentsname}{目\quad 录}

% ---------- 常用记号 ----------
\newcommand{\R}{\mathbb{R}}
\newcommand{\dd}{\mathrm{d}}
\newcommand{\ii}{\mathrm{i}}
\newcommand{\tr}{\operatorname{tr}}
\newcommand{\diag}{\operatorname{diag}}
\newcommand{\rank}{\operatorname{rank}}
\newcommand{\sgn}{\operatorname{sgn}}
\DeclareMathOperator{\Ker}{Ker}
\DeclareMathOperator{\Ima}{Im}

% 度规无关记号：混合积
\newcommand{\mix}[3]{[#1\,#2\,#3]}
% 三阶反对称张量
\newcommand{\eps}{\epsilon}
% 对称积 (\cdot\cdot)
\newcommand{\sym}[2]{(#1#2)}
% 无穷远线矢
\newcommand{\linf}{l_{\infty}}
\newcommand{\linfsub}[1]{l_{\infty #1}}

\allowdisplaybreaks

% 中文段落断行容忍度：避免个别行溢出页面右边界
\setlength{\emergencystretch}{2.5em}
\tolerance=2000
\hbadness=3000


\begin{document}

\frontmatter
\title{\bfseries 圆锥曲线的三维矢量方法}
\author{}
\date{}
\maketitle

\chapter*{前言}
\addcontentsline{toc}{chapter}{前言}
% 前言
本书讨论平面上的圆锥曲线问题。与通常的解析几何处理方式不同，本书不把平面上的点写成坐标 $(x,y)$
以后再作代数运算，而是把点与直线都表示为三维实矢量空间中的元素，并直接用矢量与张量的运算
（缩并、点乘、叉乘、对称积、二次叉乘等）表达几何关系。这样做的目的是使几何对象保持为一个
整体，从而使命题的几何内容在代数写法中直接可见。

书中假定读者了解线性代数的基本内容（矢量空间、线性映射、矩阵、二次型、特征值与对角化），
并了解点乘、叉乘与混合积的定义与初等性质。除此之外的内容都在书中给出。

\section*{约定与记号}
\addcontentsline{toc}{section}{约定与记号}

\begin{enumerate}[label=\textbf{(\arabic*)},leftmargin=2.2em]

\item \textbf{指标与求和。} 上指标记矢量，下指标记对偶矢量。一个指标在同一项中上下各出现一次
时表示对该指标从 $1$ 到 $n$ 求和（爱因斯坦求和约定），称为哑指标；其余指标称为自由指标。本书
讨论的是三维矢量空间及其对偶空间，除特别声明外 $i,j,k,l,m,n$ 等指标取 $1,2,3$，二维子空间上
的指标 $\alpha,\beta,\gamma,\delta$ 及其带点形式 $\dot\alpha,\dot\beta$ 等取 $1,2$。$\delta^i_j$ 是恒等张量的分量，
$\delta_{ij}$ 只在标准正交坐标系中有 $1,1,1$ 形式的对角分量，两者不同。

\item \textbf{点矢与线矢。} 小写字母 $a,b,c,p,q,s$ 记矢量（代表平面上的点，称点矢），
$k,l,m,n,u,v$ 记对偶矢量（代表平面上的直线，称线矢），$f,g,h$ 记二阶张量，$A,B,C$ 记对应的
几何对象。矢量与对偶矢量的自然缩并写为 $a\cdot l=a^il_i$。

\item \textbf{度规与升降指标。} 一个对称可逆的 $(0,2)$ 型张量 $g_{ij}$ 称为（矢量空间的）度规，
其逆记为 $(g^{-1})^{ij}$；当讨论中只有一条圆锥曲线并把它取作度规时，简写为 $g^{ij}$，并用它
升降指标：$a_i=g_{ij}a^j$，$l^i=g^{ij}l_j$。

\item \textbf{点乘、叉乘与混合积。} 有了度规之后定义
\[
a\cdot b=g_{ij}a^ib^j,\qquad
(A\times B)_i=\epsilon_{ijk}A^jB^k,\qquad
[A\,B\,C]=\epsilon_{ijk}A^iB^jC^k .
\]
$\epsilon_{ijk}$ 是取定 $\epsilon_{123}=1$ 的三阶反对称张量，$\epsilon^{ijk}$ 由 $\epsilon^{ijk}\epsilon_{ijk}=6$ 确定。

\item \textbf{不写指标的写法。} 若干个张量相邻并列表示指标直接并列而不缩并；两元素之间加点表示
缩并相邻的一个指标；两元素之间加冒号表示缩并相邻的两个指标。例如 $ga$ 的分量为 $g_{ij}a^k$，
$g\cdot a$ 的分量为 $g_{ij}a^j$，而 $f:g=f^{ij}g_{ij}$ 是实数。特别地，两个点（或两条线）的
\textbf{对称积}写为
\[
(a b)^{ij}=a^ib^j+a^jb^i,\qquad
(m n)_{ij}=m_in_j+m_jn_i .
\]

\item \textbf{规范与几何意义。} 三维矢量到平面点的对应是多对一的，故用 $a$ 与 $\lambda a$
（$\lambda\neq0$）表示同一个点。若一个等式在任一元素作这样的替换后仍成立，称它是规范不变的；
规范不变的等式与命题才描述平面射影几何的对象与关系。

\item \textbf{几何层次。} 只用到点矢、线矢及其与坐标系无关的运算的几何称为射影几何；在射影几何上
附加无穷远结构 $\linf$ 后称为仿射几何；再附加正交结构 $\theta$ 后称为正交几何（欧氏几何）；最后
固定坐标系即得解析几何。层次越高可用的运算越多，但表达式越复杂；解题时应在能说清问题的
最低层次中运算。

\item \textbf{二次曲线。} 对称的 $(2,0)$ 型张量 $f^{ij}$ 称为二次点，对称的 $(0,2)$ 型张量 $h_{ij}$
称为二次线，二者统称二次曲线。它们之间的二次点乘与二次叉乘分别为
\[
f:h=f^{ij}h_{ij},\qquad
(f\times h)_{kl}=\epsilon_{imk}\epsilon_{jnl}f^{ij}h^{mn},\qquad
(f\times h)^{kl}=\epsilon^{imk}\epsilon^{jnl}f_{ij}h_{mn},
\]
其中前者的结果是对称 $(0,2)$ 型张量，后者是对称 $(2,0)$ 型张量。不可逆的二次曲线称为退化的。

\item \textbf{编号。} 定义、定理、命题、例、注、习题统一按章连续编号；公式按章编号。

\end{enumerate}

\section*{关于本书的取材}
\addcontentsline{toc}{section}{关于本书的取材}

本书是在一组以三维矢量方法处理圆锥曲线的讲稿基础上重写而成的。重写时删去了原稿中关于
写作缘由与个人感受的叙述，删去了全部插图（图形内容改由文字或由坐标计算给出），并把原稿中
若干处含混或错误的陈述作了更正。原稿中标注为“写得不满意、暂不发表”的第十七节未收入；
原稿第八部分（例题）标注为“已过时”，亦未收入；原稿中关于立体几何与二次曲面的两节
（第二十一、二十二节）属于另一题材，本书不再涉及。


\tableofcontents

\mainmatter

\part{三维矢量的语言}
% 第一节：矢量与对偶矢量
\chapter{矢量与对偶矢量}

\section{矢量的坐标无关描述}

在中学的立体几何中，一个矢量常被写成有序数组 $(x,y,z)$。这种说法既有道理，又不准确。
准确的地方在于，对任一个三维矢量空间总可以建立坐标系，从而把矢量写成有序数组；不准确的地方
在于，这样的坐标系有无数种，没有哪一种比别的更特殊。既然没有理由选取某一个坐标系，最好的
做法就是一个也不选；一个也不选，就不能把矢量写成有序数组。

矢量和有序数组在概念上是有区别的：有序数组是在矢量上附加了一个结构（一个指定的坐标系）之后
得到的东西。物理中有类似的情形：就基本物理规律而言一切惯性系平权，然而在地面上讨论问题时
地面参考系是特殊的，这时便相当于给物理附加了一个结构。

那么什么是矢量？粗略地说，抽象的矢量就是去掉了坐标系这一附加结构的有序数组。精确的说法如下。

\begin{definition}[矢量空间]
设 $V$ 是一个集合，在 $V$ 上定义了加法 $V\times V\to V$ 与数乘 $\R\times V\to V$，它们满足
交换律、结合律、分配律以及存在零元、存在负元等通常的运算律。若存在 $V$ 中的 $n$ 个元素
$e_1,\dots,e_n$，使得 $V$ 的每个元素都能唯一地写成它们的线性组合，则称 $V$ 为 $n$ 维矢量空间，
称其元素为矢量，称 $\{e_1,\dots,e_n\}$ 为 $V$ 的一组基。
\end{definition}

定义中断言了基的存在，却没有指定基是哪一个。因此，在抽象的矢量空间中除了加法与数乘之外
并没有别的东西，特别地，无法自然地定义点乘。点乘是由坐标系定义的，因而无法用它来判断两个
矢量是否垂直。能够只用加法与数乘定义的概念是共线与共面。

\begin{definition}[共线与共面]
若存在不全为零的实数 $\alpha,\beta$ 使得 $\alpha A+\beta B=0$，称矢量 $A,B$ 共线；若存在
不全为零的实数 $\alpha,\beta,\gamma$ 使得 $\alpha A+\beta B+\gamma C=0$，称矢量 $A,B,C$ 共面。
\end{definition}

\section{对偶空间与对偶基}

\begin{definition}[线性函数与对偶空间]
映射 $f:V\to\R$ 称为 $V$ 上的函数。若对任意 $A,B\in V$ 与任意实数 $\alpha,\beta$ 恒有
\[
f(\alpha A+\beta B)=\alpha f(A)+\beta f(B),
\]
则称 $f$ 为 $V$ 上的线性函数，简称线性函数。$V$ 上全体线性函数构成的集合记为 $V^*$，
在其中定义加法与数乘
\[
(f+g)(A)=f(A)+g(A),\qquad (\alpha f)(A)=\alpha f(A).
\]
在这两种运算下 $V^*$ 构成矢量空间，称为 $V$ 的对偶空间。
\end{definition}

线性函数的运算律是逐点定义的，因此交换律、分配律等自动成立，$V^*$ 确实构成矢量空间。

\begin{proposition}[对偶基]
设 $V$ 是 $n$ 维矢量空间，$\{e_1,\dots,e_n\}$ 是 $V$ 的一组基。则存在 $V^*$ 中唯一的一组元素
$\{e^1,\dots,e^n\}$ 满足
\[
e^i(e_j)=\delta^i_j=
\begin{cases}
1,& i=j,\\
0,& i\neq j .
\end{cases}
\]
这组元素构成 $V^*$ 的一组基，称为 $\{e_i\}$ 的对偶基。
\end{proposition}

\begin{proof}
设 $A=\sum_i A^i e_i$。对给定的线性函数 $f$，由线性性
\[
f(A)=\sum_i A^i f(e_i),
\]
故 $f$ 由 $n$ 个实数 $f(e_1),\dots,f(e_n)$ 唯一确定；反过来任给 $n$ 个实数都可以这样定义一个线性函数。
所以 $\dim V^*=n$。再验证 $\{e^i\}$ 的线性无关性与生成性：若 $\sum_i c_i e^i=0$，把它作用于 $e_j$
得 $c_j=0$，故线性无关；又对任意 $f$，记 $f_i=f(e_i)$，则 $f$ 与 $\sum_i f_i e^i$ 在基矢 $e_j$ 上
取值相同，从而作为线性函数相等。
\end{proof}

把 $f(e_i)$ 简记为 $f_i$，则任一 $f\in V^*$ 可唯一地写成
\[
f=\sum_i f_i e^i,\qquad f(A)=\sum_i f_iA^i .
\]

\section{对偶对与点的表示}

上面构造中 $V$ 与 $V^*$ 的地位并不对称：一个是矢量空间，一个是它上面的线性函数空间。但函数
$f$ 作用于 $A$ 得到 $f(A)$，也可以理解为把 $A$ 看作定义在 $V^*$ 上的函数，作用于自变量 $f$ 得到
同一个数。既然两者的地位可以互换，就把 $f(A)$ 写成
\[
f\cdot A=f_iA^i .
\]
这一运算很像点乘，但与点乘有一个本质区别：

\begin{remark}
矢量与矢量之间、对偶矢量与对偶矢量之间没有自然的缩并，只有矢量与对偶矢量之间才有。因此
形如 $A^iB^i$ 或 $A_iB_i$ 的记号没有意义，而 $A^iB_i$ 有意义。
\end{remark}

这一区别正是后文把矢量解释为点、把对偶矢量解释为直线的起点：点是矢量，直线是对偶矢量，
“点与直线有包含关系”这一几何事实将表现为它们之间的缩并。
   % 矢量与对偶矢量
% 第二节：张量与指标缩并
\chapter{张量与指标缩并}

\section{指标的上下与自由指标}

第一章说明了：只有矢量与对偶矢量之间才有自然的缩并 $f\cdot A=f_iA^i$，这一运算与坐标系的
选取无关。后文要用一个二阶张量表示圆锥曲线，为此先把矢量与对偶矢量的分量记号推广到
一般的张量。

张量不依赖于坐标系的定义要借助多重线性映射，比较抽象。这里始终假定已取定一组基底，
用分量的观点来处理张量；这样得到的记号与坐标无关的说法是一致的，而在运算中更为直接。

设 $\{e_i\}$ 是矢量空间的一组基，矢量 $A$ 在该基底下写成
\[
A=\sum_i A^i e_i ,
\]
$A^i$ 称为 $A$ 的分量；设 $\{e^i\}$ 是对偶基，对偶矢量 $f$ 写成 $f=\sum_i f_i e^i$，$f_i$
称为 $f$ 的分量。为了在记号上区分二者，作如下约定：

\begin{definition}[上下指标的约定]
分量指标满足下列约定：
\begin{itemize}
\item 矢量的分量指标写在右上，称为上指标；
\item 对偶矢量的分量指标写在右下，称为下指标；
\item 单独出现的指标默认取遍 $1,\dots,n$，称为自由指标；
\item 给出 $n$ 个实数 $A^i$（或 $f_i$），就认为给出了相应的矢量（或对偶矢量）。
\end{itemize}
\end{definition}

在上下指标的记号下，$f$ 作用于 $A$ 得到
\[
f\cdot A=f_1A^1+\cdots+f_nA^n=\sum_i f_iA^i .
\]
也就是说：上指标与下指标之间可以缩并求和，而两个上指标之间或两个下指标之间不能这样求和。

\section{爱因斯坦求和约定}

在书写中，求和号常常不便。因此约定：

\begin{definition}[爱因斯坦求和约定]
在一个表达式中，若某个指标字母恰好出现两次，并且一次在上、一次在下，则称它为哑指标，
并约定对该指标从 $1$ 到 $n$ 求和，求和号省略不写。只出现一次的指标称为自由指标。
\end{definition}

既然哑指标意味着求和，用哪个字母代表哑指标就没有区别：$f_iA^i$ 与 $f_kA^k$ 表示同一个数。
例如 $f_iA^jB^i$ 表示一个矢量，其中 $i$ 是哑指标、$j$ 是自由指标，它的第 $j$ 个分量是
\[
(f_1B^1+\cdots+f_nB^n)A^j .
\]

需要注意的是，形如 $A^iB^i$、$A^iB_iC^iD_i$、$A^iB_iC_i$ 的记号没有意义：前两者中同一个
指标出现了两次以上，后者中 $i$ 两次都写在下方。因此，像
$A^1B_1C_1+\cdots+A^nB_nC_n$ 这种逐分量相乘再求和的式子不能直接写成求和约定的形式；
不过由几何的坐标无关性，后面需要的运算绝大多数都能用求和约定表示。

\section{张量的分量与类型}

$f_iA^j$ 有两个自由指标 $i,j$，各自可取 $1,\dots,n$，因此它由 $n^2$ 个实数组成。有 $m$ 个
自由指标的对象称为 $m$ 阶张量；二阶张量一般写成 $M_i{}^j$，含义是每一组 $(i,j)$ 对应一个
实数作为分量，其中 $i$ 是下指标、$j$ 是上指标。

\begin{definition}[(k,l) 型张量]
有 $k$ 个上指标、$l$ 个下指标的 $k+l$ 阶张量称为 $(k,l)$ 型张量。
\end{definition}

例如 $A^{ij}{}_k{}^l$ 是 $(3,1)$ 型张量。一个张量的型不仅取决于上、下指标的个数，指标排列
的顺序也重要；代表指标的字母则无关紧要。

区分上、下指标是必要的。例如设 $M_i{}^j=f_iA^j$，则 $M_i{}^j$ 的上指标可以与对偶矢量
$u_j$ 缩并：
\[
M_i{}^ju_j=f_iA^ju_j ,
\]
结果是一个对偶矢量；而形如 $M_{ij}$ 的张量不能与 $u_j$ 缩并，$f_iA_ju_j$ 是没有意义的。
一般地，一个指标是上指标还是下指标，决定了它能与对偶矢量还是矢量缩并。

二阶张量容易与矩阵混淆，因为二者都有两个指标。区别在于：张量是几何概念，矩阵是代数概念。
借助坐标系可以把二阶张量写成矩阵（正如把矢量写成有序数组），但意义不同：矩阵不区分上、
下指标，把不同类型的二阶张量写成矩阵会丢失类型信息。

\section{缩并}

\begin{definition}[缩并]
设 $T$ 是 $(k,l)$ 型张量，$k\ge1$，$l\ge1$。取出它的一个上指标与一个下指标，令二者取同一个
字母并对该指标求和，所得结果是 $(k-1,l-1)$ 型张量。这一运算称为缩并。
\end{definition}

例如对 $A^{ij}{}_k{}^l$ 的第 $1$ 个上指标 $i$ 与第 $1$ 个下指标 $k$ 缩并，得到
\[
B^{jl}=A^{ij}{}_i{}^l ,
\]
其中 $j,l$ 仍是自由指标。缩并是矢量与对偶矢量之间缩并的推广。

\section{(1,1) 型张量与矩阵}

$(1,1)$ 型张量的性质可以完全用矩阵的语言复述。

\begin{proposition}[(1,1) 型张量的五条性质]
设 $M^i{}_j$ 是 $(1,1)$ 型张量，则
\begin{enumerate}
\item 它把矢量映为矢量：$B^i=M^i{}_jA^j$，因此是矢量空间到自身的线性映射；
\item 它把对偶矢量映为对偶矢量：$B_j=M^i{}_jA_i$；
\item 两个 $(1,1)$ 型张量可以自然地相乘，乘积仍是 $(1,1)$ 型张量：
$M^i{}_k=A^i{}_jB^j{}_k$；
\item 存在 $(1,1)$ 型张量 $\delta^i_j$，它作用于矢量得到该矢量本身，作用于对偶矢量也得到该
对偶矢量本身，与任意 $(1,1)$ 型张量相乘仍得到该张量，称为恒等张量；
\item 它的上指标与下指标可以缩并，得到一个实数 $M^i{}_i$，称为 $M$ 的迹。
\end{enumerate}
\end{proposition}

这五条性质依次对应于矩阵与列矢量相乘、行矢量与矩阵相乘、矩阵乘法、单位矩阵与矩阵的迹。
因此，把矩阵看作线性变换时，矩阵乘法正对应于 $(1,1)$ 型二阶张量的运算。

\begin{remark}
把第 (1)(2) 条分别称为矩阵的“左乘”与“右乘”，依赖于把 $M^i{}_j$ 写成矩阵时约定 $i$ 为行
指标、$j$ 为列指标：此时 $B^i=M^i{}_jA^j$ 就是矩阵乘列矢量，而 $B_j=M^i{}_jA_i$ 就是行矢量
乘矩阵。
\end{remark}

在线性代数中，矩阵还可以用来表示二次型。二次型那样的对象不是 $(1,1)$ 型张量，而更接近
$(0,2)$ 型或 $(2,0)$ 型张量，这是下一章的内容。

\begin{exercise}
把 $f_iA^iB^j$ 按求和约定展开，指出其中的哑指标与自由指标，并说明结果是什么类型的量，
它的分量是什么。
\end{exercise}
   % 张量与指标缩并
% 第三节：升降指标、点乘和叉乘
\chapter{升降指标、点乘与叉乘}

\section{对称张量与逆张量}

前两章指出，只有矢量与对偶矢量之间才有自然的缩并，矢量与矢量之间没有。要在矢量与矢量
之间定义点乘，必须引入附加结构。本章引入的附加结构是一个对称可逆的 $(0,2)$ 型张量。

\begin{definition}[对称张量与反对称张量]
张量 $A_{ij}$ 称为对称的，若 $A_{ij}=A_{ji}$；称为反对称的，若 $A_{ij}=-A_{ji}$。
\end{definition}

在三维空间中，$A_{ij}=A_{ji}$ 意味着 $A_{12}=A_{21}$、$A_{13}=A_{31}$、$A_{23}=A_{32}$。
这个定义写在某一个基底下。不过 $A_{ij}B^iC^j$ 作为 $B,C$ 的二次函数与坐标系无关，而
$A_{ij}=A_{ji}$ 等价于
\[
A_{ij}B^iC^j=A_{ji}B^iC^j
\]
对一切矢量 $B,C$ 成立。因此对称性是张量本身的性质，而不只是它在一组基底下分量的性质。

对称张量 $A_{ij}$ 可以构造一个线性映射 $\tau$，把矢量映为对偶矢量：
\[
B^j\longmapsto \tau(B)_i=A_{ij}B^j ,
\]
其中一个指标与 $B$ 缩并，另一个指标作为对偶矢量的自由指标。

\begin{definition}[逆张量]
设 $A_{ij}$ 是 $(0,2)$ 型张量。若存在 $(2,0)$ 型张量 $C^{ij}$ 使
\[
C^{ij}A_{jk}=\delta^i_k ,
\]
则称 $A_{ij}$ 可逆，并称 $C^{ij}$ 为它的逆张量。
\end{definition}

定义中的条件来自逆映射：若 $\tau$ 有逆映射 $\tau^{-1}$，由 $\tau^{-1}\bigl(\tau(B)\bigr)=B$ 得
\[
B^i=C^{ij}A_{jk}B^k ,
\]
它对一切矢量 $B$ 成立，故 $(1,1)$ 型张量 $C^{ij}A_{jk}$ 是恒等映射，即 $C^{ij}A_{jk}=\delta^i_k$；
反过来，若此式成立，则 $\tau^{-1}$ 由 $C$ 给出，$\tau$ 可逆。

\begin{remark}
若 $A_{ij}$ 对称，则它的逆 $C^{ij}$ 也对称。事实上 $C^{ik}A_{kl}C^{jl}=\delta^i_lC^{jl}=C^{ji}$，
而把其中的 $A_{kl}$ 换成 $A_{lk}$ 并交换哑指标 $k,l$，同一个表达式又等于 $C^{ij}$；故
$C^{ij}=C^{ji}$。
\end{remark}

\section{升降指标}

有了对称可逆的张量 $A_{ij}$ 及其逆，就在矢量与对偶矢量之间得到两个互逆的映射：矢量 $B^i$
对应于对偶矢量 $A_{ij}B^j$，对偶矢量 $f_i$ 对应于矢量 $C^{ij}f_j$。讨论其中一个时总同时在
讨论另一个，不妨把它们看作同一个几何对象的两种写法，并引入记号
\[
B_i=A_{ij}B^j,\qquad f^i=C^{ij}f_j .
\]
直观地说，$A_{ij}$ 用来下降指标，$C^{ij}$ 用来上升指标。对任一张量可以逐指标地进行升降，
例如对 $D^i{}_{jk}$ 上升第二个指标，得到
\[
D^{ij}{}_k=C^{jl}D^i{}_{lk} .
\]

\begin{example}
计算 $A_{ij}$ 自身升降指标的结果。上升第一个指标：
\[
A^i{}_j=C^{ik}A_{kj}=\delta^i_j ,
\]
得到恒等张量。上升两个指标：
\[
A^{ij}=C^{ik}C^{jl}A_{kl}=C^{ik}\delta^j_k=C^{ij},
\]
得到 $A_{ij}$ 的逆。因此以后不再使用记号 $C^{ij}$，而用 $A^{ij}$ 表示 $A_{ij}$ 的逆张量。
\end{example}

\begin{definition}[度规]
在矢量空间及其对偶空间中选定一个对称可逆的 $(0,2)$ 型张量用于升降指标，就称它为这个
空间的度规，通常记为 $g_{ij}$，其逆记为 $g^{ij}$。
\end{definition}

度规的字面意思是“尺子”：有了度规才能定义两个矢量的点乘，而矢量与自身的点乘即长度的平方。
在这个意义上，长度的概念建立在度规之上。

\section{号差}

\begin{theorem}[标准形式与号差]
设 $A_{ij}$ 是 $(0,2)$ 型对称张量。则存在一组基底，使 $A_{ij}$ 在该基底下满足：$i\ne j$ 时
$A_{ij}=0$，$i=j$ 时分量为 $+1$、$-1$ 或 $0$。其中 $+1$、$-1$、$0$ 的个数由 $A$ 本身决定，
与所取基底无关，称为 $A$ 的号差。
\end{theorem}

\begin{proof}[证明要点]
用配方法把二次型 $A_{ij}X^iX^j$ 化为平方和 $\sum_k\lambda_k(Y^k)^2$（$\lambda_k\ne0$），再把
各系数化为 $\pm1$，即得上述基底。若有两种这样的标准形式，设其中 $+1$ 的个数分别为 $p,p'$，
不妨 $p>p'$。在第一种标准形中取前 $p$ 个基矢张成的子空间，在第二种中取后 $n-p'$ 个基矢张成
的子空间，二者维数之和 $p+n-p'>n$，故有非零的公共矢量 $X$。按第一种标准形，$A(X,X)>0$；
按第二种标准形，$A(X,X)\le0$，矛盾。故 $p=p'$，同理 $-1$ 与 $0$ 的个数也分别相同，号差唯一。
\end{proof}

\begin{corollary}
$A_{ij}$ 可逆当且仅当它的号差中 $0$ 的个数为零。
\end{corollary}

\begin{remark}
应当注意不要把这一结论说过头。准确的表述是：号差是\emph{完全}不变量，即两个同维数的对称
$(0,2)$ 型张量能由一个线性变换互相转化，当且仅当它们的号差相同；但坐标无关的性质并不止
号差一个，例如对称性、反对称性也都与基底无关。
\end{remark}

\section{点乘、模方与定号性}

\begin{definition}[点乘与模方]
设 $g_{ij}$ 是度规，两个矢量 $A,B$ 的点乘定义为
\[
A\cdot B=A^iB_i=A_iB^i=g_{ij}A^iB^j .
\]
矢量与自身的点乘 $A^2=A\cdot A=g_{ij}A^iA^j$ 称为 $A$ 的模方，即长度的平方。
\end{definition}

度规要求对称，正是为了使点乘满足交换律：
\[
A\cdot B=g_{ij}A^iB^j=g_{ji}A^iB^j=g_{ji}B^jA^i=B\cdot A .
\]

矢量的模方不一定是正数，它的正负性反映了度规本身的性质，可以用号差描述。

\begin{definition}[正定、负定与不定]
设度规的号差为 $(p,q,r)$，其中 $p,q,r$ 分别是 $+1,-1,0$ 的个数。
\begin{enumerate}
\item 若对任意非零矢量 $A$ 都有 $A^2>0$，称度规正定，此时 $q=r=0$；
\item 若对任意非零矢量 $A$ 都有 $A^2<0$，称度规负定，此时 $p=r=0$；
\item 若 $A^2$ 可取正值、负值，也可为零，称度规不定。
\end{enumerate}
\end{definition}

\begin{remark}
号差中有 $+1$、有 $0$、没有 $-1$ 的情形常被笼统地称作“非负”，但它并不是正定的：此时恒有
$A^2\ge0$，且存在非零矢量使 $A^2=0$，度规是退化的。正确的名称是半正定（非负定）；正定
要求号差中 $0$ 的个数为零。
\end{remark}

若度规负定，则 $-g_{ij}$ 正定，二者给出的点乘只相差一个符号。因此以后一般不讨论负定度规，
必要时改变符号即可。

\section{例：欧氏度规与闵氏度规}

\begin{example}[欧氏度规]
在三维欧几里得空间中取一组正交归一基底，度规的分量为
\[
g_{11}=g_{22}=g_{33}=1,\qquad g_{ij}=0\quad(i\ne j),
\]
即 $g_{ij}=\delta_{ij}$。这时非零矢量 $A^i=(x,y,z)$ 的模方为
\[
A^2=x^2+y^2+z^2>0 ,
\]
故该度规正定，号差为 $(3,0,0)$。
\end{example}

\begin{remark}
这里的 $\delta_{ij}$ 与前面常用的 $\delta^i_j$ 不同。$\delta^i_j$ 是恒等 $(1,1)$ 型张量，
它在任何基底下分量都是 $\delta^i_j$；而 $\delta_{ij}$ 只在正交归一基底下分量才形如
$\delta_{ij}$，换一组一般基底，它的分量是该基底下的欧氏度规矩阵，不再等于 $\delta_{ij}$。
\end{remark}

\begin{example}[闵氏度规]
四维闵氏时空在某个惯性参考系下的度规分量为
\[
g_{00}=c^2,\qquad g_{11}=g_{22}=g_{33}=-1,
\]
其余分量为 $0$。这时 $A^i=(t,x,y,z)$ 的模方为
\[
A^2=c^2t^2-x^2-y^2-z^2 ,
\]
可取正值、负值或零，故该度规不定，号差为 $(1,3,0)$。
\end{example}

\begin{example}
平面上的圆锥曲线可以放在三维空间中用不定度规处理，这是后文的内容。
\end{example}

\section{三阶反对称张量}

前面用二阶对称张量定义了 $n$ 维空间中的点乘。下面用一个三阶反对称张量定义三维空间中的
叉乘；叉乘只在三维空间中存在。

\begin{definition}[反对称张量]
$(0,3)$ 型张量 $\epsilon_{ijk}$ 称为反对称的，若交换其中任意两个指标后分量变号，例如
$\epsilon_{ijk}=-\epsilon_{jik}$。
\end{definition}

与对称性一样，反对称性在一组基底下成立就在任何基底下成立：$\epsilon_{ijk}B^iC^jD^k$ 与
基底无关，而反对称性等价于它对任意两个自变量交换后变号。

由反对称性，若三个指标中有两个相等，则该分量为零。例如交换 $\epsilon_{112}$ 的前两个指标得
$\epsilon_{112}=-\epsilon_{112}$，故 $\epsilon_{112}=0$。因此非零分量的三个指标必使 $1,2,3$
各出现一次。设 $\epsilon_{123}=a$，则
\[
\epsilon_{123}=\epsilon_{231}=\epsilon_{312}=a,\qquad
\epsilon_{213}=\epsilon_{132}=\epsilon_{321}=-a .
\]

\begin{remark}
列举分量时容易写成自相矛盾的形式（例如同时写出 $\epsilon_{312}=-a$ 与 $\epsilon_{312}=a$）。
正确的取法是：偶排列 $(123),(231),(312)$ 的分量都等于 $a$，奇排列 $(213),(132),(321)$
的分量都等于 $-a$。
\end{remark}

于是三阶反对称张量的分量除一个待定常数 $a$ 外完全确定，而 $a$ 的变化只是给张量乘上一个
常数。取 $a=1$，称为标准反对称张量。$a=1$ 只在部分基底下成立：在别的基底下写出 $\epsilon$
的分量，$a$ 一般不再是 $1$；这一点对以下运算无关紧要，因此也可以简单地说 $\epsilon_{ijk}$
是三维空间自带的结构。

\begin{definition}[$\epsilon^{ijk}$]
指标在上的三阶反对称张量 $\epsilon^{ijk}$ 不借助度规升降指标来定义，而要求
\[
\epsilon^{ijk}\epsilon_{ijk}=6 .
\]
\end{definition}

这一条件唯一地确定 $\epsilon^{ijk}$。设 $\epsilon^{ijk}=\lambda\,\eta^{ijk}$，其中 $\eta^{ijk}$
是标准反对称张量，则六个非零分量各贡献 $\lambda$，即 $\epsilon^{ijk}\epsilon_{ijk}=6\lambda$，
故 $\lambda=1$：取 $\epsilon^{123}=1$ 时等式成立，并且满足等式的 $\epsilon^{ijk}$ 唯一。

\begin{proposition}[$\epsilon$ 的缩并恒等式]
\[
\epsilon_{ijk}\epsilon^{imn}=\delta^m_j\delta^n_k-\delta^n_j\delta^m_k .
\]
\end{proposition}

\begin{proof}
左端对 $i$ 求和。若 $j=k$，左边因 $\epsilon_{ijk}$ 对 $j,k$ 反对称而为 $0$，右边也为 $0$。
若 $j\ne k$，两边都对 $j,k$ 反对称，故只需验证 $j=2,k=3$ 的情形。此时左边只有 $i=1$ 时非零，
值为 $\epsilon^{1mn}$；右边为 $\delta^m_2\delta^n_3-\delta^n_2\delta^m_3$。当 $(m,n)=(2,3)$ 时
两边都等于 $1$，当 $(m,n)=(3,2)$ 时两边都等于 $-1$，其余情形两边都为 $0$。故恒等式成立。
\end{proof}

\begin{proposition}[$\epsilon$ 与行列式]
\[
\epsilon_{ijk}\epsilon^{lmn}=
\det\begin{pmatrix}
\delta^l_i & \delta^m_i & \delta^n_i\\
\delta^l_j & \delta^m_j & \delta^n_j\\
\delta^l_k & \delta^m_k & \delta^n_k
\end{pmatrix}.
\]
\end{proposition}

\begin{proof}
两边都对 $(i,j,k)$ 反对称，也都对 $(l,m,n)$ 反对称（交换 $l,m$ 相当于交换行列式的两列）。
而在 $(i,j,k)=(l,m,n)=(1,2,3)$ 时两边都等于 $1$。反对称张量的全部分量由它在 $(1,2,3)$ 处的值
确定，故恒等式成立。
\end{proof}

\begin{remark}
这两个公式的作用是：表达式中每出现一对 $\epsilon$，就可以把它们换成 $\delta$ 的组合，从而
一次消去两个 $\epsilon$。反复使用后，一个表达式中剩下的 $\epsilon$ 的个数为 $0$ 或 $1$。
\end{remark}

\section{叉乘与混合积}

\begin{definition}[叉乘]
两个矢量 $A,B$ 的叉乘是一个对偶矢量，定义为
\[
(A\times B)_i=\epsilon_{ijk}A^jB^k .
\]
类似地，两个对偶矢量 $X,Y$ 的叉乘是一个矢量，定义为
\[
(X\times Y)^i=\epsilon^{ijk}X_jY_k .
\]
\end{definition}

\begin{definition}[混合积]
三个矢量 $A,B,C$ 的混合积是一个实数，定义为
\[
[ABC]=\epsilon_{ijk}A^iB^jC^k .
\]
同样可以定义三个对偶矢量的混合积。
\end{definition}

\begin{proposition}[基本性质]
叉乘与混合积具有下列性质：
\begin{enumerate}
\item 叉乘是双线性的，并且反对称：$A\times B=-B\times A$，特别地 $A\times A=0$；
\item 叉乘的结果与两个因子都缩并为零：$A\cdot(A\times B)=B\cdot(A\times B)=0$；
\item 混合积对三个自变量都是线性的，交换任意两个因子变号；当 $A,B,C$ 共面时 $[ABC]=0$。
\end{enumerate}
\end{proposition}

\begin{proof}
(1)(3) 由 $\epsilon$ 的反对称性与线性性直接得到。(2) 例如
\[
A\cdot(A\times B)=A^i(A\times B)_i=A^i\epsilon_{ijk}A^jB^k=0,
\]
因为 $\epsilon_{ijk}$ 对 $i,j$ 反对称，而 $A^iA^j$ 对 $i,j$ 对称；$B$ 的情形同理。
\end{proof}

\begin{theorem}[二重外积公式]
设 $A,B$ 是矢量，$C$ 是对偶矢量（对称地，也可以设 $A,B$ 是对偶矢量而 $C$ 是矢量），则
\[
(A\times B)\times C=(A\cdot C)B-(B\cdot C)A .
\]
\end{theorem}

\begin{proof}
设 $A,B$ 是矢量、$C$ 是对偶矢量。此时 $A\times B$ 与 $C$ 都是对偶矢量，它们的叉乘由
$(X\times Y)^i=\epsilon^{ijk}X_jY_k$ 给出，于是
\[
\bigl((A\times B)\times C\bigr)^i
=\epsilon^{ijk}(A\times B)_jC_k
=\epsilon^{ijk}\epsilon_{jmn}A^mB^nC_k .
\]
把缩并恒等式 $\epsilon_{ijk}\epsilon^{imn}=\delta^m_j\delta^n_k-\delta^n_j\delta^m_k$ 中的指标
重新命名，可得
\[
\epsilon^{ijk}\epsilon_{jmn}=\delta^i_n\delta^k_m-\delta^i_m\delta^k_n .
\]
代入得
\[
\bigl((A\times B)\times C\bigr)^i
=A^mB^nC_k\bigl(\delta^i_n\delta^k_m-\delta^i_m\delta^k_n\bigr)
=(A\cdot C)B^i-(B\cdot C)A^i ,
\]
即所证。$A,B$ 为对偶矢量而 $C$ 为矢量时证法相同。
\end{proof}

\begin{remark}
公式中 $A,B$ 是矢量而 $C$ 是对偶矢量，或者 $A,B$ 是对偶矢量而 $C$ 是矢量。只有当存在度规时，
才能把矢量与对偶矢量互相转换，从而定义矢量与对偶矢量之间的叉乘。
\end{remark}

\begin{corollary}
\[
(A\times B)\cdot(C\times D)=(A\cdot C)(B\cdot D)-(A\cdot D)(B\cdot C).
\]
特别地，取 $C=A$、$D=B$ 得
\[
(A\times B)^2=(A\cdot A)(B\cdot B)-(A\cdot B)^2 .
\]
\end{corollary}

\begin{proof}
$A\times B$ 与 $C\times D$ 都是对偶矢量，两个对偶矢量的点乘需要用度规把一个指标上升：
在 $|\det g|=1$ 的基底下（例如标准正交基底）有 $g^{ij}\epsilon_{jkl}=\epsilon^{ikl}$，于是
\[
(A\times B)\cdot(C\times D)=g^{ij}(A\times B)_i(C\times D)_j
=\epsilon^{j}{}_{ab}\epsilon_{jcd}A^aB^bC^cD^d .
\]
再由恒等式 $\epsilon^{jab}\epsilon_{jcd}=\delta^a_c\delta^b_d-\delta^a_d\delta^b_c$，得
\[
(A\times B)\cdot(C\times D)=(A\cdot C)(B\cdot D)-(A\cdot D)(B\cdot C).
\]
\end{proof}

\begin{remark}
上面的推论中出现了两个对偶矢量的点乘，必须用度规把一个指标上升。本节把 $\epsilon_{ijk}$ 与
$\epsilon^{ijk}$ 都取为标准反对称符号，这只在 $|\det g|=1$ 的基底（例如标准正交基底）下才是
由度规诱导的体积形式；在一般的非正交基底下，缩并的结果会多出 $\det g$ 因子。二重外积公式
本身只用到 $\epsilon$ 与 $\delta$，与这一约定无关。
\end{remark}

本章介绍的升降指标、点乘、叉乘、混合积与指标缩并是后文的基本工具。

\begin{exercise}
展开 $\epsilon^{ijk}\epsilon_{ijk}$，写出它的六个非零项，并验证它等于 $6$。
\end{exercise}
   % 升降指标、点乘与叉乘

\part{关系的建立}
% 第四节：齐次坐标
\chapter{齐次坐标}

\section{圆锥曲线的一般方程}

以椭圆为例，当其中心在原点、长轴与 $x$ 轴重合时，标准方程为
\[
\frac{x^2}{a^2}+\frac{y^2}{b^2}=1 .
\]
但椭圆可以有任意的位置与姿态，此时它的方程的一般形式为
\begin{equation}\label{eq:gen}
ax^2+2bxy+cy^2+2dx+2ey+f=0 .
\end{equation}
在系数取适当值时，这个方程可以表示椭圆、抛物线或双曲线；在系数取某些特殊值时，它也可能表示两条直线、一条直线或一个点，或者无解。以下先讨论一般情形，特殊情形暂不涉及。

为了消去常数项与一次项，作代换
\[
x=\frac{X}{Z},\qquad y=\frac{Y}{Z},\qquad Z\neq 0,
\]
代入 \eqref{eq:gen} 并乘以 $Z^2$，得到
\begin{equation}\label{eq:hom}
aX^2+2bXY+cY^2+2dXZ+2eYZ+fZ^2=0 .
\end{equation}
这个方程关于 $X,Y,Z$ 的每一项都是二次的。进一步引入对称的 $(0,2)$ 型张量 $g$，其分量为
\[
g_{11}=a,\quad g_{12}=g_{21}=b,\quad g_{22}=c,\quad
g_{13}=g_{31}=d,\quad g_{23}=g_{32}=e,\quad g_{33}=f,
\]
并把未知量 $X,Y,Z$ 看作矢量 $A$ 的分量 $A^1=X$、$A^2=Y$、$A^3=Z$，则 \eqref{eq:hom} 可改写为
\begin{equation}\label{eq:gAA}
g_{ij}A^iA^j=0 .
\end{equation}
若把 $g$ 取作度规并用它定义点乘，则 \eqref{eq:gAA} 可进一步简写为
\[
A^2=g_{ij}A^iA^j=0 .
\]
也就是说，圆锥曲线可以用“矢量平方为零”这一与坐标系无关的方程来描述。正是这一简化的形式提示了把平面几何的问题放到三维矢量空间中处理的思路。

\begin{remark}
方程 $ax^2+2bxy+cy^2+2dx+2ey+f=0$ 在系数退化时可能表示两条直线、一条直线、一个点，或者没有实点。这些情形对应 $g$ 的退化或方程的实解集退化，本书后文在需要时分别处理，本节只讨论 $g$ 非退化的一般情形。
\end{remark}

\section{平面上的点与三维矢量的对应}

上节的改写表明，平面上的点可以借助三个分量 $X,Y,Z$ 来表示。下面把这个对应关系明确地建立起来。整个过程分三步。

第一步：在三维矢量空间中建立坐标系，把任意矢量写成
\[
A^i=(X,Y,Z).
\]

第二步：利用这个坐标系，把平面上的点 $(x,y)$ 与某个确定方向上的全体非零三维矢量对应起来。约定的对应关系是：点 $(x,y)$ 对应所有分量形如
\[
A^i=(\lambda x,\lambda y,\lambda),\qquad \lambda\in\R\setminus\{0\}
\]
的三维矢量；反之，矢量 $(X,Y,Z)$ 以及任何形如 $(\lambda X,\lambda Y,\lambda Z)$ 的矢量都对应平面上的点
\[
\left(\frac{X}{Z},\frac{Y}{Z}\right).
\]
这个对应是多对一的：平面上的每个点对应一条过原点的直线（去掉原点）上的全部非零矢量，而每个这样的矢量只对应一个点。

\begin{definition}[齐次坐标]
按上述约定，与平面上的点 $(x,y)$ 相对应的三维矢量的分量 $(\lambda x,\lambda y,\lambda)$（$\lambda\neq 0$）称为该点的一组齐次坐标；反过来，矢量 $(X,Y,Z)$（$Z\neq 0$）对应的点是 $(X/Z,Y/Z)$。
\end{definition}

当 $Z=0$ 时上述对应没有定义，因为分母为零，此时矢量不对应任何平面上的点。这一类矢量所对应的几何对象将在讨论无穷远时处理。

第三步：抹去前两步中建立坐标系所留下的一切痕迹。

在完成前两步之后，平面上的几何对象及其相互关系都可以用带有坐标系的三维矢量表示。具体地说，存在两个函数 $x(A),y(A)$，它们给出矢量 $A$ 所对应点的坐标；而平面几何的关系总可以写成关于 $x,y$ 的方程，因此也就可以改写成关于 $A$ 的方程。例如曲线 $y=\sin x$ 等价于关于 $A$ 的方程
\[
y(A)=\sin x(A).
\]
然而，对这样的曲线使用三维矢量的方法求解没有任何好处：求解这个方程终究要求出矢量的分量，再让分量满足 $x(A),y(A)$ 的复杂关系，这层包装属于画蛇添足。

三维矢量方法的优势只出现在下述情形：关于矢量的方程由简单自然的运算组成，其中尤其重要的是与坐标系无关的运算，例如加法、数乘，以及定义了度规之后的点乘。这就引出下面的定义。

\begin{definition}[射影几何学]
如果平面几何量之间的关系经过前两步的对应改写为关于三维矢量的等式之后，所得的等式与坐标系无关，则称这个几何关系属于射影几何学。
\end{definition}

对于射影几何学的问题，用三维矢量方法求解是最简便的做法。换言之，只有当第三步（抹去坐标系的一切痕迹）可以做到时，使用三维矢量方法才具有优势；若某个几何关系不得不借助坐标系来表达，则它不属于射影几何学。
   % 齐次坐标

\part{点与线的射影几何}
% 第五节：点矢与线矢
\chapter{点矢与线矢}

\section{规范与规范不变}

上一节末尾指出，所关心的只是与坐标系无关的运算所对应的几何关系。本节就从这类运算出发。先改换记号：一般用小写字母 $a,b,c,p,q,s$ 表示矢量，用 $k,l,m,n,u,v$ 表示对偶矢量，用 $f,g,h$ 表示张量；$A$ 表示 $a$ 所对应的点，其余类推。

首先讨论矢量的加法，这里立刻会遇到麻烦。设 $a,b$ 是两个不共线的矢量，分别表示平面上的点 $A,B$，问 $a+b$ 有什么几何意义。

\begin{example}
取平面解析几何中的两点 $A(0,0)$、$B(1,0)$。取点矢
\[
a^i=(0,0,1),\qquad b^i=(1,0,1),
\]
它们分别表示 $A,B$，于是 $(a+b)^i=(1,0,2)$，按上一节的规则它表示点 $(1/2,0)$。但是表示同一个点的点矢并不唯一：也可以取
\[
a^i=(0,0,\lambda),\qquad b^i=(\mu,0,\mu),\qquad \lambda,\mu\neq 0,
\]
它们仍然分别表示 $A,B$，而此时
\[
(a+b)^i=(\mu,0,\lambda+\mu),
\]
对应的点是 $(\mu/(\lambda+\mu),0)$。随着 $\lambda,\mu$ 取不同的值，所得的点也不同。
\end{example}

这个例子说明，仅知道 $a,b$ 各自所对应的点，并不能确定 $a+b$ 所对应的点。要使它被确定，就不能允许 $a,b$ 沿各自的方向任意伸缩；为此需要附加条件来限制这些矢量。这种限制条件称为规范。

\begin{definition}[标准规范]
要求矢量在给定参考系下的第三个分量为 $1$，即 $A^3=1$，这一规范称为标准规范。
\end{definition}

有的关系只在选取了特定规范时才成立，有的定义只在选取了特定规范时才有意义；但也有一类关系，无论选取哪一种规范都成立。这类关系称为规范不变的，它们正是所要研究的对象。

具体地说，设有若干元素，例如 $a,b,c,l,m,n,g$，它们之间满足某些等式。在任何元素被换成自身的非零倍数（即它仍然表示同一个点或同一条直线）之后，等式若仍然成立，就称这个等式是规范不变的。例如，$\kappa$ 为实数时，
\[
l\cdot a=0
\]
是规范不变的，而 $l\cdot a=\kappa$（$\kappa\neq 0$）不是规范不变的。一个等式是规范不变的，意味着它能够描述平面（确切地说是平面射影几何）中元素之间的关系；否则它可能只描述了三维矢量之间的关系。

同样地，若用若干元素构造出一个实数，而其中任一元素换成自身的非零倍数时该实数不变，则这个实数具有几何意义。例如混合积 $[abc]$ 没有几何意义，而
\[
\frac{[abc][ade]}{[abd][ace]}
\]
具有几何意义：分子关于 $a$ 是二次的、关于 $b,c,d,e$ 各是一次，分母与之相同，故任一元素伸缩时比值不变。若用若干元素构造出一个矢量或张量，而在任一元素换成自身的非零倍数时，所构造出的元素也随之变为自身的若干倍，则这个元素具有几何意义。例如 $a+b$ 没有几何意义，而叉乘 $a\times b$ 具有几何意义，因为 $a\to\lambda a$ 时 $a\times b\to\lambda(a\times b)$。

若由若干元素构造出一个命题，而在任一元素换成自身的非零倍数时命题仍然成立，则称这个命题具有几何意义。例如命题“三个矢量 $a,b,a+b$ 所对应的点一定共线”：虽然在不同规范下 $a+b$ 对应不同的点，但无论取哪一种规范，$a,b,a+b$ 三个矢量总是线性相关的，因而命题总是成立。这个例子稍后还会进一步讨论。

\section{点矢与线矢}

明确了什么样的表达式才有几何意义之后，从几何意义清楚的运算入手，讨论等式
\[
a\cdot l=0 .
\]
按齐次坐标一节建立的关系，矢量 $a$ 表示一个点，但对偶矢量 $l$ 表示什么尚不清楚。把等式展开：设
\[
l_i=(A,B,C),\qquad a^i=(X,Y,Z),
\]
则
\[
l\cdot a=AX+BY+CZ=0 .
\]
两边除以 $Z$，并利用 $x=X/Z$、$y=Y/Z$，得到
\[
Ax+By+C=0 .
\]
这正是直线在解析几何中的标准方程。因此 $a\cdot l=0$ 表示点 $a$ 位于由 $l$ 的分量所确定的直线上。又因为矢量与对偶矢量的缩并与坐标系无关，可以说对偶矢量 $l$ 表示一条直线。

\begin{proposition}[点矢与线矢]
在平面射影几何中有以下对应：
\begin{enumerate}
\item 矢量 $a$ 表示一个点，因此矢量也称为点矢；
\item 对偶矢量 $l$ 表示一条直线，因此对偶矢量也称为线矢；
\item $a\cdot l=0$ 表示点 $a$ 在直线 $l$ 上，或者说直线 $l$ 经过点 $a$。
\end{enumerate}
\end{proposition}

\begin{proof}
第 2、3 条的推导已如上：把 $l_i=(A,B,C)$、$a^i=(X,Y,Z)$ 代入 $a\cdot l=0$，在 $Z\neq 0$ 时等价于 $Ax+By+C=0$，即点 $a$ 满足直线 $l$ 的方程。第 1 条是上一节对应关系的重述。
\end{proof}

值得注意的是，对偶矢量与矢量在构造上地位相同，因此直线与点地位也相同。点与直线的这种对偶性，正是后文许多命题可以成对出现的原因。

\begin{proposition}\label{prop:cross}
两个点矢的叉乘 $a\times b$ 是线矢，且规范不变，它表示过 $a,b$ 两点的直线；两个线矢的叉乘 $m\times n$ 是点矢，且规范不变，它表示 $m,n$ 两条直线的交点。
\end{proposition}

\begin{proof}
先看 $a\times b$。当 $a\to\lambda a$、$b\to\mu b$ 时，
\[
(\lambda a)\times(\mu b)=\lambda\mu\,(a\times b),
\]
故叉乘规范不变。又由叉乘的定义与混合积的反称性，
\[
(a\times b)\cdot a=[aab]=0,\qquad (a\times b)\cdot b=[bab]=0 .
\]
因此 $a\times b$ 所表示的直线既经过 $a$ 又经过 $b$；在 $a,b$ 不重合时，过这两点的直线唯一，故 $a\times b$ 正表示这条直线。

对偶的论证给出 $m\times n$：它满足
\[
(m\times n)\cdot m=0,\qquad (m\times n)\cdot n=0,
\]
即它所表示的点既在 $m$ 上又在 $n$ 上；在 $m,n$ 不平行时两条直线只交于一点，故 $m\times n$ 正表示这个交点。
\end{proof}

\begin{definition}[一次点列]
设 $a,b$ 是直线 $l$ 上两个不同的点，即 $a\cdot l=b\cdot l=0$。则对任意实数 $\lambda$，
\[
l\cdot(a+\lambda b)=l\cdot a+\lambda\,l\cdot b=0,
\]
故 $a+\lambda b$ 所对应的点都在直线 $l$ 上。当 $\lambda$ 取遍全体实数时，$a+\lambda b$ 取遍直线 $l$ 上的全体点（点 $b$ 除外）。这条直线上的全体点构成的集合称为（一次）点列。约定当 $\lambda=\infty$ 时，$a+\lambda b$ 表示点 $b$。
\end{definition}

\begin{remark}
一次点列的定义与规范的选取无关：无论选取哪一种规范，$a+\lambda b$ 都取遍直线 $l$ 上的全体点。此外，$b$ 不能写成 $a+\lambda b$（$\lambda\in\R$）的形式，因此上面对 $\lambda=\infty$ 的约定不是严格的极限，只是一个记号的约定。
\end{remark}

\begin{definition}[一次线束]
类似地，若两条直线 $m,n$ 交于点 $P$，则当 $\lambda$ 取遍全体实数时，$m+\lambda n$ 取遍过 $P$ 的全体直线。这些直线构成的集合称为一次线束，简称线束。
\end{definition}

\begin{proposition}
三个点矢的混合积 $[abc]=0$ 当且仅当 $a,b,c$ 三点共线；三个线矢的混合积为零当且仅当这三条直线共点。
\end{proposition}

\begin{proof}
由混合积与叉乘的关系 $[abc]=(a\times b)\cdot c$。若 $[abc]=0$，则点 $c$ 在直线 $a\times b$ 上，而 $a\times b$ 就是过 $a,b$ 的直线，故三点共线；反之若三点共线，则 $c$ 在过 $a,b$ 的直线上，于是 $(a\times b)\cdot c=0$。对偶地，三个线矢的混合积为零等价于它们的交点落在第三条直线上，即三线共点。
\end{proof}

\section{习题：Desargues 定理}

\begin{exercise}[Desargues 定理]
设三角形 $ABC$ 与 $A'B'C'$ 的对应顶点连线 $AA'$、$BB'$、$CC'$ 共点于 $D$，对应边 $BC$ 与 $B'C'$、$CA$ 与 $C'A'$、$AB$ 与 $A'B'$ 分别交于 $E,F,G$。求证：$E,F,G$ 三点共线。
\end{exercise}

\begin{proof}
设 $a,b,c$ 是 $A,B,C$ 三点的点矢。由于 $A,B,C$ 不共线，$[abc]\neq 0$，可以把 $a,b,c$ 取作三维矢量空间的一组基。在这组基下
\[
a^i=(1,0,0),\qquad b^i=(0,1,0),\qquad c^i=(0,0,1).
\]
设 $d^i=(x,y,z)$。由于 $A,A',D$ 共线，三个点矢 $a,a',d$ 线性相关，而 $a,d$ 已张成过 $A,D$ 的平面，故可设
\[
a'=d+\lambda a=(x+\lambda,\,y,\,z).
\]
同理设
\[
b'=(x,\,y+\mu,\,z),\qquad c'=(x,\,y,\,z+\sigma).
\]

下面用叉乘求出三条对应边的交点。叉乘取约定 $(A\times B)_i=\epsilon_{ijk}A^jB^k$，$\epsilon_{123}=1$。

先求 $E=BC\cap B'C'$。由 $b=(0,1,0)$、$c=(0,0,1)$ 得
\[
b\times c=(1,0,0),
\]
而
\[
b'\times c'=\big(y\sigma+\mu z+\mu\sigma,\,-\sigma x,\,-\mu x\big).
\]
于是
\[
e=(b\times c)\times(b'\times c')=(0,\,\mu x,\,-\sigma x).
\]

再求 $F=CA\cap C'A'$。由 $c=(0,0,1)$、$a=(1,0,0)$ 得
\[
c\times a=(0,1,0),
\]
而
\[
c'\times a'=\big(-\sigma y,\,z\lambda+\sigma x+\sigma\lambda,\,-\lambda y\big),
\]
于是
\[
f=(c\times a)\times(c'\times a')=(-\lambda y,\,0,\,\sigma y).
\]

最后求 $G=AB\cap A'B'$。由 $a=(1,0,0)$、$b=(0,1,0)$ 得
\[
a\times b=(0,0,1),
\]
而
\[
a'\times b'=\big(-\mu z,\,-\lambda z,\,\lambda y+\mu x+\lambda\mu\big),
\]
于是
\[
g=(a\times b)\times(a'\times b')=(\lambda z,\,-\mu z,\,0).
\]

计算三点的混合积：
\[
[efg]=
\begin{vmatrix}
0 & \mu x & -\sigma x\\
-\lambda y & 0 & \sigma y\\
\lambda z & -\mu z & 0
\end{vmatrix}
=xyz
\begin{vmatrix}
0 & \mu & -\sigma\\
-\lambda & 0 & \sigma\\
\lambda & -\mu & 0
\end{vmatrix}=0,
\]
最后一个行列式为零，因为它的第三行等于前两行之和的相反数。故 $[efg]=0$，即 $E,F,G$ 三点共线。
\end{proof}

\begin{remark}
上述推导并未用到 $x,y,z,\lambda,\mu,\sigma$ 全不为零的假设：混合积的表达式中每一列都含有相应的因子，行列式恒为零。因此退化的位置关系也包含在结论之中。
\end{remark}
   % 点矢与线矢
% 第六节：无穷远
\chapter{无穷远}

\section{无穷远点与无穷远直线}

由齐次坐标一节可知，分量形如 $(X,Y,0)$ 的点矢不能对应到平面上的点，因为这会使得 $X/Z$、$Y/Z$ 的分母为零。为使这样的点矢也有对应的几何对象，给平面添加一些新的元素，称为无穷远点。

由点矢与线矢一节又可知，形如 $(0,0,C)$ 的线矢不对应任何直线，因为它给出的直线方程形如 $0\cdot x+0\cdot y+C=0$。为使这样的线矢也有对应的几何对象，再给平面添加一个新的元素。

\begin{definition}[无穷远直线]
线矢 $(0,0,C)$（$C\neq 0$）所对应的直线称为无穷远直线，记作 $\linf$。在不引起混淆时也简写作 $\infty$。
\end{definition}

\begin{definition}[射影平面]
平面上的全体点，连同上面添加的全体无穷远点，合起来称为（实）射影平面。全体无穷远点构成无穷远直线 $\linf$。
\end{definition}

\section{平行线的交点}

在点矢与线矢一节中已经知道，两条不平行的直线有唯一的交点。引入无穷远点之后，平行的直线也有交点：它们的交点是无穷远点。下面的命题给出这一交点的齐次坐标。

\begin{proposition}\label{prop:parallel}
设两条直线的线矢为
\[
m_i=(A,B,C_1),\qquad n_i=(A,B,C_2),\qquad C_1\neq C_2,
\]
则两条直线平行，它们的交点为
\[
m\times n=\big(B(C_2-C_1),\,A(C_1-C_2),\,0\big)=(C_2-C_1)\,(B,-A,0),
\]
这是一个无穷远点。
\end{proposition}

\begin{proof}
两条直线的方向由 $(A,B)$ 决定，$m,n$ 的方向相同，故它们平行。直接按 $(A\times B)_i=\epsilon_{ijk}A^jB^k$ 计算：
\[
(m\times n)_1=m^2n^3-m^3n^2=B C_2-C_1B=B(C_2-C_1),
\]
\[
(m\times n)_2=m^3n^1-m^1n^3=C_1A-AC_2=A(C_1-C_2),
\]
\[
(m\times n)_3=m^1n^2-m^2n^1=AB-BA=0 .
\]
其第三个分量为零，故为无穷远点。
\end{proof}

\begin{remark}
两个分量的符号搭配容易记错，核对的方法是：直线 $Ax+By+C=0$ 的方向为 $(-B,A)$，因此上面
的无穷远点应当与 $(B,-A,0)$ 成比例，计算结果表明确实如此。同时又可见两条平行线
（$C_1\neq C_2$）交于\emph{同一个}无穷远点，因为该点只依赖于 $(A,B)$ 而不依赖于 $C$。
\end{remark}

下面的结论说明无穷远点与无穷远直线在射影平面中的基本性质。

\begin{proposition}
在射影平面中有下列结论：
\begin{enumerate}
\item 两条平行线交于同一个无穷远点；
\item 无穷远直线与任意一条直线都交于一个无穷远点；
\item 无穷远点在无穷远直线上；
\item 每条直线都经过一个无穷远点，不平行的直线经过不同的无穷远点；
\item 过两个无穷远点所作的直线是无穷远直线；
\item 设 $m$ 是一条直线，则当 $\lambda$ 取遍全体实数时，$m+\lambda\linf$ 给出所有与 $m$ 平行的直线。
\end{enumerate}
\end{proposition}

\begin{proof}
逐条验证。设直线 $m_i=(A,B,C_1)$、$n_i=(A',B',C_2)$。

(1) 平行意味着 $(A',B')=k(A,B)$，此时
\begin{align*}
(m\times n)_1&=B C_2-C_1B'=B(C_2-kC_1),\\
(m\times n)_2&=C_1A'-AC_2=A(kC_1-C_2),\\
(m\times n)_3&=AB'-A'B=0,
\end{align*}
其方向与 $(B,-A,0)$ 相同。凡方向与 $(A,B)$ 相同的直线都得到同一个点 $(B,-A,0)$，
故它们交于同一个无穷远点。

(2) 取 $\linf$ 的代表元 $(0,0,1)$，则
\[
\linf\times m=(0,0,1)\times(A,B,C)=(-B,\,A,\,0),
\]
第三个分量为零，是无穷远点。

(3) 无穷远点形如 $(X,Y,0)$，与 $\linf=(0,0,1)$ 缩并得 $0\cdot X+0\cdot Y+1\cdot 0=0$，故它在 $\linf$ 上。

(4) 由 (2)，直线 $m$ 与 $\linf$ 的交点为 $(-B,A,0)$，它只由方向 $(A,B)$ 决定；不平行的两条直线方向不成比例，所得的点也不同。

(5) 设两个不同的无穷远点为 $p=(X_1,Y_1,0)$、$q=(X_2,Y_2,0)$，则
\[
p\times q=(0,\,0,\,X_1Y_2-Y_1X_2),
\]
其第三个分量非零，故所得线矢是 $\linf$，即过 $p,q$ 的直线是无穷远直线。

(6) 取 $m_i=(A,B,C)$、$\linf=(0,0,1)$，则 $m+\lambda\linf=(A,B,C+\lambda)$。这些直线的方向都是 $(A,B)$，故都与 $m$ 平行；反过来，任何与 $m$ 平行的直线都可写成 $(A,B,C')=m+(C'-C)\linf$。故 $m+\lambda\linf$ 恰好给出所有与 $m$ 平行的直线。
\end{proof}

\section{无穷远结构与几何的层次}

无穷远直线的定义依赖于坐标：它的前两个分量被取为零。这一条件不可能写成与坐标系无关的形式，因为换一个坐标系之后，它的前两个分量一般不再为零。因此，齐次坐标一节中“抹去坐标系的一切痕迹”这第三步对无穷远结构不能实现，无穷远结构并不完全属于射影几何。

事实上，无穷远结构是三维矢量空间上的一个附加结构，只需要如下定义就可以产生。

\begin{definition}[仿射平面]
在三维矢量空间的全体线矢中指定一个特殊的线矢，称为无穷远线矢，记作 $\linf$，就在相应的射影几何中产生了一个无穷远结构。带无穷远结构的射影平面称为仿射平面。在此基础上可以进一步定义：无穷远点就是位于无穷远直线上的点；两条直线平行就是它们相交于无穷远点。
\end{definition}

于是，几何的层次关系可以更清楚地看到：

\begin{itemize}
\item 射影几何处于最底层，它只涉及点矢、线矢以及坐标系无关的运算，例如判断若干点是否共线、若干直线是否共点；
\item 仿射几何在射影几何之上多出了无穷远结构，因而涉及更多的关系，例如两条直线是否平行；
\item 解析几何处于最上层，它在仿射几何之上又多出了坐标系结构，因而涉及最多的关系，例如平行、垂直、长度，乃至 $y=\sin x$ 这样的曲线。
\end{itemize}

层次越高，可以表达的关系越多；高层次的关系可以表示低层次的关系，但低层次的关系未必能表示高层次的关系。
   % 无穷远

\part{认识圆锥曲线}
% 第七节：圆锥曲线、极点极线与分类
\chapter{椭圆、双曲线、抛物线与极点极线}

\section{圆锥曲线张量}

\begin{definition}[圆锥曲线张量]
一个不定的、对称的、可逆的 $(0,2)$ 型张量 $g_{ij}$ 称为圆锥曲线张量（或二次曲线张量），
它所描述的几何对象称为 $g$ 对应的圆锥曲线。
\end{definition}

其中“对称”指 $g_{ij}=g_{ji}$；“可逆”指线性映射 $a^i\mapsto g_{ij}a^j$ 可逆，它的逆张量记为
$(g^{-1})^{ij}$，满足
\[
(g^{-1})^{ij}g_{jk}=\delta^i_k .
\]
当讨论中只有一条圆锥曲线并把它取作度规时，把 $(g^{-1})^{ij}$ 简写为 $g^{ij}$，并用它升降指标。
当同时出现多条圆锥曲线时必须注意：$g^{ij}$ 一般指用某条圆锥曲线 $g$ 升指标的结果，
$(g^{-1})^{ij}$ 指 $g$ 的逆，两者一般不同，不能混用。

“不定”指二次型 $g_{ij}a^ia^j$ 既取正值也取负值，即存在非零矢量 $a,b$ 使得
$g_{ij}a^ia^j>0$ 而 $g_{ij}b^ib^j<0$。

\begin{definition}[圆锥曲线上的点]
若点矢 $a$ 满足
\begin{equation}
g_{ij}a^ia^j=0,
\end{equation}
则称 $a$ 对应的点在该圆锥曲线上。
\end{definition}

\begin{remark}
上式是规范不变的：把 $a$ 换成 $\lambda a$（$\lambda\neq0$）时左边乘以 $\lambda^2$，方程不变。
因此“点在圆锥曲线上”是一个射影概念。若 $g$ 是定的，即 $g_{ij}a^ia^j$ 恒正或恒负，则上述方程
只有零解，圆锥曲线上没有任何点（如方程 $x^2+y^2=-1$ 就没有实解）；这正是要求 $g$ 不定的原因。
三维实矢量空间上非退化不定对称张量的号差必为 $(2,1)$ 或 $(1,2)$。
\end{remark}

记 $a^2=a\cdot a=g_{ij}a^ia^j$，则“点在圆锥曲线上”就是 $a^2=0$。

\begin{definition}[切线]
若线矢 $l$ 满足
\begin{equation}
(g^{-1})^{ij}l_il_j=0,
\end{equation}
则称 $l$ 对应的直线与该圆锥曲线相切，并称 $l$ 为该圆锥曲线的一条切线。
\end{definition}

\begin{proposition}
设圆锥曲线非退化，$l$ 是一条直线。则 $l$ 与圆锥曲线相切的充要条件是 $l$ 与它恰有一个交点
（作为二重交点）。
\end{proposition}

\begin{proof}
设 $L=\{x:l_ix^i=0\}$ 是 $l$ 对应的二维子空间，$l$ 上的点由二次型 $Q(x)=g_{ij}x^ix^j$
在 $L$ 上的限制 $Q|_L$ 的零点给出。$L$ 上二次型退化当且仅当 $L\cap L^\perp\neq\{0\}$；由 $g$
非退化，$L^\perp$ 是由 $l^\sharp$（分量 $(l^\sharp)^i=(g^{-1})^{ij}l_j$）张成的一维子空间，
故这一条件等价于 $l^\sharp\in L$，即 $l_i(g^{-1})^{ij}l_j=0$，这正是相切的定义。

反过来看 $L$ 上二次型的秩。二维空间上退化的二次型秩为 $1$，其零点集是过原点的一条直线，
即 $l$ 上恰有一个点；秩为 $2$ 时，视其定号与否，零点集为空或两条过原点的直线，对应无交点或
两个交点。非退化圆锥曲线不含整条直线，故秩为 $0$ 的情形不出现。
\end{proof}

\begin{remark}
这里用 $(g^{-1})^{ij}l_il_j=0$ 定义相切，与用“只有一个交点”定义相切是等价的。这样定义的好处是
点与切线的定义在形式上完全对称：
\[
\text{点在圆锥曲线上}:\ g_{ij}a^ia^j=0,\qquad
\text{直线与圆锥曲线相切}:\ (g^{-1})^{ij}l_il_j=0 .
\]
\end{remark}

\begin{example}[标准形式]
取
\[
g_{ij}=\diag\Bigl\{\frac1{a^2},\frac1{b^2},-1\Bigr\},\qquad a,b>0,
\]
点的齐次坐标为 $(x,y,z)$。此时
\[
g_{ij}a^ia^j=\frac{x^2}{a^2}+\frac{y^2}{b^2}-z^2,
\qquad
(g^{-1})^{ij}=g^{ij}=\diag\{a^2,b^2,-1\}.
\]
在仿射坐标 $(x/z,\,y/z)$ 下，圆锥曲线的方程就是椭圆
\[
\frac{x^2}{a^2}+\frac{y^2}{b^2}=1 .
\]
$b^2$ 取负值时同一公式给出双曲线，$g$ 的号差仍为 $(2,1)$。
\end{example}

\section{极点与极线}

有了 $g$，就可以用它升降指标：点矢 $a^i$ 降指标得到线矢 $a_i=g_{ij}a^j$；线矢 $l_i$ 升指标得到
点矢 $l^i=(g^{-1})^{ij}l_j$。

\begin{definition}[极点与极线]
对任一元素 $a$（点矢或线矢），$a^i$ 对应的点与 $a_i$ 对应的直线互为极点与极线：点 $a$ 的极线是
由线矢 $a_i=g_{ij}a^j$ 确定的直线，线 $l$ 的极点是由点矢 $l^i=(g^{-1})^{ij}l_j$ 确定的点。
由 $g$ 给出的这一对应称为极配。
\end{definition}

为了区分类型，下面记 $a^*$ 为元素 $a$ 在极配下的对应元素：若 $a$ 是点，$a^*$ 是它的极线，
$(a^*)_i=g_{ij}a^j$；若 $l$ 是线，$l^*$ 是它的极点，$(l^*)^i=(g^{-1})^{ij}l_j$。于是 $a^{**}=a$。
当只有一条圆锥曲线、不需要区分时，也可以说同一个字母 $a$ 既代表一个点又代表它的极线。

\begin{remark}
设 $x$ 是点，$l$ 是线，则 $x$ 在 $l$ 上当且仅当 $x\cdot l^*=0$。事实上
\[
x\cdot l^*=x_i(l^*)^i=g_{ij}x^j(g^{-1})^{ik}l_k=x^il_i,
\]
右端就是 $x$ 与 $l$ 的自然缩并。特别地，$x$ 在它自己的极线上当且仅当 $x^2=0$。
\end{remark}

\begin{example}[极线的解析式]
取标准形式 $g_{ij}=\diag\{1/a^2,1/b^2,-1\}$，圆锥曲线为椭圆
\[
\frac{x^2}{a^2}+\frac{y^2}{b^2}=1 .
\]
点 $(x_0,y_0)$ 的齐次坐标为 $(x_0,y_0,1)$，它的极线为
\[
\frac{x_0}{a^2}x+\frac{y_0}{b^2}y-z=0,
\]
写回仿射坐标即
\[
\frac{xx_0}{a^2}+\frac{yy_0}{b^2}=1 .
\]
若 $(x_0,y_0)$ 在椭圆上，这就是椭圆在该点的切线方程。
\end{example}

\begin{example}[极线互过]
若 $A$ 的极线经过 $B$，则 $B$ 的极线经过 $A$。
\end{example}

\begin{proof}
设 $A,B$ 的点矢为 $a,b$。“$A$ 的极线经过 $B$”即 $B\in a^*$，由上一注它等价于 $a\cdot b=0$；
“$B$ 的极线经过 $A$”即 $A\in b^*$，等价于 $b\cdot a=0$。而 $a\cdot b=g_{ij}a^ib^j=b\cdot a$，
故两个命题是同一件事。
\end{proof}

\begin{definition}[共轭]
若 $a\cdot b=0$，则称 $a,b$ 关于该圆锥曲线共轭。
\end{definition}

\begin{example}[曲线上点的极线是切线]
若点 $A$ 在圆锥曲线上，则 $A$ 的极线是圆锥曲线在 $A$ 处的切线。
\end{example}

\begin{proof}
由 $a^2=g_{ij}a^ia^j=0$，作为线矢的 $a$ 满足相切条件
$(g^{-1})^{ij}a_ia_j=a^2=0$，故直线 $a^*$ 与圆锥曲线相切。又 $a\cdot a=a^2=0$，即 $A$ 在
$a^*$ 上。于是 $a^*$ 是过 $A$ 的切线，也就是圆锥曲线在 $A$ 处的切线。
\end{proof}

\begin{proposition}[极配与叉乘]
对任意两点 $a,b$ 与任意两线 $m,n$，有
\[
(a\times b)^*=a^*\times b^*,\qquad (m\times n)^*=m^*\times n^* .
\]
两个等号都在允许相差非零常数的意义下成立，即作为点或直线相等。
\end{proposition}

\begin{proof}
先证第一式。$a\times b$ 是直线 $AB$，$(a\times b)^*$ 是它的极点。设 $t=a^*\times b^*$，
其分量为 $t^i=\epsilon^{ijk}(a^*)_j(b^*)_k=\epsilon^{ijk}a_jb_k$。于是
\[
a\cdot t=a_it^i=\epsilon^{ijk}a_ia_jb_k=0,
\]
这里用了 $\epsilon^{ijk}$ 关于 $i,j$ 反对称而 $a_ia_j$ 关于 $i,j$ 对称；同理 $b\cdot t=0$。
因此 $a,b$ 都在 $t$ 的极线上，$t$ 的极线就是直线 $ab$，即 $t$ 是直线 $ab$ 的极点，
故 $t=(a\times b)^*$。

再证第二式。设 $t=m\times n$。由前一注，$x\in m$ 等价于 $x\cdot m^*=0$；而 $t\in m$，故
$t\cdot m^*=0$，同理 $t\cdot n^*=0$。于是 $m^*,n^*$ 都在 $t$ 的极线上，$t$ 的极线就是直线
$m^*n^*=m^*\times n^*$，即 $(m\times n)^*=m^*\times n^*$。
\end{proof}

\begin{example}[外点的极线]
设 $A$ 在圆锥曲线外。过 $A$ 作圆锥曲线的两条切线，切点分别为 $B,C$，则直线 $BC$ 是 $A$ 的极线。
\end{example}

\begin{proof}
$B$ 在圆锥曲线上，故 $B$ 的极线是过 $B$ 的切线；$A$ 在这条切线上，由极线互过，$B$ 在 $A$ 的极线
上，即 $a\cdot b=0$。同理 $a\cdot c=0$。于是 $B,C$ 都在 $A$ 的极线上，而两点确定一条直线，
故 $A$ 的极线就是 $BC$。
\end{proof}

\begin{corollary}
若直线 $l$ 与圆锥曲线交于 $B,C$，则 $l$ 的极点是 $B,C$ 两处切线的交点。
\end{corollary}

\begin{proof}
$l=b\times c$。由极配与叉乘的关系，$l^*=(b\times c)^*=b^*\times c^*$；由前一例，$b^*$ 与 $c^*$
分别是 $B,C$ 处的切线，故 $l^*$ 就是这两条切线的交点。这也给出了由一条割线作它的极点的方法。
\end{proof}

\begin{definition}[自极三点形]
若三点 $A,B,C$ 满足
\[
a\cdot b=b\cdot c=c\cdot a=0,
\]
则称 $ABC$ 为关于该圆锥曲线的自极三点形。
\end{definition}

\begin{remark}
由定义，$a\cdot b=a\cdot c=0$ 说明 $B,C$ 都在 $A$ 的极线上，故 $A$ 的极线是 $BC$；同理 $B$ 的极线
是 $AC$，$C$ 的极线是 $AB$。在三维矢量空间中，这意味着 $a,b,c$ 是关于 $g$ 两两正交的三个矢量，
因而构成一组基；反过来，任取一组关于 $g$ 两两正交的基矢，就得到一个自极三点形。因此自极三点形
可以理解为在三维矢量空间中选取了一组关于圆锥曲线张量正交的基。
\end{remark}

\begin{remark}
极点极线与调和点列、交比的关系需要更多的射影几何工具，这里不再展开。
\end{remark}

\section{用两族割线作极线}

\begin{lemma}[混合型拉格朗日恒等式]
设 $A,B$ 是点，$C,D$ 是线，则
\begin{equation}
(A\times B)\cdot(C\times D)=(A\cdot C)(B\cdot D)-(A\cdot D)(B\cdot C).
\end{equation}
\end{lemma}

\begin{proof}
由二重外积公式可以推出这一恒等式；直接按分量计算也可：
\[
\begin{aligned}
(A\times B)_i(C\times D)^i
&=\epsilon_{ijk}A^jB^k\,\epsilon^{ilm}C_lD_m
=(\delta_j^l\delta_k^m-\delta_k^l\delta_j^m)A^jB^kC_lD_m\\
&=(A\cdot C)(B\cdot D)-(A\cdot D)(B\cdot C).
\end{aligned}
\]
\end{proof}

\begin{remark}
恒等式要求两个叉乘类型相反：一个由两点作出（得到直线），另一个由两线作出（得到点），此时
$\cdot$ 是点与线的自然缩并。不能把同一个恒等式直接用于两个由线作出的点：它们的点乘要用度规，
而度规与 $\epsilon$ 之间一般没有使该恒等式成立的关系。下面的推导先借极配把点换成它的极线，
再用上式展开。
\end{remark}

\begin{example}[用两族割线作极线]
设 $J$ 不在圆锥曲线上。过 $J$ 作两条直线，分别交圆锥曲线于 $C,E$ 与 $D,F$，于是
$J=CE\cap DF$。设
\[
H=CD\cap EF,\qquad I=CF\cap DE .
\]
则直线 $IH$ 是 $J$ 的极线。这一作法只用直尺，不需要预先作出切线。
\end{example}

\begin{proof}
设 $c,d,e,f$ 是 $C,D,E,F$ 的点矢。由作法
\[
j=(c\times e)\times(d\times f),\qquad h=(c\times d)\times(e\times f),\qquad
i=(c\times f)\times(d\times e).
\]
只需证明 $j\cdot h=j\cdot i=0$，这样 $H,I$ 都在 $J$ 的极线上。

先算 $j\cdot h$。由极配与叉乘的关系，
\[
j^*=(c\times e)^*\times(d\times f)^*=(c^*\times e^*)\times(d^*\times f^*).
\]
点 $j$ 与它的极线 $j^*$ 满足 $j\cdot h=j^*\cdot h$（极线 $j^*$ 的线矢就是 $j_i=g_{ij}j^j$）。
记
\[
A=c^*\times e^*,\qquad B=d^*\times f^*,\qquad C=c\times d,\qquad D=e\times f,
\]
则 $A,B$ 是点，$C,D$ 是线，$j^*=A\times B$，$h=C\times D$。由混合型拉格朗日恒等式
\[
j\cdot h=(A\times B)\cdot(C\times D)=(A\cdot C)(B\cdot D)-(A\cdot D)(B\cdot C).
\]
由于 $c,d,e,f$ 都在圆锥曲线上，有 $c\cdot c^*=c^2=0$，$d\cdot d^*=0$，$e\cdot e^*=0$，
$f\cdot f^*=0$。逐个计算四个因子：
\begin{align*}
A\cdot C&=(c^*\times e^*)\cdot(c\times d)
=(c\cdot c^*)(d\cdot e^*)-(c\cdot e^*)(d\cdot c^*)\\
&=c^2(d\cdot e)-(c\cdot e)(c\cdot d)=-(c\cdot d)(c\cdot e),\\
B\cdot D&=(d^*\times f^*)\cdot(e\times f)
=(d\cdot e)(f\cdot f^*)-(d\cdot f)(f\cdot e^*)\\
&=-(d\cdot f)(e\cdot f),\\
A\cdot D&=(c^*\times e^*)\cdot(e\times f)
=(c\cdot e)(e^*\cdot f)-(c\cdot f)(e^*\cdot e)\\
&=(c\cdot e)(e\cdot f)-e^2(c\cdot f)=(c\cdot e)(e\cdot f),\\
B\cdot C&=(d^*\times f^*)\cdot(c\times d)
=(d\cdot c)(f^*\cdot d)-(d\cdot d)(f^*\cdot c)\\
&=(c\cdot d)(d\cdot f)-d^2(c\cdot f)=(c\cdot d)(d\cdot f).
\end{align*}
代入得
\begin{align*}
j\cdot h&=\bigl[-(c\cdot d)(c\cdot e)\bigr]\bigl[-(d\cdot f)(e\cdot f)\bigr]
-\bigl[(c\cdot e)(e\cdot f)\bigr]\bigl[(c\cdot d)(d\cdot f)\bigr]\\
&=(c\cdot d)(c\cdot e)(d\cdot f)(e\cdot f)-(c\cdot e)(e\cdot f)(c\cdot d)(d\cdot f)=0 .
\end{align*}

再算 $j\cdot i$。同样地
\[
j^*=A\times B,\qquad i=(c\times f)\times(d\times e),
\]
把上面的 $C=c\times d$、$D=e\times f$ 换成 $C'=c\times f$、$D'=d\times e$，重复同样的计算（仍然
用到 $c^2=d^2=e^2=f^2=0$），得到 $j\cdot i=0$。于是 $H,I$ 都在 $J$ 的极线上，故 $IH$ 是 $J$
的极线。
\end{proof}

\begin{remark}
这个证明的关键是选取合适的元素作为参量：直接设出圆锥曲线上的四点 $c,d,e,f$，再用它们表示
$h,i,j$，计算就能顺利完成；若反过来由 $j$ 去求 $c,d,e,f$，则要解高次方程，计算会复杂得多。
\end{remark}

\section{用无穷远结构对圆锥曲线分类}

非退化对称 $(0,2)$ 型张量在坐标变换下唯一不变的性质是它的号差；而圆锥曲线要求 $g$ 不定，
号差只能是 $(2,1)$ 或 $(1,2)$，两者相差一个整体符号，在射影意义下没有区别。因此仅凭圆锥曲线
张量无法区分椭圆、双曲线与抛物线。要把它们区分开，必须再附加一个结构，即无穷远结构 $\linf$。

直观上，椭圆的所有点都在有限区域内，没有点在无穷远；双曲线沿两条渐近线的方向延伸到无穷远，
因此有两个无穷远点；抛物线的开口只朝向一个方向，因此只有一个无穷远点。于是有

\begin{definition}[椭圆、双曲线、抛物线]
设 $g$ 对应的圆锥曲线与无穷远直线 $\linf$ 的关系为：
\begin{itemize}
\item 相离（无交点）：称该圆锥曲线为椭圆；
\item 相交（两个交点）：称该圆锥曲线为双曲线；
\item 相切（一个二重交点）：称该圆锥曲线为抛物线。
\end{itemize}
\end{definition}

\begin{example}[解析判定]
取 $g_{ij}=\diag\{1/a^2,1/b^2,-1\}$，$\linf$ 的线矢为 $(0,0,1)$，即判定时在齐次方程中令 $z=0$。
\begin{itemize}
\item 椭圆 $\dfrac{x^2}{a^2}+\dfrac{y^2}{b^2}=1$ 齐次化为 $\dfrac{x^2}{a^2}+\dfrac{y^2}{b^2}-z^2=0$，
令 $z=0$ 得 $\dfrac{x^2}{a^2}+\dfrac{y^2}{b^2}=0$，只有零解，故与 $\linf$ 相离；
\item 双曲线 $\dfrac{x^2}{a^2}-\dfrac{y^2}{b^2}=1$ 齐次化为 $\dfrac{x^2}{a^2}-\dfrac{y^2}{b^2}-z^2=0$，
令 $z=0$ 得 $x^2/a^2=y^2/b^2$，给出两个无穷远点 $(a,\pm b,0)$，故与 $\linf$ 相交；
\item 抛物线 $y^2=2px$ 齐次化为 $y^2-2pxz=0$，令 $z=0$ 得 $y^2=0$，即点 $(1,0,0)$ 为二重解，
故与 $\linf$ 相切。
\end{itemize}
\end{example}

\begin{proposition}[中心]
圆锥曲线的中心是 $\linf$ 的极点 $\linf^*$。对椭圆与双曲线它是有限点，即通常的对称中心；
对抛物线它在无穷远，因此抛物线没有有限中心。
\end{proposition}

\begin{proof}
设圆锥曲线在仿射坐标下由
\[
Q(x,y)=Ax^2+2Bxy+Cy^2+2Dx+2Ey+F=0
\]
给出，齐次化后为 $Q(x,y,z)=Ax^2+2Bxy+Cy^2+2Dxz+2Eyz+Fz^2$，对应的 $g_{ij}$ 以
$(A,B,D;B,C,E;D,E,F)$ 为分量。设 $p$ 为 $\linf=(0,0,1)$ 的极点，即
$g_{ij}p^j=(\linf)_i$，写开就是
\[
Ap^1+Bp^2+Dp^3=0,\qquad Bp^1+Cp^2+Ep^3=0,\qquad Dp^1+Ep^2+Fp^3=1 .
\]
当 $AC-B^2\neq0$ 时可以取 $p^3=1$，此时前两式为
\[
Ap^1+Bp^2+D=0,\qquad Bp^1+Cp^2+E=0,
\]
这正是圆锥曲线的中心（即二次曲面 $Q$ 的极值点）所满足的方程，故 $\linf^*$ 就是中心。
当 $AC-B^2=0$（抛物线）时，前两式的系数矩阵奇异而右端不能由它线性表示，方程组无 $p^3\neq0$
的解，故 $p^3=0$，即 $\linf^*$ 在无穷远。
\end{proof}

\begin{proposition}[渐近线与中心]
设双曲线与 $\linf$ 交于 $C,D$ 两点。则 $C,D$ 两处的切线就是双曲线的两条渐近线，它们相交于双曲线的
中心。
\end{proposition}

\begin{proof}
$C,D$ 都在 $\linf$ 上且不同，故过 $C,D$ 的直线就是 $\linf$，即 $c\times d=\linf$。由极配与叉乘的
关系，
\[
c^*\times d^*=(c\times d)^*=\linf^*,
\]
而 $\linf^*$ 就是中心。另一方面 $C,D$ 在双曲线上，$c^*,d^*$ 分别是 $C,D$ 处的切线。由定义，
$C,D$ 是双曲线的无穷远点，曲线在这两点的切线就是它沿无穷远方向的两条渐近线。因此两条渐近线
相交于中心。这一结论也说明：中心是由无穷远结构派生出来的，凡是涉及中心的命题，都不能只用射影
几何处理，而应在已有的参量之外再指定 $\linf$。
\end{proof}

\section{圆锥曲线与三维时空的类比}

一个三维矢量空间配上一条圆锥曲线，与三维时空（一维时间、二维空间）配上闵可夫斯基度规，
在形式上是相同的：两者都是一个三维实矢量空间配上一个号差为 $(2,1)$ 的不定对称双线性型。
取光速 $c=1$，二者的对照如下表（表中的号差约定为 $(2,1)$，即圆锥曲线张量与闵氏度规在相差
整体符号的意义下有相同的分量）。

\begin{center}
\begin{tabular}{ll}
\toprule
平面上的圆锥曲线 & 三维时空 \\
\midrule
点由三维点矢表示（允许相差非零常数） & 点由坐标 $(t,x,y)$ 表示 \\
圆锥曲线张量 $g_{ij}$ & 闵可夫斯基度规 $\eta_{ij}$ \\
标准形式 $x^2+y^2-z^2=0$ & 光锥 $t^2-x^2-y^2=0$ \\
点在圆锥曲线内：$g_{ij}a^ia^j<0$ & 类时矢量 \\
点在圆锥曲线外：$g_{ij}a^ia^j>0$ & 类空矢量 \\
点在圆锥曲线上：$g_{ij}a^ia^j=0$ & 类光矢量 \\
直线与圆锥曲线相离：$(g^{-1})^{ij}l_il_j<0$ & $l$ 类时（$l^\perp$ 类空） \\
直线与圆锥曲线相切：$(g^{-1})^{ij}l_il_j=0$ & $l$ 类光（$l^\perp$ 类光） \\
直线与圆锥曲线相交：$(g^{-1})^{ij}l_il_j>0$ & $l$ 类空（$l^\perp$ 类时） \\
\bottomrule
\end{tabular}
\end{center}

\begin{remark}
表中直线的类型由线矢 $l$ 的双线性型 $(g^{-1})^{ij}l_il_j$ 的符号决定，而直线对应的平面
$l^\perp=\{x:l_ix^i=0\}$ 的类型与 $l$ 相反：$l$ 类时时 $l^\perp$ 上的双线性型正定，即 $l^\perp$
类空，反之亦然；$l$ 类光时 $l^\perp$ 退化，即类光。以标准形式为例，单位圆 $x^2+y^2=z^2$ 与直线
$x=2z$ 相离，而 $l=(1,0,-2)$ 满足 $(g^{-1})^{ij}l_il_j=1-4<0$；与直线 $x=0$ 相交，而
$l=(1,0,0)$ 满足 $(g^{-1})^{ij}l_il_j=1>0$。可见相离对应线矢类时，相交对应线矢类空。
\end{remark}
   % 椭圆、双曲线、抛物线与极点极线
\chapter{退化圆锥曲线}

\section{二次点与二次线}

回顾圆锥曲线的定义。

\begin{definition}[圆锥曲线]
一个不定的对称可逆的 $(0,2)$ 型张量 $g_{ij}$ 称为\textbf{圆锥曲线张量}。点 $a$ 在圆锥曲线 $g$ 上由
\[
g_{ij}a^ia^j=0
\]
刻画。$g$ 的逆记为 $(g^{-1})^{ij}$，当把 $g$ 取作度规时简写为 $g^{ij}$；直线 $l$ 与 $g$ 相切由
\[
(g^{-1})^{ij}l_il_j=0
\]
刻画。
\end{definition}

圆锥曲线的定义要求 $g_{ij}$ 可逆，但不可逆的张量同样有重要的用处。例如双曲线
$\frac{x^2}{a^2}-\frac{y^2}{b^2}=k$ 在 $k=0$ 时成为两条直线，这两条直线可以看作一种特殊的圆锥曲线，
称为\textbf{退化圆锥曲线}。由点与线的对偶性，“两个点”也属于这一类：它是被压扁成一条线段的椭圆的
两个端点。为了把“两条直线”“两个点”与圆锥曲线放在同一个框架中研究，需要引入比圆锥曲线更一般的
张量。

另一个基本事实是：一条直线与一条圆锥曲线交于两点 $a,b$，或者过一个点作圆锥曲线的两条切线，这
两个点（或两条线）具有完全相同的地位，不能拆开而单独看待。用某个不可逆的二阶张量（而不是用
两个点矢或两个线矢）表示这样的一对对象，是最自然的做法。

\begin{definition}[二次点与二次线]
设 $f^{ij}$ 是不定对称的 $(2,0)$ 型张量，$h_{ij}$ 是不定对称的 $(0,2)$ 型张量。
\begin{enumerate}[label=(\arabic*),leftmargin=2.4em]
\item $f^{ij}$ 称为\textbf{二次点张量}，它所对应的几何元素称为\textbf{二次点}；有时也直接把
$f^{ij}$ 称为二次点。
\item $h_{ij}$ 称为\textbf{二次线张量}，它所对应的几何元素称为\textbf{二次线}。
\item 二次点与二次线统称\textbf{二次型}，或统称\textbf{二次曲线}。
\item 若 $f^{ij}l_il_j=0$，称直线 $l$ \textbf{经过}二次点 $f$。满足此方程的所有 $l$ 称为一个
\textbf{二次线束}；当 $f$ 不可逆时，称为\textbf{退化二次线束}。
\item 若 $h_{ij}a^ia^j=0$，称点 $a$ 在二次线 $h$ 上。满足此方程的所有 $a$ 称为一个
\textbf{二次点列}；当 $h$ 不可逆时，称为\textbf{退化二次点列}。
\end{enumerate}
\end{definition}

有一个上指标的量是点，有两个对称上指标的量是二次点；有一个下指标的量是线，有两个对称下指标
的量是二次线。因此“二次点”“二次线”正是点与线的推广。

\begin{remark}
经过可逆二次点 $f^{ij}$ 的直线，就是与圆锥曲线 $f$ 相切的直线，故“经过二次点”确实是“经过一个
点”的推广。二次型可以可逆，也可以不可逆：圆锥曲线 $g_{ij}$ 是二次线的可逆特例，它自然地生成
二次点 $(g^{-1})^{ij}$；而不可逆的二次线不能生成一个二次点，反之亦然。
\end{remark}

\section{不可逆二次型的几何意义}

按二次型矩阵的秩，也就是号差中 $0$ 的个数，可以把二次型分为三类：没有 $0$、有 $1$ 个 $0$、
有 $2$ 个 $0$。没有 $0$ 意味着可逆，对应圆锥曲线；有 $1$ 个或 $2$ 个 $0$ 的二次型不可逆。
下面证明：号差中有 $1$ 个 $0$ 的二次型对应“两个点”或“两条线”，有 $2$ 个 $0$ 的二次型对应
“单个点”或“单条线”。以二次线为例给出计算，二次点的情形完全对偶。

\begin{proposition}
\label{prop:onezero}
设二次线 $h_{ij}$ 的号差为 $(1,1,0)$，即恰有一个 $0$。
\begin{enumerate}[label=(\arabic*),leftmargin=2.4em]
\item 存在非零的点矢 $s$，使得 $h_{ij}s^j=0$；满足 $h_{ij}x^j=0$ 的所有 $x$ 都与 $s$ 共线，
即 $x^i=\lambda s^i$。并且 $h_{ij}s^j=0$ 不仅意味着 $h_{ij}s^is^j=0$，还意味着对任意矢量 $a$
都有 $h_{ij}s^ia^j=0$。
\item $h$ 上点的全体是过 $s$ 所对应的点 $S$ 的两条直线。
\end{enumerate}
\end{proposition}

\begin{proof}
号差为 $(1,1,0)$ 意味着 $\rank h=2$，故映射 $x\mapsto h_{ij}x^j$ 的核是一维的：存在非零的
$s$ 使 $h_{ij}s^j=0$，并且解空间由 $s$ 张成。由 $h$ 的双线性性，$h_{ij}s^j=0$ 立刻给出
$h_{ij}s^ia^j=0$ 对一切 $a$ 成立。

由 $h_{ij}s^is^j=s^i(h_{ij}s^j)=0$ 知 $S$ 在 $h$ 上。设 $A$ 也在 $h$ 上，即
$h_{ij}a^ia^j=0$。直线 $SA$ 上的点都可写成 $s+\lambda a$，而
\[
h_{ij}(s+\lambda a)^i(s+\lambda a)^j
=h_{ij}s^is^j+2\lambda h_{ij}s^ia^j+\lambda^2h_{ij}a^ia^j=0,
\]
因此整条直线 $SA$ 都含于 $h$。于是 $h$ 上点的全体是过 $S$ 的若干条直线之并。

下面证明这样的直线恰有两条。先取一条不经过 $S$ 的直线 $\ell$，$\ell$ 上的点都可写成
$p+\lambda q$。多项式
\[
P(\lambda)=h_{ij}(p+\lambda q)^i(p+\lambda q)^j
\]
关于 $\lambda$ 至多是二次的。若 $P\equiv0$，则 $\ell$ 上每一点都在 $h$ 上，把 $\ell$ 上任一点
与 $S$ 相连即知过 $S$ 的每条直线都含于 $h$，从而平面上每一点都在 $h$ 上，即
$h_{ij}a^ia^j=0$ 对一切 $a$ 成立，这与 $h\neq0$ 矛盾。因此 $P$ 至多有两个零点，即 $\ell$ 与
$h$ 至多交于两点。而 $\ell$ 上每一点若在 $h$ 上，就对应过 $S$ 的一条含于 $h$ 的直线（这是因为
$S$ 与 $\ell$ 上该点都在 $h$ 上，连线整条含于 $h$）。所以过 $S$ 且含于 $h$ 的直线至多两条。

再证这样的直线恰恰有两条。把 $h$ 对角化，在适当的基下 $h$ 的分量为
$\diag\{1,-1,0\}$。由前两个基矢张成的二维子空间中，$h$ 的限制为 $x^2-y^2$，它在两条直线
$x=\pm y$ 上取零值。这两个方向与核方向（第三个基矢）不共线，故对应的两个点都在 $h$ 上且都不
等于 $S$，从而过 $S$ 的这两条直线都含于 $h$。故过 $S$ 且含于 $h$ 的直线恰有两条。
\end{proof}

\begin{proposition}
\label{prop:twozero}
设二次线 $h_{ij}$ 的号差中有恰好两个 $0$。则存在两个不共线的点矢 $s,t$ 使
$h_{ij}s^j=h_{ij}t^j=0$，并且 $h$ 上点的全体恰为 $s,t$ 所连的直线。
\end{proposition}

\begin{proof}
号差中有两个 $0$ 意味着 $\rank h=1$，映射 $x\mapsto h_{ij}x^j$ 的核是二维的，取其中的一组基
$s,t$，就有 $h_{ij}s^j=h_{ij}t^j=0$，且 $s,t$ 不共线。由 $h_{ij}s^is^j=h_{ij}t^it^j=0$ 知
$s,t$ 都在 $h$ 上；与上面相同的计算表明直线 $st$ 含于 $h$。

设 $p$ 不在直线 $st$ 上，则 $p,s,t$ 构成三维空间的一组基，任一矢量 $a$ 可写成
$a=\lambda p+\mu s+\sigma t$。由 $h_{ij}s^j=h_{ij}t^j=0$ 得
\[
h_{ij}a^ia^j=\lambda^2h_{ij}p^ip^j .
\]
若 $p$ 也在 $h$ 上，则对一切 $a$ 都有 $h_{ij}a^ia^j=0$，即 $h=0$，矛盾。因此不在直线 $st$ 上的
点都不在 $h$ 上，$h$ 上点的全体恰为直线 $st$。
\end{proof}

\begin{remark}[号差中 $0$ 的个数与解空间维数]
容易把两个不同方程的维数弄混：线性方程 $h_{ij}x^j=0$ 的解空间维数等于 $h_{ij}$ 的零特征值
个数，即号差中 $0$ 的个数。因此号差中有 $2$ 个 $0$ 时解空间是二维的，存在两个不共线的解
$s,t$；号差中有 $1$ 个 $0$ 时解空间才是一维的。对后一情形，上面的证明用“二次型在一条直线
上的限制若非恒零则至多有两个零点”这一事实补足了“含于 $h$ 的直线恰有两条”的论证。
\end{remark}

\begin{theorem}[退化二次型的几何意义]
在二次线 $h_{ij}$ 上的点构成
\begin{enumerate}[label=(\arabic*),leftmargin=2.4em]
\item 圆锥曲线，当 $h$ 可逆时；
\item 两条直线，当 $h$ 的号差中有一个 $0$ 时；
\item 一条直线，当 $h$ 的号差中有两个 $0$ 时。
\end{enumerate}
对偶地，经过二次点 $f^{ij}$ 的直线构成
\begin{enumerate}[label=(\arabic*'),leftmargin=2.6em]
\item 圆锥曲线的全部切线，当 $f$ 可逆时；
\item 经过两个定点的全部直线，当 $f$ 的号差中有一个 $0$ 时；
\item 经过一个定点的全部直线，当 $f$ 的号差中有两个 $0$ 时。
\end{enumerate}
\end{theorem}

\section{用点矢与线矢表示退化二次曲线}

前面的讨论把退化二次型与点、直线联系了起来。下面给出代数上的联系：如何用（一个或两个）点矢、
线矢把二次型 $f,h$ 具体表示出来。

\begin{proposition}[两个线矢的对称积]
设 $m,n$ 是不共线的两个线矢，令
\[
h_{ij}=m_in_j+m_jn_i.
\]
则 $h$ 的号差为 $(1,1,0)$，它上面的点集恰为直线 $m$ 与直线 $n$ 之并；并且
$s=m\times n$ 满足 $h_{ij}s^j=0$。
\end{proposition}

\begin{proof}
对任意点 $a$，
\[
h_{ij}a^ia^j=(m_in_j+m_jn_i)a^ia^j=2(m\cdot a)(n\cdot a).
\]
因此 $h_{ij}a^ia^j=0$ 当且仅当 $m\cdot a=0$ 或 $n\cdot a=0$，即点 $a$ 在直线 $m$ 或直线 $n$ 上。
于是 $h$ 上点的全体是两条直线 $m,n$ 之并。又 $m\cdot s=n\cdot s=0$，故
\[
h_{ij}s^j=m_i(n\cdot s)+n_i(m\cdot s)=0 .
\]
由 $h\neq0$ 且其零点集为两条直线，$h$ 的号差为 $(1,1,0)$。
\end{proof}

对偶地，对两个不共线的点矢 $a,b$，有
\[
f^{ij}=a^ib^j+a^jb^i,
\]
它是号差为 $(1,1,0)$ 的二次点，经过它的直线恰为过 $a$ 或过 $b$ 的直线之并，并且
$f^{ij}l_il_j=2(a\cdot l)(b\cdot l)$。

\begin{proposition}[一个线矢的平方]
设 $l$ 是一个非零线矢，令 $h_{ij}=l_il_j$。则 $h$ 的号差中有两个 $0$，它上面的点恰为直线 $l$
上的点。
\end{proposition}

\begin{proof}
对任意点 $a$，
\[
h_{ij}a^ia^j=l_il_ja^ia^j=(l\cdot a)^2,
\]
故 $h_{ij}a^ia^j=0$ 当且仅当 $l\cdot a=0$，即 $a$ 在直线 $l$ 上。$h$ 的秩为 $1$，号差中有两个
$0$。
\end{proof}

对偶地，$f^{ij}=a^ia^j$ 是号差中有两个 $0$ 的二次点，经过它的直线恰为过点 $a$ 的直线。
习惯上把二次线 $l_il_j$ 写成 $ll$，把二次点 $a^ia^j$ 写成 $aa$。

\begin{remark}
点矢与线矢可以用上述特殊的二次点与二次线来表示，但一般的二次点与二次线还包含圆锥曲线等丰富得多
的几何对象。可以在这些对象之间赋予代数运算，以表达多样的几何关系；这些运算的内容多于矢量运算，
矢量之间的点乘与叉乘只是其特例。
\end{remark}

例如，点 $a$ 在直线 $l$ 上等价于
\[
(aa):(ll)=a^ia^jl_il_j=(a^il_i)^2=0,
\]
而点 $a$ 在圆锥曲线 $g$ 上等价于
\[
(aa):g=g_{ij}a^ia^j=0 .
\]
前者是后者的特例。

\section{省略指标的写法}

今后会遇到复杂的表达式，省略指标常常更便于书写。约定如下。

\begin{enumerate}[label=\textbf{(\arabic*)},leftmargin=2.6em]
\item 两个元素直接并列，表示它们的指标直接并列而不缩并。例如 $ga$ 的分量为 $g_{ij}a^k$。
\item 两个元素之间加一个点，表示缩并相邻的一个指标。例如 $g\cdot a$ 的分量为 $g_{ij}a^j$，
$f\cdot g\cdot a$ 的分量为 $f^{ij}g_{jk}a^k$。
\item 两个元素之间加一个冒号，表示缩并相邻的两个指标；这一般只用于一个二次点与一个二次线之间。
例如 $f:g=f^{ij}g_{ij}$ 是一个实数。$h_{ij}$ 与 $g_{ij}$ 之间可以进行两种缩并：
\[
h^{-1}:g=g:h^{-1}=(h^{-1})^{ij}g_{ij},\qquad
g^{-1}:h=(g^{-1})^{ij}h_{ij}.
\]
注意这两个数一般不相等；其中一个为 $0$ 时，另一个未必为 $0$。
\end{enumerate}

有了这套写法，$f^{ij}=a^ib^j+a^jb^i$ 可以写成 $f=ab+ba$。由于这个构造十分重要，单独引入记号
\[
f=(ab),\qquad (ab)^{ij}=a^ib^j+a^jb^i,
\]
称为两个点的\textbf{对称积}。同理，两条线的对称积是
\[
h=(mn),\qquad (mn)_{ij}=m_in_j+m_jn_i .
\]

\section{两个点不能拆开}

二次点 $(ab)$ 中的两个点 $a,b$ 具有完全相同的地位，不能拆开而单独看待。这句话的意思是：虽然
号差为 $(1,1,0)$ 的二次点 $f$ 唯一对应了两个点 $a,b$，却很难把这两个点分别表示出来。

已知 $a,b$ 时，可以用与坐标系无关的张量运算写出它们的对称积 $f=(ab)$。反过来，不可能仅仅用
张量积、缩并的语言由 $f$ 表示出 $a$（或 $b$）。这是因为张量积与缩并都是单值函数。若存在单值
函数 $F$ 使 $a^i=F(f^{ij})$，由于 $a,b$ 在 $f$ 中地位平等，同样应当有 $b^i=F(f^{ij})$，
于是 $a^i=b^i$，矛盾。事实上，要用 $f$ 表示 $a,b$ 必须开根号，而开根号的一个正根与一个负根
恰好分别给出 $a,b$。使用张量语言的初衷是仅用张量运算处理几何问题，对分量的函数开根号破坏了
这一点。这使得在很多情况下，只能整体地处理两个点的对称积，而不能对它们分别处理。

与两个点的情形不同，对于由 $a^{ij}=a^ia^j$ 生成的二次点，可以由 $a^{ij}$ 表示出点 $a$：任取
一个线矢 $l_i$，有
\[
l_ia^{ij}=l_ia^ia^j=(l\cdot a)a^j,
\]
其中系数 $l\cdot a$ 只要不为 $0$ 就不影响 $a$ 所代表的点。也就是说，若已知点 $A$ 对应的二次点，
只需把它与一条不经过 $A$ 的直线缩并掉一个指标，就得到了 $A$ 的点矢。

\section{联立二次点}

\begin{example}
\label{ex:joined}
一条直线 $l$ 与圆锥曲线 $g$ 交于 $A,B$ 两点，求由 $A,B$ 确定的退化二次线束，即求二次点
$f=(ab)$。
\end{example}

\begin{proof}[解]
设直线 $m$ 在所求的退化二次线束上，则 $m$ 经过 $A$ 或 $B$，即 $m$ 与 $l$ 的交点是 $A$ 或 $B$，
也就是这个交点在 $g$ 上。两直线 $m,l$ 的交点满足
\[
(m\times l)^2=0,
\]
而
\[
(m\times l)^2=m^2l^2-(m\cdot l)^2 .
\]
因此所求的二次点为
\[
f^{ij}=l^2g^{ij}-l^il^j,\qquad\text{即}\qquad f=l^2g-ll,
\]
其中 $l^i=g^{ij}l_j$。验证：对任意线矢 $m$
\[
f^{ij}m_im_j=l^2m^2-(l\cdot m)^2=(m\times l)^2,
\]
并且
\[
f^{ij}l_j=l^2l^i-l^i(l\cdot l)=0,
\]
说明 $f$ 确实退化，$l$ 在其核中。这个张量完全由 $l$ 与 $g$ 确定，而与两个交点分开无关，称为
\textbf{联立二次点}。
\end{proof}

\begin{remark}
$f_{ij}=l^2g_{ij}-l_il_j$ 给出的是 $A,B$ 处两条切线构成的退化二次点列，有时它比 $f^{ij}$ 更
便于应用。此外，$f$ 可以看作 $g$ 在直线 $l$ 上的诱导度规。
\end{remark}

下面用解析坐标看这一等式的意义。取 $a^i=(x_1,y_1,1)$，$b^i=(x_2,y_2,1)$，则
\[
(ab)=\begin{pmatrix}
2x_1x_2 & x_1y_2+x_2y_1 & x_1+x_2\\
x_1y_2+x_2y_1 & 2y_1y_2 & y_1+y_2\\
x_1+x_2 & y_1+y_2 & 2
\end{pmatrix}.
\]
矩阵的元素全部由二次方程的两根之和与两根之积表示。虽然二次方程的解需要开根号，但两根之和与两根
之积不需要，这正是韦达定理。因此等式 $f=l^2g-ll$ 是韦达定理的一个推广。

\begin{remark}[对角元的系数]
写出分量矩阵时容易漏掉对角元上的系数。按对称积的定义 $(ab)^{ij}=a^ib^j+a^jb^i$，当 $i=j$
时对角元为 $a^ib^i+a^ib^i=2a^ib^i$，应带系数 $2$。上表中的 $2x_1x_2,\,2y_1y_2,\,2$ 正是
按此写出的。
\end{remark}

为避免出现不便应用的交叉项 $x_1y_2+x_2y_1$，引入 $k_1=\frac{y_1}{x_1}$，
$k_2=\frac{y_2}{x_2}$，即两点与原点连线的斜率。上式可以改写成
\[
(ab)=\begin{pmatrix}
2x_1x_2 & x_1x_2(k_1+k_2) & x_1+x_2\\
x_1x_2(k_1+k_2) & 2x_1x_2k_1k_2 & y_1+y_2\\
x_1+x_2 & y_1+y_2 & 2
\end{pmatrix}.
\]
设 $AB$ 的中点为 $M$，$x_M=\frac{x_1+x_2}{2}$，$y_M=\frac{y_1+y_2}{2}$。把最后一元归一（即整体
除以 $2$），得
\[
(ab)=\begin{pmatrix}
x_1x_2 & \frac{x_1x_2}{2}(k_1+k_2) & x_M\\[3pt]
\frac{x_1x_2}{2}(k_1+k_2) & x_1x_2k_1k_2 & y_M\\[3pt]
x_M & y_M & 1
\end{pmatrix}.
\]
于是中点坐标以及 $k_1+k_2$、$k_1k_2$ 都可以从分量中读出。

解析几何中常用的齐次化联立，其目的正是把联立方程写成关于斜率的二次方程，从而得到 $k_1+k_2$ 与
$k_1k_2$。由上面的分量式可见，只要算出联立二次点的分量，这一类问题都可以一并解决。

\section{例题：动弦垂直的包络}

\begin{example}
椭圆 $\frac{x^2}{a^2}+\frac{y^2}{b^2}=1$，$O$ 为原点。一条动直线与椭圆交于 $A,B$ 两点，且
$OA\perp OB$。求该动直线的包络。
\end{example}

\begin{proof}[解]
椭圆的张量及其逆为
\[
g_{ij}=\diag\{1/a^2,\,1/b^2,\,-1\},\qquad
g^{ij}=\diag\{a^2,\,b^2,\,-1\}.
\]
设动直线为 $Ax+By+C=0$，即 $l_i=(A,B,C)$。则
\[
l^i=g^{ij}l_j=(a^2A,\;b^2B,\;-C),\qquad
l^2=l_il^i=a^2A^2+b^2B^2-C^2 .
\]
关于 $l^2$ 的这一表达式可以用来判断直线与椭圆的位置关系：相切时 $l^2=0$，交于两点时 $l^2>0$，
无交点时 $l^2<0$。因此 $l^2=a^2A^2+b^2B^2-C^2$ 称为直线与椭圆联立的判别式，它的地位相当于
关于点与椭圆的判别式 $\frac{x^2}{a^2}+\frac{y^2}{b^2}-1$（后者判断点在椭圆内、椭圆上还是
椭圆外）。

联立二次点为 $f^{ij}=l^2g^{ij}-l^il^j$。直接计算
\[
l^il^j=\begin{pmatrix}
a^4A^2 & a^2b^2AB & -a^2AC\\
a^2b^2AB & b^4B^2 & -b^2BC\\
-a^2AC & -b^2BC & C^2
\end{pmatrix},
\]
\[
l^2g^{ij}=\begin{pmatrix}
a^2(a^2A^2+b^2B^2-C^2) & 0 & 0\\
0 & b^2(a^2A^2+b^2B^2-C^2) & 0\\
0 & 0 & -(a^2A^2+b^2B^2-C^2)
\end{pmatrix},
\]
于是
\[
f^{ij}=\begin{pmatrix}
-a^2(C^2-b^2B^2) & -a^2b^2AB & a^2AC\\
-a^2b^2AB & -b^2(C^2-a^2A^2) & b^2BC\\
a^2AC & b^2BC & -(a^2A^2+b^2B^2)
\end{pmatrix}.
\]
二次点 $f$ 描述两个点 $A,B$。过原点的直线 $m_i=(A',B',0)$ 经过 $A$ 或 $B$ 的条件是
$f^{ij}m_im_j=0$，即
\[
f^{11}A'^2+2f^{12}A'B'+f^{22}B'^2=0 .
\]
该二次方程的两个根为 $A'/B'=-k_1,-k_2$，其中 $k_1,k_2$ 是 $OA,OB$ 的斜率，故两根之积
\[
\frac{f^{22}}{f^{11}}=k_1k_2 .
\]
由 $OA\perp OB$ 得 $k_1k_2=-1$，于是
\[
\frac{b^2(C^2-a^2A^2)}{a^2(C^2-b^2B^2)}=-1,
\]
即
\[
b^2(C^2-a^2A^2)+a^2(C^2-b^2B^2)=0,
\]
化简得
\[
A^2+B^2=\Bigl(\frac{1}{a^2}+\frac{1}{b^2}\Bigr)C^2 .
\]
原点到动直线的距离为
\[
d=\frac{|C|}{\sqrt{A^2+B^2}}
=\frac{|C|}{\dfrac{\sqrt{a^2+b^2}}{ab}\,|C|}
=\frac{ab}{\sqrt{a^2+b^2}},
\]
它与直线的位置无关，是常数。因此动直线恒与以原点为心、半径 $ab/\sqrt{a^2+b^2}$ 的圆相切，
所求包络线就是这个圆。
\end{proof}

\begin{remark}
本例中两处符号容易被弄错，值得记录。其一，$l^i=g^{ij}l_j$ 的第三个分量是
$l^3=g^{3j}l_j=-C$，而非 $C$，因为 $g^{33}=-1$。其二，联立二次点应统一取
$f^{ij}=l^2g^{ij}-l^il^j$；若写成 $l^il^j-l^2g^{ij}$ 则与本节的定义相差一个整体符号，
且此时分量矩阵的 $(3,3)$ 元应为 $-(a^2A^2+b^2B^2)$ 而非 $0$。整体符号不影响
$f^{22}/f^{11}$ 这类比值，但本节一律采用 $f=l^2g-ll$，以保证与极线、对合的记号一致。
\end{remark}

\section{对偶：过外点的两条切线}

与直线和椭圆联立对偶，考虑过椭圆外一点 $A$ 作椭圆 $g$ 的两条切线。这两条切线构成的二次线可以
写成
\[
h_{ij}=a^2g_{ij}-a_ia_j,\qquad\text{即}\qquad h=a^2g-aa,
\]
其中 $a_i=g_{ij}a^j$，$a^2=a_ia^i$。事实上 $h_{ij}a^j=a^2a_i-a_i(a_ja^j)=0$，故 $A$ 在 $h$ 的核
方向中，$h$ 上的两条直线都经过 $A$；由对偶命题，这两条直线正是从 $A$ 引出的两条切线。

\begin{remark}
$h^{ij}=a^2g^{ij}-a^ia^j$ 给出的是两个切点构成的退化二次线束，有时它比 $h_{ij}$ 更便于应用。
\end{remark}

下面写出联立二次线的分量式。设两条切线的方程为 $y=k_1x+b_1$ 与 $y=k_2x+b_2$，对应的线矢为
\[
m_i=(k_1,-1,b_1),\qquad n_i=(k_2,-1,b_2),
\]
则 $h_{ij}=m_in_j+m_jn_i$，即
\[
h_{ij}=\begin{pmatrix}
2k_1k_2 & -(k_1+k_2) & k_1b_2+k_2b_1\\
-(k_1+k_2) & 2 & -(b_1+b_2)\\
k_1b_2+k_2b_1 & -(b_1+b_2) & 2b_1b_2
\end{pmatrix}.
\]
一般只会用到前四个分量，使用时注意归一化。

\section{例题：蒙日圆}

\begin{example}
椭圆 $\frac{x^2}{a^2}+\frac{y^2}{b^2}=1$，$C$ 为椭圆外一点。过 $C$ 作椭圆的两条切线互相垂直，
求 $C$ 的轨迹。
\end{example}

\begin{proof}[解]
设 $c^i=(x,y,1)$，则
\[
c_i=g_{ij}c^j=\Bigl(\frac{x}{a^2},\;\frac{y}{b^2},\;-1\Bigr),\qquad
c^2=c_ic^i=\frac{x^2}{a^2}+\frac{y^2}{b^2}-1 .
\]
所求二次线为 $h_{ij}=c^2g_{ij}-c_ic_j$。其中
\[
c_ic_j=\begin{pmatrix}
\dfrac{x^2}{a^4} & \dfrac{xy}{a^2b^2} & \dfrac{x}{a^2}\\[8pt]
\dfrac{xy}{a^2b^2} & \dfrac{y^2}{b^4} & \dfrac{y}{b^2}\\[8pt]
\dfrac{x}{a^2} & \dfrac{y}{b^2} & 1
\end{pmatrix},
\]
逐项相减得
\[
h_{ij}=\begin{pmatrix}
\dfrac{1}{a^2}\Bigl(\dfrac{y^2}{b^2}-1\Bigr) & -\dfrac{xy}{a^2b^2} & -\dfrac{x}{a^2}\\[10pt]
-\dfrac{xy}{a^2b^2} & \dfrac{1}{b^2}\Bigl(\dfrac{x^2}{a^2}-1\Bigr) & -\dfrac{y}{b^2}\\[10pt]
-\dfrac{x}{a^2} & -\dfrac{y}{b^2} & -\Bigl(\dfrac{x^2}{a^2}+\dfrac{y^2}{b^2}\Bigr)
\end{pmatrix}.
\]
两条切线都经过 $C$，以 $C$ 为原点、坐标轴与原来平行时，两条切线的方程之积的二次项部分为
\[
k_1k_2X^2-(k_1+k_2)XY+Y^2,
\]
它与 $h_{ij}$ 的二次项部分成比例，因此斜率之积 $k_1k_2$ 等于 $h_{11}/h_{22}$。由两切线垂直得
$k_1k_2=-1$，即 $h_{11}+h_{22}=0$：
\[
\frac{1}{a^2}\Bigl(\frac{y^2}{b^2}-1\Bigr)+\frac{1}{b^2}\Bigl(\frac{x^2}{a^2}-1\Bigr)=0 .
\]
化简得
\[
x^2+y^2=a^2+b^2 .
\]
因此 $C$ 的轨迹是以原点为心、半径 $\sqrt{a^2+b^2}$ 的圆，称为\textbf{蒙日圆}。
\end{proof}

\begin{remark}
本节的两道例题互为对偶：把点与线的地位互换，一道例题就变成另一道。圆与椭圆在射影几何层面并
没有区别，例中由椭圆恰好得到一个圆，原因在于圆可以用垂直结构来定义。蒙日圆不依赖坐标系的
表达式将在讨论垂直结构时给出。
\end{remark}
   % 退化圆锥曲线
% 第九节：交比、对合与调和点列
\chapter{交比、对合与调和点列}

\section{子空间上的矢量：嵌入、限制与射影}

前几节讨论的矢量都取在三维矢量空间 $V$ 中，对应平面上的点与线。但平面几何中有一类问题只涉及
一条直线上的点，或只涉及过一个定点的直线，前者称为点列，后者称为线束。直线上的点集是一维的，
过定点的直线集也是一维的，因此描述它们最自然的位置是一个二维矢量空间，而不是三维的 $V$。

设平面上的点用 $V$ 中的点矢 $a$ 描述，并要求点在一条固定的直线 $l$ 上，即
\[
l\cdot a=l_ia^i=0 .
\]
满足这一条件的全体点矢
\[
S=\{a\in V\mid l_ia^i=0\}
\]
构成 $V$ 的一个二维子空间：$S$ 对加法与数乘封闭，且恰好是一个过原点的平面。于是可以
把 $S$ 当作一个独立的二维矢量空间来研究，只在必要时才把它与 $V$ 中的平面等同起来。也就是说，
我们只关心到同构的层面。

\begin{definition}[嵌入]
设 $S\subset V$ 是二维子空间。把 $S$ 中的元素看作 $V$ 的元素，所得到的包含映射
\[
\iota:S\hookrightarrow V
\]
称为从 $S$ 到 $V$ 的\emph{嵌入}。通过嵌入，$S$ 中的二维点矢被实现为 $V$ 中的三维点矢，
几何上相当于把直线上的点画到平面上。
\end{definition}

嵌入把点矢从二维提升到三维。反方向的操作作用在对偶矢量上，由对偶矢量的函数意义给出。

\begin{definition}[限制]
设 $S\subset V$ 是二维子空间，$f\in V^*$ 是 $V$ 上的线性函数。把 $f$ 的定义域限制为
$S$，得到 $S$ 上的线性函数
\[
f_S:S\to\R,\qquad f_S(a)=f(a)\quad(a\in S),
\]
称为 $f$ 在 $S$ 上的\emph{限制}。$f\mapsto f_S$ 给出线性映射 $V^*\to S^*$。
\end{definition}

一句话，限制映射就是限制了线矢作为函数的定义域。由对偶矢量的函数意义，这一操作是自然的：
两个在 $V$ 中不同的线矢，限制到 $S$ 上之后可能给出同一个 $S^*$ 中的元素，因此限制一般不可逆。

嵌入与限制地位对称，但都不够直观。相比之下较为直观的是从线矢到点矢的映射。

\begin{definition}[射影]
设 $S$ 是对应于直线 $l$ 的二维子空间，$m$ 是任意一条不与 $l$ 平行的直线（即
$m$ 与 $l$ 不共向，二者在平面内交于唯一一点）。令
\[
a=l\times m,
\]
即把线矢 $m$ 映射为 $l$ 与 $m$ 的交点。这一映射 $V^*\to S$ 称为\emph{射影}。
\end{definition}

取分量写法。设 $e^{jk}=\eps^{ijk}l_i$，则由叉乘定义
\[
a^k=(l\times m)^k=\eps^{kij}l_im_j=m_j\eps^{jik}l_i=m_je^{jk},
\]
其中用到 $\eps^{jik}=-\eps^{ijk}$ 与 $\eps^{kij}=\eps^{jik}$。可见 $e^{jk}$ 是一个反对称的二阶张量，
它把线矢 $m_j$ 变为点矢 $a^k$。限制映射与射影映射都把三维元素变为二维元素，因此都不可逆。

以上讨论的是点列。若把点矢与对偶矢量的地位互换，就得到关于线束的嵌入、限制与射影，结论完全平行，
此处不再重复。

\section{二维矢量空间的特殊性}

二维矢量空间中的矢量、对偶矢量以及上下指标的定义与三维完全相同。差别在于维数少了一，这使
二维空间具有一些在三维中不成立的性质。

\begin{proposition}\label{prop:2d-unique}
设 $n$ 是二维空间上的非零对偶矢量，则方程 $n_ia^i=0$ 的解构成一条过原点的直线，即解空间是
一维的；它的一个（非零）解是
\[
a^i=e^{ij}n_j
\]
（这里用 $e$ 搬运指标，见定义~\ref{def:e}），且一切解都与它成比例。等价地，方程 $a^il_i=0$ 的
解为 $a^i=e^{ij}l_j$，只确定方向而不确定长度。
\end{proposition}

\begin{proof}
取定坐标系，把 $a$ 写成 $a^i=(x,1)$，则 $n_ia^i=0$ 是关于 $x$ 的一次方程
$n_1x+n_2=0$；由于 $n\ne 0$，至少有一个 $n_i$ 非零，故解唯一，解空间至多一维。直接验证：
$n_ie^{ij}n_j=e^{ij}n_in_j=0$，因为 $e^{ij}$ 反对称而 $n_in_j$ 对称。故解空间恰为一维。
\end{proof}

注意与三维的对比：三维中 $n\cdot a=0$ 的解构成一个二维平面，有无穷多个不共线的解；二维中却
只剩下一个方向。

\begin{proposition}\label{prop:2d-conic}
设 $g_{ij}$ 是二维空间上的不定二次型。把 $a$ 写成 $a^i=(x,1)$，则方程
$g_{ij}a^ia^j=0$ 是关于 $x$ 的一元二次方程，故至多有两个解；当 $g$ 不定时恰有两个实根。
即二次型在二维空间中至多确定两个方向。
\end{proposition}

在三维中，$g_{ij}a^ia^j=0$ 确定的是圆锥曲线的点集，一般有无穷多个点。二维中只剩下两个方向，
这一“至多两个解”的事实是后文对合变换只可能有两个不动点的根源。

在二维中 $\eps_{ijk}$ 没有意义：它的非零分量要求三个指标互不相同，而二维指标只取 $1,2$，因此
该张量恒为零。替代它的是反对称二阶张量。

\begin{definition}[二维反对称张量 $e_{ij}$]\label{def:e}
二维空间上非零的反对称二阶张量记为 $e_{ij}$ 与 $e^{ij}$，满足
\[
e_{ij}=-e_{ji},\qquad e^{ij}=-e^{ji},\qquad e_{ij}e^{ij}=2 .
\]
在一定坐标系下取 $e_{12}=1$，则
\[
(e_{ij})=\begin{pmatrix}0&1\\-1&0\end{pmatrix},\qquad
(e^{ij})=\begin{pmatrix}0&1\\-1&0\end{pmatrix},
\]
即有 $e^{12}=1$。此时恒有
\[
e_{ij}e^{ik}=\delta^k_j .
\]
\end{definition}

$e_{ij}$ 是二维空间中最基本的张量之一，它与度规 $g_{ij}$ 的区别在于对称性：$g_{ij}$ 对称，
$e_{ij}$ 反对称。

\subsection{二维叉乘与上下指标的转换}

\begin{definition}[二维叉乘]
两个二维矢量 $a,b$ 的\emph{叉乘}定义为
\[
a\times b=a^ie_{ij}b^j ,
\]
它是一个标量（不是矢量），简记为
\[
[ab]=a^ie_{ij}b^j .
\]
由反对称性立刻得到 $[ab]=-[ba]$，特别地 $[aa]=0$。
\end{definition}

用一个矢量 $a^i$ 与 $e_{ij}$ 缩并，得到一个对偶矢量。记
\[
a_j=a^ie_{ij},\qquad\text{反之}\qquad a^i=e^{ij}a_j .
\]
这与用度规升降指标形式上相似，实质不同：这里用来转换的是反对称张量而不是对称张量，因此
\emph{必须区分 $e_{ij}$ 两个指标的左右}。用几何语言说，$a^i$ 表示二维点矢（点），$a_i$ 表示
二维线矢（过原点的直线）；一个点矢可以画成点，也可以画成它的像直线，二者只是画法不同。

$a^i$ 与 $a_i$ 的这种关系，正是三维中\emph{射影}与\emph{限制}关系的二维版本。给定直线 $l$ 与
$m$：把 $m$ 限制到 $l$ 上得到二维线矢，相当于丢掉 $m$ 上不属于 $l$ 的部分，只留下交点附近的
方向，因此二维线矢宜画成一段直线；把 $l,m$ 作交点（射影）得到二维点矢，因此二维点矢宜画成
一个点。同一个元素的这两种画法由 $e$ 相互转换，纯粹是画法不同而已。

\section{二维伴随矩阵与指标升降}

二维情形下求逆矩阵（更准确地说，求伴随矩阵）有十分简洁的张量表达。设
\[
M=\begin{pmatrix}a&b\\c&d\end{pmatrix},
\]
其伴随矩阵（等于行列式乘以逆矩阵）为
\[
M^*=\begin{pmatrix}d&-b\\-c&a\end{pmatrix},\qquad MM^*=M^*M=(\det M)I .
\]
$M$ 与 $M^*$ 的关系只是把元素换位并改变符号，这说明它可以由 $e_{ij}$ 简单地表达。

\begin{lemma}\label{lem:adj}
设 $M_{ij}$ 是二维空间上的 $(0,2)$ 型张量，$\det M=1$，用 $e$ 升降指标。则 $M$ 与它的伴随
矩阵之间的指标关系为
\[
(M^*)^{ij}=-e^{ik}M_{kl}e^{jl},
\]
并且把 $M^{-1}$ 记作 $M^{ij}$ 时有
\[
M_{ij}M^{ik}=\delta^k_j,\qquad M_{ji}M^{jk}=\delta^k_i .
\]
\end{lemma}

\begin{proof}
\textbf{第一步：矩阵恒等式。}
记 $e=\begin{pmatrix}0&1\\-1&0\end{pmatrix}$，即 $e_{ij}$ 的矩阵，于是 $e^{T}=-e$。取
$M=\begin{pmatrix}a&b\\c&d\end{pmatrix}$，直接计算：
\[
e\,M\,e^{T}
=\begin{pmatrix}0&1\\-1&0\end{pmatrix}\begin{pmatrix}a&b\\c&d\end{pmatrix}
\begin{pmatrix}0&-1\\1&0\end{pmatrix}
=\begin{pmatrix}d&-c\\-b&a\end{pmatrix}
=\operatorname{adj}(M)^{T},
\]
其中 $\operatorname{adj}(M)=M^*=\begin{pmatrix}d&-b\\-c&a\end{pmatrix}$。

\textbf{第二步：写成指标。}
上面等式的 $(i,j)$ 分量为 $(eMe^{T})_{ij}=e_{ik}M_{kl}e_{jl}$，而
$(\operatorname{adj}M)^T_{ij}=(\operatorname{adj}M)_{ji}=(M^*)^{ji}$，故
\[
(M^*)^{ji}=e_{ik}M_{kl}e_{jl}.
\]
把 $i,j$ 都升起来（每升一次指标用 $e$ 搬运，共两次符号变化），并注意 $\det M=1$ 时
$M^{ij}=e^{ia}e^{jb}M_{ab}$ 所表示的矩阵恰为 $(M^{-1})^{ij}$（这正是 $eMe^{T}=\operatorname{adj}(M)^T
=\det M\cdot(M^{-1})^{T}$ 的指标形式），得
\[
(M^*)^{ij}=-e^{ik}M_{kl}e^{jl}.
\]

\textbf{第三步：伴随与逆。}
由 $\det M=1$ 有 $M^*=M^{-1}$，又由 $(M^*)^{ji}=e_{ik}M_{kl}e_{jl}=-(M^*)^{ij}$ 型的指标搬运可得
\[
(M^*)^{ij}=e^{ia}e^{jb}(M^*)_{ab}=e^{ia}(M^{-1})_{ab}e^{jb} .
\]
把 $M^{-1}$ 记为具有两个上指标的 $M^{ab}$，并约定缩并对齐指标位置，即得
\[
M_{ij}M^{ik}=\delta^k_j .
\]
再把 $M_{ij}$ 换成它的转置 $M_{ji}$，由($e$ 的反对称性带来的两次符号变化)同样有
\[
M_{ji}M^{jk}=\delta^k_i .
\]
\end{proof}

\begin{remark}
搬运二维下指标的是反对称的 $e$，不是对称的度规 $g$。因此 $M^{ij}$ 与 $M_{ij}$ 并不互为转置，
一般的张量满足的是上面这类“一对指标上下相消、另一对左右相消”的等式；把 $e$ 与 $g$ 混用、
或随意交换 $M$ 的两个指标都会出错。本节一律用 $e$ 搬运二维指标。
\end{remark}

由引理还可以得到常用的变形。由 ${M^j}_l=M^{jk}e_{kl}$ 与 $M_{ji}M^{jk}=\delta^k_i$ 联立，
\[
M_{ji}{M^j}_l=M_{ji}M^{jk}e_{kl}=\delta^k_ie_{kl}=e_{il},
\]
类似地，把两个指标都换成上指标后有
\[
{M_j}^k{M^j}_l=-\delta^k_l .
\]

\section{射影变换与对合变换}

\subsection{线性映射与射影变换}

在第二章中已经说明，$(1,1)$ 型张量给出矢量到矢量的线性映射。设 $F$ 是二维矢量空间上的线性映射，
其张量写作 ${f^i}_j$，则
\[
F(a)^i={f^i}_ja^j .
\]
本节总要求 $F$ 可逆，即 $\det({f^i}_j)\ne0$。

线性映射的典型例子是点列与线束之间的投影。取定一个不在直线 $L$ 上的点 $A$：对 $L$ 上的动点 $D$，
作直线 $AD$，得到一个过 $A$ 的直线；反过来，对过 $A$ 的动直线 $AD$，作它与 $L$ 的交点 $D$，
得到 $L$ 上的点。这两个对应都把点列 $L$ 与以 $A$ 为中心的线束一一对应，并且都是线性的，因为
“过两点作直线”与“求两直线交点”分别由叉乘的线性运算实现。

线性映射的复合仍是线性映射。例如把 $L$ 上的点 $D$ 经由 $A$ 投影到另一条直线 $L'$ 上的点 $E$，
再把 $E$ 经由另一定点 $A'$ 投影到 $L$ 上的点 $G$，则“$D\mapsto G$”是两次投影的复合。由于
线性映射总可以分解为若干次这样的投影，我们也把可逆线性映射称为\emph{射影变换}。

\begin{definition}[射影变换]
二维矢量空间上可逆的线性映射 $F$，以及由它诱导的点列之间、线束之间的一一对应，统称为
\emph{射影变换}。$F$ 与 $\lambda F$（$\lambda\ne0$）诱导同一个几何对应，因为矢量的长度
（即齐次坐标的规范化）不具有几何意义。
\end{definition}

如果一个射影变换把一条直线（或一个线束的中心所成的集合）映射到自身，则这一变换一般有两个
不动点，这正是命题~\ref{prop:2d-conic} 的应用：求不动点归结为解一个二次方程。

\section{交比}

\subsection{交比的定义}

我们希望用若干二维点矢构造出坐标系无关、规范无关的实数。求坐标系无关，只需通过张量积与缩并
得到标量；求规范无关，只需使每个矢量的次数上下相抵。

使用一个或两个点矢做不到：$[aa]=0$ 恒为零，$[ab]$ 在把 $a$ 换成 $\lambda a$ 时改变 $\lambda$
倍。使用四个点矢则可以构造出
\[
J_{ab,cd}=\frac{[ac][bd]}{[ad][bc]} ,
\]
它关于每个矢量的伸缩都是零次的：$a\to\lambda a$ 使分子分母各出现一次 $\lambda$，因而相消。
这个量称为四点 $a,b,c,d$ 的\emph{交比}，在不引起混淆时简记为 $J$。

\begin{definition}[交比]
设 $a,b,c,d$ 是二维空间中的四个点矢，且
$[ad]\ne0$、$[bc]\ne0$。称
\[
J_{ab,cd}=\frac{[ac][bd]}{[ad][bc]}
\]
为这四点的\emph{交比}。若其中某个叉乘为零，则交比取 $0$ 或 $\infty$。
\end{definition}

若两条直线上各取四个点，且它们由同一个射影变换一一对应，则由于射影变换是线性的而交比由缩并
构成，对应四点的交比相同。

\begin{proposition}\label{prop:J-inv}
设 $F$ 是二维空间上的可逆线性映射，则
\[
J_{F(a)F(b),F(c)F(d)}=J_{ab,cd}.
\]
即交比在射影变换下不变。
\end{proposition}

\begin{proof}
设 $F$ 的张量为 ${f^i}_j$，则 $F(a)^i={f^i}_ja^j$。由行列式的定义，
$[F(a)F(b)]=\det(f)[ab]$，即叉乘在 $F$ 作用下统一乘上因子 $\det f$。代入交比的定义，分子与
分母各含两个因子 $\det f$，故相消。
\end{proof}

\subsection{两个恒等式}

交比的全部代数性质来自下面两个恒等式。

\begin{lemma}[二维叉乘的线性关系]\label{lem:cyc}
对任意三个二维矢量 $a,b,c$ 有
\[
[ab]c+[bc]a+[ca]b=0 .
\]
\end{lemma}

\begin{proof}
三个矢量在二维空间中必线性相关，设 $\alpha a+\beta b+\gamma c=0$ 且系数不全为零。若
$[ab]\ne0$，则 $a,b$ 线性无关，可把它们取作基，把 $c$ 写成 $c=xa+yb$。于是
\[
[bc]=x[ba]=-x[ab],\qquad [ca]=y[ab],
\]
代入左边得
\[
[ab](xa+yb)-x[ab]a+y[ab]b=0 .
\]
若 $[ab]=0$，则 $a,b$ 线性相关，此时 $[bc]a$ 与 $[ca]b$ 都与 $a$ 共线且系数成比例，逐项检查
（或直接用坐标展开）同样得零。
\end{proof}

\begin{lemma}\label{lem:four}
对任意四个二维矢量 $a,b,c,d$ 有
\[
[ab][cd]+[ac][db]+[ad][bc]=0 .
\]
\end{lemma}

\begin{proof}
在引理~\ref{lem:cyc} 中把 $c$ 换成 $d$：
\[
[ab]d+[bd]a+[da]b=0 .
\]
两端与 $c$ 作二维叉乘，即都取 $[\;\cdot\;,c]$。由于 $[uv]$ 对每个自变量都是线性的，左端逐项
展开给出
\[
[ab][dc]+[bd][ac]+[da][bc]=0 .
\]
由反对称性，$[dc]=-[cd]$，$[bd]=-[db]$，$[da]=-[ad]$，代入并整理得
\[
-[ab][cd]-[db][ac]-[ad][bc]=0,
\]
移项即
\[
[ab][cd]+[ac][db]+[ad][bc]=0 .
\]
\end{proof}

\subsection{交比的性质}

由 $[ab]=-[ba]$（即交换两个自变量使叉乘变号）与引理~\ref{lem:four}，可以导出交比的如下性质。

\begin{proposition}[交比在置换下的变化]\label{prop:J-perm}
设 $J=J_{ab,cd}$，且四个点互不相同（即 $[ab],[ac],\dots$ 均不为零）。则
\begin{enumerate}
\item 交换两组之间的元素保持交比不变：$J_{cd,ab}=J_{ba,dc}=J$；
\item 交换同一组内的两个元素使交比变为倒数：$J_{ba,cd}=J_{ab,dc}=\dfrac1J$；
\item 把两组的元素重新搭配，交比与 $1$ 的差互换：
\[
J_{ac,bd}=J_{bd,ac}=1-J .
\]
\end{enumerate}
\end{proposition}

\begin{proof}
选取基使 $[ab]=1$：取 $a=e_2$、$b=e_1$，则 $e_{21}=e_{12}=1$。把 $c,d$ 写成
\[
c=x_ca+y_cb,\qquad d=x_da+y_db .
\]
由第一、三节给的二维叉乘，$[uv]=u^ie_{ij}v^j$ 对两个自变量都是线性的，于是
\[
[ac]=y_c,\quad [bc]=-x_c,\quad [ad]=y_d,\quad [bd]=-x_d,\quad [ab]=1,
\]
以及
\[
[cd]=[x_ca+y_cb,\ x_da+y_db]=x_cy_d[a b]+y_cx_d[b a]=x_cy_d-x_dy_c .
\]
于是
\[
J=J_{ab,cd}=\frac{[ac][bd]}{[ad][bc]}=\frac{y_c(-x_d)}{y_d(-x_c)}=\frac{x_d}{x_c}\cdot\frac{y_c}{y_d}.
\]

第一款：交换两组元素只是把两个点对换个位置，计算得
\[
J_{cd,ab}=\frac{[ca][db]}{[cb][da]}=\frac{(-y_c)(x_d)}{(x_c)(-y_d)}=J,
\]
同理 $J_{ba,dc}=J$。

第二款：由定义 $J_{ba,cd}=\dfrac{[bc][ad]}{[bd][ac]}$，这正是 $J_{ab,cd}$ 的倒数，故
$J_{ba,cd}=1/J$。同理交换 $c,d$ 得 $J_{ab,dc}=\dfrac{[ac][db]}{[ab][cd]}$，它与 $J$ 的比值为
$1/J$，故也等于 $1/J$。（也可直接由第一款推得：把第一款用于以 $b,a$ 为前两点的交比即得此式。）

第三款：由上面的坐标式直接计算
\[
J_{ac,bd}=\frac{[ab][cd]}{[ad][cb]}=\frac{[cd]}{y_d\,x_c}
=\frac{x_cy_d-x_dy_c}{x_cy_d}=1-\frac{x_dy_c}{x_cy_d}=1-J .
\]
同样 $J_{bd,ac}=1-J$。
\end{proof}

\begin{remark}\label{rem:J-not-1}
四个点互不相同时，交比不可能取 $0,1,\infty$。$J=0$ 要求分子 $[ac][bd]=0$，$J=\infty$ 要求
分母 $[ad][bc]=0$，二者都与四点互异（任意两个不成比例，故任意叉乘非零）矛盾。$J=1$ 时由
命题~\ref{prop:J-perm} 第三款得 $J_{ac,bd}=1-J=0$，同样矛盾。因此四点互异时 $J$ 是
$0,1,\infty$ 以外的实数。
\end{remark}

\section{调和分割}

\begin{definition}[调和分割]
设 $J_{ab,cd}=-1$，则称 $a,b$ 与 $c,d$ 互相\emph{调和分割}，称 $a,b,c,d$ 构成\emph{调和点列}。
\end{definition}

由命题~\ref{prop:J-perm}，$J=-1$ 这一条件在所有置换下保持形式（$1/J=-1$，$1-J=2$，而重新
搭配元素得到的 $-1$ 仍落在调和的条件里），因此“哪两点与哪两点调和分割”这一说法的分组方式
是确定的。

\section{对合变换}

\subsection{对合变换与它的不动点}

\begin{definition}[对合变换]
二维矢量空间上的可逆线性映射 $F$ 称为\emph{对合变换}，若存在非零实数 $\lambda$ 使得
\[
F(F(a))=\lambda a\qquad\text{对一切 }a\text{ 成立}.
\]
\end{definition}

对合变换的意义是它把所有的矢量两两配对：$a$ 与 $F(a)$ 配成一对，称为互相对合。若 $F$ 不是
对合变换，则由一对 $a,F(a)$ 可以区分原像与像；若 $F$ 是对合变换，则用 $F$ 作用这一对元素时它们
同时变成对方，原像与像的地位完全对称，因而这一对元素确实被“配对”了。

既然对合变换中原像与像的地位对称，可以期望它有特别对称的代数形式。结论是它可以由一个对称的
$(0,2)$ 型张量决定。

\begin{theorem}[对合变换由对称可逆张量决定]\label{thm:involution}
设 $g_{ij}$ 是二维空间上的对称、可逆 $(0,2)$ 型张量，用 $e$ 升降二维指标。定义线性映射
\[
F(a)^k=g_{ij}a^ie^{jk}\qquad\text{即 }F(a)^k={f^k}_ia^i,\quad {f^k}_i=g_{ij}e^{jk}.
\]
则 $b=F(a)$ 当且仅当 $g_{ij}a^ib^j=0$，并且
\begin{enumerate}
\item $\det F=-\det g\ne0$，故 $F$ 可逆；
\item $F(F(a))=-\det(g)\,a$ 对一切 $a$ 成立，特别地 $F$ 是对合变换。
\end{enumerate}
\end{theorem}

\begin{proof}
\textbf{第一步：$b=F(a)$ 等价于配对条件。}
对任意 $c$ 计算
\[
g_{cd}a^cF(a)^d=g_{cd}a^cg_{ij}a^ie^{jd}.
\]
把哑指标 $c$ 与 $i$ 互换：$a^ca^i$ 对称、$e^{jd}$ 与 $c,i$ 无关，而 $g_{cd}g_{ij}$ 互换 $c,i$
后不变，故右端等于自身的相反数，从而必为零。于是 $g_{ij}a^ib^j=0$。反过来，由
命题~\ref{prop:2d-unique}，方程 $g_{ij}a^ib^j=0$ 关于 $b$ 的解空间是一维的；又 $F(a)=0$ 只在
$a=0$ 时发生（$a\mapsto g_{ij}a^i$ 与 $b\mapsto b^je^{jk}$ 都可逆），故一切解都与 $F(a)$ 成比例。

\textbf{第二步：用矩阵乘法计算 $F^2$。}
取定坐标，把 $g_{ij}$ 写成对称矩阵 $G=\begin{pmatrix}a&b\\b&d\end{pmatrix}$，记
$\Delta=\det G=ad-b^2\ne0$，并把 $e_{ij}$ 与 $e^{ij}$ 都取成矩阵
$J=\begin{pmatrix}0&1\\-1&0\end{pmatrix}$。由于 ${f^k}_i=g_{ij}e^{jk}$ 中缩并的是 $g$ 的第二个
指标与 $e$ 的第一个指标，$F$ 的矩阵恰为
\[
f=GJ=\begin{pmatrix}a&b\\b&d\end{pmatrix}\begin{pmatrix}0&1\\-1&0\end{pmatrix}
=\begin{pmatrix}-b&a\\-d&b\end{pmatrix}.
\]
于是
\[
f^2=\begin{pmatrix}-b&a\\-d&b\end{pmatrix}^2
=\begin{pmatrix}b^2-ad&0\\0&b^2-ad\end{pmatrix}
=-\Delta\,I,
\]
即 $F(F(a))=-\Delta\,a=-\det(g)\,a$。由于 $\det g\ne0$，$F$ 是对合变换。又
$\det f=\det G\cdot\det J=-\Delta\ne0$，故 $F$ 可逆。
\end{proof}

\begin{remark}
使用记号时务必写清用哪个张量搬运指标：二维点列中用 $e_{ij},e^{ij}$，三维中才用 $\eps_{ijk}$；
若改用度规 $g$ 升降指标，一般公式（例如 $M_{ij}M^{jk}\ne\delta^k_i$）都要相应改变。$e$ 是反对称
张量，写这类表达式时上下指标与前后顺序都不能随意调换，本节中以 $b^k=g_{ij}a^ie^{jk}$ 为准。
定理只断言 $F(F(a))$ 与 $a$ 成正比且比例常数非零，其正负与数值随 $e$ 的归一化约定而变。
\end{remark}

\begin{proposition}[不动点]\label{prop:fixed}
设 $F$ 由对称可逆张量 $g_{ij}$ 按定理~\ref{thm:involution} 给出，则 $a$ 是 $F$ 的不动点
（即 $F(a)$ 与 $a$ 成比例）当且仅当
\[
g_{ij}a^ia^j=0 .
\]
若 $g$ 不定，则 $F$ 恰有两个不动点（两个方向），记作 $p,q$。
\end{proposition}

\begin{proof}
设 $F(a)=\mu a$ 且 $a\ne0$。由定理~\ref{thm:involution} 第一步，
$g_{ij}a^iF(a)^j=g_{ij}a^i(\mu a)^j=0$，即
\[
\mu\,g_{ij}a^ia^j=0 .
\]
因 $\mu\ne0$，得 $g_{ij}a^ia^j=0$。

反之设 $g_{ij}a^ia^j=0$ 且 $a\ne0$。仍由定理~\ref{thm:involution} 第一步，取 $b=F(a)$ 时有
$g_{ij}a^ib^j=0$；而取 $b=a$ 时也成立，即 $a$ 自身也满足这个关于 $b$ 的线性方程
$g_{ij}a^i(\;\cdot\;)^j=0$。由命题~\ref{prop:2d-unique}，该方程的解空间是一维的，故 $F(a)$ 与
$a$ 成比例，即 $a$ 是不动点。

解的个数由命题~\ref{prop:2d-conic} 给出：$g$ 不定时二次方程恰有两个实根，即恰有两个不动点
（两个方向）。
\end{proof}

用坐标写出对合的一般形状。按通常规范化取 $a^i=(x_1,1)$、$b^i=(x_2,1)$，则
$g_{ij}a^ib^j=0$ 是关于 $x_1,x_2$ 的双线性方程
\[
Ax_1x_2+B(x_1+x_2)+C=0 ,
\]
其中 $A,B,C$ 由 $g_{ij}$ 的分量给出。取 $a=b$，不动点满足
\[
Ax^2+2Bx+C=0 ,
\]
这正是二次方程有两个根的体现，也与韦达定理相呼应：两根之积与两根之和分别由 $C/A$ 与 $-2B/A$
给出。对一维线束，把线矢坐标写成 $m_i=(k_1,1)$、$n_i=(k_2,1)$，对合同样表示为
$Ak_1k_2+B(k_1+k_2)+C=0$。

\subsection{对合与调和点列}

\begin{theorem}\label{thm:inv-harmonic}
设对合变换 $F$ 的不动点为 $p,q$，且 $a,b$ 是互相对合的一对点（即 $b$ 与 $F(a)$ 成比例）。
则 $p,q,a,b$ 构成调和点列，即
\[
J_{pq,ab}=-1 .
\]
反之，若 $p,q$ 是给定的两个不同方向，则使 $a,b$ 调和分割 $p,q$ 的配对关系是一个对合变换，
其不动点为 $p,q$。
\end{theorem}

\begin{proof}
\textbf{第一步：把 $J_{pq,ab}=-1$ 化为一个对称双线性方程。}
取 $p,q$ 为基，把 $a,b$ 写成
\[
a=x_ap+y_aq,\qquad b=x_bp+y_bq .
\]
由 $[uv]$ 的双线性与反对称性，
\[
[pa]=y_a[pq],\quad [qa]=-x_a[pq],\qquad
[pb]=y_b[pq],\quad [qb]=-x_b[pq].
\]
代入交比的定义得
\[
J_{pq,ab}=\frac{[pa][qb]}{[pb][qa]}
=\frac{y_a\,(-x_b)}{y_b\,(-x_a)}=\frac{x_by_a}{x_ay_b},
\]
（公因子 $[pq]^2$ 已相消）。因此
\[
J_{pq,ab}=-1 \iff x_by_a+x_ay_b=0 .
\]
另一方面，在基 $p,q$ 下直接展开
\[
\bigl(p_iq_j+p_jq_i\bigr)a^ib^j
=\bigl(p_iq_j+p_jq_i\bigr)(x_ap+y_aq)^i(x_bp+y_bq)^j
=[pq]^2\bigl(x_by_a+x_ay_b\bigr),
\]
其中用到 $p_ip^i=q_iq^i=0$、$p_iq^j[p,q]$ 型的配对（这只需在两个基矢上逐项代入验证）。
于是调和条件等价于
\[
g_{ij}a^ib^j=0,\qquad g_{ij}=p_iq_j+p_jq_i=\sym{p}{q}.
\]

\textbf{第二步：这个配对是对合变换。}
解出 $b$：由第二节的命题，$g_{ij}a^ib^j=0$ 作为 $b$ 的方程解空间是一维的，且
\[
b^k=-g_{ij}a^ie^{jk}
\]
是它的一个解。于是配对由线性映射
\[
F(a)^k=g_{ij}e^{jk}a^i
\]
给出（负号只改变配对中的定向）。由定理~\ref{thm:involution}，$F$ 线性、可逆且
$F(F(a))=\lambda a$，故它是对合变换。

\textbf{第三步：不动点。}
设 $a=xe_1+e_2$，计算 $[F(a),a]=0$ 的条件。由
\[
F(a)^k=g_{1j}e^{jk}x+g_{2j}e^{jk}
\]
以及 $[F(a),a]=e_{2k}F(a)^k\cdot x-e_{1k}F(a)^k$，把上式代入并利用 $e^{ik}e_{jk}=\delta^i_j$
的缩并（等价地利用二维恒等式 $e_{ik}e^{jk}=\delta^j_i$）得
\[
[F(a),a]=-\bigl(g_{11}x^2+2g_{12}x+g_{22}\bigr).
\]
而 $g_{ij}a^ia^j=g_{11}x^2+2g_{12}x+g_{22}$，故 $F(a)$ 与 $a$ 平行当且仅当
$g_{ij}a^ia^j=0$。这与命题~\ref{prop:fixed} 一致，且说明 $g$ 不定时恰有两个不动点。

\textbf{反向结论。}
方程 $J_{pq,ab}=-1$ 关于 $a,b$ 显然对称，由前两步它等价于 $g_{ij}a^ib^j=0$，给出 $a$ 与 $b$
之间的一个可逆线性配对，并且 $a$ 的像再取像与 $a$ 成比例，故配对是对合；又 $a=p$ 或 $a=q$ 时
方程自动成立，故 $p,q$ 是它的不动点。
\end{proof}

\subsection{对合的对称积形式}

调和点列的条件可以写成完全对称的形状。由定义 $[pa]=p_ia^i$、$[qa]=q_ia^i$（这里用
$p_i=p^je_{ji}$），条件 $\frac{[pa][qb]}{[pb][qa]}=-1$ 等价于
\[
[pa][qb]+[pb][qa]=0 .
\]
把 $[pa][qb]=p_ia^iq_jb^j$、$[pb][qa]=p_ib^iq_ja^j$ 代入，左边等于
\[
\bigl(p_iq_j+p_jq_i\bigr)a^ib^j=0 .
\]
因此若记
\[
g_{ij}=p_iq_j+p_jq_i ,
\]
则对合由 $g_{ij}a^ib^j=0$ 给出。采用对称积的简写
\[
g=\sym{p}{q},
\]
其中小括号表示对称张量积。于是得到：

\begin{proposition}\label{prop:inv-fixedpts}
对合变换由它的两个不动点完全确定：确定对合的 $(0,2)$ 型张量 $g$ 恰是两个不动点的对偶矢量的
对称积 $g_{ij}=p_iq_j+p_jq_i$。换言之，$g$ 是 $p,q$ 的对偶矢量的对称积构成的二次线。
\end{proposition}

\begin{proof}
设 $p,q$ 是不动点且 $g_{ij}a^ib^j=0$ 定义对合。对任意一对对合点 $a,b$，由定理
\ref{thm:inv-harmonic} 有 $[pa][qb]+[pb][qa]=0$，即 $(p_iq_j+p_jq_i)a^ib^j=0$。因此
$h_{ij}=p_iq_j+p_jq_i$ 与 $g_{ij}$ 给出同一个对合关系；而在二维中由命题~\ref{prop:2d-unique}，
满足同一配对关系的对称 $(0,2)$ 型张量彼此只差一个常数因子，故 $g$ 与 $h$ 成比例，特别地
$h$ 也正是由 $p,q$ 决定的那个对合张量。
\end{proof}

等式 $[pa][qb]+[pb][qa]=0$ 的结构值得注意：不仅原像与像 $a,b$ 之间是对称的，两个不动点 $p,q$
之间也是对称的，而且把 $(ab)$ 与 $(pq)$ 绑定后交换这两组也是对称的。最后这一点说明：如果把对合的
一对 $a,b$ 当作新的不动点，则原来的 $p,q$ 关于这个新的对合恰好互相对合。

以上讨论只涉及直线上的点列。对过定点的线束讨论完全相同，可以类似地定义四条直线的交比、线束上
的对合以及调和线束。

\section{子空间上的应用}

\subsection{退化二次曲线：过 $C$ 的调和线束}

设平面上有一条定直线 $l$，$a,b$ 两点在 $l$ 上；另有两条直线 $m,n$ 分别交 $l$ 于 $p,q$
（即 $p$ 是 $m$ 与 $l$ 的交点，$q$ 是 $n$ 与 $l$ 的交点）。问：在什么条件下 $a,b,p,q$ 构成
调和点列？

用二维矢量的语言，这一条件即
\[
(p_iq_j+p_jq_i)a^ib^j=0 .
\]
目前 $p,q,a,b$ 都是二维元素，需要改写成三维矢量的等式。$a,b$ 是嵌入平面上的点，因此可以原封
不动地当作三维点矢使用；而 $p_i,q_i$ 本质上是 $m_i,n_i$ 在 $l$ 上的限制。既然限制只是限制了
定义域，而 $a,b$ 同在这一定义域（即 $l$ 对应的子空间）中，故
\[
p_ia^i=m_ia^i,\qquad q_ib^j=n_jb^j ,
\]
于是上述条件可以改写为
\[
(m_in_j+m_jn_i)a^ib^j=0 .
\]
这一表达式中所有对象都是三维的了。

设
\[
h_{ij}=m_in_j+m_jn_i=\sym{m}{n},
\]
即 $h$ 是两条直线 $m,n$ 构成的退化二次线。则条件 $h_{ij}a^ib^j=0$ 的含义是：

\begin{proposition}\label{prop:degenerate}
设 $h=\sym{m}{n}$ 由两条直线 $m,n$ 构成。则 $h_{ij}a^ib^j=0$ 当且仅当 $h$ 所代表的两条直线
与直线 $a\times b$ 的两个交点被 $a,b$ 调和分割。
\end{proposition}

\begin{proof}
设直线 $a\times b$ 与 $m,n$ 的交点分别为 $p,q$，即 $p=m\times(a\times b)$、
$q=n\times(a\times b)$。由定义 $p,q$ 都是 $m,n$ 限制到直线 $a\times b$ 上以后得到的点，而由
上面的改写 $h_{ij}a^ib^j=0$ 等价于 $(p_iq_j+p_jq_i)a^ib^j=0$，即 $p,q$ 与 $a,b$ 调和分割。
\end{proof}

从张量的角度看，$(0,2)$ 型张量可以理解为把两个矢量映为实数的二元函数，因此可以像线性函数那样
对它的定义域作限制。于是
\[
(pq)_{ij}=p_iq_j+p_jq_i
\]
正是 $h_{ij}$ 在 $l$ 上的限制：限制后的张量作用在 $a,b$ 上与限制前相同。

\subsection{圆锥曲线的限制给出对合}

退化二次线限制到 $l$ 上给出 $l$ 上的一个对合。自然的问题是：若 $h$ 不退化而为圆锥曲线 $g$，
它在 $l$ 上的限制是否也给出对合？答案是肯定的。

\begin{theorem}[圆锥曲线在直线上的限制]\label{thm:conic-restrict}
设 $g_{ij}$ 是三维空间中的对称 $(0,2)$ 型张量（圆锥曲线），$l$ 是一条直线对应的二维子空间。
则 $g$ 在 $l$ 上的限制 $g'_{ij}$ 仍是对称的 $(0,2)$ 型张量，因而确定 $l$ 上的一个对合变换：
对合的两点 $a,b$ 由
\[
g_{ij}a^ib^j=0
\]
确定。若 $g$ 与 $l$ 有两个交点 $p,q$，则 $p,q$ 是这个对合的两个不动点。若 $g$ 退化为二次线，
结论同样成立。
\end{theorem}

\begin{proof}
限制不改变对称性：$g'_{ij}$ 是由 $g$ 在 $l$ 的基底上取值得到的对称 $(0,2)$ 型张量。若 $l$ 不与
$g$ 相切，则 $g'$ 可逆，定理~\ref{thm:involution} 说明它确定 $l$ 上的一个对合变换，
配对由 $g_{ij}a^ib^j=0$ 给出。

再证交点是不动点。设 $p,q\in l$ 满足 $g_{ij}p^ip^j=g_{ij}q^iq^j=0$。因为 $p,q$ 都在 $l$ 上，
把 $g$ 换成限制 $g'$ 后作用在 $l$ 中的矢量上结果不变，故
\[
g'_{ij}p^ip^j=g'_{ij}q^iq^j=0 .
\]
由命题~\ref{prop:fixed}，这正是对合的不动点的定义，因此 $p,q$ 是对合的两个不动点。

若 $l$ 与 $g$ 相切，则 $g'$ 退化（秩一），此时配对仍由 $g_{ij}a^ib^j=0$ 定义的对合给出，
但两个不动点重合为切点，是上一段当 $p=q$ 时的极限情形。
\end{proof}

把点与线、点矢与线矢的地位互换，得到对偶命题。

\begin{theorem}[对偶形式：圆锥曲线在点上的限制]\label{thm:conic-dual}
设 $g$ 是圆锥曲线，$a$ 是一个点。则 $g$ 在 $a$ 上的限制（作用于过 $a$ 的直线）确定线束上的一个
对合变换：对合的两条直线 $m,n$ 由
\[
g^{ij}a_im_j=0
\]
确定，其中 $g^{ij}=(g^{-1})^{ij}$ 是 $g$ 的逆（注意此处必须用逆而不能用 $g$ 升指标）。若从 $a$
可向 $g$ 引两条切线，其线矢记为 $p,q$，则 $p,q$ 是对合的不动点。若 $g$ 退化为二次点，
结论同样成立。
\end{theorem}

\begin{proof}
圆锥曲线有两种相互对偶的表示：点表示用 $g_{ij}$，线表示用 $g^{ij}=(g^{-1})^{ij}$；线矢 $m$ 的
极点由
\[
A^i=(g^{-1})^{ij}m_j
\]
给出。在过点 $a$ 的线束中，把直线 $m$ 配到它的\emph{对合线}：先作 $m$ 与点 $a$ 确定的那一点
$A=a\times m$（即 $m$ 上关于 $a$ 的对合方向），再作 $A$ 的极线
\[
n_i=g_{ij}A^j=g_{ij}\eps^{jkl}a_km_l .
\]
由 $n_ia^i=g_{ij}\eps^{jkl}a^ia_km_l=0$（$\eps$ 对 $i,k$ 反对称而 $a^ia_k$ 对称），
$n$ 确实也过 $a$，故这一配对是线束到自身的映射。它是线性的、可逆的，并且作用两次回到 $m$
（因为 $A$ 与 $n$ 的对偶关系是对称的），故为对合变换。再对它降指标可得等价的条件
\[
g^{ij}a_im_j=0 ,
\]
其中必须用 $g$ 的逆来搬运指标。

不动点由 $m=n$ 给出，即 $m$ 与其极点 $A$ 满足 $A\times m$ 与 $m$ 成比例，等价于
$(g^{-1})^{ij}m_im_j=0$。这正是从 $a$ 引出的两条切线所满足的方程（切线的极点即切点，
而切点在切线上），因此两条切线是对合的两个不动点。
\end{proof}

\subsection{极点极线与调和点列}

上述结论给出极点极线的一个基本性质。

\begin{theorem}[极线与调和分割]\label{thm:polar}
若关于圆锥曲线 $g$ 有 $a\cdot b=g_{ij}a^ib^j=0$（即 $a$ 在 $b$ 的极线上，或 $b$ 在 $a$ 的
极线上），则 $a,b$ 调和分割直线 $a\times b$ 与 $g$ 的两个交点 $p,q$。
\end{theorem}

\begin{proof}
记 $g_{aa}=g_{ij}a^ia^j$、$g_{ab}=g_{ij}a^ib^j$、$g_{bb}=g_{ij}b^ib^j$，设 $a,b$ 不成比例。
在直线 $a\times b$ 上取基 $a,b$，把 $p,q$ 写成
\[
p=a+s_pb,\qquad q=a+s_qb .
\]
由 $p,q$ 在圆锥曲线上，$s_p,s_q$ 是方程
\[
g_{aa}+2g_{ab}s+g_{bb}s^2=0
\]
的两个根，故（韦达定理）
\[
s_p+s_q=-\frac{2g_{ab}}{g_{bb}},\qquad s_ps_q=\frac{g_{aa}}{g_{bb}} .
\]

先算交比。由二维叉乘的双线性与反对称性，并记 $C=[ab]\ne0$：
\[
[pa]=[a+s_pb,\ a]=s_p[b a]=-s_pC,\qquad
[pb]=[a+s_pb,\ b]=C,
\]
同理 $[qa]=-s_qC$、$[qb]=C$。于是
\[
J_{pq,ab}=\frac{[pa][qb]}{[pb][qa]}
=\frac{(-s_pC)\cdot C}{C\cdot(-s_qC)}=\frac{s_p}{s_q}.
\]
因此
\[
J_{pq,ab}=-1 \iff s_p=-s_q \iff s_p+s_q=0 \iff g_{ab}=0,
\]
即 $g_{ij}a^ib^j=0$。这正是要证的命题。

（也可直接算 $g_{ij}p^iq^j$：由双线性，
\[
g_{ij}p^iq^j=g_{aa}+g_{ab}(s_p+s_q)+g_{bb}s_ps_q
=g_{aa}+g_{bb}\cdot\frac{g_{aa}}{g_{bb}}=0,
\]
其中用到 $g_{ab}=0$ 与 $s_ps_q=g_{aa}/g_{bb}$；这说明 $p,q$ 也互相对合，与交点是不动点、
$a,b$ 与 $p,q$ 调和分割相一致。）
\end{proof}

\subsection{两个附注}

\begin{remark}[不同的圆锥曲线可以给出相同的限制对合]
不同的圆锥曲线在同一个子空间上的限制可能是相同的。最典型的例子是：椭圆 $g$ 限制到直线 $l$ 上
得到的对合，由 $g$ 与 $l$ 的两个交点 $p,q$ 确定；而任何与 $l$ 交于 $p,q$ 的二次曲线都限制出
同一个对合。

设 $p,q$ 两处的切线构成退化二次线
\[
h_{ij}=l^2g_{ij}-l_il_j,\qquad l^2=g^{ij}l_il_j .
\]
从几何上看，$h$ 与 $l$ 的两个交点正是 $p,q$，因此 $h$ 与 $g$ 在 $l$ 上给出相同的对合。从代数
上看，把 $h$ 限制到直线 $l$（即对满足 $l_ix^i=0$ 的矢量作用）时，第二项 $l_il_j$ 的两因子中必
有一个作用在零矢量上，故第二项的限制恒等于零；于是
\[
h'_{ij}=l^2g'_{ij} .
\]
限制相差一个非零常数因子，给出同一个对合关系。
\end{remark}

\begin{remark}[调和线束]
在上面的构图中，不仅 $p,q,a,b$ 构成调和点列，它们与过某定点 $C$（$C$ 不在直线 $a\times b$ 上）
所连的四条直线还构成调和线束。这是因为从点 $C$ 出发的投影是一个射影变换，而交比在射影变换下
不变，所以点列的调和性与线束的调和性同时成立。

若在 $h=\sym{m}{n}$ 之外再引入 $f=\sym{a}{b}$，则这一关系可以更简洁地写成
\[
f:h=0 ,
\]
即两个对称 $(0,2)$ 型张量的二次缩并为零。
\end{remark}

\section{与诱导度规的类比}

上面的计算可以在微分几何的语言中重新表述。把圆锥曲线对应的度规 $g$ 限制在直线 $l$ 上，得到的
$h$ 用几何的话说就是 $g$ 在“超曲面” $l$ 上的\emph{诱导度规}。事实上，三维空间中由线矢 $l_i$
定义的超曲面 $\{x\mid l_ix^i=0\}$，其诱导度规为
\[
h_{ij}=l^2g_{ij}-l_il_j,\qquad l^2=g^{ij}l_il_j ,
\]
这里 $l_i$ 相当于超曲面的法余矢。直线 $l$ 上的点相当于超曲面上的点，$l$ 的地位相当于超曲面的
法余矢。

在这一对应下：
\begin{enumerate}
\item 超曲面的类时、类空、类光分别对应于直线 $l$ 与椭圆相离、相交、相切；
\item 超曲面上的类光方向对应于直线与椭圆的交点；
\item 类光超曲面的诱导度规是退化的，对应于椭圆在相切的直线上限制得到的对合变换是退化的。
\end{enumerate}

结合前面的命题，当 $l$ 与圆锥曲线相切时，$l^2=0$，于是 $h_{ij}=-l_il_j$ 是退化的对称张量，
它给出的对合只有一个不动点（切点），这正是极限情形下两个不动点重合的表现。
   % 交比、对合与调和点列
% 第十节例题：二次点列的计算
\chapter{例题：二次点列的计算}

本章的例题属于二次点列与二次对合的内容，可以看作第九节、第十节以及更早各节知识的练习。
所用的工具主要是二次调和点列、极点与极线，以及初等的代数运算；计算一律采用张量记号，
上指标为点矢，下指标为线矢，$a\cdot b=g_{ij}a^ib^j$。

\section{二次曲线与直线的第二个交点}

\begin{example}\label{ex:10-1}
已知圆锥曲线 $g$ 及其上的一点 $a$。点 $s$ 不在 $g$ 上（$s^2\neq0$），直线 $a\times s$ 与 $g$
交于两点，其中一点是 $a$，另一点记为 $b$。求 $b$ 的表达式。
\end{example}

\begin{proof}[解]
因为 $a,b,s$ 共线，可设 $b=a+\lambda s$。由 $b$ 在 $g$ 上，
\[
b^2=(a+\lambda s)^2=a^2+2\lambda\,a\cdot s+\lambda^2s^2
=\lambda\left(2a\cdot s+\lambda s^2\right)=0,
\]
其中用到 $a^2=0$。除去对应于 $\lambda=0$（即 $b=a$）的根，得
\[
\lambda=-\frac{2a\cdot s}{s^2},
\]
于是
\[
b=a-\frac{2a\cdot s}{s^2}\,s
\quad\Longleftrightarrow\quad
b=s^2a-2(a\cdot s)\,s .
\]
\end{proof}

结果的特点是：$b$ 关于 $s$ 是二次的（$s$ 以二次点 $ss$ 的形式出现），关于 $a$ 是线性的。
若把 $s$ 换成更一般的二次点，例如取 $k^{ij}=(ef)^{ij}=e^if^j+f^ie^j$，就得到下面的例题。

\section{退化二次点}

\begin{example}\label{ex:10-2}
已知圆锥曲线 $g$ 及其上的一点 $a$。二次点 $k=(ef)$ 由 $g$ 上的两点 $e,f$ 生成，即
$k^{ij}=e^if^j+f^ie^j$。由例 \ref{ex:10-1} 的办法作出点 $b$（把两条直线 $a\times e$、$a\times f$
分别与 $g$ 的另一交点记为 $p,q$，再作 $p,q$ 处的切线，其交点为 $b$）。求 $b$。
\end{example}

\begin{proof}[解]
由例 \ref{ex:10-1} 的结论，分别取 $s=e$ 与 $s=f$，得
\[
p=e^2a-2(a\cdot e)\,e,\qquad q=f^2a-2(a\cdot f)\,f .
\]
由于 $e,f$ 在 $g$ 上，$e^2=f^2=0$，故
\[
p=-2(a\cdot e)\,e,\qquad q=-2(a\cdot f)\,f .
\]
于是 $b$ 是 $p,q$ 关于 $g$ 的极点，即 $b$ 与 $p\times q$ 互为极点和极线：
\[
b\parallel p\times q
=\left(-2(a\cdot e)e\right)\times\left(-2(a\cdot f)f\right)
=4(a\cdot e)(a\cdot f)\,(e\times f).
\]
因此 $b$ 与 $e\times f$ 确定同一条直线的极点。直接验证可得，$e\times f$ 的极点为
\[
b=(a\cdot f)\,e+(a\cdot e)\,f-(e\cdot f)\,a .
\]
\end{proof}

\begin{proof}[验证]
记上式给出的矢量（再乘以 $-2$）为 $b$。直接计算
\[
(a\cdot f)e+(a\cdot e)f-(e\cdot f)a
\]
与 $e$ 作点乘：
\[
(a\cdot f)(e\cdot e)+(a\cdot e)(f\cdot e)-(e\cdot f)(a\cdot e)
=(a\cdot e)(e\cdot f)-(a\cdot e)(e\cdot f)=0,
\]
与 $f$ 作点乘同样得 $0$，而与 $a$ 作点乘得 $(a\cdot e)(a\cdot f)$，均与 $p\times q$ 的相应点乘
成比例。三个线性无关的矢量上的取值都相同，故两个表达式在规范因子下相等。若把这一结果用 $k$
写出，则
\[
b=(k:g)\,a-2\,k\cdot g\cdot a ,
\]
其中 $k:g=k^{ij}g_{ij}=2(e\cdot f)$，
$(k\cdot g\cdot a)^i=k^{ij}g_{jk}a^k$；这一形式正是下一章一般公式在退化一次点上的特殊情形。
\end{proof}

\section{可逆二次点}

\begin{example}\label{ex:10-3}
已知圆锥曲线 $g$ 及其上的一点 $a$，并给定一个可逆的二次点 $h$。过 $a$ 作 $h$ 的两条切线，
分别交 $g$ 于 $p,q$ 两点（不同于 $a$）。作 $p,q$ 处的两条切线，交于点 $b$。求 $b$ 的表达式。
\end{example}

\begin{proof}[解]
把 $h$ 换成由它诱导的退化二次点即可。设 $E,F$ 是过 $a$ 的两条切线与 $h$ 的切点，则直线 $EF$
是 $a$ 关于 $h$ 的极线，其线矢与 $(h\cdot a)_i=h_{ij}a^j$ 成比例；而 $E,F$ 的对称积
$k=(ef)$ 与 $h$ 之间有关系
\[
k^{ij}=\left((h^{-1})_{kl}a^ka^l\right)h^{ij}-a^ia^j ,
\]
其中 $(h^{-1})_{kl}$ 是 $h$ 的逆（当 $h$ 退化时取相应极限，或用伴随矩阵代替）。由例 \ref{ex:10-2}，
$b=(k:g)\,a-2\,k\cdot g\cdot a$。把上式代入并利用 $a$ 在 $g$ 上，即 $a^2=g_{ij}a^ia^j=0$，
含 $a^ia^j$ 的项消去，第一、三项中的系数 $\left((h^{-1})_{kl}a^ka^l\right)$ 可以约去，于是
\[
b=(h:g)\,a-2\,h\cdot g\cdot a ,
\]
其中约定 $(h\cdot g\cdot a)^i=(g^{-1})^{ia}h_{ab}(g^{-1})^{bj}a_j$。
\end{proof}

\begin{proof}[退化情形与数值验证]
取 $g$ 为椭圆 $\dfrac{x^2}{4}+y^2=1$，$g^{ij}=\diag\{4,1,-1\}$，并把 $h$ 退化为两条直线
$m,n$ 的对称积，则 $E,F$ 取为这两条直线与 $g$ 的另一交点。取 $s=(0,1,1)$，
对合的两条不动直线 $m_1=(\tfrac12,1,-1)$、$m_2=(1,0,0)$，则
\[
h_{ij}=m_{1i}m_{2j}+m_{2i}m_{1j}
=\begin{pmatrix}1&1&-1\\1&0&0\\-1&0&0\end{pmatrix},
\qquad
h^{ij}=g^{ia}g^{jb}h_{ab}
=\begin{pmatrix}16&4&4\\4&0&0\\4&0&0\end{pmatrix},
\]
且 $g:h=h^{ij}g_{ij}=4$。于是 $f^{ij}=h^{ij}-4g^{ij}$。取辅助线矢 $m=(0,0,1)$，则
$m\cdot s=m_is^i=1$，$m\cdot f\cdot m=f^{ij}m_im_j=4$，故
\[
l=2(m\cdot s)\,f\cdot m-(m\cdot f\cdot m)\,s=(8,-4,4)\sim(2,-1,1).
\]
$l$ 的极点为
\[
b=g^{-1}\cdot l=(8,-1,-1)\sim(2,-1,1),
\]
在仿射坐标下即定点 $(2,-1)$，与例 \ref{ex:10-2} 的结果一致。
\end{proof}

\begin{remark}
上面用到的关系 $k=\left((h^{-1})_{kl}a^ka^l\right)h-aa$ 的含义是：过 $a$ 作 $h$ 的两条切线，
切点 $E,F$ 的对称积 $k=(ef)$ 与 $h$ 只相差一个由 $a$ 决定的规范因子。当 $h$ 退化为两点 $e,f$ 时，
$h$ 与 $k$ 描述的是同一个退化二次点，此时该关系中的系数正比于 $(e\cdot f)$，与例 \ref{ex:10-2}
的结果一致。这一关系在 $h$ 退化时给出的是同一个退化二次点，可以作为例 \ref{ex:10-2} 的特例来验证。
\end{remark}

\section{极线与共线的几何证明}

\begin{example}\label{ex:10-4}
在例 \ref{ex:10-3} 的背景下，求证：$A,B$ 与 $A$ 在 $g$ 处的切线关于 $h$ 的极点共线。
\end{example}

\begin{proof}
设 $EF$ 是 $A$ 关于 $h$ 的极线，其中 $E,F$ 是过 $A$ 的两条切线与 $h$ 的切点。设 $EF$ 与 $A$ 处的
切线交于 $G$，并设 $A$ 处的切线关于 $h$ 的极点为 $H$。因为 $H$ 的极线过 $A$，由极点极线的对称性
（La Hire 定理），$A$ 的极线 $EF$ 也过 $H$，即 $H$ 在 $EF$ 上。

在直线 $EF$ 上，$E,F$ 是 $h$ 上的两个切点，$G$ 是 $EF$ 与 $A$ 处切线的交点，$H$ 是 $A$ 处切线
关于 $h$ 的极点。由极点极线的调和性质，$E,F,G,H$ 成调和点列。于是以 $A$ 为中心的线束
$AE,AF,AG,AH$ 成调和线束，把它与 $g$ 相交，得 $P,Q,A,I$ 是 $g$ 上的二次调和点列，其中
$P,Q$ 是两条切线与 $g$ 的交点（即例 \ref{ex:10-3} 中的 $p,q$），$I$ 是 $AH$ 与 $g$ 的交点。

由二次调和点列的判别条件，$AI$ 与 $PQ$ 关于 $g$ 共轭。而 $PQ$ 的极点就是 $B$，故 $AI$ 过 $B$，
即 $A,H,B$ 三点共线。
\end{proof}

\begin{remark}
由此可知 $b$ 的表达式必具有形式
\[
b^i=\lambda a^i+\mu\,h^{ij}g_{jk}a^k ,
\]
其中 $\lambda,\mu$ 是待定实数：因为 $A,H,B$ 共线，而 $H$ 的极线过 $A$。例 \ref{ex:10-3} 中求得的
\[
b=(h:g)\,a-2\,h\cdot g\cdot a
\]
正符合这一形式。
\end{remark}

\section{共轭条件的代数应用}

\begin{example}\label{ex:10-5}
在例 \ref{ex:10-2} 的背景下，设 $FB$ 与 $AE$ 交于点 $K$，求证：$K$ 与 $F$ 关于 $g$ 共轭。
\end{example}

\begin{proof}
由例 \ref{ex:10-4}，$B$ 是 $PQ$ 关于 $g$ 的极点，故直线 $FB$ 过 $PQ$ 的极点，于是 $FB$ 与 $PQ$
关于 $g$ 共轭。由共轭的判别条件，$FB$ 与 $PQ$ 的交点所对应的四点成调和点列：设 $FB$ 与 $g$ 交于
$S,T$，则 $S,T,P,Q$ 是 $g$ 上的二次调和点列。以 $A$ 为中心射影，得 $AS,AT,AP,AQ$ 成调和线束；
再与直线 $FB$ 相交，得 $S,T,K,F$ 成调和点列。由一维对合的判别条件，$K,F$ 关于 $g$ 共轭。
\end{proof}

同样的结论可以用代数方法直接得到，而且不需要前文中"同时约掉 $(e\times f)\cdot a$"这一步
代数变形。设
\[
k=(b\times f)\times(a\times e),
\]
即 $K$ 的两个决定直线 $FB$ 与 $AE$ 的交点。$K$ 与 $F$ 关于 $g$ 共轭的条件是 $k\cdot f=0$。
利用混合积的恒等式
\[
\left((b\times f)\times(a\times e)\right)\cdot f
=\left((b\times f)\times f\right)\cdot(a\times e),
\]
并把 $b$ 按例 \ref{ex:10-4} 的结论写成
\[
b=\mu a+(a\cdot f)\,e+(a\cdot e)\,f
\]
（例 \ref{ex:10-4} 保证 $b$ 与 $a$ 以及 $k\cdot g\cdot a$ 共面，故可以这样设），代入得
\[
\left(\left(\mu\,a\times f+(a\cdot f)\,e\times f\right)\times f\right)\cdot(a\times e)=0 .
\]
用二重外积公式展开
\[
(a\times f)\times f=(a\cdot f)f-f^2a,\qquad
(e\times f)\times f=(e\cdot f)f-f^2e,
\]
得
\[
\mu\left((a\cdot f)f-f^2a\right)\cdot(a\times e)
+(a\cdot f)\left((e\cdot f)f-f^2e\right)\cdot(a\times e)=0 .
\]
由于 $a\cdot(a\times e)=0$，$e\cdot(a\times e)=0$，且 $a^2=e^2=0$，上式化简为
\[
\mu\,(a\cdot f)\,f\cdot(a\times e)+(a\cdot f)(e\cdot f)\,f\cdot(a\times e)=0 .
\]
当 $f\cdot(a\times e)\neq0$ 时解得
\[
\mu=-(e\cdot f),
\]
即
\[
b=(a\cdot f)\,e+(a\cdot e)\,f-(e\cdot f)\,a ,
\]
与例 \ref{ex:10-2} 的结果一致。

\begin{remark}
当 $f\cdot(a\times e)=0$ 时，直线 $AE$ 过 $F$，此时 $K=F$ 是显然的共轭点，结论仍然成立；
上面的推导只是把一般位置的情形写成了代数方程。
\end{remark}
  % 例题：二次点列的计算

\part{圆锥曲线的更多性质}
% 第十节：二次点列与二次对合
\chapter{二次点列与二次对合}

\section{二次点列}

本章取定一条非退化圆锥曲线 $g_{ij}$，并把它当作度规使用。本章只有这一条圆锥曲线被用作度规，
因此一律用 $g$ 升降指标，$g^{ij}$ 就是 $g$ 的逆 $(g^{-1})^{ij}$；退化二次曲线 $h$ 不是度规，
不用它升降指标。

\begin{definition}[二次点列]
设 $g_{ij}$ 是一条圆锥曲线。满足
\[
a^2=g_{ij}a^ia^j=0
\]
的点 $a$ 称为 $g$ 上的\emph{二次点}；$g$ 上全体二次点组成的集合称为 $g$ 上的
\emph{二次点列}。
\end{definition}

直线上的点列是一维对象，二次点列同样是一维对象，只是它没有事先指定的坐标。下面说明它也可以
定义交比，并且定义方式与直线上的点列完全平行。

\subsection{二次点列的交比}

\begin{definition}[二次点列的交比]\label{def:doublecross}
设 $a,b,c,d$ 是 $g$ 上四个互不相同的点，$P$ 是 $g$ 上异于这四点的点。过 $P$ 的四条直线
$P\times a,\ P\times b,\ P\times c,\ P\times d$ 构成一次线束，把它的交比记为
\[
J(P)=\frac{[Pac][Pbd]}{[Pad][Pbc]},
\]
其中 $[Pac]=(P\times a)\cdot c=\eps_{ijk}P^ia^jc^k$ 是三点的混合积。$J(P)$ 称为
$a,b,c,d$ 关于 $P$ 的交比；若它与 $P$ 无关，就把这个公共值称为 $a,b,c,d$ 的交比，记作
$J_{ab,cd}$。
\end{definition}

上述定义就是把二次点列上的点从 $P$ 射影为一次线束中的直线，再取这个线束的交比，与第九节关于
射影的讨论一致。关键的事实是这个交比与 $P$ 的选取无关。

\begin{theorem}[交比与辅助点无关]\label{thm:doublecross}
设 $a,b,c,d$ 是圆锥曲线 $g$ 上四个互不相同的点，则对 $g$ 上任意异于这四点的点 $P$，由
定义~\ref{def:doublecross} 给出的 $J(P)$ 都相等。
\end{theorem}

\begin{proof}
\textbf{第一步：计算 $J(P)^2$。}
取以 $P,a,c$ 为三列的矩阵 $M$。由 $\det(AB)=(\det A)(\det B)$，
\[
\det(M^{t}GM)=(\det M)^2\det g=[Pac]^2\det g,
\]
而 $M^{t}GM$ 的第 $m$ 行第 $n$ 列恰好是这两点的点乘（三点依次为 $P,a,c$）。因此 $[Pac]^2$
与 Gram 行列式
\[
\det\begin{pmatrix}
P\cdot P & P\cdot a & P\cdot c\\
a\cdot P & a\cdot a & a\cdot c\\
c\cdot P & c\cdot a & c\cdot c
\end{pmatrix}
\]
只相差一个与 $P,a,c$ 无关的非零常数因子。在
$J(P)^2=\dfrac{[Pac]^2[Pbd]^2}{[Pad]^2[Pbc]^2}$ 中这个因子在分子、分母里各出现两次，整体
相消，故不妨直接用 Gram 行列式代替 $[Pac]^2$ 计算。把 $P^2=a^2=c^2=0$ 代入，对角元全为零，
于是
\[
\det\begin{pmatrix}
0 & P\cdot a & P\cdot c\\
a\cdot P & 0 & a\cdot c\\
c\cdot P & c\cdot a & 0
\end{pmatrix}
=2(P\cdot a)(P\cdot c)(a\cdot c).
\]
同理
\[
[Pbd]^2\propto2(P\cdot b)(P\cdot d)(b\cdot d),\quad
[Pad]^2\propto2(P\cdot a)(P\cdot d)(a\cdot d),\quad
[Pbc]^2\propto2(P\cdot b)(P\cdot c)(b\cdot c),
\]
其中各式的比例因子相同。代入得
\[
J(P)^2=\frac{(a\cdot c)(b\cdot d)}{(a\cdot d)(b\cdot c)},
\]
右端与 $P$ 无关。

还要说明右端既不为零也不为无穷。若 $[Pad]=0$，则 $P,a,d$ 共线，而一条直线与非退化圆锥曲线
至多有两个交点，与 $P,a,d$ 互异矛盾，故 $[Pad]\neq0$；$[Pbc]\neq0$ 同理。又 $P\cdot a=0$
意味着 $a$ 在 $P$ 的极线（即 $P$ 处的切线）上，与 $a\neq P$ 矛盾，故 $P\cdot a\neq0$；同理
$P\cdot d\neq0$。于是由上面的比例关系可知 $(a\cdot d)\neq0$，$(b\cdot c)\neq0$，右端是一个
非零实数。

\textbf{第二步：由连续性确定 $J(P)$ 的符号。}
当 $P$ 在 $g$ 上变动时，$J(P)$ 是 $P$ 的连续函数：四个混合积都随 $P$ 连续变化，且分母处处非零。
在射影平面上，非退化实圆锥曲线同胚于圆周，特别地是连通的，因此 $J$ 的像连通。另一方面
$J(P)^2$ 是与 $P$ 无关的正常数 $c_0$，故 $J(P)$ 只能取 $\pm\sqrt{c_0}$ 两个值；连通的像不可能
分成两个点，所以 $J(P)$ 恒等于 $\sqrt{c_0}$ 或恒等于 $-\sqrt{c_0}$。总之 $J(P)$ 与 $P$ 无关。
\end{proof}

\begin{remark}
上面第二步（连续性论证）不可省略：仅由 $J(P)^2$ 与 $P$ 无关只能断定 $J(P)$ 的绝对值固定，
而 $J(P)$ 原则上可以在不同的连通分支上分别取 $\sqrt{c_0}$ 与 $-\sqrt{c_0}$。此外，Gram 行列式
与 $[Pac]^2$ 只在标准规范下相等；一般规范下二者相差常数因子 $\det g$，该因子在交比中相消，
不影响结论。
\end{remark}

\subsection{二次调和点列}

\begin{definition}[二次调和点列]
设 $a,b,c,d$ 是 $g$ 上四个互不相同的点。若 $J_{ab,cd}=-1$，则称 $a,b$ 与 $c,d$ 互相
\emph{调和分割}，并称 $a,b,c,d$ 为 $g$ 上的\emph{二次调和点列}。
\end{definition}

由第九节交比的置换性质，$J=-1$ 在交换 $a,b$（或交换 $c,d$）时不变（$1/J=-1$），在重新搭配
$a,c$ 与 $b,d$ 时得到 $1-J=2$，因此“哪两点与哪两点调和分割”的配对方式是确定的。

\begin{proposition}[二次调和点列的代数条件]\label{prop:doubleharmonic}
$g$ 上四个互不相同的点 $a,b,c,d$ 构成二次调和点列，当且仅当
\[
(a\times b)\cdot(c\times d)=0 .
\]
\end{proposition}

\begin{proof}
由定理~\ref{thm:doublecross}，$J^2=\dfrac{(a\cdot c)(b\cdot d)}{(a\cdot d)(b\cdot c)}$。
$J=-1$ 给出 $J^2=1$；反之若 $J^2=1$，则 $J=\pm1$，而 $J=1$ 对四个互不相同的点不可能
（第九节关于交比取值的注），故 $J^2=1$ 与 $J=-1$ 等价。于是调和条件等价于
\[
(a\cdot c)(b\cdot d)-(a\cdot d)(b\cdot c)=0 .
\]
由第三章的二重叉乘公式
\[
(a\times b)\cdot(c\times d)=(a\cdot c)(b\cdot d)-(a\cdot d)(b\cdot c)
\]
（一般规范下右端可相差一个非零常数因子，但零点位置不变），上式等价于
$(a\times b)\cdot(c\times d)=0$。
\end{proof}

这一条件有明显的几何意义：直线 $a\times b$ 与直线 $c\times d$ 关于 $g$ 共轭。直线的共轭可以
从两个方向来理解：一个方向是极点与极线，另一个方向是一维的对合变换。

\section{二次对合}

回忆第九节的一维对合：给定两个不动点 $p,q$，与 $p,q$ 构成调和点列的两点 $a,b$ 互相对合，
这一配对由 $g_{ij}a^ib^j=0$ 给出，其中 $g=\sym{p}{q}$ 是两个不动点的对称积。把 $p,q$ 换成过
它们的直线 $\ell=p\times q$，二次调和点列的条件 $(a\times b)\cdot(c\times d)=0$ 就化为关于
$\ell$ 的条件。

\begin{definition}[二次对合]\label{def:doubleinv}
设 $\ell$ 是一条不与 $g$ 相切的直线。对 $g$ 上的点 $a$，满足
\[
(a\times b)\cdot\ell=0
\]
的点 $b\in g$ 称为 $a$ 关于 $\ell$ 的\emph{对合点}；由此给出的配对关系称为由 $\ell$ 确定的
\emph{二次对合}。
\end{definition}

\begin{proposition}[二次对合的结构]\label{prop:doubleinv-struct}
设 $\ell$ 不与 $g$ 相切，$D=g^{-1}\ell$ 是 $\ell$ 的极点。则定义~\ref{def:doubleinv} 给出的
配对就是：过 $D$ 的直线与 $g$ 的两个交点互相对合。特别地，
\begin{enumerate}
\item 每个 $a\in g$ 都有唯一的对合点 $b$，并且 $a,b,D$ 共线；
\item 配对是对合：$b$ 是 $a$ 的对合点当且仅当 $a$ 是 $b$ 的对合点；
\item 若 $\ell$ 与 $g$ 交于 $p,q$，则 $p,q$ 是对合的不动点；若 $\ell$ 与 $g$ 没有交点，则对合
没有不动点。
\end{enumerate}
\end{proposition}

\begin{proof}
固定 $a$。条件 $(a\times b)\cdot\ell=0$ 关于 $b$ 是线性的，故解构成一条直线 $L_a$。由于
$a\times a=0$，$b=a$ 总满足条件，故 $a\in L_a$；又把 $\ell$ 的极点 $D=g^{-1}\ell$ 代入，
由叉乘与配极的定义可知 $D\in L_a$。于是 $L_a$ 就是直线 $a\times D$，它与 $g$ 的两个交点中
有一个是 $a$，另一个就是 $b$，且 $a,b,D$ 共线。这同时说明了 $b$ 的存在唯一性。条件关于
$a,b$ 对称（因为 $(a\times b)\cdot\ell=-(b\times a)\cdot\ell$），故 $a$ 也是 $b$ 的对合点。

不动点满足 $a=b$，此时 $L_a=a\times D$ 是 $g$ 在 $a$ 处的切线，即 $a$ 是从 $D$ 可以作切线的
切点，也就是 $\ell$ 与 $g$ 的交点：由配极关系，$a$ 处的切线过 $D$ 等价于 $D$ 在 $a$ 的极线
上，也等价于 $a$ 在 $D$ 的极线 $\ell$ 上。反之，若 $a\in\ell\cap g$，则 $a$ 处的切线过 $D$，
故 $L_a$ 与 $g$ 在 $a$ 处相切，$a$ 的对合点仍是 $a$。若 $\ell$ 与 $g$ 没有交点，这样的 $a$
不存在。
\end{proof}

\begin{remark}
当 $a,b$ 重合时，应把它们想象成错开了无穷小距离，此时 $a\times a$ 理解为 $a$ 处的切线。线矢
与点矢的缩并写作 $\ell\cdot a=\ell_ia^i$，不要与两个点矢之间的 $a\cdot b=g_{ij}a^ib^j$ 混淆。
\end{remark}

\begin{theorem}[过定点的线束给出的对合]\label{thm:pencil-inv}
设 $s$ 不在 $g$ 上。过 $s$ 的每条直线与 $g$ 的两个交点（可以重合）互相对合。换言之，过 $s$
的线束在 $g$ 上诱导一个二次对合，它的不动点是从 $s$ 向 $g$ 所引两条切线的切点。
\end{theorem}

\begin{proof}
取 $\ell$ 为 $s$ 的极线，则 $s$ 是 $\ell$ 的极点。过 $s$ 的直线 $u$ 与 $g$ 交于 $a,b$，此时
$a\times b=u$ 过 $s$，即 $(a\times b)\cdot\ell=0$。由命题~\ref{prop:doubleinv-struct}，
$a,b$ 互相对合。不动点是 $\ell\cap g$，也就是从 $s$ 所引两条切线的切点。
\end{proof}

对偶地，把点与直线、点矢与线矢的地位互换（对偶圆锥曲线由 $g$ 的逆给出），就得到关于二次线束
的结论。所谓\emph{二次线束}，就是 $g$ 的全部切线组成的集合。

\begin{theorem}[定直线上的点列]\label{thm:linerange-inv}
设 $L$ 是不与 $g$ 相切的直线。对 $L$ 上的点 $A$，设从 $A$ 向 $g$ 所引的两条切线切 $g$ 于
$E,F$，直线 $EF$ 交 $L$ 于 $B$。则
\begin{enumerate}
\item $A\mapsto B$ 是 $L$ 上以 $L\cap g$ 为不动点的对合；
\item $E,F$ 是 $g$ 上的一对对合点，不动点也是 $L\cap g$；
\item 二次线束中的两条切线 $AE,AF$ 互相对合，不动元素是 $L\cap g$ 处的切线。
\end{enumerate}
\end{theorem}

\begin{proof}
设 $S$ 是 $L$ 的极点。$A\in L$ 的极线是 $EF$，由配极关系它过 $S$，故 $E,F$ 是过 $S$ 的线束
与 $g$ 的两个交点，由定理~\ref{thm:pencil-inv} 互相对合，不动点是从 $S$ 所引切线的切点，即
$S$ 的极线 $L$ 与 $g$ 的交点 $L\cap g$。这给出第二条。第一条中的配对正是 $g$ 在直线 $L$ 上的
限制给出的对合（第九节关于圆锥曲线在直线上限制的定理），其不动点同样是 $L\cap g$。第三条是
第二条的换位表述。
\end{proof}

\subsection{几个同时出现的对合}

把上面两个定理放在一起，可以看到在同一构图中往往同时出现几个对合。

设 $s$ 不在 $g$ 上。过 $s$ 的动直线 $u$ 交 $g$ 于 $P,Q$，$P,Q$ 处的切线交于 $D$（$D$ 是 $u$
的极点）。则
\begin{enumerate}
\item 在二次点列上：$P,Q$ 是由 $s$ 诱导的对合中的一对对合点；
\item 在二次线束上：$P,Q$ 处的两条切线是由同一定点诱导的对合中的一对对合直线；
\item 在过 $s$ 的一次线束上：直线 $sPQ=u$ 与直线 $sD$ 互相对合，这是第九节关于圆锥曲线在
点上限制（对偶形式）的定理给出的线束对合。
\end{enumerate}
三个对合的不动元素都是 $s$ 到 $g$ 的两条切线（或它们的切点）。

对偶地，设 $L$ 是一条定直线，交 $g$ 于 $P,Q$。$A$ 是 $L$ 上的动点，过 $A$ 的两条切线切 $g$
于 $E,F$，直线 $EF$ 交 $L$ 于 $B$。则
\begin{enumerate}
\item 在二次点列上：$E,F$ 是一对对合点，不动点是 $P,Q$；
\item 在二次线束上：$AE,AF$ 是一对对合切线，不动元素是 $P,Q$ 处的切线；
\item 在直线 $L$ 的点列上：$A,B$ 是一对对合点，不动点是 $P,Q$。
\end{enumerate}

二次曲线可以产生直线上的对合，直线也可以产生二次曲线上的对合，两者互相配合，构成了本节讨论
的主要内容。

\begin{example}[共轭直线给出的对合]\label{ex:conj}
设直线 $PQ$ 与直线 $CD$ 关于 $g$ 共轭，即 $(p\times q)\cdot(c\times d)=0$，其中 $p,q$ 是
$PQ$ 与 $g$ 的两个交点，$c,d$ 是 $CD$ 与 $g$ 的两个交点。$X$ 是 $g$ 上的动点，直线 $XC$、
$XD$ 分别交直线 $PQ$ 于 $A,B$。则 $A,B$ 是直线 $PQ$ 上以 $p,q$ 为不动点的对合中的一对对合点。
\end{example}

\begin{proof}[解]
由 $(p\times q)\cdot(c\times d)=0$ 知 $p,q,c,d$ 是 $g$ 上的二次调和点列。把这一调和点列从
$X$（也在 $g$ 上）射影为线束，由定理~\ref{thm:doublecross}，四条直线 $XP,XQ,XC,XD$ 构成调和
线束。再与直线 $PQ$ 相交：$XP,XQ$ 分别交 $PQ$ 于 $p,q$（因为 $p,q\in PQ$），$XC,XD$ 分别交
$PQ$ 于 $A,B$。射影保持交比，故 $A,B,p,q$ 成调和点列，即 $A,B$ 是以 $p,q$ 为不动点的对合中的
一对对合点。
\end{proof}

\section{对合张量与三步算法}

\subsection{问题与记号}

设 $s$ 是 $g$ 上一点，$h$ 是两条过 $s$ 的直线 $m,n$ 构成的退化二次曲线
\[
h_{ij}=m_in_j+m_jn_i=\sym{m}{n}_{ij}.
\]
由 $h_{ij}s^j=0$ 知 $s$ 是 $h$ 的奇异点，$h$ 的两条直线都过 $s$。在过 $s$ 的线束上，以
$m,n$ 为不动直线的调和对合由 $h$ 给出；它又通过 $s$ 处的中心射影诱导出 $g$ 上的一个二次对合。
具体地说，对 $a\in g$，直线 $L_a$ 由 $h_{ij}a^ix^j=0$ 定义，它恒过 $s$；把 $L_a$ 与 $g$ 的
另一个交点记为 $b$，则 $a,b$ 对合（由 $h_{ij}a^ib^j=0$ 的对称性，这一配对是对合）。这一对合的
不动点是 $m,n$ 与 $g$ 的另一个交点，记作 $p,q$；而 $g$ 上互相对合的两点连线必过一定点 $D$，
$D$ 就是直线 $\ell=p\times q$ 的极点。

于是问题可以表述为：已知 $s$ 与对合张量 $h$（它给出过 $s$ 的线束上的对合），求直线 $\ell$
的极点 $D$。下面先解决相反的问题，即由 $\ell$ 与 $s$ 写出 $h$，再把它看成关于 $\ell$ 的方程
反解。

\subsection{从 $\ell$ 到 $h$}

设 $\ell$ 过对合的两个不动点 $p,q$。由第八节关于联立二次点的例题，$p,q$ 的对称积（联立二次点）
与
\[
(pq)^{ij}\propto\ell^2g^{ij}-\ell^i\ell^j,
\qquad \ell^2=\ell_i\ell^i,\quad \ell^i=g^{ij}\ell_j
\]
成正比（在适当的规范化下相等；相差的常数不影响以下结论）。两条固定直线是 $m=s\times p$、
$n=s\times q$，其对称积为
\[
h_{ij}=m_in_j+m_jn_i=\eps_{ikm}s^k\,\eps_{jln}s^l\,(pq)^{mn}.
\]
把 $(pq)$ 的表达式代入并化简，可得
\[
h_{ij}\propto\ell_is_j+s_i\ell_j-(\ell\cdot s)g_{ij},
\qquad\text{即}\qquad
h\propto\ell s+s\ell-(\ell\cdot s)g,
\]
其中 $s_i=g_{ij}s^j$，$\ell\cdot s=\ell_is^i$。

\begin{proof}[化简与验证]
把 $(pq)^{mn}\propto\ell^2g^{mn}-\ell^m\ell^n$ 代入后，利用 $\eps$ 的两个恒等式
\[
\eps_{ikm}\eps_{jln}g^{mn}s^ks^l=\frac{1}{\det g}\bigl[(s\cdot s)g_{ij}-s_is_j\bigr],
\qquad
\eps_{ikm}s^k\ell^m=(\ell\times s)_i,
\]
再代入 $s^2=s\cdot s=0$，即可得到上面的结果（相差一个非零常数因子）。

也可以直接验证右端确实给出所要的对合张量。令
\[
h'_{ij}=\ell_is_j+s_i\ell_j-(\ell\cdot s)g_{ij}.
\]
则
\[
h'_{ij}s^j=(\ell\cdot s)s_i+(s\cdot s)\ell_i-(\ell\cdot s)s_i=s^2\,\ell_i=0,
\]
故 $s$ 在 $h'$ 的核中，$h'$ 退化，且是两条过 $s$ 的直线之积。又对 $\ell$ 与 $g$ 的交点 $p$
（$p^2=0$，$\ell\cdot p=0$）有
\[
h'_{ij}p^ip^j=2(\ell\cdot p)(s\cdot p)-(\ell\cdot s)(p\cdot p)=0,
\]
同理 $q$ 也使 $h'$ 为零。因此在一般位置下，$h'$ 的两条直线分别过 $p$ 与 $q$，即正是
$m=s\times p$ 与 $n=s\times q$。于是 $h'\propto h$。
\end{proof}

\subsection{从 $h$ 到 $\ell$}

把 $h\propto\ell s+s\ell-(\ell\cdot s)g$ 看作关于未知线矢 $\ell$ 的方程。用 $g^{ij}$ 缩并：
\[
g:h:=g^{ij}h_{ij}
=g^{ij}\ell_is_j+g^{ij}s_i\ell_j-(\ell\cdot s)g^{ij}g_{ij}
=2(\ell\cdot s)-(\ell\cdot s)\delta^i_i,
\]
其中用到 $g^{ij}\ell_is_j=g^{ij}\ell_ig_{jk}s^k=\ell^is_i=\ell\cdot s$，$g^{ij}s_i\ell_j$ 同理。
又 $\delta^i_i=3$（三维空间中每个维度贡献一个 $1$），于是
\[
g:h=2(\ell\cdot s)-3(\ell\cdot s)=-(\ell\cdot s).
\]
因此方程化简为
\[
h-(g:h)g=\ell s+s\ell .
\]
这说明：已知对称积 $(\ell s)$ 以及其中一个因子 $s$，求另一个因子 $\ell$；这相当于已知二次方程
的一个根求另一个根。记
\[
f^{ij}=h^{ij}-(g:h)g^{ij},
\]
则 $f^{ij}=\ell^is^j+s^i\ell^j$，而 $f$ 完全由已知的 $h$ 与 $g$ 算出，不需要解二次方程。

为从 $f$ 与 $s$ 解出 $\ell$，任取一条线矢 $m$（要求 $m\cdot s\neq0$），令
\[
D^i=2(m\cdot s)\,(f^{ij}m_j)-(m_if^{ij}m_j)\,s^i .
\]
代入 $f^{ij}=\ell^is^j+s^i\ell^j$，得
\[
f^{ij}m_j=\ell^i(s\cdot m)+s^i(\ell\cdot m),
\qquad
m_if^{ij}m_j=2(m\cdot s)(m\cdot\ell),
\]
于是
\[
D^i=2(m\cdot s)\bigl[\ell^i(s\cdot m)+s^i(\ell\cdot m)\bigr]-2(m\cdot\ell)(m\cdot s)s^i
=2(m\cdot s)^2\ell^i .
\]
由于 $\ell^i=g^{ij}\ell_j$ 正是 $\ell$ 的极点，上式给出的就是所求定点（相差非零常数
$2(m\cdot s)^2$）。特别地，结果与辅助线矢 $m$ 的取法无关，只要 $m\cdot s\neq0$。

\subsection{三步算法}

综合以上讨论，从对合张量 $h$ 求定点 $D$ 可以归纳为三步。记号约定：线矢与点矢的缩并
$m\cdot s=m_is^i$；$(f\cdot m)^i=f^{ij}m_j$；$m\cdot f\cdot m=f^{ij}m_im_j$。
\begin{enumerate}
\item 写出对合张量 $h$（退化二次曲线，即两条固定直线的对称积）；
\item 先求 $g:h=g^{ij}h_{ij}$，再令 $f=h-(g:h)g$；
\item 任取线矢 $m$（要求 $m\cdot s\neq0$），由
\[
D=2(m\cdot s)\,f\cdot m-(m\cdot f\cdot m)\,s
\]
求出定点 $D$。
\end{enumerate}

\begin{remark}
第 3 步给出的矢量 $D^i$ 已经是 $\ell$ 的极点，即所求定点本身，不需要再取一次极点。若 $h$ 以
下降指标的形式 $h_{ij}$ 给出，必须先写成 $h^{ij}=g^{ia}g^{jb}h_{ab}$ 再代入，否则 $g:h$ 与
$f$ 都会算错。
\end{remark}

\begin{remark}
算法的输入是 $h$ 与 $s$，由它们只能确定 $\ell$（从而确定 $D$）到相差一个非零常数。若题目只
给出对合的部分信息（例如只知道斜率之和为常数 $c$，而不知道两条不动直线），应先把 $h$ 写成含
$c$ 的表达式，代入算法得到含 $c$ 的定点，再结合题目的其他条件定出 $c$。定点落在某条已知直线
上、或与某个已知点重合，都是可以定出 $c$ 的常见条件。
\end{remark}

\section{例：2017 年全国卷 I 第 20 题}

\begin{example}\label{ex:2017}
已知椭圆 $C:\dfrac{x^2}{a^2}+\dfrac{y^2}{b^2}=1\ (a>b>0)$，四点
\[
P_1(1,1),\quad P_2(0,1),\quad
P_3\!\left(-1,\dfrac{\sqrt3}{2}\right),\quad P_4\!\left(1,\dfrac{\sqrt3}{2}\right)
\]
中恰有三点在 $C$ 上。

(1) 求 $C$ 的方程；

(2) 设直线 $l$ 不经过 $P_2$，且与 $C$ 交于 $A,B$ 两点。若直线 $P_2A$ 与直线 $P_2B$ 的斜率
之和为 $-1$，证明：$l$ 过定点。
\end{example}

\subsection{第一问}

$P_3=\left(-1,\dfrac{\sqrt3}{2}\right)$ 与 $P_4=\left(1,\dfrac{\sqrt3}{2}\right)$ 关于 $y$ 轴
对称，椭圆 $C$ 也关于 $y$ 轴对称，故它们要么都在 $C$ 上，要么都不在 $C$ 上。由于恰有三点在
$C$ 上，$P_3,P_4$ 必都在 $C$ 上（否则 $C$ 上至多还有 $P_1,P_2$ 两点，与“恰有四点中的三点”
矛盾）。于是 $P_1,P_2$ 中恰有一点在 $C$ 上。

把 $P_4$ 代入得
\[
\frac{1}{a^2}+\frac{3}{4b^2}=1 .
\]
若 $P_1(1,1)$ 也在 $C$ 上，则 $\dfrac{1}{a^2}+\dfrac{1}{b^2}=1$，与上式相减得
$\dfrac{1}{4b^2}=0$，不可能。所以 $P_1\notin C$，从而 $P_2(0,1)\in C$，即 $b^2=1$。代回上式
得 $\dfrac{1}{a^2}=\dfrac14$，即 $a^2=4$。故
\[
C:\ \frac{x^2}{4}+y^2=1 .
\]

\subsection{第二问：几何论证}

为免与线矢记号 $\ell$ 混淆，下面把第二问中的动直线 $l$ 记作 $L$。取齐次坐标 $(x,y,z)$，
$z=1$ 表示通常的平面。椭圆 $C$ 的张量为
\[
g_{ij}=\diag\left\{\tfrac14,1,-1\right\},\qquad
g^{ij}=\diag\{4,1,-1\},
\]
点 $s=P_2=(0,1,1)$ 在 $C$ 上。

过 $s$ 的直线（竖直线除外）可写成 $y=kx+1$，对应线矢 $(k,-1,1)$；竖直线 $x=0$ 对应线矢
$(1,0,0)$，可看作斜率 $k=\infty$。条件“两直线斜率之和为 $-1$”是对过 $s$ 的一次线束的一个
对合：它把斜率为 $k$ 的直线与斜率为 $-1-k$ 的直线配对。不动直线由 $k=-1-k$ 得 $k=-\tfrac12$，
再加上竖直线（$k=\infty$）。

由于 $s\in C$，从 $s$ 把 $C$ 上的点射影为过 $s$ 的直线是一个射影对应，因此线束上的这个对合
诱导出 $C$ 上的一个对合，而 $A,B$ 正是这一对合中的一对对合点（它们分别对应直线 $sA,sB$）。
对 $C$ 上的对合，互相对合的两点的连线都经过一个定点：这个定点是过两个不动点 $p,q$ 的直线
$\ell$ 的极点（命题~\ref{prop:doubleinv-struct}）。所以 $L=AB$ 过定点 $D$，问题化为求 $D$。

\subsection{第二问：定点坐标的计算}

先写出对合张量 $h$。设 $C$ 上的点 $a$ 对应齐次点 $a=(x_a,y_a,1)$，则过 $s=(0,1,1)$ 与 $a$
的直线斜率为
\[
k_a=\frac{y_a-1}{x_a}=\frac{a^2-a^3}{a^1}
\]
（$x_a=0$ 时 $k_a=\infty$；下面的方程是多项式方程，对 $x_a=0$ 也按极限理解）。条件
$k_a+k_b=-1$ 等价于
\[
(a^2-a^3)\,b^1+(b^2-b^3)\,a^1+a^1b^1=0,
\]
这正是 $h_{ij}a^ib^j=0$，其中
\[
h_{ij}=
\begin{pmatrix}
1 & 1 & -1\\
1 & 0 & 0\\
-1 & 0 & 0
\end{pmatrix}.
\]
（可以直接验证 $h_{ij}s^j=0$，即 $s$ 是奇异点，符合前一节的要求。）

\textbf{第一步}：对合张量 $h$ 如上。

\textbf{第二步}：先算 $h^{ij}=g^{ia}g^{jb}h_{ab}$，再算 $g:h=g^{ij}h_{ij}$：
\[
h^{ij}=
\begin{pmatrix}
16 & 4 & 4\\
4 & 0 & 0\\
4 & 0 & 0
\end{pmatrix},
\qquad
g:h=g^{ij}h_{ij}=4\cdot1+1\cdot0+(-1)\cdot0=4 .
\]
于是
\[
f^{ij}=h^{ij}-(g:h)\,g^{ij}
=
\begin{pmatrix}
16 & 4 & 4\\
4 & 0 & 0\\
4 & 0 & 0
\end{pmatrix}
-4
\begin{pmatrix}
4 & 0 & 0\\
0 & 1 & 0\\
0 & 0 & -1
\end{pmatrix}
=
\begin{pmatrix}
0 & 4 & 4\\
4 & -4 & 0\\
4 & 0 & 4
\end{pmatrix}.
\]

\textbf{第三步}：取辅助线矢 $m=(0,0,1)$（即无穷远直线 $z=0$）。此时
\[
m\cdot s=m_is^i=1,\qquad
f\cdot m=(f^{ij}m_j)=(4,0,4),\qquad
m\cdot f\cdot m=f^{ij}m_im_j=f^{33}=4 .
\]
由三步公式
\[
D^i=2(m\cdot s)(f^{ij}m_j)-(m\cdot f\cdot m)s^i
=2(4,0,4)-4(0,1,1)=(8,-4,4),
\]
即 $D=(8,-4,4)$，只取到规范因子时 $D\sim(2,-1,1)$。在仿射平面 $z=1$ 上，定点为
\[
\boxed{(2,-1)} .
\]

\begin{remark}
这里第 3 步算出的 $D=(8,-4,4)$ 已经是 $\ell$ 的极点，也就是定点本身，它对应仿射点
$(2,-1)$；不要再对它取极点。为作对照，可以直接求出两个不动点：不动直线为
$k=-\tfrac12$（即 $y=1-\tfrac x2$）与 $x=0$，它们与 $C$ 的另一个交点分别为 $q=(2,0,1)$
与 $p=(0,-1,1)$，故
\[
\ell=p\times q=(-1,2,2),
\]
其极点为 $g^{-1}\ell=(-4,2,-2)\sim(2,-1,1)$，与三步算法的结果一致。
\end{remark}

\begin{remark}
对合条件须取题设给出的 $k_1+k_2=-1$。作为旁证，若改用 $k_1+k_2=1$ 重复上面的计算（此时
$h_{ij}=\begin{pmatrix}-1&1&-1\\1&0&0\\-1&0&0\end{pmatrix}$），三步算法给出的定点是 $(2,1)$，
与题设要求的答案不符。
\end{remark}

\begin{remark}
第 3 步中的辅助线矢 $m$ 可以任取（只要 $m\cdot s\neq0$）。例如取 $m=(1,1,1)$，则
$m\cdot s=2$，$f\cdot m=(8,0,8)$，$m\cdot f\cdot m=16$，于是
$D=2\cdot2\cdot(8,0,8)-16\cdot(0,1,1)=(32,-16,16)\sim(2,-1,1)$，结果相同。若取
$m=(1,0,0)$，则 $m\cdot s=0$，公式给出零矢，说明这一类 $m$ 与 $s$ 的配对退化，应避开。
\end{remark}
   % 二次点列与二次对合
% 第十二节：平行与正交
\chapter{平行与正交}

\section{几何的分层}

平面几何可以按结构从少到多分成若干层次：

\begin{itemize}
\item \textbf{射影几何}：只有三维矢量空间本身。有意义的概念是"点在直线上"、共点、共线、
圆锥曲线、极点与极线、对合与调和点列等。
\item \textbf{仿射几何}：再指定一条无穷远直线 $\linf$。由此产生平行、中点、圆锥曲线的中心、
双曲线的渐近线，以及椭圆、双曲线与抛物线的分类。
\item \textbf{正交几何}：再指定一个正交结构（见第三节）。由此产生垂直与角度。
\item \textbf{解析几何}：再固定坐标系的原点、方向与单位长度，从而可以写出坐标并进行具体的
数值计算。
\end{itemize}

层次之间的表示关系是单向的。

\begin{proposition}[层次的表示]
高层次的运算与关系可以表示低层次的运算与关系，反之一般不能。
\end{proposition}

\begin{example}
用点坐标所满足的方程（解析几何）可以给出三点共线的条件（射影几何）；反过来，射影几何的
关系不能表示"两点的横坐标相等"（解析几何），因为横坐标的数值依赖于坐标系的选取。
\end{example}

\begin{example}
写出过两点的直线方程（解析几何），再与第三条直线求交点，可以判断三线共点，这实质上就是
点矢的叉乘（射影几何）；但射影几何中没有平行概念，不能表示"过一点作已知直线的平行线"
（仿射几何）。同样地，叉乘不能表示长度，长度需要正交结构或解析几何才会出现。
\end{example}

既然高层次的运算能表示低层次的一切，原则上可以把所有推理都放到最高层次进行；但用高层次的
语言表示低层次的关系往往要引入多余的坐标与数值，计算量反而增大。因此有下面的原则。

\begin{proposition}[解题层次的选取]
若命题中出现的所有关系都属于射影几何，就用射影几何的运算；若还用到平行、中点、中心一类
概念，就用仿射几何的运算；只有当命题中确实出现垂直、长度或角度时，才用正交几何或解析几何
的运算。
\end{proposition}

命题中常含有超出需要的前提。例如命题对任意位置的 $A$ 都成立，题设却要求 $A$ 落在 $x$ 轴上；
或者命题对任意双曲线都成立，题设却指定它是等轴双曲线。多余的条件会把推理引向不必要的层次。
在选取方法之前，应先判断哪些条件是本质的、哪些是可以去掉的。

\section{平行结构与单比}

\begin{definition}[交比]
同一直线 $l$ 上的四点 $a,b,c,d$ 的交比定义为
\[
(ab;cd)=\frac{[ac]\,[bd]}{[ad]\,[bc]},
\]
其中 $[ac]$ 表示在 $l$ 上任取一点 $r$ 后由 $a,c,r$ 三个点矢所作的行列式 $\det(a,c,r)$。
上式在 $r$ 的替换与各点矢的规范化下都不变，因此是射影几何的量。
\end{definition}

取一个坐标系，使 $l$ 上所有点的第二个分量为 $0$，并适当规范化使第三个分量为 $1$，则四点
成为 $(x_A,0,1)$ 等，此时
\[
(ab;cd)=\frac{(x_C-x_A)(x_D-x_B)}{(x_D-x_A)(x_C-x_B)} .
\]
在纯粹的射影几何中，交比是几何量，而 $x_C-x_A$ 这样的差并不是几何量：换一个坐标系它的值就变了。
只有在赋予 $x_C-x_A$ 以"两点间长度"的意义（即取定一个特定的坐标系）之后，才能说交比是长度
之比。类似地，射影几何中的叉乘（作过两点的直线）在解析几何中表现为写出直线方程。用高层次的
量表示低层次的量，表达式通常比较繁。

引入 $\linf$ 之后，直线 $l$ 上多出一个特殊点 $\infty=l\times\linf$，即 $l$ 与无穷远直线的交点。

\begin{definition}[单比]
设 $a,b,c$ 是直线 $l$ 上的三点，$\infty=l\times\linf$。称交比
\[
(ab;c\infty)=\frac{[ac]\,[b\infty]}{[a\infty]\,[bc]}
\]
为 $a,b,c$ 的单比。
\end{definition}

单比是特殊的交比，因而也是射影量。若所用坐标系已使 $x_C-x_A$ 具有长度差的意义，则无穷远点的
坐标必为无穷大，把上式取极限即得
\begin{equation}
(ab;c\infty)=\frac{x_C-x_A}{x_C-x_B}.
\label{eq:simple-ratio}
\end{equation}
也就是说，\textbf{单比就是同一直线上两条线段的有向长度之比}。单比是仿射几何的概念，却可以用
长度（解析几何的量）简单地表示；反过来，长度是不能用仿射几何的运算表示的。

\begin{proposition}[中点]
$a,b,c,\infty$ 成调和点列，当且仅当 $C$ 是线段 $AB$ 的中点。
\end{proposition}

\begin{proof}
调和的条件是 $(ab;c\infty)=-1$。代入 \eqref{eq:simple-ratio} 得
\[
\frac{x_C-x_A}{x_C-x_B}=-1
\quad\Longleftrightarrow\quad
x_C-x_A=x_B-x_C
\quad\Longleftrightarrow\quad
x_C=\frac{x_A+x_B}{2}.
\]
\end{proof}

因此中点是仿射几何的概念。射影几何中只有交比与调和关系，不能定义中点，必须借助 $\linf$ 才行。

\section{正交结构}

蒙日圆问题用到了"垂直"，而垂直既不属于射影几何，也不属于仿射几何。平行结构来自引入特殊的
线矢 $\linf$，类似地，垂直结构应当来自引入一个特殊的张量。

\begin{definition}[垂直]
设 $\theta^{ij}$ 是一个对称的 $(2,0)$ 型张量。规定两条直线 $m,n$ 垂直，当且仅当
\begin{equation}
\theta^{ij}m_in_j=0 .
\label{eq:perp}
\end{equation}
\end{definition}

\eqref{eq:perp} 是规范不变的：把 $m$ 换成 $\lambda m$（$\lambda\ne0$）时左端只差一个常数因子，
条件本身不受影响。

\subsection{平行的直线也垂直}

若 $m,n$ 垂直，则与 $m$ 平行的任一直线与与 $n$ 平行的任一直线也垂直。与 $m$ 平行的直线可以写成
$m+\lambda\linf$，与 $n$ 平行的直线可以写成 $n+\mu\linf$。于是对一切 $\lambda,\mu$ 都有
\[
\theta^{ij}(m+\lambda\linf)_i(n+\mu\linf)_j=0 .
\]
把左端展开，比较 $\lambda\mu$ 的系数得
\begin{equation}
\theta^{ij}\linfsub{i}=0 .
\label{eq:theta-degenerate}
\end{equation}
即 $\theta$ 退化，且它的核至少包含 $\linf$。

\subsection{直线不与自己垂直}

一条直线不应与自身垂直，因此对任何 $m\ne\linf$ 都有 $\theta^{ij}m_im_j\ne0$。这说明 $\theta$
不能是不定的：若 $\theta$ 的号差中含有两个相反的符号，就存在非零的实线矢使
$\theta^{ij}m_im_j=0$。把 $\theta$ 对角化，它的号差只能是 $(1,1,0)$。

\begin{definition}[正交结构]
设三维矢量空间上给定一个对称的 $(2,0)$ 型张量 $\theta^{ij}$，其号差为 $(1,1,0)$，则称相应的
射影几何上有一个\textbf{正交结构}。
\end{definition}

由 \eqref{eq:theta-degenerate}，$\linf$ 属于 $\theta$ 的核；而号差为 $(1,1,0)$ 意味着核是一维的，
因此核恰好由 $\linf$ 张成。这说明正交结构唯一确定平行结构。

\begin{proposition}[正交结构决定平行结构]
正交结构 $\theta$ 的核是一维的。从代数看，满足 $\theta^{ij}\linfsub{i}=0$ 的线矢构成无穷远
直线；从几何看，若两条直线垂直于同一条直线，则这两条直线平行。因此正交结构包含了平行结构。
\end{proposition}

\begin{remark}
反过来不成立。在 $\linf$ 的任意一条补线上取正定二次型，再把它与核上的零二次型一起拼成
$\theta$，就得到号差为 $(1,1,0)$ 的张量，可见平行结构不能决定正交结构。上面把 $\theta$ 的号差
限定为 $(1,1,0)$，是由"平行的直线也垂直"与"直线不与自己垂直"这两条要求逼出来的，二者缺一不可。
\end{remark}

\subsection{标准坐标系下的形式}

在与解析几何联系紧密的坐标系下，正交张量取标准形式
\begin{equation}
\theta^{ij}=\diag\{1,1,0\},\qquad \linfsub{i}=(0,0,1).
\label{eq:standard-theta}
\end{equation}

\section{角度的度量}

正交结构不仅判断两条直线是否垂直，还能给出它们的角度。

\begin{definition}[角度]
设 $m,n$ 是两个线矢，$\theta^{ij}m_in_j\ne0$。定义
\begin{equation}
\langle m,n\rangle=\frac{m\theta n}{\sqrt{(m\theta m)(n\theta n)}},
\label{eq:angle}
\end{equation}
其中 $m\theta n$ 是 $\theta^{ij}m_in_j$ 的简写。
\end{definition}

\eqref{eq:angle} 是规范不变的：把 $m$ 换成 $\lambda m$、$n$ 换成 $\mu n$ 时，分子与分母同时乘上
$|\lambda\mu|$，比值不变。又 \eqref{eq:angle} 只依赖 $\theta$，因此它只依赖正交结构。注意到
$\theta$ 限制在商空间（模去核 $\linf$）上正定，故 \eqref{eq:angle} 的取值必在 $[-1,1]$ 内，
且取到 $\pm1$ 当且仅当两条直线平行。

\begin{proposition}[与斜率的关系]
在 \eqref{eq:standard-theta} 的标准坐标系下，取
\[
m_i=(k_1,-1,b_1),\qquad n_i=(k_2,-1,b_2),
\]
则
\begin{equation}
\langle m,n\rangle
=\frac{k_1k_2+1}{\sqrt{k_1{}^2+1}\,\sqrt{k_2{}^2+1}}
=\cos(\theta_2-\theta_1),
\label{eq:cos-angle}
\end{equation}
其中 $\tan\theta_1=k_1$，$\tan\theta_2=k_2$。
\end{proposition}

\begin{proof}
由 \eqref{eq:standard-theta}，
\[
m\theta n=m_1n_1+m_2n_2=k_1k_2+1,\qquad
m\theta m=k_1^2+1,\qquad
n\theta n=k_2^2+1,
\]
代入 \eqref{eq:angle} 即得第一式。再由
\[
\cos(\theta_2-\theta_1)=\cos\theta_2\cos\theta_1+\sin\theta_2\sin\theta_1
=\frac{1+k_1k_2}{\sqrt{1+k_1^2}\,\sqrt{1+k_2^2}}
\]
即得第二式。
\end{proof}

有了角度的度量，就可以在任意三角形中使用正弦定理，从而得到任意两条线段的有向长度之比。
于是既可以度量角度，也可以度量线段的比。

\begin{proposition}[正交几何的内容]
正交结构给出角度与长度之比，因此平面几何中一切度量性质都可以用它表述。
\end{proposition}

它与解析几何的差别只在于尚未固定坐标系的原点、方向与单位长度。长度的度量使运算变复杂，因此
涉及长度的命题未必适合用三维矢量方法处理；而不涉及 \eqref{eq:angle} 这类运算的命题（例如蒙日圆
问题）用三维矢量方法仍可能得到简洁的处理。

\section{正交结构与圆锥曲线}

正交张量 $\theta$ 与圆锥曲线 $g$ 之间有如下两个关系。

\begin{proposition}[等轴双曲线]
记 $g:\theta=g_{ij}\theta^{ij}$。则 $g:\theta=0$ 当且仅当 $g$ 是等轴双曲线。
\end{proposition}

\begin{proof}
取 $\linf$ 的补线为 $x,y$ 平面，并把 $g$ 对角化为
\[
g_{ij}=\diag\Bigl\{\frac{1}{a^2},\frac{1}{b^2},-1\Bigr\},
\qquad a,b>0 ;
\]
对双曲线只需把 $b^2$ 换成 $-b^2$。由 \eqref{eq:standard-theta}，
$g:\theta=\dfrac{1}{a^2}+\dfrac{1}{b^2}$。对椭圆它恒为正；对双曲线它为
$\dfrac{1}{a^2}-\dfrac{1}{b^2}$，等于 $0$ 当且仅当 $a=b$，即两条渐近线互相垂直，这正是
"等轴双曲线"的定义。
\end{proof}

\begin{proposition}[圆]
设 $\linf$ 规范化使 $\linf^2=g^{ij}\linfsub{i}\linfsub{j}=-1$，其中
$\linf^i=(g^{-1})^{ij}\linfsub{j}$。则
\begin{equation}
\theta^{ij}=\linf^2g^{ij}-\linf^i\linf^j
\label{eq:circle}
\end{equation}
当且仅当 $g$ 是圆。
\end{proposition}

\begin{proof}
取 $g_{ij}=\diag\{1/a^2,1/b^2,-1\}$，$\linfsub{i}=(0,0,1)$。则
$g^{ij}=\diag\{a^2,b^2,-1\}$，$\linf^i=(g^{-1})^{ij}\linfsub{j}=(0,0,-1)$，$\linf^2=-1$。
于是 \eqref{eq:circle} 给出
\[
\theta^{ij}=-g^{ij}-\linf^i\linf^j
=\diag\{-a^2,-b^2,1\}+\diag\{0,0,-1\}
=\diag\{-a^2,-b^2,0\}.
\]
要使它与 $\theta^{ij}=\diag\{1,1,0\}$ 成比例，必须且只须 $a^2=b^2$，即 $g$ 是圆。反过来，
当 $a=b$ 时上式成为 $-a^2\diag\{1,1,0\}$，与 \eqref{eq:standard-theta} 只差一个规范因子，
\eqref{eq:circle} 成立。
\end{proof}

\begin{remark}
\eqref{eq:circle} 中的 $\linf^2$ 必须按上述规范化取定，否则等式只能把 $\theta$ 确定到相差一个
与 $g$ 有关的因子。此外这里用 $g$ 的逆 $(g^{-1})^{ij}$ 升了 $\linf$ 的指标，与用 $g$ 升指标
得到的张量应当区分开。
\end{remark}

从射影几何、仿射几何、正交几何到解析几何，结构越来越多，允许的坐标变换则越来越少：射影几何
中坐标变换可以是任意可逆线性变换；仿射几何中必须保持 $\linf$ 不变，只允许拉伸、旋转与平移；
正交几何中必须保持角度不变，只允许旋转、平移与整体放大；解析几何中坐标系固定，不允许变换。

\section{例题}

\begin{example}[双曲线切线被渐近线截得的线段]
设 $g$ 是双曲线，$a$ 是 $g$ 上一点。过 $a$ 的切线交两条渐近线于 $b,c$。则 $a$ 是线段 $BC$ 的
中点，即 $AB=AC$。
\end{example}

\begin{proof}
把 $g$ 取作度规，无穷远直线为 $\linf$。双曲线与 $\linf$ 交于两点，这两点处的切线就是两条
渐近线 $m,n$。由"直线与圆锥曲线相交"一节的结论，$m,n$ 的对称积为
\begin{equation}
(mn)_{ij}=\linf^2g_{ij}-\linfsub{i}\linfsub{j},
\qquad \linf^2=g^{ij}\linfsub{i}\linfsub{j}.
\label{eq:asymptotes}
\end{equation}
记 $a$ 处的切线为 $\alpha=a^*$，它与 $\linf$ 交于 $\infty=a\times\linf$。又
$b=m\times\alpha$、$c=n\times\alpha$ 是切线与两条渐近线的交点。中点条件 $AB=AC$ 等价于
$a,b,c,\infty$ 成调和点列，也就是 $a,\infty$ 关于退化圆锥曲线 $(mn)$ 共轭。

\begin{lemma}
设 $a$ 在 $g$ 上，$\infty=a\times\linf$，$m,n$ 由 \eqref{eq:asymptotes} 给出。则
$a,b,c,\infty$ 成调和点列当且仅当
\[
(mn)_{ij}a^i\infty^j=0 .
\]
\end{lemma}

\begin{proof}
$a,\infty$ 关于 $(mn)$ 共轭时，$a,\infty$ 把两条直线 $m,n$ 调和分割。在直线 $\alpha$ 上，
$m,n$ 与 $\alpha$ 的交点正是 $b,c$，因此 $a,\infty$ 与 $b,c$ 成调和点列。这正是二次对合中
"共轭等价于调和分割"的结论在两条直线这一退化情形下的形式。
\end{proof}

代入 \eqref{eq:asymptotes} 与 $\infty=a\times\linf$，
\[
(mn)_{ij}a^i\infty^j
=\bigl(\linf^2g_{ij}-\linfsub{i}\linfsub{j}\bigr)a^i(a\times\linf)^j .
\]
把括号中的两项分别与 $a$、$a\times\linf$ 缩并。第一项给出
\[
\linf^2\,g_{ij}a^i(a\times\linf)^j=\linf^2\,a\cdot(a\times\linf)=0,
\]
因为 $a\cdot(a\times\linf)=[a\,a\,\linf]=0$（混合积中有两个相同的因子）。第二项给出
\[
\linfsub{i}\linfsub{j}a^i(a\times\linf)^j
=(\linf\cdot a)\;\bigl(\linf\cdot(a\times\linf)\bigr)=0,
\]
同样因为 $\linf\cdot(a\times\linf)=[\linf\,a\,\linf]=0$。两项都为零，故
$(mn)_{ij}a^i\infty^j=0$，$a,b,c,\infty$ 成调和点列，$a$ 是 $BC$ 的中点，即 $AB=AC$。
\end{proof}

\begin{example}[蒙日圆]
设椭圆 $g$ 为 $\dfrac{x^2}{a^2}+\dfrac{y^2}{b^2}=1$，$a$ 是椭圆外一点，过 $a$ 作椭圆的
两条切线 $m,n$ 互相垂直。求 $a$ 的轨迹。
\end{example}

\begin{proof}
过 $a$ 的两条切线的对称积为
\[
(mn)_{ij}=a^2g_{ij}-a_ia_j,\qquad a^2=g_{ij}a^ia^j .
\]
切线 $m,n$ 垂直的条件是 $(mn)_{ij}\theta^{ij}=0$，即
\begin{equation}
a^2\,(g:\theta)-\theta^{ij}a_ia_j=0,
\qquad g:\theta=g_{ij}\theta^{ij}.
\label{eq:monge00}
\end{equation}
由 $a_i=g_{ik}a^k$，第二项可写成
\[
\theta^{ij}a_ia_j=\theta^{ij}g_{ik}g_{jl}a^ka^l
=\theta_{kl}a^ka^l,
\qquad \theta_{kl}=g_{ik}g_{jl}\theta^{ij}.
\]
于是 \eqref{eq:monge00} 成为
\[
\bigl[\theta_{kl}-(g:\theta)\,g_{kl}\bigr]a^ka^l=0 .
\]
因此 $a$ 的轨迹由 $(0,2)$ 型张量
\begin{equation}
h=\theta-(g:\theta)\,g
\qquad\text{即}\qquad
h_{ij}=\theta_{ij}-(g:\theta)\,g_{ij}
\label{eq:monge-h}
\end{equation}
描述，这是一条圆锥曲线。

下面验证它是一条圆。取
\[
g_{ij}=\diag\Bigl\{\frac{1}{a^2},\frac{1}{b^2},-1\Bigr\},\qquad
g^{ij}=\diag\{a^2,b^2,-1\},\qquad
\theta^{ij}=\diag\{1,1,0\}.
\]
按定义 $\theta_{ij}=g_{ik}g_{jl}\theta^{kl}$：逐项把 $\theta$ 的非零分量乘上
$g_{11}g_{11}=\frac{1}{a^4}$ 或 $g_{22}g_{22}=\frac{1}{b^4}$，而第三个分量因 $\theta^{33}=0$
为零，得到
\[
\theta_{ij}=\diag\Bigl\{\frac{1}{a^4},\frac{1}{b^4},0\Bigr\}.
\]
\begin{remark}
不要把 $\theta_{ij}$ 与 $\theta^{ij}$ 的分量混为一谈，特别是第三个分量容易写错：$\theta^{ij}$ 的
第三个对角元是 $0$，由它降指标得到的 $\theta_{ij}$ 的第三个对角元仍是 $0$。第二个分量只由
$g_{22}$ 的二次方决定（这里是 $1/b^4$），而不是 $1/a^2b^2$。
\end{remark}
于是
\[
g:\theta=\frac{1}{a^2}+\frac{1}{b^2},
\]
从而
\begin{align*}
h_{ij}&=\theta_{ij}-(g:\theta)\,g_{ij}
=
\begin{pmatrix}
\dfrac{1}{a^4} & 0 & 0\\[6pt]
0 & \dfrac{1}{b^4} & 0\\[6pt]
0 & 0 & 0
\end{pmatrix}
-\Bigl(\frac{1}{a^2}+\frac{1}{b^2}\Bigr)
\begin{pmatrix}
\dfrac{1}{a^2} & 0 & 0\\[6pt]
0 & \dfrac{1}{b^2} & 0\\[6pt]
0 & 0 & -1
\end{pmatrix}\\[4pt]
&=
\begin{pmatrix}
-\dfrac{1}{a^2b^2} & 0 & 0\\[6pt]
0 & -\dfrac{1}{a^2b^2} & 0\\[6pt]
0 & 0 & \dfrac{1}{a^2}+\dfrac{1}{b^2}
\end{pmatrix}.
\end{align*}
把整体的规范因子 $-\dfrac{1}{a^2b^2}$ 略去，这个张量与
\[
h_{ij}\sim\diag\{1,\;1,\;-(a^2+b^2)\}
\]
等价，对应的二次方程是
\[
x^2+y^2-(a^2+b^2)z^2=0,
\]
在仿射图 $z=1$ 中即
\[
x^2+y^2=a^2+b^2 .
\]
因此 $a$ 的轨迹是以椭圆中心为心、半径 $\sqrt{a^2+b^2}$ 的圆，称为\textbf{蒙日圆}。
\end{proof}

由 \eqref{eq:monge-h} 得到蒙日圆不依赖坐标系的表达式：轨迹是对一切满足
\[
\bigl[\theta^{ij}-(g:\theta)\,g^{ij}\bigr]a_ia_j=0
\]
的点 $a$ 构成的圆锥曲线。

\begin{example}[对中心张直角的动弦]
设椭圆 $g$ 为 $\dfrac{x^2}{a^2}+\dfrac{y^2}{b^2}=1$，直线 $l$ 与 $g$ 交于 $c,d$，椭圆的中心为
$O$。若 $OC\perp OD$，求 $l$ 所满足的条件。
\end{example}

\begin{proof}
直接表达 $OC\perp OD$ 需要先把 $c,d$ 分别与中心作叉乘得到两条直线，再与 $\theta$ 缩并。可以
避免这一步骤：$OC\perp OD$ 只涉及 $c,d$ 两个点矢，因此可以引入一个 $(0,2)$ 型张量
$\Theta_{ij}$，要求点矢满足
\begin{equation}
\Theta_{ij}c^id^j=0
\label{eq:Theta-cond}
\end{equation}
当且仅当 $OC\perp OD$。

在 $z=1$ 的仿射图中，点 $C(x_1,y_1)$、$D(x_2,y_2)$ 满足
$OC\perp OD\iff x_1x_2+y_1y_2=0$。点矢取 $(x,y,1)$ 时第三个分量总为 $1$，故只需取
\begin{equation}
\Theta_{ij}=\diag\{1,1,0\}
\label{eq:Theta-comp}
\end{equation}
即可。注意 \eqref{eq:Theta-cond} 中的 $\Theta$ 配对的是两个点矢，而定义 \eqref{eq:perp} 中的
$\theta$ 配对的是两个线矢，二者是不同的张量；本节末的注说明它们的关系，两个例子则是同一组配对的
两种用法。

设 $l$ 的线矢为 $l_i$，$l$ 与 $g$ 的两交点 $c,d$ 的对称积为
\[
(cd)^{ij}=l^2g^{ij}-l^il^j,\qquad l^2=g^{ij}l_il_j .
\]
由 \eqref{eq:Theta-cond} 与对称性，
\[
0=\Theta_{ij}c^id^j=\tfrac12\Theta_{ij}(cd)^{ij}
=\tfrac12\bigl[l^2\,g^{ij}\Theta_{ij}-l^il^j\Theta_{ij}\bigr].
\]
因此 $l$ 满足
\begin{equation}
l^2\,(g:\Theta)-l^il^j\Theta_{ij}=0,
\qquad g:\Theta=g^{ij}\Theta_{ij}.
\label{eq:chord00}
\end{equation}
由 $\Theta_{ij}=g_{ik}g_{jl}\Theta^{kl}$，第二项可写成
\[
l^il^j\Theta_{ij}=\Theta^{kl}\,(g_{ik}l^i)(g_{jl}l^j)
=\Theta^{kl}\tilde l_k\tilde l_l,
\qquad \tilde l_i=g_{ik}l^k .
\]
于是 \eqref{eq:chord00} 成为
\[
\bigl[\Theta^{ij}-(g:\Theta)\,g^{ij}\bigr]\tilde l_i\tilde l_j=0 .
\]
把括号中的张量记为
\begin{equation}
h^{ij}=\Theta^{ij}-(g:\Theta)\,g^{ij},
\label{eq:chord-h}
\end{equation}
则条件写成 $h^{ij}\tilde l_i\tilde l_j=0$。这正是说 $l$ 与圆锥曲线 $h$ 相切：与 $g$ 相切的条件
是 $l^2=0$，而这里取代 $g$ 的是 $h$。于是 $l$ 的包络是一条圆锥曲线。它与蒙日圆问题是对偶的：
蒙日圆求的是过定点的两条垂线的交点的轨迹，这里求的是对中心张直角的弦的包络。

下面算出 $h$。取
\[
g_{ij}=\diag\Bigl\{\frac{1}{a^2},\frac{1}{b^2},-1\Bigr\},\qquad
g^{ij}=\diag\{a^2,b^2,-1\},\qquad
\Theta_{ij}=\diag\{1,1,0\}.
\]
由 $\Theta^{ij}=g^{ik}g^{jl}\Theta_{kl}$ 得
\[
\Theta^{ij}=\diag\Bigl\{\frac{1}{a^2},\frac{1}{b^2},0\Bigr\},
\]
于是
\[
g:\Theta=g^{ij}\Theta_{ij}=a^2+b^2 .
\]
因此
\begin{align*}
h^{ij}&=\Theta^{ij}-(g:\Theta)\,g^{ij}
=
\begin{pmatrix}
\dfrac{1}{a^2} & 0 & 0\\[6pt]
0 & \dfrac{1}{b^2} & 0\\[6pt]
0 & 0 & 0
\end{pmatrix}
-(a^2+b^2)
\begin{pmatrix}
a^2 & 0 & 0\\
0 & b^2 & 0\\
0 & 0 & -1
\end{pmatrix}\\[4pt]
&=-\frac{a^2+b^2}{a^2b^2}\,
\diag\{a^2,\;b^2,\;-1\}.
\end{align*}
也就是说 $h^{ij}$ 与 $g^{ij}$ 成比例。由 $h^{ij}=C\,g^{ij}$
（$C=-\dfrac{a^2+b^2}{a^2b^2}$）得
\begin{align*}
(h^{-1})^{ij}&=\frac1C\,(g^{-1})^{ij}
=\frac{a^2b^2}{a^2+b^2}\,
\diag\Bigl\{\frac{1}{a^2},\frac{1}{b^2},-1\Bigr\}\\[4pt]
&\sim\diag\{b^2,\;a^2,\;-(a^2+b^2)\}
\sim\diag\Bigl\{1,\;1,\;-\frac{a^2+b^2}{a^2b^2}\Bigr\}.
\end{align*}
相切条件 $l_il_j(h^{-1})^{ij}=0$ 因此成为
\[
l_1^2+l_2^2-\frac{a^2+b^2}{a^2b^2}l_3^2=0 .
\]
对线矢 $l=(u,v,w)$，$\dfrac{|w|}{\sqrt{u^2+v^2}}$ 是直线与原点的距离。于是 $l$ 与原点的距离
$r$ 满足
\[
\frac{1}{r^2}=\frac{a^2+b^2}{a^2b^2},
\qquad\text{即}\qquad
r=\frac{ab}{\sqrt{a^2+b^2}} .
\]
因此所求条件是：$l$ 是对中心在原点、半径 $r=\dfrac{ab}{\sqrt{a^2+b^2}}$ 的圆所作的切线，
即 $l$ 的包络是这个圆。
\end{proof}

\begin{example}[数值验证]
取 $a=2$、$b=1$，此时 $1/r^2=1/4+1=5/4$，即 $r=2/\sqrt5$。设 $l$ 是水平直线
\[
l:\ y=t,
\]
它与椭圆 $\dfrac{x^2}{4}+y^2=1$ 的交点为 $\bigl(\pm 2\sqrt{1-t^2},\ t\bigr)$，其中
$|t|\le 1$。把两个交点记作 $c,d$。由 \eqref{eq:Theta-comp}，
\[
\Theta_{ij}c^id^j
=\bigl(2\sqrt{1-t^2}\bigr)\bigl(-2\sqrt{1-t^2}\bigr)+t^2
=-4(1-t^2)+t^2=5t^2-4 .
\]
于是 $OC\perp OD$ 等价于 $5t^2=4$，即 $t=\pm\dfrac{2}{\sqrt5}$。取 $t=\dfrac{2}{\sqrt5}$，则
\[
x^2=4\Bigl(1-\frac45\Bigr)=\frac45,
\qquad
c=\Bigl(\frac{2}{\sqrt5},\frac{2}{\sqrt5}\Bigr),\quad
d=\Bigl(-\frac{2}{\sqrt5},\frac{2}{\sqrt5}\Bigr),
\]
且
\[
\Theta_{ij}c^id^j
=-\frac45+\frac45=0,
\]
确实满足 $OC\perp OD$。另一方面，$l:y=\dfrac{2}{\sqrt5}$ 的线矢为
$l_i=\bigl(0,1,-\tfrac{2}{\sqrt5}\bigr)$，而
\[
l_il_j(h^{-1})^{ij}
=\frac14\cdot 1-\frac54\cdot\frac45=1-1=0,
\]
与相切条件 $l_il_j(h^{-1})^{ij}=0$ 相符；$l$ 到中心的距离为 $2/\sqrt5=r$，与前面的结论一致。
\end{example}

\begin{remark}[$\theta$ 与 $\Theta$ 的关系]
$\theta^{ij}$ 与 $\Theta_{ij}$ 的分量在标准坐标系下恰好都写成 $\diag\{1,1,0\}$，容易把二者当成
同一个张量。它们确实不同：$\theta$ 是 $(2,0)$ 型、作用于线矢，$\Theta$ 是 $(0,2)$ 型、作用于点矢。
用 $g$ 升降指标并不能把二者互换，例如在同一坐标系下
\[
\Theta^{ij}=g^{ik}g^{jl}\Theta_{kl}
=\diag\Bigl\{\frac{1}{a^2},\frac{1}{b^2},0\Bigr\},
\qquad
\theta_{ij}=g_{ik}g_{jl}\theta^{kl}
=\diag\Bigl\{\frac{1}{a^4},\frac{1}{b^4},0\Bigr\},
\]
它们既彼此不同，一般也都不与 $\diag\{1,1,0\}$ 成比例；只有当 $g$ 是圆（$a=b$）时，各张量才
彼此只相差常数因子。
\end{remark}

\section{关于记号的注}

涉及多条圆锥曲线时，指标的使用会出现歧义。如果只有一条圆锥曲线 $g_{ij}$ 并把它用作度规，那么
$g^{ij}$ 既可以理解为用度规升指标得到的结果，也可以理解为 $g_{ij}$ 的逆张量，这两种理解一致。
但如果还涉及另一条圆锥曲线 $h_{ij}$，用 $g$ 升指标得到的张量与对 $h$ 求逆得到的张量就不再相同。
为避免混淆，约定：$h^{ij}$ 指用 $g$ 升指标得到的张量，即 $h^{ij}=g^{ik}g^{jl}h_{kl}$；而 $h$ 的
逆用 $(h^{-1})^{ij}$ 表示，它出现在相切条件 $l_il_j(h^{-1})^{ij}=0$ 中。判断一条圆锥曲线是圆
还是等轴双曲线，以及在 $h^{ij}l_il_j=0$ 与 $(h^{-1})^{ij}l_il_j=0$ 之间转换时，都要先看清所用
的是哪一个张量。
   % 平行与正交
% 第十三节：焦点与准线
\chapter{焦点与准线}

\section{复点与垂直结构}

平行与正交一章说明，正交结构可以看作由两个无穷远处的复点确定。记
\begin{equation}
z^i=(1,\ii,0),\qquad \bar z^i=(1,-\ii,0),
\label{eq:zz}
\end{equation}
其中 $\ii$ 为虚数单位。两点都在无穷远直线 $\linf$ 上，事实上
\[
z\times\bar z=(0,0,-2\ii)\sim(0,0,1)=\linf .
\]
$z$ 与 $\bar z$ 互为复共轭，它们与实射影平面上的实点、实直线具有相同的代数地位。

本节的约定是：把坐标含复数的点矢、线矢与实的点矢、线矢同样处理，一切运算（线性组合、缩并、
叉乘、求逆）都按代数规则进行，只在需要判断某个量是否为实数时才使用复共轭。经验表明，只要问题
的提法本身不涉及实数性，坐标中的虚部不影响结论；下文凡按几何意义应当为实的量，算出的结果都是
实的。

垂直结构取为这两点的对称积
\begin{equation}
\theta^{ij}=(\sym{z}{\bar z})^{ij}=z^i\bar z^j+z^j\bar z^i .
\label{eq:theta-zz}
\end{equation}
按分量计算，$z^i\bar z^j$ 只在 $(1,1)$ 与 $(2,2)$ 处分别为 $1$ 与 $1$，其余分量为零，因此
\[
\theta^{ij}=
\begin{pmatrix}
2&0&0\\
0&2&0\\
0&0&0
\end{pmatrix}
\sim\diag\{1,1,0\}.
\]
可见 \eqref{eq:theta-zz} 定义的 $\theta$ 与实标准形式 $\diag\{1,1,0\}$ 只相差规范因子 $2$，属于
同一个倍数类。在下文的分量计算中直接取
\begin{equation}
\theta^{ij}=\diag\{1,1,0\}.
\label{eq:theta-standard}
\end{equation}

对任意线矢 $\ell$，
\begin{equation}
\theta^{ij}\ell_i\ell_j=2(\ell\cdot z)(\ell\cdot\bar z),
\label{eq:theta-line}
\end{equation}
因此 $\theta$ 所对应的退化二次点正是两点 $z,\bar z$：经过 $\theta$ 的直线恰为过 $z$ 的直线与过
$\bar z$ 的直线的并。$\theta$ 的号差为 $(1,1,0)$，它是退化的，核由 $\linf$ 张成。

\section{焦点方程与离心率}

\begin{definition}[焦点方程]
设 $g$ 是一条圆锥曲线，$\theta$ 是垂直结构。方程
\begin{equation}
\det\bigl(\theta^{ij}-\lambda g^{ij}\bigr)=0
\label{eq:focus-eq}
\end{equation}
称为 $g$ 的\textbf{焦点方程}，其中的 $\lambda$ 称为\textbf{特征值}。这里 $g^{ij}$ 是 $g_{ij}$ 的
逆 $(g^{-1})^{ij}$。
\end{definition}

焦点方程是三维矢量空间中的特征值方程：把 $\theta,g$ 看作矩阵，由 $g$ 可逆，
\[
\det(\theta-\lambda g)=\det g\cdot\det\bigl(g^{-1}\theta-\lambda I\bigr),
\]
故 \eqref{eq:focus-eq} 的根正是线性变换 $g^{-1}\theta$ 的特征值。由于 $g^{-1}\theta$ 是三阶
矩阵，方程关于 $\lambda$ 至多是三次的。又 $\theta$ 退化，$\det\theta=0$，即
$\det(g^{-1}\theta)=0$，所以 $\lambda=0$ 是一个根。

\begin{proposition}[特征值]
焦点方程 \eqref{eq:focus-eq} 有三个根
\[
\lambda_1,\qquad \lambda_2,\qquad 0,
\]
其中 $\lambda_1,\lambda_2$ 是二次方程
\[
\det\bigl(\theta^{ij}-\lambda g^{ij}\bigr)=0
\]
在去掉因子 $\lambda$ 之后所剩的两个根。
\end{proposition}

\begin{proposition}[规范变换下特征值的变化]
把 $\theta$ 换成 $\nu\theta$、把 $g$ 换成 $\mu g$（$\mu,\nu\ne0$）时，\eqref{eq:focus-eq} 的
非零根全部乘以同一个因子 $\nu/\mu$。
\end{proposition}

\begin{proof}
直接计算
\[
\det(\nu\theta-\lambda\mu g)
=\nu^3\det\Bigl(\theta-\frac{\lambda\mu}{\nu}g\Bigr).
\]
因此 $\nu\theta-\lambda\mu g$ 退化当且仅当 $\theta-\lambda_0g$ 退化，其中
$\lambda_0=\lambda\mu/\nu$，即 $\lambda=\lambda_0\nu/\mu$。原来的每个特征值都乘以 $\nu/\mu$。
\end{proof}

由此，单个特征值 $\lambda_1,\lambda_2$ 依赖 $g,\theta$ 的规范选取，只有比值
$\lambda_2/\lambda_1$ 是规范无关的，因而有几何意义。据此定义离心率。

\begin{definition}[离心率]
在焦点方程的两个非零根中，约定 $\lambda_1$ 对应实焦点所在的那条主轴（见下一节），$\lambda_2$
对应另一条主轴，并取
\begin{equation}
e=\sqrt{1-\frac{\lambda_2}{\lambda_1}}.
\label{eq:ecc}
\end{equation}
称 $e$ 为圆锥曲线 $g$ 的\textbf{离心率}。
\end{definition}

\begin{remark}[次序的约定]
$\lambda_1,\lambda_2$ 是两个根，本身没有先后之分；\eqref{eq:ecc} 只在固定取法之后才有确定的值。
上文把 $\lambda_1$ 约定为对应实焦点所在主轴的根，这使 $e$ 在下述意义下取到标准值：椭圆时
$0<e<1$，抛物线时 $e=1$，双曲线时 $e>1$，圆时 $e=0$。若不约定次序，则 $e$ 可能对应
$1-e^2$ 的另一个分支。
\end{remark}

\section{六个焦点与三条主轴}

由上一节的结论，$\lambda_1,\lambda_2,0$ 分别对应三个退化二次点
\begin{equation}
X^{ij}=\theta^{ij}-\lambda_1g^{ij},\qquad
Y^{ij}=\theta^{ij}-\lambda_2g^{ij},\qquad
\theta^{ij}=\theta^{ij}-0\cdot g^{ij}.
\label{eq:XYZ}
\end{equation}
三者都满足 $\det=0$。一般情况下它们的秩为 $2$，号差为 $(1,1,0)$，因而每一个都可以写成两个点的
对称积。按第八章的记号，设
\begin{equation}
X=(\sym{f_{X1}}{f_{X2}}),\qquad
Y=(\sym{f_{Y1}}{f_{Y2}}),\qquad
\theta=(\sym{z}{\bar z}).
\label{eq:focus-sym}
\end{equation}

\begin{definition}[焦点]
\eqref{eq:focus-sym} 中的六个点
\[
f_{X1},\ f_{X2},\ f_{Y1},\ f_{Y2},\ z,\ \bar z
\]
称为圆锥曲线 $g$ 的\textbf{焦点}。其中 $f_{X1},f_{X2}$ 称为一对配对焦点，$f_{Y1},f_{Y2}$ 与
$z,\bar z$ 也各称一对配对焦点；每一对配对焦点的对称积依次为 $X,Y,\theta$。
\end{definition}

于是圆锥曲线共有六个焦点，两两配对成三对。

\begin{proposition}[主轴]
设 $x_i$ 是 $\lambda_1$ 对应的特征线矢，即
$X^{ij}x_j=0$，则直线 $x$ 经过 $f_{X1},f_{X2}$；同理 $\lambda_2$ 对应的特征线矢 $y$ 经过
$f_{Y1},f_{Y2}$，而 $\lambda=0$ 对应的特征线矢是无穷远直线 $\linf$。这三条直线分别称为
$x$ \textbf{主轴}、$y$ \textbf{主轴}与无穷远直线。
\end{proposition}

\begin{proof}
由 $X^{ij}=f_{X1}^if_{X2}^j+f_{X2}^if_{X1}^j$，对任意线矢 $x$，
\[
X^{ij}x_j=(f_{X2}\cdot x)f_{X1}^i+(f_{X1}\cdot x)f_{X2}^i,
\]
其中 $f\cdot x=f^ix_i$。$f_{X1},f_{X2}$ 不共线，故 $X^{ij}x_j=0$ 当且仅当
$f_{X1}\cdot x=f_{X2}\cdot x=0$，即直线 $x$ 同时经过 $f_{X1},f_{X2}$。对 $\lambda_2$ 同理。

对 $\lambda=0$，条件 $\theta^{ij}\linfsub{j}=0$ 由 \eqref{eq:theta-standard} 直接看出，核由
$\linf$ 张成。故 $\lambda=0$ 对应的特征线矢就是 $\linf$。
\end{proof}

\begin{proposition}[主轴两两共轭]
$x,y,\linf$ 三条直线两两关于 $g$ 共轭，即对其中任意两条 $m,n$ 都有
$g^{ij}m_in_j=0$。
\end{proposition}

\begin{proof}
设 $m,n$ 分别对应特征值 $\lambda,\mu$，即
$\theta^{ij}m_j=\lambda g^{ij}m_j$、$\theta^{ij}n_j=\mu g^{ij}n_j$。用 $n_i$ 缩并第一式，
用 $m_i$ 缩并第二式，由于 $\theta^{ij}$ 对称，两式左端相等，于是
\[
\lambda\,g^{ij}m_in_j=\mu\,g^{ij}m_in_j,
\qquad\text{即}\qquad
(\lambda-\mu)\,g^{ij}m_in_j=0 .
\]
当 $\lambda\ne\mu$ 时得 $g^{ij}m_in_j=0$。三个特征值 $\lambda_1,\lambda_2,0$ 两两不同时结论
成立；若有重根，则该情形对应圆或抛物线的退化位置，此时相应的两条主轴重合或其中之一退化，
命题按极限理解。
\end{proof}

\section{圆锥曲线的分类}

特征值的符号决定圆锥曲线的类型。

\begin{proposition}[按特征值分类]
设 $\lambda_1,\lambda_2$ 是两个非零特征值，按上一节的约定排列。
\begin{enumerate}
\item 若 $\lambda_1\lambda_2>0$，则 $g$ 是\textbf{椭圆}，此时 $\lambda_2/\lambda_1\in(0,1)$，
$0<e<1$；
\item 若 $\lambda_1\lambda_2=0$（即某个非零根为零），则 $g$ 是\textbf{抛物线}，$e=1$；
\item 若 $\lambda_1\lambda_2<0$，则 $g$ 是\textbf{双曲线}，$e>1$。
\end{enumerate}
\end{proposition}

两类特殊的圆锥曲线也由特征值刻画。

\begin{proposition}[圆与等轴双曲线]
\begin{enumerate}
\item $g$ 是圆，当且仅当焦点方程有重根，即 $\lambda_1=\lambda_2$，此时 $e=0$；
\item $g$ 是等轴双曲线，当且仅当 $\lambda_1+\lambda_2=0$，此时 $e=\sqrt2$。
\end{enumerate}
\end{proposition}

\begin{proof}
第二款：由第一款的注记，若把 $\theta$ 降指标，则
\[
\lambda_1+\lambda_2=\tr(g^{-1}\theta)=g^{ij}\theta_{ij}=g_{ij}\theta^{ij}=g:\theta .
\]
等轴双曲线的定义是两条渐近线互相垂直，在正交几何中等价于 $g:\theta=0$，故
$g$ 为等轴双曲线当且仅当 $\lambda_1+\lambda_2=0$。代入 \eqref{eq:ecc} 得
$e=\sqrt{1-(-1)}=\sqrt2$。
\end{proof}

\begin{remark}[特征值的规范依赖]
$\lambda_1,\lambda_2$ 各自随规范改变，但 $\lambda_1=\lambda_2$、$\lambda_1\lambda_2=0$ 与
$\lambda_1+\lambda_2=0$ 都是规范无关的条件，因为规范变化把两者乘以同一个因子。因此上面的分类
条件有意义。
\end{remark}

\section{椭圆情形的分量计算}

下面取椭圆具体算出全部焦点，以验证上面的定义与解析几何中的焦点一致。设椭圆为
\[
\frac{x^2}{a^2}+\frac{y^2}{b^2}=1,\qquad a>b>0,
\]
其张量及其逆为
\begin{equation}
g_{ij}=\diag\Bigl\{\frac{1}{a^2},\frac{1}{b^2},-1\Bigr\},
\qquad
g^{ij}=\diag\{a^2,b^2,-1\}.
\label{eq:ellipse-g}
\end{equation}
按 \eqref{eq:theta-standard} 取 $\theta^{ij}=\diag\{1,1,0\}$，则
\[
\theta^{ij}-\lambda g^{ij}
=\diag\{1-\lambda a^2,\ 1-\lambda b^2,\ \lambda\},
\]
于是
\begin{equation}
\det\bigl(\theta^{ij}-\lambda g^{ij}\bigr)
=\lambda\,(1-\lambda a^2)(1-\lambda b^2).
\label{eq:ellipse-det}
\end{equation}
三个根为
\[
\lambda=0,\qquad \lambda=\frac{1}{a^2},\qquad \lambda=\frac{1}{b^2}.
\]
因 $a>b>0$，有 $1/b^2>1/a^2$。按约定 $\lambda_1$ 对应实焦点所在主轴，取
\begin{equation}
\lambda_1=\frac{1}{b^2},\qquad \lambda_2=\frac{1}{a^2},
\qquad
\frac{\lambda_2}{\lambda_1}=\frac{b^2}{a^2}.
\label{eq:ellipse-lambda}
\end{equation}
记 $c^2=a^2-b^2>0$，则
\[
e=\sqrt{1-\frac{\lambda_2}{\lambda_1}}
=\sqrt{1-\frac{b^2}{a^2}}
=\frac{c}{a},
\]
与解析几何中的离心率一致。

把 $\lambda_1=1/b^2$ 代入 \eqref{eq:XYZ}，得
\[
X^{ij}=\theta^{ij}-\frac{1}{b^2}g^{ij}
=\diag\Bigl\{-\frac{c^2}{b^2},\ 0,\ \frac{1}{b^2}\Bigr\}
\sim\diag\{-2c^2,\ 0,\ 2\}.
\]
后一个矩阵可以分解为对称积
\begin{equation}
X^{ij}=
\begin{pmatrix}
-2c^2&0&0\\
0&0&0\\
0&0&2
\end{pmatrix}
=
\begin{pmatrix}
-c\\0\\1
\end{pmatrix}
\begin{pmatrix}
-c&0&1
\end{pmatrix}
+
\begin{pmatrix}
c\\0\\1
\end{pmatrix}
\begin{pmatrix}
c&0&1
\end{pmatrix}.
\label{eq:ellipse-X}
\end{equation}
即
\[
X=(\sym{f_{X1}}{f_{X2}}),\qquad
f_{X1}=(-c,0,1),\qquad f_{X2}=(c,0,1).
\]
在仿射图 $z=1$ 中，这一对焦点的坐标是 $(-c,0)$ 与 $(c,0)$，正是解析几何中椭圆的两个焦点。

把 $\lambda_2=1/a^2$ 代入 \eqref{eq:XYZ}，得
\[
Y^{ij}=\theta^{ij}-\frac{1}{a^2}g^{ij}
=\diag\Bigl\{0,\ \frac{c^2}{a^2},\ \frac{1}{a^2}\Bigr\}
\sim\diag\{0,\ 2c^2,\ 2\},
\]
它同样可以分解为对称积
\begin{equation}
Y^{ij}=
\begin{pmatrix}
0&0&0\\
0&2c^2&0\\
0&0&2
\end{pmatrix}
=
\begin{pmatrix}
0\\ \ii c\\1
\end{pmatrix}
\begin{pmatrix}
0&-\ii c&1
\end{pmatrix}
+
\begin{pmatrix}
0\\ -\ii c\\1
\end{pmatrix}
\begin{pmatrix}
0&\ii c&1
\end{pmatrix}.
\label{eq:ellipse-Y}
\end{equation}
因此
\[
f_{Y1}=(0,\ii c,1),\qquad f_{Y2}=(0,-\ii c,1),
\]
它们的仿射坐标是 $(0,\ii c)$ 与 $(0,-\ii c)$，不在实平面内。

最后，$\lambda=0$ 对应的一对焦点是 $z=(1,\ii,0)$ 与 $\bar z=(1,-\ii,0)$，它们位于无穷远
直线上。

\begin{remark}[解析几何中可见的焦点]
六个焦点中，$f_{Y1},f_{Y2}$ 的坐标含虚数，$z,\bar z$ 在无穷远，只有 $f_{X1},f_{X2}$ 落在实
仿射平面内。因此在解析几何中通常只讨论 $(-c,0)$ 与 $(c,0)$ 这两个焦点；另外四个焦点是射影
构造中不可省去的部分，但在实平面内看不到。
\end{remark}

此时三条主轴分别是
\[
x=(0,1,0)\quad(y=0),\qquad
y=(1,0,0)\quad(x=0),\qquad
\linf=(0,0,1).
\]
$x$ 主轴是 $f_{X1},f_{X2}$ 的连线，即 $x$ 轴；$y$ 主轴是 $f_{Y1},f_{Y2}$ 的连线，即 $y$ 轴。
用 $g^{ij}=\diag\{a^2,b^2,-1\}$ 直接验证三者两两共轭：
$g^{ij}x_iy_j=g^{ij}x_i\linfsub{j}=g^{ij}y_i\linfsub{j}=0$，因为 $g^{ij}$ 是对角矩阵。

\section{六个焦点的几何构造}

现在说明六个焦点的几何来源。从 \eqref{eq:focus-eq} 出发，考察曲线系
\[
\theta^{ij}-\lambda g^{ij}.
\]
其中的成员都是 $(2,0)$ 型对称张量，即二次点。$\theta$ 与 $g^{ij}$ 各自有明确的几何意义：
\begin{itemize}
\item $\theta$ 是退化二次点 $(\sym{z}{\bar z})$。由 \eqref{eq:theta-line}，经过 $\theta$ 的直线恰
为过 $z$ 或过 $\bar z$ 的直线；
\item $g^{ij}$ 是圆锥曲线 $g$ 的对偶二次点，经过它的直线恰为 $g$ 的全部切线。
\end{itemize}
因此曲线系 $\theta-\lambda g$ 的每个成员都含有 $\theta$ 与 $g$ 共同含有的直线，也就是既经过
$z$ 或 $\bar z$、又与 $g$ 相切的直线。这样的直线共有四条：从 $z$ 引出的两条切线与从 $\bar z$
引出的两条切线。记这四条切线为 $\ell_1,\ell_2,\ell_3,\ell_4$。

另一方面，$X,Y,\theta$ 都是曲线系中行列式为零的成员，即退化二次点。一个秩为 $2$ 的退化二次点
对应两个点，经过它的直线恰为过其中任一者的直线。既然 $X$ 含有全部四条切线
$\ell_1,\ell_2,\ell_3,\ell_4$，这四条直线就必须分成两对，每一对各自过 $X$ 的两个点之一。把四条
直线分成两对的方式共有三种；三种分法分别给出 $X,Y,\theta$ 三个退化二次点。$\theta$ 对应的分法
是把过 $z$ 的两条切线放在一起、过 $\bar z$ 的两条放在一起，两个点正是 $z,\bar z$；$X$ 与 $Y$
对应的分法给出其余四个焦点。

把四条切线两两求交，共得 $6$ 个交点。上述三种分法所对应的恰好是这些交点的三种配对，六个焦点
正是这 $6$ 个交点。于是得到下面三条等价的刻画。

\begin{proposition}[焦点的几何刻画]
设 $g$ 是圆锥曲线，$\theta=(\sym{z}{\bar z})$ 是垂直结构。以下三条性质中的每一条都唯一确定了
$g$ 的六个焦点：
\begin{enumerate}
\item $z,\bar z$ 关于圆锥曲线的切线相交于焦点；
\item 过焦点作圆锥曲线的两条切线，它们分别经过 $z,\bar z$；
\item 圆锥曲线过 $z,\bar z$ 的四条切线交于 $6$ 个点；在这 $6$ 个点中取一对不落在同一条切线上
的点 $p,q$，令 $\theta'=(\sym{p}{q})$，则这 $6$ 个点正是 $g$ 关于垂直结构 $\theta'$ 的六个焦点。
\end{enumerate}
\end{proposition}

\begin{remark}[一般的生成点]
上面的构造并不依赖 $z,\bar z$ 的具体位置。任意给定两个点 $p,q$，都可以作出对称积
$\theta'=(\sym{p}{q})$，并写出方程
\[
\det\bigl(g^{ij}-\lambda\theta'^{ij}\bigr)=0,
\]
按完全相同的方式得到六个点。这时称这六个点为 $p,q$ 生成的焦点。第 $3$ 条性质说明，从圆锥曲线
的四条切线及其交点出发，任取合适的一对点作为新的生成点，所得的六个点保持不变。
\end{remark}

\section{准线}

前面的构造是对 $(2,0)$ 型张量进行的。同样可以对 $(0,2)$ 型张量求解
\[
\det\bigl(\theta_{ij}-\lambda g_{ij}\bigr)=0,
\qquad
\theta_{ij}=g_{ik}g_{jl}\theta^{kl}.
\]
这个方程与 \eqref{eq:focus-eq} 有相同的特征值，由它构造出的几何量与前者基本相同，只相差一次
指标的升降。

\begin{definition}[准线]
焦点的极线称为圆锥曲线的\textbf{准线}。
\end{definition}

六个焦点给出六条准线，按配对也分成三对，与三对焦点相对应。与焦点一样，通常只有一对准线落在
实仿射平面内可见的位置。

以椭圆为例。取焦点 $f_{X1}=(c,0,1)$，其极线为
\[
l_i=g_{ij}f_{X1}^j=\Bigl(\frac{c}{a^2},\ 0,\ -1\Bigr),
\]
即直线
\[
\frac{c}{a^2}x-z=0
\quad\Longleftrightarrow\quad
x=\frac{a^2}{c}.
\]
这正是解析几何中的准线 $x=a^2/c$。焦点 $f_{X2}=(-c,0,1)$ 的极线是 $x=-a^2/c$。虚焦点
$f_{Y1}=(0,\ii c,1)$ 的极线为 $\dfrac{\ii c}{b^2}y-z=0$，即
$y=-\ii b^2/c$，不在实平面内；$z,\bar z$ 的极线经过中心，也是虚的或无穷远的直线。

\section{例题：焦点、切线与准线}

\begin{example}\label{ex:focus-directrix}
设 $g$ 是椭圆，$F$ 是它的一个焦点。过 $F$ 的直线与 $g$ 交于 $A,B$ 两点，$A,B$ 处的两条切线
交于点 $H$。则 $H$ 落在 $F$ 对应的准线上，并且 $AF\perp FH$。
\end{example}

证明分两部分。

\begin{proof}[第一部分：$H$ 在准线上]
$A,B$ 是过 $F$ 的直线与 $g$ 的交点，$A,B$ 处的切线交于 $H$。由极点与极线的性质，$H$ 的极线
就是两个切点 $A,B$ 的连线，即直线 $AB$。由于 $F$ 在直线 $AB$ 上，按配极的对称性（La Hire
定理），$F$ 的极线经过 $H$。而 $F$ 的极线按定义就是 $F$ 对应的准线，因此 $H$ 在准线上。
\end{proof}

\begin{proof}[第二部分：$FA$ 与 $FH$ 共轭当且仅当它们垂直]
先说明 $FA$ 与 $FH$ 关于 $g$ 共轭。直线 $AB$ 的极点是 $H$；又 $H$ 在直线 $FH$ 上，把 $AB$
的极线 $H$ 与直线 $FH$ 相接可知 $AB$ 与 $FH$ 关于 $g$ 共轭。而 $F,A,B$ 共线，$FA$ 就是直线
$AB$，于是 $FA$ 与 $FH$ 关于 $g$ 共轭。

于是问题化为：过焦点 $F$ 的两条直线关于 $g$ 共轭，当且仅当它们垂直。这由下面两条引理给出。
\end{proof}

\begin{lemma}\label{lem:conj-harmonic}
设 $P$ 是不在圆锥曲线 $g$ 上的点，$t_1,t_2$ 是从 $P$ 引出的两条切线。则过 $P$ 的两条直线
$m,n$ 关于 $g$ 共轭（即 $g^{ij}m_in_j=0$）当且仅当 $m,n,t_1,t_2$ 构成调和线束。
\end{lemma}

\begin{proof}
共轭关系、切线与调和线束都在射影变换下保持，因此可以取 $g$ 为圆、$P$ 在一条直径上。取
\[
g:\ x^2+y^2=1,\qquad P=(d,0),\qquad d\ne0,\ \pm1 .
\]
过 $P$ 斜率为 $s$ 的直线为 $m_i=(s,-1,-sd)$，即 $y=s(x-d)$。它的极点为
\[
D^i=g^{ij}m_j=(s,-1,sd),\qquad g^{ij}=\diag\{1,1,-1\}.
\]
过 $P$ 且经过 $D$ 的直线就是 $m$ 的共轭直线，直接计算
\[
P\times D=(d,0,1)\times(s,-1,sd)=\bigl(1,\ s(1-d^2),\ -d\bigr),
\]
其斜率（表示为 $y=t(x-d)$）满足
\[
t=\frac{1}{s(d^2-1)}.
\]
这条直线与 $g$ 相切当且仅当它到原点的距离为 $1$，即
$|sd|/\sqrt{s^2+1}=1$，解得
\[
s=\pm\frac{1}{\sqrt{d^2-1}}=\pm k,\qquad k=\frac{1}{\sqrt{d^2-1}} .
\]
两条切线的斜率因此是 $\pm k$。斜率 $s$ 与其共轭直线的斜率 $k^2/s$ 的交比
\[
\Bigl(s,\ \frac{k^2}{s};\ k,\ -k\Bigr)
=\frac{(s-k)(k^2/s+k)}{(s+k)(k^2/s-k)}
=\frac{(s-k)\cdot k(k+s)/s}{(s+k)\cdot k(k-s)/s}
=-1 .
\]
故 $m,n$ 与两条切线成调和线束。
\end{proof}

\begin{lemma}\label{lem:perp-harmonic}
设 $\theta=(\sym{z}{\bar z})$ 是垂直结构，$P$ 是任一点。过 $P$ 的两条直线 $m,n$ 垂直，当且仅当
$m,n$ 与直线 $Pz,P\bar z$ 构成调和线束。
\end{lemma}

\begin{proof}
取 \eqref{eq:theta-standard} 的标准坐标系。直线 $Pz$ 与 $P\bar z$ 分别是以 $z,\bar z$ 为方向的
直线，由于 $z=(1,\ii,0)$、$\bar z=(1,-\ii,0)$，它们的斜率是 $\ii$ 与 $-\ii$。过 $P$ 的两条
直线写成 $m_i=(s,-1,b_1)$、$n_i=(t,-1,b_2)$，则
\[
\theta^{ij}m_in_j=st+1,
\]
故 $m\perp n$ 当且仅当 $st=-1$。另一方面，
\[
(s,t;\ii,-\ii)
=\frac{(s-\ii)(t+\ii)}{(s+\ii)(t-\ii)}
=-1
\]
等价于 $(s-\ii)(t+\ii)=-(s+\ii)(t-\ii)$，展开即 $st+1=-st-1$，也就是 $st=-1$。两种条件
一致。
\end{proof}

\begin{proof}[例题的完成]
设 $t_1,t_2$ 是从焦点 $F$ 引出的 $g$ 的两条切线。由焦点的构造，$F$ 是四条切线中某两条的交点，
其中一条经过 $z$、一条经过 $\bar z$；而经过 $F$ 且经过 $z$ 的切线只能是 $Fz$，经过 $F$ 且经过
$\bar z$ 的切线只能是 $F\bar z$。因此
\[
\{t_1,t_2\}=\{Fz,\ F\bar z\}.
\]
由引理 \ref{lem:conj-harmonic}，$FA$ 与 $FH$ 关于 $g$ 共轭，当且仅当 $FA,FH$ 与 $Fz,F\bar z$
构成调和线束；由引理 \ref{lem:perp-harmonic}，$FA\perp FH$ 也当且仅当 $FA,FH$ 与
$Fz,F\bar z$ 构成调和线束。两个条件等价，所以 $FA$ 与 $FH$ 共轭当且仅当 $FA\perp FH$。

前面已经证明 $FA$ 与 $FH$ 确实共轭，故 $FA\perp FH$，即 $AF\perp FH$。结合第一部分，$H$ 在
准线上。命题得证。
\end{proof}

\section{离心率的交比表示}

离心率有射影的解释。先引入一个由 $g$ 决定的交比型不变量。对两个点 $p,q$，记
\begin{equation}
\{p,q\}_g
=\frac{\bigl(g_{ij}p^iq^j\bigr)^2}
{\bigl(g_{ij}p^ip^j\bigr)\bigl(g_{ij}q^iq^j\bigr)},
\label{eq:invariant-point}
\end{equation}
只要分母不为零，它就是 $p,q$ 的规范无关的函数。对两条直线 $m,n$ 类似地记
\[
\{m,n\}_g
=\frac{\bigl(g^{ij}m_in_j\bigr)^2}
{\bigl(g^{ij}m_im_j\bigr)\bigl(g^{ij}n_in_j\bigr)} .
\]
若 $m=p^*=g_{ij}p^j$、$n=q^*=g_{ij}q^j$ 是 $p,q$ 的极线，则由 $g^{ij}m_in_j=g_{ij}p^iq^j$
等等式得 $\{p,q\}_g=\{m,n\}_g$。

\begin{lemma}\label{lem:invariant-J}
设直线 $m,n$ 分别与圆锥曲线 $g$ 交于 $a,b$ 与 $c,d$ 四个互不相同的点，$J_{ab,cd}$ 是这四点在
$g$ 上的交比，则
\begin{equation}
\{m,n\}_g=\left(\frac{J_{ab,cd}+1}{J_{ab,cd}-1}\right)^2 .
\label{eq:invariant-J}
\end{equation}
特别地，若 $p,q$ 是两点，$a,b$ 与 $c,d$ 分别是从 $p,q$ 引出的切线与 $g$ 的切点，则
$\{p,q\}_g$ 也由 \eqref{eq:invariant-J} 给出。
\end{lemma}

\begin{proof}
不妨设 $m=a\times b$、$n=c\times d$。由叉乘的恒等式
\[
(a\times b)\cdot(c\times d)=(a\cdot c)(b\cdot d)-(a\cdot d)(b\cdot c).
\]
又 $a,b$ 在 $g$ 上，$a^2=b^2=0$，故
\[
(a\times b)^2=a^2b^2-(a\cdot b)^2=-(a\cdot b)^2,
\]
同理 $(c\times d)^2=-(c\cdot d)^2$。于是
\[
\{m,n\}_g
=\frac{\bigl[(a\cdot c)(b\cdot d)-(a\cdot d)(b\cdot c)\bigr]^2}{(a\cdot b)^2(c\cdot d)^2}.
\]
由二次点列交比的公式与交比在交换下标时的变化规律，有
\[
(a\cdot c)(b\cdot d)=J^2\,(a\cdot d)(b\cdot c),\qquad
(a\cdot b)(c\cdot d)=(1-J)^2(a\cdot d)(b\cdot c),
\]
其中 $J=J_{ab,cd}$（第一式是交比平方公式，第二式来自 $J_{db,ca}=1-J_{ab,cd}$ 的平方）。代入得
\[
\{m,n\}_g
=\frac{\bigl[(J^2-1)(a\cdot d)(b\cdot c)\bigr]^2}
{\bigl[(1-J)^2(a\cdot d)(b\cdot c)\bigr]^2}
=\frac{(J+1)^2(J-1)^2}{(J-1)^4}
=\left(\frac{J+1}{J-1}\right)^2 .
\]
把上式用于极线 $m=p^*$、$n=q^*$，即得关于 $p,q$ 的结论。
\end{proof}

\begin{remark}
\eqref{eq:invariant-J} 的右端在 $J\mapsto1/J$ 下不变，因此给定 $\{p,q\}_g$ 只能把 $J$ 确定到
$J$ 与 $1/J$ 两种取法。具体取哪一个由几何连续性（例如圆时应取 $J_\theta=1$）确定。
\end{remark}

现在对椭圆 $g^{ij}=\diag\{a^2,b^2,-1\}$、$g_{ij}=\diag\{1/a^2,1/b^2,-1\}$ 逐一计算。由
\eqref{eq:zz}，
\[
z\cdot z=\frac{1}{a^2}-\frac{1}{b^2},\qquad
\bar z\cdot\bar z=\frac{1}{a^2}-\frac{1}{b^2},\qquad
z\cdot\bar z=\frac{1}{a^2}+\frac{1}{b^2},
\]
其中 $z\cdot\bar z=g_{ij}z^i\bar z^j$。因此
\begin{equation}
\{z,\bar z\}_g
=\left(\frac{1/a^2+1/b^2}{1/a^2-1/b^2}\right)^2
=\left(\frac{a^2+b^2}{c^2}\right)^2 .
\label{eq:inv-z}
\end{equation}
由 \eqref{eq:ellipse-X}，$f_{X1}=(-c,0,1)$、$f_{X2}=(c,0,1)$，于是
\[
f_{X1}\cdot f_{X2}=-\frac{c^2}{a^2}-1=-\frac{a^2+c^2}{a^2},\qquad
f_{X1}\cdot f_{X1}=f_{X2}\cdot f_{X2}=\frac{c^2}{a^2}-1=-\frac{b^2}{a^2},
\]
所以
\begin{equation}
\{f_{X1},f_{X2}\}_g
=\left(\frac{a^2+c^2}{b^2}\right)^2 .
\label{eq:inv-X}
\end{equation}
由 \eqref{eq:ellipse-Y}，$f_{Y1}=(0,\ii c,1)$、$f_{Y2}=(0,-\ii c,1)$，于是
\[
f_{Y1}\cdot f_{Y2}=\frac{c^2}{b^2}-1=\frac{c^2-b^2}{b^2},\qquad
f_{Y1}\cdot f_{Y1}=f_{Y2}\cdot f_{Y2}=-\frac{c^2}{b^2}-1=-\frac{b^2+c^2}{b^2},
\]
所以
\begin{equation}
\{f_{Y1},f_{Y2}\}_g
=\left(\frac{c^2-b^2}{c^2+b^2}\right)^2
=\left(\frac{b^2-c^2}{b^2+c^2}\right)^2 .
\label{eq:inv-Y}
\end{equation}

把 \eqref{eq:inv-z}--\eqref{eq:inv-Y} 代入 \eqref{eq:invariant-J}，并利用
$c^2=a^2-b^2$，解得相应的交比
\begin{equation}
J_\theta=\frac{b^2}{a^2}=1-e^2=\frac{\lambda_2}{\lambda_1},\qquad
J_X=e^2,\qquad
J_Y=1-e^{-2}.
\label{eq:J}
\end{equation}
例如第一个：\eqref{eq:inv-z} 给出 $\{z,\bar z\}_g=((a^2+b^2)/c^2)^2$，而
$(J+1)/(J-1)$ 在 $J=b^2/a^2$ 时等于 $(a^2+b^2)/(b^2-a^2)=-(a^2+b^2)/c^2$，平方后与
\eqref{eq:inv-z} 一致，故 $J_\theta=b^2/a^2=1-e^2$。另外两式同法可得，其中 $J_X=e^2$ 由
$(e^2+1)/(e^2-1)=-(a^2+c^2)/b^2$ 得到，$J_Y=1-e^{-2}$ 由
$(J_Y+1)/(J_Y-1)=(b^2-c^2)/(b^2+c^2)$ 得到。这里 $J_\theta$ 是过 $z,\bar z$ 的四条切线的
四个切点所对应的交比，$J_X$ 与 $J_Y$ 是把 $z,\bar z$ 分别换成 $f_{X1},f_{X2}$ 与
$f_{Y1},f_{Y2}$ 后所得的同类交比。于是离心率完全由这三个交比决定。

\subsection{特殊情形与连续变化}

由 \eqref{eq:J} 可以读出几个特殊位置。
\begin{itemize}
\item $J_\theta=-1$ 时四线成调和线束，此时 $1-e^2=-1$，$e=\sqrt2$，圆锥曲线是等轴双曲线。
\item $J_\theta=0$ 时成立 $b^2=0$，$e=1$。此时 $\{z,\bar z\}_g=1$ 所对应的四条切线中有两条
重合，有一个焦点落在圆锥曲线上，圆锥曲线与无穷远直线相切，即抛物线。
\item $e=0$ 时 $\lambda_1=\lambda_2$，圆锥曲线是圆。此时 $z,\bar z$ 都落在圆锥曲线上，从它们
引出的切线各退化为一条；四个有限焦点同时重合到圆锥曲线的中心，也就是无穷远直线的极点。
\end{itemize}
相应地，$J_X=e^2$ 在圆时为 $0$，在抛物线时为 $1$，在等轴双曲线时为 $2$；$J_Y=1-e^{-2}$ 在
圆时为无穷（即该对焦点的连线退化），在抛物线时为 $0$，在等轴双曲线时为 $1/2$。

椭圆的离心率变化时，六个焦点的位置可以连续地描述。设 $a>b>0$，四个有限焦点是
$(\pm c,0)$ 与 $(0,\pm\ii c)$，其中 $c^2=a^2-b^2$。当 $a$ 与 $b$ 逐渐接近时，$c$ 减小，
$x$ 轴上的两个实焦点与 $y$ 轴上的两个虚焦点同时向中心靠拢；当 $a=b$ 时 $c=0$，四个焦点重合
于中心，此时圆锥曲线是圆，而 $z,\bar z$ 落在圆锥曲线上。若继续压缩，使 $b>a$，则原来的两个
实焦点变为虚焦点、原来的两个虚焦点变为实焦点，四个焦点重新分开：此时长轴转到 $y$ 轴，实焦点
位于 $y$ 轴。整个过程中两个无穷远焦点 $z,\bar z$ 的位置不变。

\section{注：共焦点圆锥曲线族}

最后补充关于共焦点圆锥曲线族的一个观察。

\begin{remark}[共焦点族]
设 $g,h$ 是两条圆锥曲线。方程
\[
\det(g-\lambda h)=0
\]
的根给出与 $g,h$ 的焦点有关的射影数据；当 $g,h$ 共焦点时，这个方程出现重根。

另一方面，焦点由四条切线决定。设 $g$ 及其过 $z,\bar z$ 的四条切线给定，则所有与这四条切线都
相切的圆锥曲线 $h$ 都与 $g$ 有相同的焦点，它们构成一个单参数的圆锥曲线族，即共焦点圆锥曲线
族。因此共焦点族既可以用方程 $\det(g-\lambda h)=0$ 从代数上刻画，也可以用四条公切线的相切条件
从几何上刻画。这两种刻画互相等价，可作为处理共焦点圆锥曲线的一种方法。
\end{remark}
   % 焦点与准线
% 第十四节：曲线系、笛沙格对合定理与圆锥曲线的光学性质
\chapter{曲线系、笛沙格对合定理与圆锥曲线的光学性质}

本章讨论由两条圆锥曲线生成的曲线系。曲线系是处理共点、共切、共焦等问题的一条统一途径：
它把"经过某些固定点"或"切于某些固定直线"这类几何条件翻译为线性条件，
因而可以用线性代数的方法处理。作为应用，我们将证明笛沙格（Desargues）对合定理，
并由此推出共焦曲线系、圆锥曲线在焦点处的切线以及椭圆、双曲线的光学性质。

\section{二次点系与二次线系}

\begin{definition}[二次点系与二次线系]
设 $g,h$ 是两个不同的二次点（即两个不同的 $(2,0)$ 型对称张量）。所有形如
\[
g+\lambda h\qquad(\lambda\in\R)
\]
的二次点构成的集合，再添上 $h$ 本身（即 $\lambda=\infty$ 的极限情形），称为 $g,h$ 生成的
\textbf{二次点系}。

设 $g,h$ 是两个不同的二次线（即两个不同的 $(0,2)$ 型对称张量）。所有形如
\[
g+\lambda h\qquad(\lambda\in\R)
\]
的二次线构成的集合，再添上 $h$ 本身，称为 $g,h$ 生成的\textbf{二次线系}。

二次点系与二次线系统称\textbf{二次曲线系}。
\end{definition}

定义中两个元素都按比例视为同一个对象，故 $\lambda$ 取实数与 $\infty$ 恰好给出由 $g,h$
张成的二维线性空间中的全部元素。必须注意，只有具有上述形式的集合才叫曲线系；
任意一堆圆锥曲线放在一起不构成曲线系。

\subsection{两类曲线系的区别}

二次点系与二次线系是两个不同的概念，因为取逆不是线性运算：
\[
(g+\lambda h)^{-1}\neq g^{-1}+\lambda h^{-1}.
\]
同样，用 $g$ 升指标与用 $g,h$ 的线性组合升指标也不相同。两者的几何差异如下。

若 $m$ 是 $g,h$ 的公切线，即
\[
g^{ij}m_im_j=h^{ij}m_im_j=0 ,
\]
则对任意实数 $\lambda$ 有
\[
(g^{ij}+\lambda h^{ij})m_im_j=0 .
\]
由定义，$g^{ij}+\lambda h^{ij}$ 是二次点系元素的系数张量，因此该直线与点系中每一条曲线都相切。
故：\textbf{二次点系中的曲线具有公共切线}。

若 $a$ 是 $g,h$ 的公共点，即
\[
g_{ij}a^ia^j=h_{ij}a^ia^j=0 ,
\]
则对任意实数 $\lambda$ 有
\[
(g_{ij}+\lambda h_{ij})a^ia^j=0 ,
\]
即点 $a$ 在二次线系的每一条曲线上。故：\textbf{二次线系中的曲线经过公共点}。

这两种性质在几何上显然不同。用曲线的语言叙述就是：在二次点系中，若有两条曲线相切于同一条直线，
则点系中所有曲线都与这条直线相切；在二次线系中，若有两条曲线经过同一点，
则线系中所有曲线都经过这个点。

"切于固定直线"与"经过固定点"给出了曲线系的必要条件。反过来，它们也给出了充分条件，
这正是下一小节的内容。

\subsection{充分性}

\begin{proposition}[相切于四条定直线的二次点系]\label{prop:tangent-family}
设 $l_1,l_2,l_3,l_4$ 是平面上四条直线，其中任意两条不平行、任意三条不共点。
则与这四条直线都相切的全体二次点构成一个二次点系。
\end{proposition}

\begin{proof}
\textbf{第一步：相切条件的计算。} 二次点（即二次曲线的齐次方程）的系数构成六维线性空间
$W\cong\R^6$，于是二次曲线作为射影对象构成五维射影空间 $\mathbb{P}W$。设 $l$ 是一条直线，
把二次曲线限制到 $l$ 上得到 $l$ 上的一个二次型。曲线与 $l$ 相切，等价于这一限制有重根，
即其判别式为零；判别式是系数的二次多项式，但在齐次形式下它是系数的二次型。
因此"与 $l$ 相切"在 $\mathbb{P}W$ 中是一个二次超曲面，特别地它是 $W$ 中一个非零代数方程，
其解集是 $W$ 中的一个超曲面。

\textbf{第二步：至少是二维的。} 四条直线两两相交，记 $t_{ij}=l_i\times l_j$。
对两两配对的四条直线，任取一种配对方式，例如 $l_1,l_2$ 为一对、$l_3,l_4$ 为一对，
则退化二次点
\[
f_1=\sym{l_1}{l_2},\qquad f_2=\sym{l_3}{l_4}
\]
都是两条直线之并，因而与 $l_1,l_2,l_3,l_4$ 全部相切。它们线性无关，
故满足条件的二次点至少张成一个二维空间。

\textbf{第三步：不可能超过二维。} 由第二步，解集在 $W$ 中至少是二维的。
下面在一个具体模型上算出它恰为二维。

取
\[
l_1:\ x=0,\qquad l_2:\ y=0,\qquad l_3:\ z=0,\qquad l_4:\ x+y+z=0 .
\]
这四条直线两两不平行、任意三条不共点。设二次点为
\[
F=ax^2+by^2+cz^2+dxy+eyz+fzx .
\]
与 $l_3:z=0$ 相切，相当于 $ax^2+dxy+by^2=0$ 有重根，即
\[
d^2-4ab=0 .
\]
与 $l_2:y=0$ 相切，相当于 $ax^2+fzx+cz^2=0$ 有重根，即
\[
f^2-4ac=0 .
\]
与 $l_1:x=0$ 相切，相当于 $by^2+eyz+cz^2=0$ 有重根，即
\[
e^2-4bc=0 .
\]
最后与 $l_4:x+y+z=0$ 相切，相当于把 $x=-y-z$ 代入后所得二次型有重根。代入得
\[
F(-y-z,\;y,\;z)
=(a+b+c)y^2+2(a+b+d+e+f)\,yz+(a+b+c+2d+2e+2f)z^2 ,
\]
它有重根当且仅当
\[
(d+e+f)^2-4(a+b+c)(a+b+c+d+e+f)=0 .
\]
关键的一点是：由前三个条件，$d,e,f$ 的平方分别被 $a,b,c$ 决定，故它们的取值只有有限种组合；
第四个条件再从中筛出一条射影直线。为了看出这一点，作代换
\[
(a,b,c)=(u^2,v^2,w^2),
\]
则前三个条件化为
\[
d=\pm 2uv,\qquad e=\pm 2vw,\qquad f=\pm 2wu ,
\]
这是有限多种符号选取。对每一种符号选取，第四个条件成为 $(u,v,w)$ 的一个非平凡齐次方程，
它在射影直线 $(u:v:w)$ 上至多给出有限个点。
另一方面，由第二步，它的解集至少是一条射影直线（因为
$\sym{l_1}{l_3}$ 与 $\sym{l_2}{l_4}$ 所张成的二维空间中的元素都满足四个条件）。
有限集与整条射影直线只能相等，故解集恰为一条射影直线，
即满足四个相切条件的二次点在 $W$ 中构成二维线性子空间。

\textbf{第四步：二维子空间必为二次点系。} 已证满足四个相切条件的二次点构成二维线性子空间
$D\subset W$。取其中两个线性无关的元素 $G,H$（例如第二步中的 $f_1,f_2$），
则 $D=\{G+\lambda H:\lambda\in\R\}\cup\{H\}$，这正是二次点系的定义。
\end{proof}

\begin{proposition}[经过四个定点的二次线系]\label{prop:point-family}
设 $a,b,c,d$ 是平面上四点，其中任意三点不共线。则经过这四点的全体二次线构成一个二次线系。
\end{proposition}

\begin{proof}
\textbf{第一步：至少是二维的。} 注意"两条直线的乘积"给出的退化二次曲线是两个一次因子的乘积，
因此它经过两条直线上的所有点。于是退化二次线
\[
g=\sym{(a\times d)}{(b\times c)},\qquad
h=\sym{(a\times c)}{(d\times b)}
\]
都经过 $a,b,c,d$ 四点。

它们线性无关。为看出这一点，取坐标
\[
a=(1,0,0),\qquad b=(0,1,0),\qquad c=(0,0,1),\qquad d=(1,1,1),
\]
（三点不共线使这一取法可行）并把系数按 $(x^2,y^2,z^2,xy,yz,zx)$ 的顺序写成向量。由
\[
(a\times d)=(-1,0,1),\quad (b\times c)=(1,0,0),\quad
(a\times c)=(0,1,0),\quad (d\times b)=(-1,0,-1)
\]
得
\[
g\sim(0,0,0,0,0,2),\qquad
h\sim(0,2,0,0,2,2),
\]
这两个向量显然不成比例，故 $g,h$ 线性无关，满足条件的二次线至少张成一个二维空间。

\textbf{第二步：不可能超过二维。} 保持上面的坐标。设经过 $a,b,c,d$ 的二次线为
\[
G=\alpha x^2+\beta y^2+\gamma z^2+\delta xy+\eps yz+\zeta zx .
\]
"经过 $a$"给出 $\alpha=0$；"经过 $b$"给出 $\beta=0$；"经过 $c$"给出 $\gamma=0$；
"经过 $d$"给出
\[
(\alpha+\beta+\gamma)+(\delta+\eps+\zeta)=0 .
\]
于是约束条件为
\[
\alpha=\beta=\gamma=0,\qquad \delta+\eps+\zeta=0 .
\]
这四个条件中，前三个分别只涉及 $\alpha,\beta,\gamma$ 中的一个，第四个只涉及 $\delta,\eps,\zeta$，
故它们线性无关，约束矩阵的秩恰为 $4$。因此解空间维数为 $6-4=2$。

\textbf{第三步：二维解空间必为二次线系。} 由第二步，解空间 $D$ 是二维的；
由第一步，$g,h\in D$ 且线性无关，故
\[
D=\{g+\lambda h:\lambda\in\R\}\cup\{h\},
\]
即 $D$ 是由 $g,h$ 生成的二次线系。
\end{proof}

\begin{remark}
命题~\ref{prop:tangent-family} 与命题~\ref{prop:point-family} 互为对偶：
把点与线互换，四定直线变成四个对偶点，"相切于四条定直线"变成"经过四个对偶点"。
在应用中，"二次线系"比"二次点系"更常用。
\end{remark}

\subsection{二次线系的解析形式}

设两个圆锥曲线的方程分别为
\[
G(x,y)=0,\qquad H(x,y)=0 ,
\]
其中 $G,H$ 是二次多项式。当 $\lambda$ 取遍实数与 $\infty$ 时，方程
\[
G(x,y)+\lambda H(x,y)=0
\]
描述的正是 $g,h$ 生成的二次线系。反之，任给二次线系，
把其中一个元素作为 $G$、另一个元素关于同一坐标系的二次多项式作为 $H$，
就得到这种解析形式。因此在解析几何的语境中，求过若干定点的圆锥曲线的问题，
可以直接写成含参数 $\lambda$ 的方程再代入条件求解。

相切于四条定直线的二次点系在解析几何中不太直观，因此不常直接使用。
它最重要的应用是下一节的共焦曲线系。

\section{共焦曲线系}

回忆焦点与准线的射影定义：设 $g$ 是圆锥曲线，$\theta^{ij}$ 是正交张量，则
\[
\det\bigl(\theta^{ij}-\lambda g^{ij}\bigr)=0
\]
称为焦点方程。这是一个三次方程，$\lambda=0$ 是它的一个根（因为 $\theta$ 退化），
另外两个根记为 $\lambda_1,\lambda_2$。与 $\lambda_1,\lambda_2$ 对应的两个 $(2,0)$ 型退化二次点记为
\[
X^{ij}=\theta^{ij}-\lambda_1 g^{ij},\qquad Y^{ij}=\theta^{ij}-\lambda_2 g^{ij},
\]
它们分别可写成一对焦点的对称积；$\lambda=0$ 对应的 $\theta=(z\bar z)$ 也同样是一对复共轭点的对称积。
与 $X$ 对应的焦点记为 $f_1,f_2$，即
\[
X=\sym{f_1}{f_2},\qquad \theta=\sym{z}{\bar z}.
\]

\begin{proposition}[共焦曲线系]\label{prop:confocal}
设 $f_1,f_2$ 是两个焦点，$X=\sym{f_1}{f_2}$，$\theta=\sym{z}{\bar z}$ 是正交结构。
则以 $f_1,f_2$ 为焦点的全体圆锥曲线构成二次点系，称为\textbf{共焦曲线系}；
它的元素可以写成
\[
g^{ij}=X^{ij}+\lambda\theta^{ij},
\]
$\lambda$ 取遍实数与 $\infty$。其中用 $g$ 升降指标，$X,\theta$ 都是 $(2,0)$ 型张量。
\end{proposition}

\begin{proof}
只需说明共焦曲线系中所有曲线有公共切线。设直线 $m$ 同时经过 $f_1$ 与 $f_2$，
即 $m\cdot f_1=m\cdot f_2=0$。由 $X=\sym{f_1}{f_2}$ 得
\[
X^{ij}m_im_j=2(m\cdot f_1)(m\cdot f_2)=0 .
\]
又 $z,\bar z$ 是无穷远处的两个复共轭点，任何直线都经过无穷远直线，
故任何有限直线的线矢 $m$ 都满足
\[
\theta^{ij}m_im_j=2(m\cdot z)(m\cdot\bar z)=0 .
\]
于是对任意 $\lambda$ 有
\[
X^{ij}m_im_j+\lambda\theta^{ij}m_im_j=0 ,
\]
即 $m$ 与共焦曲线系中每一条曲线都相切。族中所有曲线都以过 $f_1,f_2$ 的直线为公共切线，
故它是二次点系；又 $X,\theta$ 线性无关，故 $\{X^{ij}+\lambda\theta^{ij}\}$ 是它的一般元素。
\end{proof}

\begin{remark}
$\theta=(z\bar z)$ 作为退化二次点，表示两条过 $z,\bar z$ 的直线之并，
这两条直线是无穷远处的两条"迷向直线"。因此 $\theta^{ij}m_im_j=0$ 对任意有限直线自动成立，
这一点在计算中要注意：不能因为 $\theta$ 退化为零就丢掉这一项。
另外，$\lambda$ 的具体数值依赖于 $X,\theta$ 的规范化选取，但曲线族
$\{X^{ij}+\lambda\theta^{ij}\}$ 本身与规范化无关。
\end{remark}

\section{例：椭圆上四点共圆的充要条件}

\begin{example}[椭圆上四点共圆的斜率条件]\label{ex:concyclic}
设 $\Gamma:\frac{x^2}{a^2}+\frac{y^2}{b^2}=1$ 是椭圆，$A,B,C,D$ 是它上面四个互不相同的点。
则 $A,B,C,D$ 共圆的充要条件是
\[
k_{AB}+k_{CD}=0 ,
\]
其中 $k_{AB},k_{CD}$ 分别是弦 $AB,CD$ 的斜率；当弦为竖直线时约定斜率为无穷，
相应的条件按极限理解。
\end{example}

\begin{proof}
先建立斜率条件的代数形式。用齐次坐标 $(x,y,z)$，点与线分别是三维点矢与线矢。
点 $a,b$ 的连线为 $m=a\times b$，其线矢写为 $m_i=(m_1,m_2,m_3)$，
则直线的斜率是 $k=-m_1/m_2$（当 $m_2=0$ 即竖直线时斜率为无穷）。

设 $x_\infty=(1,0,0)$ 与 $y_\infty=(0,1,0)$ 分别是 $x$ 轴方向与 $y$ 轴方向的无穷远点，
$\sym{x_\infty}{y_\infty}$ 是它们的对称积。由
\[
\bigl(\sym{x_\infty}{y_\infty}\bigr)^{ij}m_in_j
=x_\infty^im_in_jy_\infty^j+x_\infty^jm_jn_iy_\infty^i
=m_1n_2+m_2n_1
\]
可知，直线 $m,n$ 满足
\[
\bigl(\sym{x_\infty}{y_\infty}\bigr)^{ij}m_in_j=0
\]
当且仅当 $m_1n_2+m_2n_1=0$。把 $m_1=-m_2k_{AB}$、$n_1=-n_2k_{CD}$ 代入，
即得 $k_{AB}+k_{CD}=0$。因此斜率条件等价于
\begin{equation}\label{eq:concyclic}
\bigl(\sym{x_\infty}{y_\infty}\bigr)^{ij}(a\times b)_i(c\times d)_j=0 .
\end{equation}

再计算椭圆与圆在这两项上的取值。椭圆的系数张量为
$g_{ij}=\diag\{1/a^2,\,1/b^2,\,-1\}$，于是
\[
g_{ij}x_\infty^iy_\infty^j=g_{12}=0 .
\]
半径为 $r$、中心为 $(p,q)$ 的圆方程为
\[
(x-p)^2+(y-q)^2-r^2=0 ,
\]
其二次部分的系数矩阵为 $\diag\{1,1,0\}$，故同样有
\[
h_{ij}x_\infty^iy_\infty^j=h_{12}=0 .
\]
也就是说，椭圆与圆都满足式~\eqref{eq:concyclic} 所描述的条件
（等价地，$g$ 与 $h$ 都经过 $x_\infty,y_\infty$ 这一对对偶点，即都切于直线
$\sym{x_\infty}{y_\infty}$ 所代表的两条"沿坐标方向的退化直线"）。

\textbf{充分性。} 设 $A,B,C,D$ 在 $\Gamma$ 上且共圆，圆记为 $\omega$，其系数张量为 $h$。
由命题~\ref{prop:point-family}，过 $A,B,C,D$ 的全体二次线构成由 $g,h$ 生成的二次线系。
弦 $AB,CD$ 的并是退化的二次曲线，它经过 $A,B,C,D$，故其系数张量可以写成
\[
f=g+\lambda h
\]
（$\lambda$ 取实数或 $\infty$；$f$ 与弦对的系数张量至多相差一个非零常数因子，
这不影响下面的齐次方程）。于是
\[
f_{ij}x_\infty^iy_\infty^j
=g_{ij}x_\infty^iy_\infty^j+\lambda h_{ij}x_\infty^iy_\infty^j=0 .
\]
又 $f$ 是两条直线 $AB,CD$ 之并的系数张量，故 $(AB)$ 与 $(CD)$ 满足式~\eqref{eq:concyclic}，
于是 $k_{AB}+k_{CD}=0$。

\textbf{必要性。} 设 $A,B,C,D$ 在 $\Gamma$ 上且 $k_{AB}+k_{CD}=0$。由于四点都在 $\Gamma$ 上，
过这四点的曲线中有一个标准的圆：任取其中三点所定的圆 $\omega$，由圆的方程知
$h_{ij}x_\infty^iy_\infty^j=0$，又 $\omega$ 与 $\Gamma$ 在 $A,B,C,D$ 四点相交，
故 $\omega$ 也在过这四点的二次线系中，即 $h=g+\lambda h_0$ 型。于是与充分性同样的计算表明
$f_{ij}x_\infty^iy_\infty^j=0$，反过来给出 $k_{AB}+k_{CD}=0$。
更直接地：由 $k_{AB}+k_{CD}=0$，弦对 $(AB)(CD)$ 满足式~\eqref{eq:concyclic}；
过 $A,B,C$ 的圆 $\omega$ 也满足该式；两者都在过 $A,B,C$ 的二维线性空间中，
而在四点 $\omega\cap\Gamma$ 处取值的线性泛函
$f\mapsto f_{ij}x_\infty^iy_\infty^j$ 在它们上同时为零，
故在它们的公共点 $D$ 上，$\omega$ 与弦对 $CD$ 满足同一线性关系，从而 $D$ 也在 $\omega$ 上。
\end{proof}

\begin{remark}
式~\eqref{eq:concyclic} 的几何含义是：$x_\infty,y_\infty$ 两点关于直线 $AB,CD$ 成调和点列。
这一点可以用交比直接验证：四条直线 $AB,CD$ 与过 $x_\infty,y_\infty$ 的两条直线的交比等于
$\dfrac{k_{AB}k_{CD}-1}{k_{AB}k_{CD}+1}$，当且仅当 $k_{AB}+k_{CD}=0$ 时它等于 $-1$。
因此"斜率之和为零"与"无穷远方向的对称点关于两弦调和"是同一件事。
这里使用的是两个实无穷远点 $x_\infty,y_\infty$；由于椭圆与圆在坐标轴方向上都没有切点，
$g$ 与 $h$ 在 $x_\infty,y_\infty$ 上的取值均为零，这正是所需的条件。
\end{remark}

\section{笛沙格对合定理}

\subsection{定理的陈述与改写}

\begin{theorem}[笛沙格对合定理]\label{thm:desargues}
设 $C,D,E,F$ 是四个点（任意三点不共线），$l$ 是一条定直线。
考虑经过 $C,D,E,F$ 四点的全体圆锥曲线；它们与 $l$ 交于一对动点。
则 $l$ 上存在一个对合变换 $f$，使得对每一条这样的圆锥曲线，
它与 $l$ 的两个交点都关于 $f$ 互相对合。若某条圆锥曲线与 $l$ 相切，
则切点是这个对合的不动点。
\end{theorem}

\textbf{证明的思路。} 第一步把几何命题代数化：由命题~\ref{prop:point-family}，
过四个定点的全体圆锥曲线正是二次线系 $\{g+\lambda h\}$，其中 $g,h$ 是两条退化的二次线
（四点两两连线之并）。这样就把待定元素从"四个点加一条动曲线"减少到"两个张量加一个参数"。
第二步把二次线系限制到定直线 $l$ 上。限制映射是线性的，且圆锥曲线在 $l$ 上的限制正是它与 $l$
两个交点的对称积。于是平面上的问题化为二维空间上的一个命题。
第三步利用限制映射的线性性：线系在 $l$ 上的像仍是二维线性族，而二维空间中的对合条件
是关于参数 $\lambda$ 的线性等式，因此无数个等式等价于两个等式，解出它们即可。

\subsection{限制到定直线}

设二次线系为 $g+\lambda h$。关于限制映射，回忆第九节：三维空间中 $(0,2)$ 型对称张量可以限制为
子空间上的 $(0,2)$ 型对称张量，且这一映射是线性的。记 $g,h$ 限制到 $l$ 上的像分别为 $g',h'$，
则由线性性，
\[
(g+\lambda h)\big|_l=g'+\lambda h'.
\]
设这条圆锥曲线与 $l$ 交于两点 $a,b$。圆锥曲线在 $l$ 上的限制是其与 $l$ 的两个交点的对称积，
即
\[
g'+\lambda h'=a_ib_j+a_jb_i=\sym{a}{b},
\]
其中 $a_i=e_{ij}a^j$、$b_i=e_{ij}b^j$，$e$ 是直线 $l$ 上的反对称张量（二维空间中的 $\eps$）。
因此笛沙格对合定理等价于下面的二维命题，其中 $g',h'$ 都是二维矢量空间上的 $(0,2)$ 型对称张量，
$g'+\lambda h'$ 表示它们的线性组合。

\subsection{二维命题}

\begin{proposition}[二维对合命题]\label{prop:2d-involution}
设 $g',h'$ 是二维实矢量空间上两个线性无关的 $(0,2)$ 型对称张量。
设点对 $a,b$ 使得存在实数 $\lambda$ 满足
\[
g'+\lambda h'=\sym{a}{b} ,
\]
则 $a,b$ 满足一个对合关系，即存在一个与 $\lambda$ 无关的对合变换 $f$，
使得每一对满足上式（对不同 $\lambda$）的 $a,b$ 都关于 $f$ 互相对合。
\end{proposition}

\begin{proof}
\textbf{第一步：化为标准坐标。} 先设 $g'$ 非退化。由 Sylvester 定理，
可以取 $l$ 上的坐标使
\[
g'=\begin{pmatrix}0&1\\1&0\end{pmatrix},
\]
即 $g'(x,y)=x_1y_2+x_2y_1$。之所以能这样取，是因为 $g'$ 的作用只差一个可逆线性变换：
任何非退化对称 $(0,2)$ 型张量都可以通过可逆线性变换化为上述形式。
在新的坐标下，$h'$ 仍是对称张量，设
\[
h'=\begin{pmatrix}\alpha&\delta\\\delta&\gamma\end{pmatrix}.
\]

把 $l$ 上的点写成齐次坐标 $(1,t)$，即 $x=t e_1+e_2$（$t=\infty$ 对应 $e_1$）。
记 $a$ 对应参数 $t$、$b$ 对应参数 $s$，计算
\begin{align*}
C_\lambda(s,t)
&=(g'+\lambda h')\bigl((1,s),(1,t)\bigr)\\
&=g'\bigl((1,s),(1,t)\bigr)+\lambda h'\bigl((1,s),(1,t)\bigr)\\
&=2(s+t)+\lambda\bigl(\alpha+\gamma st+\delta(s+t)\bigr).
\end{align*}
条件 $g'+\lambda h'=\sym{a}{b}$ 意味着对合点对 $a,b$ 满足 $C_\lambda(s,t)=0$。
（对称性可见 $C_\lambda(s,t)$ 关于 $s,t$ 对称，与 $\sym{a}{b}$ 的对称性一致。）

由 $C_\lambda=0$ 解出
\begin{equation}\label{eq:lambda}
\lambda=-\frac{2(s+t)}{\alpha+\gamma st+\delta(s+t)} .
\end{equation}

\textbf{第二步：消去 $\lambda$。} 由第一步，$\det(g'+\lambda h')$ 是 $\lambda$ 的二次多项式，
直接计算得
\[
\det(g'+\lambda h')
=1+\lambda(\alpha+\gamma)+\lambda^2\bigl(\alpha\gamma-4\delta^2\bigr).
\]
而 $g'+\lambda h'=\sym{a}{b}$ 是退化张量，故
\begin{equation}\label{eq:detzero}
\det(g'+\lambda h')=0 .
\end{equation}
把式~\eqref{eq:lambda} 代入式~\eqref{eq:detzero}，并乘以分母的平方，得到
\[
P+Q(s+t)+R\,st=0 ,
\]
其中
\[
P=4\delta^2+2\gamma(\alpha-\gamma),\qquad
Q=-2\delta(\alpha+\gamma),\qquad
R=2\alpha(\alpha-\gamma).
\]
推导要点如下：式~\eqref{eq:detzero} 等价于
\[
(\alpha+\gamma st+\delta(s+t))^2
-2(s+t)(\alpha+\gamma)(\alpha+\gamma st+\delta(s+t))
+4(s+t)^2(\alpha\gamma-4\delta^2)=0 .
\]
把左端按 $1,s+t,st$ 展开，$st$ 的二次项与 $(s+t)^2$ 相消，
$(s+t)^2st$ 的二次项也相消，剩下的恰为 $P+Q(s+t)+R\,st$。

\textbf{第三步：这一步定义一个对合。} 上式是关于 $s,t$ 的对称双线性方程，
对每个 $t$ 关于 $s$ 是一次方程，其解
\[
s=-\frac{P+Qt}{Rt+Q}
\]
给出一个分式线性变换，且由对称性，它的逆是它自己，故它是对合变换。
于是对每一对满足 $g'+\lambda h'=\sym{a}{b}$ 的 $a,b$，它们都满足这同一个对合关系。

\textbf{第四步：如 $g'$ 退化。} 此时限制映射把 $g'$ 变为秩一的张量，
式~\eqref{eq:detzero} 对 $\lambda$ 最多有一个解，对应的对合有一个不动点（切点），
结论仍成立；也可以由非退化情形取极限得到。
\end{proof}

\subsection{对合的张量形式}

上面的对合可以用一个 $(2,0)$ 型张量 $f$ 简洁地刻划。

\begin{proposition}[对合张量]\label{prop:involution-tensor}
设 $g',h'$ 是二维矢量空间上两个线性无关的 $(0,2)$ 型对称张量，$e^{ij}$ 是二维反对称张量。
则存在 $(2,0)$ 型张量 $f$，使得
\[
f^{ij}a_ib_j=0
\]
对每一对满足 $g'+\lambda h'=\sym{a}{b}$ 的点 $a,b$ 都成立；
这样的 $f$ 由 $f:g'=f:h'=0$ 确定（相差一个非零常数因子），
并且具体可以取
\[
f_{il}=g'_{ij}e^{jk}h'_{kl}+g'_{lj}e^{jk}h'_{ki}
\]
（即 $f=(g'\cdot e\cdot h')$，小括号表示对称化）。
\end{proposition}

\begin{proof}
\textbf{第一步：两个条件的等价性。} 记 $a,b$ 对应的参数为 $t,s$。
"$a,b$ 关于 $f$ 互相对合"的含义是 $f^{ij}a_ib_j=0$，即
\[
f:\sym{a}{b}=0 .
\]
把 $\sym{a}{b}=g'+\lambda h'$ 代入，得
\begin{equation}\label{eq:f-lambda}
f:g'+\lambda\,f:h'=0 .
\end{equation}
式~\eqref{eq:f-lambda} 对满足条件的每一对 $a,b$（对应不同的 $\lambda$）都成立。
当 $\lambda$ 取两个不同的实数时，即得
\begin{equation}\label{eq:twocond}
f:g'=0,\qquad f:h'=0 .
\end{equation}
反过来，若式~\eqref{eq:twocond} 成立，则由式~\eqref{eq:f-lambda} 对任意 $\lambda$ 成立，
故所有对合点对 $a,b$ 都满足 $f^{ij}a_ib_j=0$。
因此对合由式~\eqref{eq:twocond} 确定。这是关于 $f$ 的三个分量的两个齐次线性方程，
解空间至少一维。

\textbf{第二步：验证所给的 $f$ 满足两个条件。} 记 $f^{(1)}_{il}=g'_{ij}e^{jk}h'_{kl}$，
$f^{(2)}_{il}=g'_{lj}e^{jk}h'_{ki}$。在这个写法下，$e$ 作用于中间指标，
对称化以后 $f$ 对称。计算 $f:g'$：
\begin{align*}
f:g'
&=f^{(1)}_{il}g'^{il}+f^{(2)}_{il}g'^{il}\\
&=g'_{ij}e^{jk}h'_{kl}g'^{il}+g'_{lj}e^{jk}h'_{ki}g'^{il} .
\end{align*}
对第一项，把与 $g'$ 有关的部分缩并（等价地，把 $e$ 写在后两项之间并对 $g'$ 缩并）：
\[
g'_{lj}e^{jk}g'_{ik}=g'_{lj}e^{jk}g'_{ki}.
\]
交换哑指标 $j,k$ 并利用 $g'_{ki}=g'_{ik}$、$e^{jk}=-e^{kj}$ 得
$g'_{lj}e^{jk}g'_{ki}=-g'_{lk}e^{kj}g'_{ji}=-g'_{lj}e^{jk}g'_{ki}$，故这一项为零。
于是
\[
f:g'=g'_{ij}e^{jk}h'_{kl}g'^{il}
=\bigl(g'_{ij}e^{jk}g'_{il}\bigr)h'_{kl}.
\]
利用二维恒等式
\[
g'_{ij}e^{jk}g'_{il}=\det(g')\,e_{il},
\]
得到
\[
f:g'=\det(g')\,e_{il}h'^{il}=0 ,
\]
最后一步用到 $e$ 反对称而 $h'$ 对称：$e_{il}h'^{il}=0$。
完全类似地交换 $g'$ 与 $h'$ 的角色得 $f:h'=0$。这正是式~\eqref{eq:twocond}。

\textbf{第三步：在标准坐标下求出 $f$。} 取坐标使
$g'=\begin{pmatrix}0&1\\1&0\end{pmatrix}$ 且
$h'=\begin{pmatrix}\alpha&0\\0&\gamma\end{pmatrix}$（由第二步中 $f:g'=f:h'=0$
以及 $e_{12}=1$，可以直接算出 $f$ 的形式）。此时
\[
e^{jk}=\begin{pmatrix}0&1\\-1&0\end{pmatrix},\qquad
g'e=\begin{pmatrix}-1&0\\0&1\end{pmatrix},
\]
于是
\[
f=g'e h'+h'e g'
=\begin{pmatrix}-\alpha&0\\0&\gamma\end{pmatrix}
+\begin{pmatrix}-\alpha&0\\0&\gamma\end{pmatrix}
=-2\alpha\begin{pmatrix}1&0\\0&0\end{pmatrix}
+2\gamma\begin{pmatrix}0&0\\0&1\end{pmatrix}.
\]
这与命题陈述中的一般表达式一致（相差一个常数因子），特别地它与
\[
\begin{pmatrix}1&-1\\-1&1\end{pmatrix}
\]
成比例当且仅当 $\alpha=\gamma$ 的倍数关系；在第四步中我们会看到这一点。
更直接地，把 $f_{il}$ 的分量写开：$f_{11}=0$，$f_{22}=0$，而
\[
f_{12}=f_{21}=g'_{1j}e^{jk}h'_{k2}+g'_{2j}e^{jk}h'_{k1}=\gamma-\alpha,
\]
故
\[
f=(\gamma-\alpha)\begin{pmatrix}0&1\\1&0\end{pmatrix}.
\]
这说明 $f$ 也可以写成 $\begin{pmatrix}1&-1\\-1&1\end{pmatrix}$ 的倍数：
注意在坐标变换
\[
u=\frac{x+y}{2},\qquad v=\frac{x-y}{2}
\]
下，$g'=2uv$，而 $f$ 的上述形式恰对应于该坐标系中的 $\begin{pmatrix}1&-1\\-1&1\end{pmatrix}$ 的倍数。
\end{proof}

\begin{remark}
命题~\ref{prop:involution-tensor} 的第二、三两步给出了两种等价的表达：
在标准坐标下，$f$ 是 $\begin{pmatrix}1&-1\\-1&1\end{pmatrix}$ 的倍数（当 $\alpha\neq\gamma$ 时），
而在一般的坐标下它由 $f=g'e h'+h'e g'$ 给出。
由 $e_{ij}e^{ij}=2$ 与 $e_{ij}e^{ik}=\delta^k_j$ 可以验证
$e^{ij}$、$e_{ij}$ 的作用恰好把 $\sym{a}{b}$ 变成与 $a,b$ 垂直的方向，
这正是对合的几何意义。
\end{remark}

\subsection{证明结构小结}

定理~\ref{thm:desargues} 的证明用到了三步。

\begin{enumerate}
\item \textbf{几何命题的代数化。} 用曲线系的性质，把"过四点的圆锥曲线族"改写为二次线系
$g+\lambda h$，把几何图形的待定元素转化为两个张量与一个参数。
\item \textbf{限制到直线。} 用限制映射把平面上的二次线系转化为直线上的二维"线系"
$g'+\lambda h'$，消去了与该直线无关的几何元素。
\item \textbf{用线性性把无穷多个等式化为两个等式。} 直线上的对合条件是关于 $\lambda$ 的线性等式，
它要对所有 $\lambda$ 成立，等价于 $f:g'=f:h'=0$ 两个等式；解出 $f$ 就得到了对合。
\end{enumerate}

\subsection{推论}

\begin{corollary}[直线上的对合]\label{cor:line-involution}
设 $c,d,s,t$ 是直线 $l$ 上的四点。则
\begin{enumerate}
\item 存在 $l$ 上的一个对合变换 $f$，使得 $c,d$ 互相对合且 $s,t$ 互相对合；
\item 对 $l$ 上任意两点 $a,b$，它们关于 $f$ 互相对合，当且仅当存在实数（或 $\infty$）$\lambda$ 使
\[
\sym{a}{b}=\sym{c}{d}+\lambda\sym{s}{t} ;
\]
\item 特别地，设 $p,q$ 是 $l$ 上两点，$f$ 是以它们为不动点的对合，则对 $l$ 上任意两点 $a,b$，
它们关于 $f$ 对合，当且仅当存在 $\lambda$ 使
\[
\sym{a}{b}=\sym{p}{p}+\lambda\sym{q}{q} .
\]
\end{enumerate}
\end{corollary}

\begin{proof}
(1) 由命题~\ref{prop:2d-involution}，$c,d$ 满足一个对合；由同一命题，
存在另一个对合使 $s,t$ 对合。但直线上的对合构成的族是一维的，
而由命题~\ref{prop:2d-involution} 的证明，对合由 $f:g'=f:h'=0$ 确定，
其中 $g'=\sym{c}{d}$、$h'=\sym{s}{t}$。因此这两个条件确定的对合是同一个。
具体地，$(cd)$ 与 $(st)$ 张成的线性族中每个元素都是一个对合张量，
而 $\sym{c}{d}$ 与 $\sym{s}{t}$ 的任意线性组合都同时满足
$f:\sym{c}{d}=0$ 与 $f:\sym{s}{t}=0$ 中的各自一个，故族中元素对应的对合都包含 $c,d$ 与 $s,t$。

(2) 设 $f$ 是 (1) 中的对合，其张量记为 $F$。由 $F$ 满足
$F:\sym{c}{d}=F:\sym{s}{t}=0$ 知 $F$ 与 $\sym{c}{d}$、$\sym{s}{t}$ 都正交。
反之，任何与它们都正交的 $F$ 都给出同一个对合。现在
$\sym{a}{b}$ 关于 $f$ 对合，即 $\sym{a}{b}$ 与 $F$ 正交。
由二维空间的维数，正交于 $F$ 的空间是二维的；而 $\sym{c}{d}$ 与 $\sym{s}{t}$
也在这个空间中且线性无关，故
\[
\sym{a}{b}=\lambda_1\sym{c}{d}+\lambda_2\sym{s}{t}.
\]
两边同时除以 $\lambda_2$（若 $\lambda_2=0$ 则对应 $\lambda=\infty$）即得所需形式。

(3) 以 $p,q$ 为不动点的对合，其张量是 $\sym{p}{q}$；而 $\sym{p}{q}$ 与 $\sym{p}{p}$、$\sym{q}{q}$
都正交。于是与 (2) 同样的论证给出 $\sym{a}{b}$ 落在 $\sym{p}{p}$ 与 $\sym{q}{q}$ 张成的空间中。
\end{proof}

\section{对偶命题与共焦曲线系的结论}

\subsection{对偶命题}

把定理~\ref{thm:desargues} 中的点与线互换，得到它的对偶命题。

\begin{proposition}[对偶命题]\label{prop:dual-desargues}
设 $l_1,l_2,l_3,l_4$ 是四条定直线，$A$ 是一个定点。
则以 $l_1,l_2,l_3,l_4$ 为公共切线的全体圆锥曲线（即一个二次点系），
从 $A$ 引出的两条动切线关于一个对合变换互相对合。
若这条圆锥曲线经过 $A$，则此时的两条切线重合，给出该对合的不动点。
\end{proposition}

\begin{proof}
命题~\ref{prop:tangent-family} 已证与四条定直线相切的圆锥曲线构成二次点系。
把定理~\ref{thm:desargues} 的证明中的点矢与线矢互换、限制映射换成对偶的限制，
即得本命题。具体地，二次点系在点 $A$ 处的"取值"（即过 $A$ 的切线）给出对偶空间中的二维限制，
由命题~\ref{prop:2d-involution} 得到对合；曲线经过 $A$ 时，两条切线重合为一条，
它是对合的不动点。
\end{proof}

\begin{proposition}[二次点系的切线对合]\label{prop:dual-tangent}
定点到二次点系的两条动切线是对合的；若点系中某条圆锥曲线经过该定点，
则该点处的切线是对合的不动点。
\end{proposition}

命题~\ref{prop:dual-tangent} 只是把命题~\ref{prop:dual-desargues} 中的四条定直线
换成二次点系的一般表述：二次点系由两条曲线生成，因而它的所有元素都切于四条定直线，
命题~\ref{prop:dual-desargues} 即可应用。

\subsection{共焦曲线系}

\begin{corollary}[共焦曲线系的切线对合]\label{cor:confocal-tangent}
定点到共焦曲线系的两条动切线是对合的。若共焦曲线系中某条圆锥曲线经过该定点，
则该点处的切线是对合的不动点。
\end{corollary}

共焦曲线系是二次点系的重要特例，故本推论直接来自命题~\ref{prop:dual-tangent}。

\subsection{过定点的共焦曲线与切线垂直}

现在取定一个定点 $A$ 与一对焦点 $F_1,F_2$，考虑以 $F_1,F_2$ 为焦点的共焦曲线系
\[
g=X+\lambda\theta,\qquad X=\sym{F_1}{F_2},\quad \theta=\sym{z}{\bar z}.
\]

\begin{proposition}[过定点的两条共焦曲线]\label{prop:two-conics}
设 $A$ 是平面上一点（不在直线 $F_1F_2$ 上，且不满足退化情形）。
则以 $F_1,F_2$ 为焦点的圆锥曲线中，恰有一条椭圆和一条双曲线经过 $A$，
且它们在 $A$ 点的两条切线互相垂直。
\end{proposition}

\begin{proof}
\textbf{第一步：恰有两条曲线经过 $A$。} 由推论~\ref{cor:confocal-tangent}，
从 $A$ 到共焦曲线系的动切线构成一个对合变换，而对合必有两个不动点。
不动点对应的正是经过 $A$ 的曲线在该点的切线（两条切线重合），
故恰有两条曲线经过 $A$（计重数；若不退化则两条不同）。

\textbf{第二步：两条曲线在 $A$ 点的切线互相垂直。}
设过 $A$ 的两条曲线在 $A$ 点的切线分别为 $p,q$，它们是对合的不动点。
对任意 $\lambda$，设从 $A$ 到曲线 $g_\lambda=X+\lambda\theta$ 的两条切线为 $m,n$，
则 $m,n$ 互相对合。而对合的两点被对合的不动点调和分割，故 $p,q,m,n$ 成调和线束。
这一条件等价于
\begin{equation}\label{eq:conjugate}
g_\lambda:(mn)=0
\end{equation}
（即 $m,n$ 关于曲线 $g_\lambda$ 共轭）。

现在取 $\lambda=0$，即 $g_0=X=\sym{F_1}{F_2}$。它是两条直线的并，
这两条直线分别过 $F_1$ 与 $F_2$，故从 $A$ 到它的两条切线就是 $AF_1$ 与 $AF_2$。
代入式~\eqref{eq:conjugate}，得 $X:(pq)=0$。
再取 $\lambda=\infty$，即 $g_\infty=\theta$。曲线 $\theta$ 的两条切线是两条迷向直线，
它们互相垂直，且都经过 $A$（因为 $\theta$ 是退化二次点，即无穷远处的一对复共轭点，
从 $A$ 引出的两条切线就是过 $A$ 的两条迷向直线）。由正交张量的定义，
$m,n$ 垂直的条件为 $\theta^{ij}m_in_j=0$，即 $\theta:(mn)=0$。于是
\[
\theta:(pq)=0 .
\]
结合 $X:(pq)=0$ 得
\[
g_\lambda:(pq)=\bigl(X+\lambda\theta\bigr):(pq)=0
\]
对一切 $\lambda$ 成立，特别地对经过 $A$ 的这两条曲线成立。
特别地取其中的椭圆，设它在 $A$ 点的切线为 $p,q$，则
$X:(pq)=\theta:(pq)=0$。

\textbf{第三步：由 $X:(pq)=0$ 与 $\theta:(pq)=0$ 读出垂直。}
把点 $A$ 处的线束看作二维点空间：过 $A$ 的直线用它与无穷远直线的交点表示，
于是 $p,q$ 是这一二维空间中的两个点，而 $X,\theta$ 限制到该空间上成为对称 $(0,2)$ 型张量。
条件
\[
X:(pq)=0,\qquad \theta:(pq)=0
\]
说明 $p,q$ 是 $X$ 与 $\theta$ 在 $A$ 处限制出的两条二次曲线的公共的一对对合点。

由 $X=\sym{F_1}{F_2}$，把 $X$ 限制到 $A$ 处的线束上，$F_1,F_2$ 给出两条固定直线
$AF_1,AF_2$ 的方向点，而 $X|_A$ 正是这两个方向点的对称积，即
\[
X|_A=\sym{AF_1}{AF_2}.
\]
于是 $X:(pq)=0$ 就是 $X|_A:(pq)=0$，即 $p,q$ 与 $AF_1,AF_2$ 成调和线束；
而由 $\theta$ 的定义，$\theta^{ij}p_iq_j=0$ 正是 $p\perp q$。
因此 $p,q$ 是一对互相垂直的直线，且与 $AF_1,AF_2$ 成调和线束。
当 $p,q$ 与 $AF_1,AF_2$ 不重合时，把 $(pq)$ 的两个因子记为 $P_1,P_2$，上述两式写开就是
\[
(AF_1\cdot P_1)(AF_2\cdot P_2)+(AF_1\cdot P_2)(AF_2\cdot P_1)=0,
\qquad \theta^{ij}P_{1i}P_{2j}=0 .
\]
第一式是说 $p,q$ 与 $AF_1,AF_2$ 调和，第二式是说 $p\perp q$；
其余情形由连续性得到。

\textbf{第四步：一条椭圆一条双曲线。} 共焦曲线系 $X+\lambda\theta$ 中，
$\lambda=0$ 给出退化的点对 $(f_1f_2)$，
$\lambda$ 从 $0$ 增大到 $\infty$ 时，曲线的二次型系数从退化变为退化，
其号差在中间必然经过两种情形：一种为定号（椭圆），一种为不定号（双曲线）。
过 $A$ 的两条曲线分别落在号差变化的两侧，故一条是椭圆、一条是双曲线。
\end{proof}

\begin{remark}
命题~\ref{prop:two-conics} 的第二步中用了"对合的两点被对合的不动点调和分割"这一事实。
它的代数形式是：若对合由张量 $F$ 给出，不动点为 $p,q$，则对任意一对对合点 $m,n$，
$p,q,m,n$ 为调和线束，等价于 $F:(mn)=0$；
而 $F$ 可取为 $X$ 与 $\theta$ 的线性组合，故 $X:(mn)=\theta:(mn)=0$ 同时成立。
特别地，$\theta:(mn)=0$ 就是 $m\perp n$。
\end{remark}

\subsection{圆锥曲线在焦点处的光学性质}

下面把上述结论翻译成关于焦点与法线的欧氏几何命题。
先叙述一个经典的欧氏几何事实。

\begin{proposition}[垂直平分线对的调和判定]\label{prop:bisector}
设 $p,q,r,s$ 是过同一点的四条直线，且 $p\perp q$。
则 $p,q,r,s$ 成调和线束，当且仅当 $p,q$ 都是直线 $r,s$ 的角平分线。
\end{proposition}

\begin{proof}
\textbf{充分性。} 不妨设 $p,q$ 为坐标轴。取 $p$ 为 $x$ 轴、$q$ 为 $y$ 轴，
设 $r,s$ 的倾角分别为 $\alpha$ 与 $-\alpha$（即 $r,s$ 关于 $x$ 轴对称，$p,q$ 是它们的角平分线）。
用斜率表示四条直线：$p$ 对应参数 $0$，$q$ 对应参数 $\infty$，
$r,s$ 分别对应斜率 $\tan\alpha$ 与 $-\tan\alpha$。线束上的交比为
\[
(p,q;r,s)
=\frac{\tan\alpha-0}{\tan\alpha-\infty}\cdot
\frac{-\tan\alpha-\infty}{-\tan\alpha-0}
=\frac{\tan\alpha}{1}\cdot\frac{1}{-\tan\alpha}=-1 ,
\]
故 $p,q,r,s$ 成调和线束。

\textbf{必要性。} 设 $p\perp q$ 且 $p,q,r,s$ 成调和线束。不妨设 $p$ 为 $x$ 轴（参数 $0$），
$q$ 为 $y$ 轴（参数 $\infty$），$r$ 的斜率为 $k\neq0$，$s$ 的斜率为 $k'$。
计算交比
\[
(p,q;r,s)=\frac{k-0}{k-\infty}\cdot\frac{k'-\infty}{k'-0}=-\frac{k}{k'} .
\]
由 $(p,q;r,s)=-1$ 得 $k'=-k$。于是 $r\) 与 $s$ 关于 $x$ 轴对称，
而 $q$ 是 $y$ 轴，故 $p,q$ 都是 $r,s$ 的角平分线。充分性中的计算是可逆的，
故两者等价。
\end{proof}

\begin{theorem}[椭圆与双曲线的光学性质]\label{thm:optical}
设 $A$ 是以 $F_1,F_2$ 为焦点的圆锥曲线上的一点。
\begin{enumerate}
\item 若 $A$ 在椭圆上，则 $A$ 点处的法线平分 $\angle F_1AF_2$；
\item 若 $A$ 在双曲线上，则 $A$ 点处的法线平分 $\angle F_1AF_2$ 的补角。
\end{enumerate}
\end{theorem}

\begin{proof}
设 $A$ 处的切线为 $m$，法线为 $n$，则 $m\perp n$。
由于 $A$ 在曲线上，$A$ 到曲线的两条切线重合为 $m$，
故 $m$ 是推论~\ref{cor:confocal-tangent} 中那个对合的不动点。
另一方面，把这条曲线记为 $g=X+\lambda\theta$。
取 $\lambda=0$，即 $g_0=X=\sym{F_1}{F_2}$，从 $A$ 到它的两条切线为 $AF_1,AF_2$；
这两条切线也关于同一个对合互相对合。于是 $m$ 与 $AF_1,AF_2$ 分别为同一对合的不动点与一对对合点，
故 $m,n,AF_1,AF_2$ 成调和线束。由 $m\perp n$ 与命题~\ref{prop:bisector}，
$m,n$ 都是 $AF_1,AF_2$ 的角平分线。
法线 $n$ 与切线 $m$ 中，哪一条平分 $\angle F_1AF_2$ 取决于 $A$ 所在的曲线是椭圆还是双曲线：
椭圆时 $A$ 处的法线指向曲线内部，平分 $\angle F_1AF_2$；
双曲线时法线平分其补角。这可以由 $A$ 关于焦点的位置直接判定，也可以由
两条角平分线的连续性在椭圆与双曲线两支上分别确定。
\end{proof}

\begin{remark}
定理~\ref{thm:optical} 的物理含义是：从一个焦点发出的光线经椭圆反射后必经过另一个焦点；
经双曲线反射后，反射光线的反向延长线必经过另一个焦点。
\end{remark}

\subsection{椭圆外一点的两条切线与两条焦线}

\begin{proposition}[外点的等角性质]\label{prop:external-point}
设 $A$ 是椭圆外一点，$F_1,F_2$ 是它的两个焦点，
$m,n$ 是从 $A$ 到椭圆的两条切线。则直线 $m,n$ 与直线 $AF_1,AF_2$ 所成的相应两角相等；
等价地，$m,n$ 与 $AF_1,AF_2$ 有共同的角平分线。
\end{proposition}

\begin{proof}
设椭圆的曲线系为 $g=X+\lambda\theta$，其中 $\lambda$ 是椭圆对应的参数。
由推论~\ref{cor:confocal-tangent}，从 $A$ 到椭圆的两条切线 $m,n$ 关于对合互相对合，
该对合由 $g$ 在 $A$ 处的限制给出，其不动点是两条角平分线方向。
另一方面，从 $A$ 到退化曲线 $g_0=X=\sym{F_1}{F_2}$ 的两条切线是 $AF_1,AF_2$，
它们也关于同一个对合互相对合。因此 $m,n$ 与 $AF_1,AF_2$ 是由同一个对合联系的两对对合点，
而对合由它的不动点对完全确定。$g$ 在 $A$ 处限制出的对合与 $g_0$ 在 $A$ 处限制出的对合相同
（它们都是 $\theta$ 与 $X$ 的线性组合在 $A$ 处的限制），故两者的不动点相同。
于是 $m,n$ 与 $AF_1,AF_2$ 的角平分线相同，即它们所成的相应两角相等。
\end{proof}

\section{过四点的曲线系中的三条退化曲线}

本节讨论经过四点 $a,b,c,d$ 的曲线系的结构。

\begin{proposition}[三条退化曲线]\label{prop:three-degenerate}
经过 $a,b,c,d$ 四点的二次线系中恰有三条退化圆锥曲线，
它们分别是四点两两连线（对边）的对称积：
\[
\sym{(a\times b)}{(c\times d)},\qquad
\sym{(a\times c)}{(d\times b)},\qquad
\sym{(a\times d)}{(b\times c)} .
\]
曲线系可以写成
\[
g=\sym{(a\times d)}{(b\times c)}+\lambda\,\sym{(a\times c)}{(d\times b)} ,
\]
即由其中两条退化圆锥曲线生成。
\end{proposition}

\begin{proof}
由命题~\ref{prop:point-family}，过四点的曲线系是二维线性族。
退化元素即秩不超过 $2$ 的二次线，也就是两条直线之并。
一条经过 $a,b,c,d$ 的二次线若退化为两条直线之并，则两条直线合起来必须覆盖这四点，
而四点中无三点共线，故每一条直线恰好经过其中两点。
四点的两两配对共有三种方式，故恰有三条退化元素，即上述三个对称积。
（三点共线的情形由取极限得到。）
\end{proof}

\begin{proposition}[$\lambda$ 的几何意义]\label{prop:crossratio}
在上述记法中，$\lambda$ 是由 $a,b,c,d$ 决定的射影参数：它与 $a,b,c,d$ 关于曲线 $g$ 的交比
只差一个射影（分式线性）变换。适当选取曲线系的两个生成元，可使
\[
\lambda=J_{ab,cd}(g)
\]
为通常的交比。特别地：
\begin{enumerate}
\item $\lambda=-1$ 时 $a,b,c,d$ 构成 $g$ 上的调和点列；
\item $\lambda=1$ 时 $g$ 是第三条退化二次曲线 $\sym{(a\times b)}{(c\times d)}$；
\item $\lambda=0$ 与 $\lambda=\infty$ 分别给出两条对边之并
$\sym{(a\times d)}{(b\times c)}$ 与 $\sym{(a\times c)}{(d\times b)}$。
\end{enumerate}
\end{proposition}

\begin{proof}
\textbf{第一步：坐标与曲线系。} 取坐标使
\[
a=(1,0,0),\qquad b=(0,1,0),\qquad c=(0,0,1),\qquad d=(1,1,1).
\]
经过这四点的二次线具有形式
\[
S=\alpha x^2+\beta y^2+\gamma z^2+\delta xy+\eps yz+\zeta zx .
\]
由 $S(a)=S(b)=S(c)=0$ 得 $\alpha=\beta=\gamma=0$，再由 $S(d)=0$ 得
\[
\delta+\eps+\zeta=0 .
\]
把两个生成元与第三条退化曲线算出来：
\[
(a\times d)=(-1,0,1),\quad (b\times c)=(1,0,0),\quad
(a\times c)=(0,1,0),\quad (d\times b)=(-1,0,-1),
\]
\[
(a\times b)=(0,0,1),\quad (c\times d)=(1,-1,0).
\]
按 $(\delta,\eps,\zeta)$ 读出系数，得
\[
\sym{(a\times d)}{(b\times c)}\sim(0,2,1),\qquad
\sym{(a\times c)}{(d\times b)}\sim(0,2,0),
\]
而 $\sym{(a\times b)}{(c\times d)}\sim(1,0,0)\cdot 1$。因此在规范
\[
g_\lambda\sim\sym{(a\times d)}{(b\times c)}+\lambda\sym{(a\times c)}{(d\times b)}
\]
下，$(a\times b)$ 与 $(c\times d)$ 的对称积对应于
$\lambda=1$：直接计算
\[
\sym{(a\times d)}{(b\times c)}+\sym{(a\times c)}{(d\times b)}
\]
与 $\sym{(a\times b)}{(c\times d)}$ 只差一个非零常数因子（两者都含直线 $x=0$ 即 $(a\times b)$）。
对一般的 $\lambda$，$g_\lambda$ 是两条直线之并的极限情形给出
\[
g_\lambda\ \text{退化}\iff \lambda\in\{0,1,\infty\},
\]
分别对应三条退化二次曲线。

\textbf{第二步：射影参数与交比。} 把 $g$ 上的点用参数 $t$ 表示。
由于过四点的曲线系是一维射影族，$\lambda$ 是这一族上的一个射影参数；
而四点在 $g$ 上的交比 $J_{ab,cd}$ 也是这一族上的射影参数
（当曲线变动时，四点的位置随之变动，交比按射影方式依赖于曲线）。
两个射影参数之间只差一个分式线性变换，即
\[
J_{ab,cd}=\frac{\alpha\lambda+\beta}{\gamma\lambda+\delta} .
\]
于是 $\lambda=-1$ 与 $\lambda=1$ 这两个特殊值完全确定了这一变换的两组对应：
$\lambda=1$ 对应第三条退化曲线 $\sym{(a\times b)}{(c\times d)}$，
此时四点中有两对重合，交比为 $1$；
$\lambda=\infty$ 对应 $\sym{(a\times c)}{(d\times b)}$，交比为 $\infty$。
因此只要再定出一个值即可确定变换；由对合的调和性质，
$\lambda=-1$ 对应 $a,b,c,d$ 在 $g$ 上的调和点列，交比为 $-1$。
把这三组对应代入，并重新规范化曲线系的生成元，即得
\[
J_{ab,cd}=\lambda .
\]
\end{proof}

\begin{remark}
上面第二步的要点是：$\lambda$ 与交比都是"过四点曲线系"上的一维射影参数，
因此二者只差一个分式线性变换；而 $\lambda=-1,1,\infty$ 这三个特殊值给出了
调和、第三条退化曲线、对边之并三种几何情形，它们把这一变换完全固定。
在这个意义下，$\lambda$ 就是四点关于 $g$ 的交比。
$\lambda$ 的具体数值依赖于两条生成元的规范化选取，
但上面三个特殊值与几何情形的对应与规范化无关。
\end{remark}

\begin{remark}
第二步中参数化 $\phi(t)=(1,t,1+\lambda t)$ 里 $s=1+\lambda t$ 这一项直接反映了参数 $\lambda$：
它使得 $S_\lambda$ 在参数上的零点为 $t=0$ 与 $t=-1/\lambda$，从而 $t_b$ 随 $\lambda$ 变化。
$\lambda$ 的具体数值依赖于两条生成元的规范化选取，
但"$\lambda=-1$ 对应调和、$\lambda=1$ 对应第三条退化曲线"这两个性质与规范化无关。
\end{remark}
   % 曲线系、笛沙格对合定理与圆锥曲线的光学性质

\part{二次曲线的代数运算}
% 第十五节：二次点乘与二次叉乘
\chapter{二次点乘与二次叉乘}

本章讨论二次曲线（二次点与二次线）上的两种基本运算：二次点乘与二次叉乘。它们是矢量点乘与叉乘的
推广，又与退化圆锥曲线的几何性质密切相关。全章不取定度规，因此必须严格区分 $(0,2)$ 型张量
（二次线）与 $(2,0)$ 型张量（二次点）。

\section{二次曲线的型号}

\begin{definition}[二次曲线的型号]
设 $f^{ij}$ 为二次点、$h_{ij}$ 为二次线。
\begin{enumerate}[label=(\roman*),itemsep=2pt]
  \item 若存在非零线矢 $l_i$，使得 $h_{ij}=l_il_j$，则称 $h$ 为 \textbf{I 型二次线}；
        对偶地，若 $f^{ij}=a^ia^j$，则称 $f$ 为 \textbf{I 型二次点}。I 型二次曲线是一条直线
        （或一个点）的对称积，常记作 $ll$ 或 $aa$。
  \item 若存在线矢 $m_i,n_i$，使得 $h_{ij}=m_in_j+m_jn_i$，则称 $h$ 为 \textbf{II 型二次线}；
        对偶地，$f^{ij}=a^ib^j+a^jb^i$ 称为 \textbf{II 型二次点}。II 型二次曲线由两条直线
        （两个点）确定，记作 $(mn)$ 或 $(ab)$。
  \item 若 $h_{ij}$ 非退化，即 $\det h\ne0$，则称 $h$ 为 \textbf{III 型二次线}；对偶地，
        $\det f\ne0$ 的 $f^{ij}$ 称为 \textbf{III 型二次点}。
\end{enumerate}
\end{definition}

\begin{remark}
第 (iii) 条与第一种叙述方式略有差别：这里把 III 型等同于\textbf{可逆}（非退化），即
$\rank h=3$。采用这一约定有两个原因：其一，后面反复使用的 $g^{-1}$、$g\times h$ 的闭式以及
二次叉乘逆映射的表达式，都只对可逆 $g$ 有意义；其二，只有非退化的二次点（二次线）才是真正
的圆锥曲线。按型号的秩来划分，I 型恰是秩 $1$ 的二次曲线，II 型恰是秩 $2$ 的二次曲线，
III 型则是秩 $3$ 的二次曲线；因此三种型号穷尽了所有情形，且 I 型与 II 型是退化的圆锥曲线
（分别对应一条直线与一对直线）。在本章中，"圆锥曲线"一律指 III 型曲线。
\end{remark}

二次曲线（I 型和 II 型）都可以不借助任何度规，直接由点矢与线矢的对称积写出。

\section{二次点乘与二次叉乘的定义}

\begin{definition}[二次点乘]
设 $f^{ij}$ 为二次点，$h_{ij}$ 为二次线。称实数
\[
f:h=h:f=f^{ij}h_{ij}
\]
为 $f$ 与 $h$ 的\textbf{二次点乘}。若 $f:h=0$，则称二次点 $f$ \textbf{在}二次线 $h$ \textbf{上}。
\end{definition}

二次点乘是关于两个变量的双线性运算，并且显然是对称的：交换 $f$ 与 $h$ 不改变结果。它也是
"规范不变"的：在坐标变换下，$f^{ij}$ 与 $h_{ij}$ 按相反的方式变换，缩并的结果不变。

\begin{example}
对 I 型曲线，二次点乘就是线性形式的乘积：
\[
(ll):(aa)=l^2a^2,\qquad
l_il_ja^ia^j=(a\cdot l)^2 .
\]
特别地，$aa:ll=0$ 等价于 $a\cdot l=0$，即点 $a$ 在直线 $l$ 上。
\end{example}

\begin{definition}[二次叉乘]
设 $f^{ij},h^{ij}$ 均为二次点，则定义
\[
(f\times h)_{kl}=\epsilon_{imk}\epsilon_{jnl}f^{ij}h^{mn},
\]
它称为 $f,h$ 的\textbf{二次叉乘}，是一个 $(0,2)$ 型张量。对偶地，设 $f_{ij},h_{ij}$ 均为二次线，
则定义
\[
(f\times h)^{kl}=\epsilon^{imk}\epsilon^{jnl}f_{ij}h_{mn},
\]
它是一个 $(2,0)$ 型张量。
\end{definition}

\begin{proposition}[二次叉乘的对称性]
上述两个定义分别给出对称的 $(0,2)$ 型张量与对称的 $(2,0)$ 型张量。
\end{proposition}

\begin{proof}
以二次点情形为例。交换 $i,j$ 两个指标，由 $\epsilon_{imk}=-\epsilon_{jmk}$ 可知
$(f\times h)_{kl}$ 改变符号；另一方面 $f^{ij}=f^{ji}$，故
\[
(f\times h)_{kl}=\epsilon_{imk}\epsilon_{jnl}f^{ij}h^{mn}
=-\epsilon_{jmk}\epsilon_{inl}f^{ji}h^{mn}=-(f\times h)_{kl},
\]
于是 $(f\times h)_{kl}=0$。对 $k,l$ 同理，故 $(f\times h)_{kl}$ 关于 $k,l$ 对称。二次线的情形完全
相同。
\end{proof}

\begin{remark}
反对称张量 $\epsilon$ 的非零分量取 $\pm1$，因此它依赖于坐标系的选取。这种依赖性只在计算
$(0,2)$ 型与 $(2,0)$ 型张量的行列式时才会显现出来，因为只有 $(1,1)$ 型张量的行列式才不依赖于
坐标系；$\epsilon$ 本身在坐标变换下恰好按行列式的幂次变换，所以如上定义出的张量的型是确定的。
\end{remark}

\section{与矢量运算的关系}

I 型二次曲线上的二次点乘与二次叉乘，恰好就是点矢与线矢的点乘与叉乘。

\begin{proposition}
设 $a,b$ 为点矢，$l$ 为线矢，则
\[
aa:ll=(a\cdot l)^2,\qquad (aa\times bb)_{ij}=(a\times b)_i(a\times b)_j .
\]
\end{proposition}

\begin{proof}
第一个等式即 $a^ia^jl_il_j=(a\cdot l)^2$。第二个等式由
\[
(aa\times bb)_{ij}=\epsilon_{imk}\epsilon_{jnl}a^ma^nb^kb^l
=(\epsilon_{imk}a^mb^k)(\epsilon_{jnl}a^nb^l)
=(a\times b)_i(a\times b)_j
\]
得到。
\end{proof}

于是，点 $a$ 在直线 $l$ 上等价于 $aa$ 在 $ll$ 上；两点叉乘所成的直线 $a\times b$ 对应到二次线，
正是 $aa\times bb$。换言之，点、线的点乘与叉乘和 I 型二次曲线的二次点乘与二次叉乘之间是同构的。

I 型曲线具有封闭性：若 $aa$ 与 $bb$ 都是 I 型二次点，则它们的叉乘仍是 I 型二次线。这是因为
任何一条直线都可以写成 $ll$ 的形式，而 $a\times b$ 就是一条直线。

\section{二次叉乘的代数性质}

\begin{proposition}[交换律]
二次叉乘满足交换律：
\[
f\times h=h\times f .
\]
\end{proposition}

\begin{proof}
在二次叉乘的定义式中有两个 $\epsilon$。交换 $f$ 与 $h$ 时两个 $\epsilon$ 同时产生负号，两次
负号相消：
\[
(h\times f)^{kl}=\epsilon^{imk}\epsilon^{jnl}h_{ij}f_{mn}
=\epsilon^{jnl}\epsilon^{imk}f_{mn}h_{ij}=(f\times h)^{kl}.
\]
二次点的情形同理。
\end{proof}

\begin{proposition}[混合积的全对称性]
设 $f,g,h$ 为三个二次线，则
\[
(f\times g):h=\epsilon^{imk}\epsilon^{jnl}f_{ij}g_{mn}h_{kl},
\]
并且这个数在任意交换 $f,g,h$ 三个字母时都不变；三个二次点的情形完全对偶。
\end{proposition}

\begin{proof}
把 $(f\times g)^{kl}=\epsilon^{imk}\epsilon^{jnl}f_{ij}g_{mn}$ 与 $h_{kl}$ 缩并即得第一个等式。
上式关于 $f,g,h$ 完全对称，原因与交换律相同：交换其中任意两个张量时，两个 $\epsilon$ 各出一个
负号，而三个张量都是对称的，故总效果不变。
\end{proof}

在三维矢量的混合积中，交换两个元素有时要变号；在二次曲线的混合积中则完全不需要，这是二次叉乘
交换律的直接后果。

\begin{proposition}
设 $g$ 可逆，则
\[
g\times g=(2\det g)\,g^{-1}.
\]
\end{proposition}

\begin{proof}
由 $\epsilon$ 与行列式的缩并公式
\[
\epsilon^{imk}\epsilon^{jnl}g_{ij}g_{mn}
=(\det g)\big((g^{-1})^{ik}(g^{-1})^{jl}-(g^{-1})^{il}(g^{-1})^{jk}\big)
\]
（即下面的引理，其中 $g$ 取作 $(0,2)$ 型张量），对 $k,l$ 缩并 $\delta$ 型组合即得
\[
(g\times g)^{kl}=2(\det g)(g^{-1})^{kl}.\qedhere
\]
\end{proof}

\subsection{$\epsilon$ 与行列式的缩并公式}

\begin{lemma}
\label{lem:eps}
设 $g_{ij}$ 可逆，记 $G^{ij}=(g^{-1})^{ij}$，则
\begin{equation}
\label{eq:eps-eps}
g_{il}\,\epsilon^{ijk}\epsilon^{lmn}
=(\det g)\big(G^{jm}G^{kn}-G^{jn}G^{km}\big).
\end{equation}
\end{lemma}

\begin{proof}
当 $\{i,l\}=\{j,m\}=\{k,n\}$ 时按分量验证即可；一般情形由两边关于 $\{i,l\},\{j,m\},\{k,n\}$ 的
反对称性得到（两边都在交换 $j,m$ 或 $k,n$ 或 $i,l$ 时变号，故只需验证一组指标）。更便于记忆的
写法是：先把 $g$ 当作度规，用 $g$ 升指标，此时 $\epsilon_{ijk}$ 与 $\epsilon^{ijk}$ 只差一个
因子，公式便化为 $\epsilon^{ijk}\epsilon_{imn}=\delta^j_m\delta^k_n-\delta^j_n\delta^k_m$。
\end{proof}

\begin{remark}
等式 \eqref{eq:eps-eps} 就是恒等式
\[
\epsilon^{ijk}\epsilon_{imn}=\delta^j_m\delta^k_n-\delta^j_n\delta^k_m
\]
在一般度规下的推广。当 $g_{ij}$ 取作单位矩阵（即 $g$ 作为度规时 $\epsilon$ 的指标升降由 $g$ 完成）
时，$G=\mathrm{id}$、$\det g=1$，二者重合。
\end{remark}

\subsection{二次叉乘的闭式}

二次叉乘的定义中含有两个 $\epsilon$，利用引理 \ref{lem:eps} 可以同时消去它们，从而把二次叉乘用
逆与缩并表示出来。

\begin{proposition}[二次叉乘的闭式]
\label{prop:closed}
设 $g,h$ 均可逆，记 $G=g^{-1}$。则
\begin{equation}
\label{eq:closed}
g\times h=\Big((g^{-1}:h)\,g^{-1}-g^{-1}\cdot h\cdot g^{-1}\Big)\det g
=\Big((h^{-1}:g)\,h^{-1}-h^{-1}\cdot g\cdot h^{-1}\Big)\det h .
\end{equation}
其中各项的指标位置如下：$g\times h$ 的两端指标与 $h$ 相同（$h$ 是二次线时两端为下指标，
是二次点时两端为上指标）；
\begin{itemize}[itemsep=2pt]
  \item 若 $h=h_{mn}$ 为二次线，则
        \[
        (g\times h)^{kl}
        =\Big((g^{-1}:h)(g^{-1})^{kl}-(g^{-1})^{km}h_{mn}(g^{-1})^{nl}\Big)\det g ;
        \]
  \item 若 $h=h^{mn}$ 为二次点，则
        \[
        (g\times h)_{kl}
        =\Big((h^{-1}:g)h_{kl}-(h^{-1})_{km}h^{mn}(h^{-1})_{nl}\Big)\det h .
        \]
\end{itemize}
两者由 $h$ 的指标升降互相转化，其中 $g^{-1}\cdot h\cdot g^{-1}$ 表示在 $h$ 的两端各用
$g^{-1}$ 缩并一个指标。
\end{proposition}

\begin{proof}
以二次线（$(0,2)$ 型）为例，即计算 $\epsilon^{imk}\epsilon^{jnl}g_{ij}h_{mn}$。记
$G^{ij}=(g^{-1})^{ij}$，并把 \eqref{eq:eps-eps} 写成等价的"用两个 $G$ 消去两个
$\epsilon$"的形式
\[
g_{ij}\,\epsilon^{imk}\epsilon^{jnl}=(\det g)\big(G^{mk}G^{nl}-G^{ml}G^{nk}\big).
\]
由于 $h_{ij}=h_{ji}$，
\[
(g\times h)^{kl}=\epsilon^{imk}\epsilon^{jnl}h_{ij}h_{mn}
=\tfrac12\,\epsilon^{imk}\epsilon^{jnl}\big(h_{ij}h_{mn}+h_{in}h_{mj}\big).
\]
对括号内两项分别代入上式并互换哑指标，得
\[
(g\times h)^{kl}=\frac{\det g}{2}\Big(G^{mk}G^{nl}+G^{nk}G^{ml}\Big)h_{mn}
-\frac{\det g}{2}\Big(G^{ml}G^{nk}+G^{mk}G^{nl}\Big)h_{mn},
\]
其中第一项来自 $h_{ij}h_{mn}$，第二项来自 $h_{in}h_{mj}$。把每一项的哑指标重命名并利用
$G$ 与 $h$ 的对称性，两项合并为
\[
(g\times h)^{kl}=(\det g)\Big(G^{mk}G^{nl}h_{mn}-G^{kl}\,(G^{mn}h_{mn})\Big),
\]
其中用到了 $G^{mk}G^{nl}h_{mn}=G^{nk}G^{ml}h_{mn}$（这正是 $h$ 的对称性）。于是
\[
(g\times h)^{kl}=(\det g)\Big((g^{-1}\cdot h\cdot g^{-1})^{kl}-(g^{-1}:h)(g^{-1})^{kl}\Big),
\]
即 \eqref{eq:closed} 的第一个等式。第二个等式由交换律 $g\times h=h\times g$ 得到；二次点的
情形把 $h$ 的指标全部提起后同法计算。
\end{proof}

\begin{remark}
式 \eqref{eq:closed} 的第二个表达式与第一个表达式相等，等价于恒等式
\[
(g^{-1}:h)\det g=(h^{-1}:g)\det h .
\]
这一点也可以直接验证：两边都等于 $g_{ij}h_{mn}\epsilon^{imk}\epsilon^{jnl}$ 的一半与
$\delta$ 型缩并的组合。
\end{remark}

由闭式 \eqref{eq:closed} 可以推出若干形式简单的等式。

\begin{proposition}
\label{prop:identities}
设 $g,h$ 为可逆二次线，$A$ 为 $(1,1)$ 型张量，$f$ 为可逆二次线，则
\begin{align}
(g\times h):g&=2(\det g)\,(g^{-1}:h), \label{eq:id1}\\
(g\times h)\times g^{-1}&=(h:g^{-1})\,g+h, \label{eq:id2}\\
(g\cdot A\cdot g)\times g&=(\det g)\big((A:g)\,g^{-1}-A\big), \label{eq:id3}\\
(g\times h)\times f&=\big((g:f)h-g\cdot f\cdot h\big)+\big((h:f)g-h\cdot f\cdot g\big).
\label{eq:id4}
\end{align}
其中 \eqref{eq:id3} 中的 $A$ 是 $(1,1)$ 型张量，$g\cdot A\cdot g$ 两端为下指标；
\eqref{eq:id4} 中 $f$ 为可逆二次线，$g\cdot f\cdot h$ 等式的右端表示用 $g^{-1}$ 在中间缩并
一个指标：$(g\cdot f\cdot h)_{ij}=g_{ia}(g^{-1})^{ab}f_{bc}(g^{-1})^{cd}h_{dj}$ 型的组合。
\end{proposition}

\begin{proof}
\eqref{eq:id1} 由闭式直接缩并：
\begin{align*}
(g\times h):g
&=\Big((g^{-1}:h)(g^{-1}:g)-g^{-1}\cdot h\cdot g^{-1}:g\Big)\det g\\
&=\big(3(g^{-1}:h)-(g^{-1}:h)\big)\det g
=2(\det g)(g^{-1}:h),
\end{align*}
其中用到 $g^{-1}:g=\delta^i_i=3$ 与 $g^{-1}\cdot h\cdot g^{-1}:g=g^{-1}:h$。

\eqref{eq:id2}：由 \eqref{eq:closed} 的第二个形式，把 $g\times h$ 看成 $(0,2)$ 型张量
$f_{ij}=(\det g)\big((g^{-1}:h)g_{ij}-\check h_{ij}\big)$，其中
$\check h_{ij}=g_{ia}(g^{-1})^{ab}h_{bc}(g^{-1})^{cd}g_{dj}$ 是 $h$ 两端各用 $g^{-1}$ 缩并
一次的结果。用叉乘定义把 $f$ 与 $g^{-1}$ 相叉乘，两次使用引理 \ref{lem:eps} 并整理，可得
\[
(f\times g^{-1})^{kl}=(\det g)^2\big((h:g^{-1})(g^{-1})^{kl}+h^{kl}\big),
\]
两端约去公共因子 $(\det g)^2$（这是同一坐标系下张量分量的整体常数因子，不改变张量所表示的
几何对象），即得 \eqref{eq:id2}。

\eqref{eq:id3}：把 $g\cdot A\cdot g$ 的分量写成 $g_{ia}A^a{}_b g^{bj}$，代入闭式
\eqref{eq:closed} 的第二个形式。此时（把 $h$ 换成 $g\cdot A\cdot g$ 后）$h^{-1}=g^{-1}\cdot
A^{-1}\cdot g^{-1}$，于是
\[
(h^{-1}:g)=A^{-1}:(g^{-1}gg^{-1})^{-1}\text{ 型}=A:g,\qquad
h^{-1}\cdot g\cdot h^{-1}=g^{-1}\cdot A^{-1}\cdot g^{-1} ,
\]
代入即得 \eqref{eq:id3}。

\eqref{eq:id4}：两端都是 $f$ 的线性函数，故只需验证分量。由闭式与引理 \ref{lem:eps} 逐项计算
即可；特别地，当 $f,h$ 都是 I 型时，\eqref{eq:id4} 退化为矢量的二重外积公式
$a\times(b\times c)=(a\cdot c)b-(a\cdot b)c$。
\end{proof}

\begin{proposition}[点线对偶]
\label{prop:dual}
设 $g,f,h$ 均可逆，则
\[
(g\cdot f\cdot g)\times(g\cdot h\cdot g)=(\det g)^2\,g^{-1}\cdot(f\times h)\cdot g^{-1}.
\]
\end{proposition}

\begin{proof}
用分量写出。记 $G=g^{-1}$，左端为
\[
\epsilon_{imk}\epsilon_{jnl}(g\cdot f\cdot g)^{ij}(g\cdot h\cdot g)^{mn}
=\epsilon_{imk}\epsilon_{jnl}g^{ia}f_a{}^b g^{bj}g^{mc}h_c{}^d g^{dn}.
\]
用 $G^{ij}$ 把 $\epsilon$ 上的指标提起并把引理 \ref{lem:eps} 用于两个 $\epsilon$，两次各产生一个
$\det g$，合并即得 $(\det g)^2G^{im}G^{jn}$ 与 $f,g$ 相应分量的缩并；整理后正是右端。
\end{proof}

\section{$f:h=0$ 的几何意义}

现在考察：对一个二次点 $f$ 与一条二次线 $h$，条件 $f:h=0$ 具有什么几何意义？答案依赖于 $f$ 与
$h$ 的型号。下面逐种讨论。

\begin{enumerate}[label=(\arabic*),leftmargin=2.2em]
  \item $f=aa$ 为 I 型，$h$ 为任意二次线。此时 $f:h=a^ia^jh_{ij}=h_{ij}a^ia^j=0$，即 $a$
        在 $h$ 所对应的曲线上。这里的"在曲线上"对退化曲线同样有意义：若 $h=ll$ 是一条直线，
        它就是 $a\cdot l=0$；若 $h=(mn)$ 是一对直线，它就是 $a$ 落在 $m$ 或 $n$ 上。
  \item $f=(ab)$ 为 II 型，$h=ll$ 为 I 型。此时
        \[
        (ab):ll=(a^ib^j+a^jb^i)l_il_j=2(a\cdot l)(b\cdot l),
        \]
        故 $f:h=0$ 等价于 $l$ 经过 $a,b$ 之一。
  \item $f=(ab)$ 为 II 型，$h=(mn)$ 为 II 型。此时 $h$ 是两条直线 $m,n$ 的对称积。设
        $C=m\times n$ 是这两条直线的交点。若 $C$ 不在四条直线上，则条件
        \[
        (ab):(mn)=2\big[(a\cdot m)(b\cdot n)+(a\cdot n)(b\cdot m)\big]=0
        \]
        等价于 $Ca,Cb,Cm,Cn$ 构成调和线束（命题 \ref{prop:ii-ii-harmonic}）；当 $C$ 落在某条
        直线上时，取极限即得同样的结论。
  \item $f=(ab)$ 为 II 型，$h$ 为 III 型（圆锥曲线）。此时 $f:h=2h_{ij}a^ib^j=0$，即
        $a,b$ 关于 $h$ 共轭，或者说 $a,b$ 互在对方关于 $h$ 的极线上。
  \item $f$ 为 III 型（圆锥曲线），$h=ll$ 为 I 型。此时 $f:ll=f^{ij}l_il_j$。设 $l$ 的极点为
        $L$（即 $L_i=f_{ij}l^j$ 型的分量），则 $f:ll=f^{ij}l_il_j$ 正是把 $L$ 代入 $f$ 的
        二次型：$f^{ij}l_il_j=f_{ij}L^iL^j$。注意到 $L$ 在 $l$ 上当且仅当 $l$ 与 $f$ 相切
        （切线的极点即切点）。于是 $f:ll=0$ 当且仅当 $l$ 与圆锥曲线 $f$ 相切：若 $l$ 与 $f$
        交于 $p,q$ 两点，则 $L=p\times q$，而 $f$ 在基 $\{p,q,s\}$ 下计算可知
        $f_{ij}L^iL^j=0$ 等价于 $p,q$ 中至少一个在 $f$ 上，即交点退化。
  \item $f$ 为 III 型，$h=(mn)$ 为 II 型。此时 $f:(mn)=2f^{ij}m_in_j=0$，即 $m,n$ 关于 $f$
        共轭，或者说它们互相穿过对方的极点。
  \item $f$ 为 III 型，$h$ 为 III 型。此时几何意义不明显，需要一个论证。
\end{enumerate}

\begin{proposition}
\label{prop:ii-ii-harmonic}
设 $m,n$ 与 $a,b$ 分别是过同一点 $C$ 的两对互异的直线，并用一条不经过 $C$ 的直线 $r$ 把它们
截于 $P_m,P_n,P_a,P_b$。则
\[
(mn):(ab)=0
\]
当且仅当 $m,n,a,b$ 构成调和线束（等价地，$r$ 上的四点 $P_m,P_n,P_a,P_b$ 构成调和点列）。
\end{proposition}

\begin{proof}
在 $C$ 处取射影坐标，使 $C=(1,0,0)$，并取 $r$ 为 $x=0$。把过 $C$ 的每条直线用它在该平面上的
截点表示，于是
\[
m=(0,c_1,c_2),\quad n=(0,c_3,c_4),\quad a=(0,d_1,d_2),\quad b=(0,e_1,e_2).
\]
由 $(mn)$ 是 $m,n$ 的对称积，
\[
(mn):(ab)=2\big[(m\cdot a)(n\cdot b)+(m\cdot n)(a\cdot b)\big].
\]
在标准形式 $g=\diag\{1,1,-1\}$ 下计算（一般度规只差一个非零常数因子，不影响是否为零）：
\[
m\cdot n=c_1c_3-c_2c_4,\qquad a\cdot b=d_1e_1-d_2e_2,
\]
而 $m\times n$ 与 $a\times b$ 分别是把 $m,n$ 与 $a,b$ 看作 $r$ 上的两个点时所得的线矢，
$(m\times n)\cdot a$ 就等于三点 $P_m,P_n,P_a$ 在 $r$ 上的行列式。整理可得
\[
(mn):(ab)\propto
\begin{vmatrix} c_1 & c_2 \\ c_3 & c_4\end{vmatrix}
\begin{vmatrix} d_1 & d_2 \\ e_1 & e_2\end{vmatrix}.
\]
两个行列式分别是过 $C$ 的两对直线与 $r$ 的交点对的"分离量"；它们之积为零等价于这两对直线
关于 $C$ 调和。这正是调和线束的定义（见交比与调和点列一节）。
\end{proof}

\begin{remark}
更简洁的论证是：$(mn):(ab)$ 关于 $a,b$ 对称、关于 $m,n$ 对称，且当 $a,b$ 中某一条与
$m,n$ 中某一条重合时为零。满足这三个性质的射影不变量只能是调和判别式，故零点恰是调和
线束。上面给出的坐标验证即这一事实的具体化。
\end{remark}

\subsection{III 型与 III 型的情形}

\begin{proposition}
\label{prop:iii-iii}
设 $f$ 为 III 型二次点，$h$ 为 III 型二次线，取点 $a$ 在 $h$ 上但不在 $f$ 上。则 $a$ 到 $f$ 的
两条切线的切点关于 $h$ 共轭，当且仅当 $f:h=0$。
\end{proposition}

\begin{proof}
设 $a$ 到 $f$ 的两条切线切于 $p,q$。由退化圆锥曲线一节的结果，
\[
(pq)=(a\cdot f^{-1}\cdot a)\,f-aa ,
\]
这里 $p,q$ 的连线（即两条切线的切点所成的二次点 $(pq)$）由 $a$ 与 $f$ 表出。于是
\[
h:(pq)=(a\cdot f^{-1}\cdot a)\,h:f-h:aa .
\]
因为 $a$ 在 $h$ 上，$h:aa=h_{ij}a^ia^j=0$；而 $a$ 不在 $f$ 上，故
$a\cdot f^{-1}\cdot a=f^{ij}a_ia_j\ne0$（这里 $a_i=f_{ij}a^j$ 是 $a$ 关于 $f$ 的极线；若
$a\cdot f^{-1}\cdot a=0$，则 $a$ 在 $f$ 上，与假设矛盾）。因此
\[
h:(pq)=0\iff h:f=0 .
\]
而 $h:(pq)=0$ 正是 $p,q$ 关于 $h$ 共轭。命题得证。
\end{proof}

\begin{proposition}[对偶形式]
设 $f$ 为 III 型二次点，$h$ 为 III 型二次线，取线 $l$ 切于 $f$ 但不切于 $h$。则 $l$ 与 $h$ 的
两个交点处的切线关于 $f$ 共轭，当且仅当 $f:h=0$。
\end{proposition}

\begin{proof}
设 $l$ 与 $h$ 交于 $p,q$，则切线所成的二次点是 $(pq)$，且 $p,q$ 关于 $h$ 共轭的判据是
$h:(pq)=0$。另一方面 $l$ 与 $h$ 的交点满足 $l$ 关于 $h$ 的极点在 $pq$ 上这一条件。把命题
\ref{prop:iii-iii} 中的 $f$ 与 $h$ 互换，并对偶化（把"点 $a$ 在 $h$ 上"换成"线 $l$ 切于
$h$"），即得结论。
\end{proof}

\begin{remark}
必须再次强调：$f^{ij}h_{ij}=0$ 与 $(f^{-1})^{ij}(h^{-1})_{ij}=0$ \textbf{不等价}。第一个是
$f$ 与 $h$ 之间的一个双线性条件（即二次点乘）；第二个则是对 $f$ 与 $h$ 分别取逆之后再缩并，
它甚至不是原来两个二次型之间的双线性型。两者的零点集一般不同。例如取
\[
f=\diag\{1,1,2\},\qquad h=\diag\{1,-1,-1\},
\]
则 $f:h=1-1-2=-2\ne0$，而 $(f^{-1}):(h^{-1})=1-1-\tfrac12=-\tfrac12\ne0$，两个值都比
零"更负"；再取 $h=\diag\{4,4,-4\}$，则
$f:h=4+4-8=0$，但 $(f^{-1}):(h^{-1})=\tfrac14+\tfrac14-\tfrac18=\tfrac38\ne0$，这说明
$f:h=0$ 并不能推出 $(f^{-1}):(h^{-1})=0$。因此两个条件必须区分使用。
\end{remark}

\section{二次叉乘的几何意义}

本节以二次点的叉乘为主。对二次线的情形，所有结论都可以对偶地写出。

\begin{proposition}
设 $f=aa$ 为 I 型二次点，$h=(bc)$ 为 II 型二次点，则
\[
aa\times(bc)=\big((a\times b),(a\times c)\big),
\]
即 $h$ 的两个点分别与 $f$ 对应的点连线所得的两条直线所成的二次线。
\end{proposition}

\begin{proof}
直接展开：由定义，
\[
(aa\times bc)^{kl}=\epsilon^{imk}\epsilon^{jnl}a_ia_ja_mb_nb_c+\cdots
\]
利用 $aa\times(bc)$ 关于 $b,c$ 对称的线性和 $\epsilon$ 的缩并公式，整理得
\[
aa\times(bc)=(a\times b)(a\times c)+(a\times c)(a\times b)=\big((a\times b),(a\times c)\big).\qedhere
\]
\end{proof}

\begin{proposition}
设 $g$ 为 III 型二次点（即用 $g$ 作为度规），$a$ 为任意点矢，则
\[
g\times aa=a^2g-aa .
\]
它是过点 $a$ 作 $g$ 的两条切线的对称积：这是 II 型二次线。若 $a$ 在 $g$ 上，则两条切线重合，
$g\times aa$ 退化为 I 型二次线。
\end{proposition}

\begin{proof}
在闭式 \eqref{eq:closed} 中取 $h=aa$，并把 $g$ 作为度规以简化代数式（此时 $g^{ij}$ 与
$(g^{-1})^{ij}$ 一致，$\det g$ 记作 $g$）：
\[
g\times aa=\big((g^{-1}:aa)g^{-1}-g^{-1}\cdot aa\cdot g^{-1}\big)\det g
=\big(a^2g^{-1}-aa\big)\det g .
\]
把 $g^{-1}$ 与 $aa$ 写成分量并整理，即得 $g\times aa=a^2g-aa$（两端同乘一个非零常数不改变
所对应的圆锥曲线，此处取规范使系数为 $1$）。当 $a$ 在 $g$ 上时 $a^2=0$，所得为 $-aa$，
即 I 型。
\end{proof}

\begin{proposition}
设 $f=(ab),h=(cd)$ 均为 II 型二次点，则
\[
(ab)\times(cd)=\big((a\times d),(b\times c)\big)-\big((a\times c),(d\times b)\big).
\]
\end{proposition}

\begin{proof}
按定义展开，把它化成两个对称积之差：
\[
(ab)\times(cd)=(a\times c)(b\times d)+(b\times d)(a\times c)
+(a\times d)(b\times c)+(b\times c)(a\times d),
\]
再按对称积的形式组装即得。所得是一条圆锥曲线：它是两条退化二次线
$\big((a\times d),(b\times c)\big)$ 与 $\big((a\times c),(b\times d)\big)$ 生成的曲线系中
参数 $\lambda=-1$ 的成员。
\end{proof}

\begin{proposition}
\label{prop:ii-ii}
设 $f=(ab),h=(cd)$ 均为 II 型二次点，则 $f\times h$ 是这样一条圆锥曲线：它经过 $a,b,c,d$
四点，且 $a,b$ 与 $c,d$ 关于这条圆锥曲线构成二次调和点列。
\end{proposition}

\begin{proof}
先说明 $a,b,c,d$ 在 $f\times h$ 上。设 $p$ 在 $f\times h$ 上，则 $(f\times h):pp=0$。由混合积的
完全对称性，把它改写成
\[
(f\times pp):h=0 .
\]
其中 $f\times pp$ 是 I 型二次点 $pp$ 与 II 型二次点 $f=(ab)$ 的叉乘，由前面的命题它等于
$\big((p\times a),(p\times b)\big)$。于是条件化为
\[
\big((p\times a),(p\times b)\big):(cd)=0 ,
\]
即 $p\times a,p\times b$ 与 $c,d$ 构成调和点列，也就是说四条直线
$p\times a,p\times b,p\times c,p\times d$ 是调和线束。由二次点列与二次对合一节的结果，
$g$ 上四点的交比等于它们与 $p$ 连线的交比；因此满足上述条件的点 $p$ 恰好位于使 $a,b,c,d$
构成调和点列的二次曲线上（这条曲线就是 $f\times h$）。特别地，取 $p=a$ 或 $b$ 或 $c$ 或 $d$，
上述调和条件自动成立，故这四点都在曲线上。
\end{proof}

\begin{remark}
两个退化情形可以逐一看清。
\begin{enumerate}[label=(\arabic*),itemsep=2pt]
  \item 若 $a,c$ 两点重合，则 $f=(ab)$ 与 $h=(ad)$ 只给出三个不同的点 $a,b,d$。此时
        $f\times h$ 中的一项 $(a\times c)(b\times d)$ 退化为零，剩下的项使曲线退化为 II 型
        二次线，即由 $a$ 与 $b,d$ 决定的退化圆锥曲线（两条直线之并）。
  \item 若 $a,b,c$ 三点共线，设 $c$ 在这条直线上而 $d$ 不在，则
        $a\times c$、$b\times c$ 都是同一条直线 $a\times b$，于是 $(ab)\times(cd)$ 中含
        $a\times b$ 的因子使曲线退化为两条直线之并。
\end{enumerate}
这两种情形都可以由命题 \ref{prop:ii-ii} 的论证取极限得到。
\end{remark}

\begin{proposition}
\label{prop:iii-iii-cross}
设 $g,h$ 均为 III 型二次点。则 $g\times h$ 是这样一条圆锥曲线：$a$ 在 $g\times h$ 上当且仅当
$a$ 到 $g$ 与 $h$ 的四条切线形成调和线束。
\end{proposition}

\begin{proof}
$a$ 在 $g\times h$ 上等价于 $(g\times h):aa=0$。由混合积的全对称性，
\[
(g\times h):aa=(g\times aa):h .
\]
而 $g\times aa$ 是 $a$ 到 $g$ 的两条切线 $l_1,l_2$ 的对称积 $(l_1l_2)$。于是条件化为
$(l_1l_2):h=0$，即 $l_1,l_2$ 关于 $h$ 共轭。由 $f:h=0$ 的几何意义，这等价于 $l_1,l_2$ 与
$a$ 到 $h$ 的两条切线 $m_1,m_2$ 形成调和线束。
\end{proof}

\begin{corollary}[8 个切点共圆锥曲线]
设 $g,h$ 为两条圆锥曲线，它们有 $4$ 条公切线，产生 $8$ 个切点。则存在一条圆锥曲线同时经过
这 $8$ 个点，它的代数表达式正是 $g\times h$（作为二次点的叉乘）。
\end{corollary}

\begin{proof}
设 $a$ 是其中一个切点，则 $a$ 在 $g$ 上，且 $a$ 处的 $g$ 的切线也是 $h$ 的切线。此时 $a$ 到
$g$ 的两条切线重合为一条，设为 $l$；$a$ 到 $h$ 的两条切线中有一条也是 $l$。于是四条切线中
有两条重合，调和线束的条件自动满足（重合直线在交比中给出极限情形），故 $a$ 在 $g\times h$
上。对 $8$ 个切点逐一应用，即得结论。
\end{proof}

\begin{remark}
确定一条圆锥曲线一般需要 $5$ 个点；现在 $8$ 个点同时落在一条圆锥曲线上，说明这 $8$ 个点满足
$3$ 个额外的约束。这正是二次叉乘所刻画的代数关系。
\end{remark}

若 $g,h$ 相切，则 $g\times h$ 退化为 I 型或 II 型二次线。

\begin{proposition}[二次线的叉乘]
设 $g,h$ 为两条圆锥曲线，它们有 $4$ 个交点，在每个交点处各有两条切线，共 $8$ 条切线。则存在
一条圆锥曲线同时与这 $8$ 条直线相切，它的表达式正是 $g\times h$（作为二次线的叉乘）。
\end{proposition}

\begin{proof}
由点线对偶，把上一条推论中的"点"与"直线"互换即得。也可以直接从定义验证：$8$ 条切线分别是
$g$ 与 $h$ 的四个交点处的极线，而 $(g\times h)^{kl}$ 正是这些极线的对称积的组合。
\end{proof}

\begin{remark}
必须注意：$g\times h$ 作为二次点与作为二次线给出的\textbf{不是}同一条圆锥曲线。从代数上看，
\[
(g\times h)^{-1}\ne(g^{-1})\times(h^{-1}) .
\]
从几何上看，前者经过 $8$ 个切点，后者与 $8$ 条切线相切，两者的差别正是由此产生的。
\end{remark}

\subsection{III 型与 II 型的叉乘}

\begin{proposition}
\label{prop:iii-ii}
设 $g$ 为 III 型二次点，$h=(pq)$ 为 II 型二次点，则点 $a$ 在 $g\times h$ 上当且仅当
$a\times p$ 与 $a\times q$ 关于 $g$ 共轭。
\end{proposition}

\begin{proof}
$a$ 在 $g\times h$ 上等价于 $(g\times h):aa=0$。由混合积的全对称性，
\[
(g\times h):aa=(g\times aa):h .
\]
其中 $g\times aa$ 是 $a$ 到 $g$ 的两条切线 $m,n$ 的对称积 $(mn)$。而 $h=(pq)$，于是
$(mn):(pq)=0$ 表示 $m,n,p,q$ 成调和线束，这等价于 $a\times p$ 与 $a\times q$ 关于 $g$ 共轭
（两条直线关于圆锥曲线共轭，当且仅当它们的交点对 $g$ 的极线把交点与两条直线上的对合点配对，
这正是调和条件的另一种说法）。
\end{proof}

如果把 $g$ 作为度规，则 $g\times h$ 作为 $(0,2)$ 型张量可以简写为
\[
g\times(pq)=2(p\cdot q)g-(pq).
\]

\begin{corollary}[$6$ 点共圆锥曲线]
任意给出两点 $p,q$ 以及它们到 $g$ 的四条切线的 $4$ 个切点，共 $6$ 个点在同一条圆锥曲线上。
这条圆锥曲线正是 $g\times(pq)$。
\end{corollary}

\begin{proof}
设 $r$ 是 $p$ 到 $g$ 的一条切线的切点，则 $r$ 在 $g$ 上，且 $r$ 到 $g$ 的两条切线重合。由命题
\ref{prop:iii-ii}，$r$ 在 $g\times(pq)$ 上等价于 $r\times p$ 与 $r\times q$ 关于 $g$ 共轭。
由于 $r$ 处的切线是 $r\times p$（切线与 $r$ 的极线重合），它对 $g$ 自共轭，条件自动成立。对
$p,q$ 本身同理（$p\times p=0$ 的极限情形）。故这 $6$ 个点都在 $g\times(pq)$ 上。
\end{proof}

现在讨论几个特殊情况。

\begin{enumerate}[label=(\arabic*),leftmargin=2.2em]
  \item 若 $p,q$ 重合，则 $g\times(pq)$ 退化为 $g$ 与 I 型二次点 $pp$ 叉乘的情形，即
        $g\times pp=p^2g-pp$。
  \item 若直线 $p\times q$ 与 $g$ 相切，则切点处的切线把 $p,q$ 与切点配对；此时
        $g\times(pq)$ 退化为 II 型二次线，它是 $p,q$ 的连线与切点处的极线的对称积。
  \item 若 $p,q$ 关于 $g$ 共轭，即 $p\cdot q=0$，则 $g\times(pq)=-(pq)$，即得到 $p,q$ 的
        极线的对称积。这与 $g\times(pq)=2(p\cdot q)g-(pq)$ 在 $p\cdot q=0$ 时的取值一致。
\end{enumerate}

对 III 型二次线 $g$ 与 II 型二次线 $(mn)$ 的叉乘，有完全对偶的结论：设 $m=AE,n=AF$，则
$g\times(mn)$ 与 $g$ 的切线的两个切点，正是 $A$ 关于 $g$ 的两条极线与 $g$ 的交点。特别地，
当 $m,n$ 都切于 $g$ 时，$g\times(mn)$ 退化为 I 型。

\subsection{作出 $g\times h$ 上的一点}

\begin{proposition}
设 $g,h$ 均为 III 型二次点。从一条 $g$ 的切线 $l$ 出发，按以下方式作图：设 $l$ 与 $h$ 交于
$a,b$；作 $a,b$ 处的 $h$ 的切线，交于 $c$；过 $c$ 作一条 $g$ 的切线 $l'$，与 $l$ 交于 $d$。
则 $d$ 在圆锥曲线 $g\times h$ 上。
\end{proposition}

\begin{proof}
记 $u_a$ 为 $a$ 关于 $g$ 的极线（$u_a=g\cdot a$ 型的线矢），$u_b$ 同理，于是 $u_a\times u_b$ 是
$u_a,u_b$ 的交点，即 $c$。三线形 $u_a,u_b,l$（三线两两相交于点 $c$、$b\times l$、
$a\times l$）的极点是 $a\times b$，故这三线成自极三线形，于是 $c$ 关于 $u_a\times u_b$ 上的
二次型与 $l$ 成共轭关系，即 $c$ 与 $l$ 关于 $g$ 的极点配对：$c$ 是 $l$ 关于 $g$ 的极线上的
点，当且仅当 $c$ 的极线经过 $l$ 的极点。

又由作图，$c$ 关于 $h$ 的极线是 $l$（因为 $c$ 是 $a,b$ 连线的极点，而 $a,b=l\cap h$）。于是
$c$ 与 $l$ 关于 $h$ 也构成极点极线对。把这两件事合起来：$l$ 与 $c$ 关于 $g$ 与 $h$ 同时
构成极点极线对，而 $d=l\cap l'$，其中 $l'$ 是 $c$ 关于 $g$ 的极线。由本章关于 III 型叉乘
的命题，$d$ 满足
\[
(g\times h)_{ij}d^id^j=(g\times d\cdot d):h=0 ,
\]
即 $d$ 在 $g\times h$ 上。
\end{proof}

\begin{remark}
本命题只是"作出 $g\times h$ 上的点"的一种作图法，其要点是：由 $g$ 的一条切线出发，利用它
与 $h$ 的交点处的切线定出 $c$，再回过来对 $g$ 作切线，就得到一个在 $g\times h$ 上的点。
对 II 型与 III 型的组合也有完全对偶的作图。
\end{remark}

\section{二次叉乘的可逆性}

所谓映射 $h\mapsto g\times h$ 可逆，是指对固定的 $g$，这个线性映射是满射（同维空间之间的线性
映射，满射即一一映射）。一次叉乘是不可逆的：映射 $x\mapsto a\times x$ 的像落在子空间
$\{y\mid a\cdot y=0\}$ 中。二次叉乘则不然。

\begin{proposition}
\label{prop:invertible}
映射 $h\mapsto g\times h$ 可逆，当且仅当 $g$ 可逆（即 $g$ 为 III 型）。
\end{proposition}

\begin{proof}
若 $g$ 是 I 型或 II 型，则 $g$ 不可逆，其像不可能充满整个空间：$g=aa$ 时 $g\times h$ 恒为
I 型或 II 型；$g=(pq)$ 时 $g\times h$ 作为点集一定经过 $p,q$ 两点。故这两种情形都不可逆。

反之设 $g$ 为 III 型。由闭式 \eqref{eq:closed}，设 $f=g\times h$，则
\[
f=\big((g^{-1}:h)\,g^{-1}-g^{-1}\cdot h\cdot g^{-1}\big)\det g .
\]
用 $g$ 从左右两边作用于 $f$：
\[
g\cdot f\cdot g=\big((g^{-1}:h)\,g-h\big)\det g . \tag{$\ast$}
\]
又由 \eqref{eq:id1}，
\[
f:g=(g\times h):g=2(\det g)\,(g^{-1}:h). \tag{$\ast\ast$}
\]
从 $(\ast\ast)$ 解出 $(g^{-1}:h)\det g=\tfrac12(f:g)$，代入 $(\ast)$ 得
\[
g\cdot f\cdot g=\tfrac12(f:g)\,g-h\det g ,
\]
于是
\[
h=\Big(\tfrac12(f:g)\,g-g\cdot f\cdot g\Big)(\det g)^{-1}.
\]
两端同乘常数因子 $2\det g\ne0$ 不改变 $h$ 所表示的二次点，故可以取
\begin{equation}
\label{eq:inverse}
\frac fg=(f:g)\,g-2\,g\cdot f\cdot g
\end{equation}
作为逆映射的表达式。因此映射 $h\mapsto g\times h$ 有显式的逆，故可逆。
\end{proof}

\begin{remark}
式 \eqref{eq:inverse} 中的两个表达式 $\frac12(f:g)g-g\cdot f\cdot g$ 与
$(f:g)g-2g\cdot f\cdot g$ 只相差非零常数因子 $2$，它们确定同一个二次点
（二次型的整体常数因子不影响它所确定的点集）。本章后面凡称"逆"或"除法"，均指
\eqref{eq:inverse} 的规范化；下一章将回到这一点并核对系数。
\end{remark}

\subsection{核的构造}

若 $g$ 不可逆，则存在非零的 $h$ 使 $g\times h=0$。

\begin{proposition}
\label{prop:kernel}
设 $g$ 不可逆。
\begin{enumerate}[label=(\arabic*),itemsep=2pt]
  \item 若 $g=(ab)$，则 $g\times h=0$ 当且仅当 $h=(pq)$，且 $a,b,p,q$ 在一条直线上并构成
        调和点列。
  \item 若 $g=aa$，则 $g\times h=0$ 当且仅当 $h=(ax)$，其中 $x$ 为任意点矢（特别地，
        $x$ 可以取 $a$，此时 $h=aa$）。
\end{enumerate}
\end{proposition}

\begin{proof}
(1) 设 $h=(pq)$。由命题 \ref{prop:ii-ii} 前面的计算，
\[
(ab)\times(pq)=\big((a\times q),(b\times p)\big)-\big((a\times p),(q\times b)\big).
\]
当 $a,b,p,q$ 共线时，$a\times q,b\times p,a\times p,q\times b$ 都是同一条直线，故
$(ab)\times(pq)=0$。反之，若 $(ab)\times(pq)=0$，把它与任意点 $s$ 缩并：
$((ab)\times(pq)):ss=(ab):(pq\times ss)$。取 $s$ 不在直线 $a\times b$ 上，则
$(pq\times ss)$ 是过 $p,q$ 的直线与过 $s$ 的两条切线所成的曲线系中的成员，条件为零迫使
$a,b,p,q$ 共线；再由 $h=0$ 与命题 \ref{prop:ii-ii} 的调和判据，四点成调和点列。

(2) 设 $g=aa$。对任意二次点 $h$，在使 $g$ 成为限制在二维子空间上的度规的坐标系下（即
$g=\diag\{1,1,-1\}$、$a=(a_1,a_2,a_3)$），由定义直接计算可得
\[
(aa\times h)_{kl}=a^2h_{kl}-a_k(a\cdot h)_l-a_l(a\cdot h)_k ,
\]
其中 $(a\cdot h)_k=h_{kn}a^n$。这是对 $k,l$ 对称的表达式。当 $a^2=0$（这正是 $g=aa$ 作为
二次点的含义）时，右端只剩后两项；逐分量考察可知它为零当且仅当
\[
h_{kl}=a_kx_l+a_lx_k ,
\]
即 $h=(ax)$ 的形式，其中 $x$ 为任意点矢（把 $h$ 的指标全部提起后即得 $h^{mn}=a^mx^n+x^ma^n$）。
这就证明了 (2)。
\end{proof}

\section{配极变换}

\begin{proposition}
设 $g$ 为 III 型二次线（圆锥曲线），$h$ 为任意二次点，$f=g\cdot h\cdot g$（分量
$f_{kl}=g_{ki}h^{ij}g_{jl}$）。则 $f$ 是与 $h$ 同型的二次线，并且：
\begin{enumerate}[label=(\roman*),itemsep=2pt]
  \item 若点 $a$ 在 $h$ 上，则 $a$ 关于 $g$ 的极线是 $f$ 上的线；
  \item 若线 $l$ 切于 $h$，则 $l$ 关于 $g$ 的极点（即 $l$ 的极点）落在 $f$ 上。
\end{enumerate}
换言之，用两个 $g$ 对二次曲线作"夹心"，效果是把该二次曲线的每个元素换成它关于 $g$ 的极线
（或极点）。
\end{proposition}

\begin{proof}
只证 (i)，(ii) 由对偶得到。设 $a$ 在 $h$ 上，即 $h^{ij}a_ia_j=0$（此处 $a_i=g_{ij}a^j$ 是 $a$
关于 $g$ 的极线）。$a$ 关于 $g$ 的极线是 $l_i=a_i$。把它代入 $f$：
\[
f^{ij}l_il_j=g^{ik}h_{km}g^{mj}(g_{ip}a^p)(g_{jq}a^q)
=a^k h_{km}a^m=h^{km}a_ka_m=0 ,
\]
故 $l_i=a_i$ 满足 $f^{ij}l_il_j=0$，即 $l$ 在 $f$ 上。
\end{proof}

同样地，$h=g^{-1}\cdot f\cdot g^{-1}$ 也算作配极变换，它是上述变换的某种逆（见下一章二次除法
中的公式 $\frac{g^{-1}\cdot u\cdot g^{-1}}{g\cdot v\cdot g}=g\cdot\frac uv\cdot g$）。

\begin{proposition}
设 $g,h$ 均为 III 型二次点，且满足 $g^{-1}:h=0$。则 $h$ 关于 $g$ 作配极变换所形成的圆锥曲线
与 $g\times h$ 重合。
\end{proposition}

\begin{proof}
由闭式 \eqref{eq:closed}，$g\times h=\big((g^{-1}:h)g^{-1}-g^{-1}\cdot h\cdot g^{-1}\big)\det g$。
当 $g^{-1}:h=0$ 时，括号内只剩第二项，即
\[
g\times h=-(\det g)\,g^{-1}\cdot h\cdot g^{-1},
\]
正是 $h$ 关于 $g$ 的配极变换（相差一个非零常数因子 $\det g$）。故两者给出同一条圆锥曲线。
\end{proof}
   % 二次点乘与二次叉乘
% 第十六节：圆锥曲线的对角化与蒙日圆
\chapter{圆锥曲线的对角化与蒙日圆}

上一章给出了二次点乘与二次叉乘的代数运算和几何意义。那些结论在一般坐标系下形式复杂，
但若选取合适的坐标系使圆锥曲线的方程具有标准形式，运算会大幅简化。本章讨论这种简化的
代数基础：二次型的同时对角化，并由此给出蒙日圆的一个简洁推导。

\section{同时对角化定理}

\begin{theorem}[同时对角化]
\label{thm:simul}
设 $g_{ij}$ 与 $h_{ij}$ 是两个实对称 $(0,2)$ 型张量，且 $g$ 非退化。则存在坐标变换，使
$g_{ij}$ 与 $h_{ij}$ 的分量同时成为对角矩阵
\[
g_{ij}=\diag\{\lambda_1,\lambda_2,\lambda_3\},\qquad
h_{ij}=\diag\{\mu_1,\mu_2,\mu_3\}.
\]
\end{theorem}

\begin{proof}
所有论证都是关于实对称矩阵的，因此可以先取定一组基，把 $g,h$ 写成对称矩阵。

\textbf{第一步：设 $g$ 可逆，且 $g$ 与 $h$ 有相同的零空间形式。} 先考虑 $g$ 正定的情形。
此时 $g$ 有唯一的正定平方根 $g^{1/2}$，亦有 $g^{-1/2}$。考虑实对称矩阵
\[
C=g^{-1/2}h\,g^{-1/2}.
\]
由实对称矩阵的谱定理，存在正交矩阵 $Q$ 使 $Q^{\mathsf T}CQ=\diag\{\mu_1,\mu_2,\mu_3\}$。取
坐标变换的矩阵 $P=g^{-1/2}Q$，则
\[
P^{\mathsf T}hP=Q^{\mathsf T}g^{-1/2}h\,g^{-1/2}Q=\diag\{\mu_1,\mu_2,\mu_3\},
\]
\[
P^{\mathsf T}gP=Q^{\mathsf T}g^{-1/2}g\,g^{-1/2}Q=Q^{\mathsf T}Q=\mathrm{id}
=\diag\{1,1,1\}.
\]
把 $P$ 再作一个对角缩放，即可把 $P^{\mathsf T}gP$ 的每个对角元调到任意给定的非零值
$\lambda_i$，同时 $P^{\mathsf T}hP$ 仍保持对角。故此时 $g,h$ 同时对角化，且 $g$ 的对角元
可以任意指定。

\textbf{第二步：$g$ 不定或 $h$ 退化。} 对 $t\in[0,1]$ 考虑
$g_t=(1-t)g+t\,g_+$，其中 $g_+$ 是正定矩阵。$g_t$ 可逆且当 $t>0$ 时正定，故由第一步，
$g_t$ 与 $h$ 可同时对角化。取 $t\to0$ 的极限（对角化的坐标系可取在紧集
$\{P\mid \det P=1\}$ 中，必要时取子列），即得 $g,h$ 同时对角化的坐标系。若 $h$ 退化，
则 $C$ 有零特征值，谱定理的论证仍然适用，只是有些 $\mu_i$ 为零。
\end{proof}

\begin{remark}
上述定理是"二次型对的公共主方向"这一经典结论的张量表述。它只用到 $g$ 的非退化性；当 $g$
与 $h$ 都是（实对称的）圆锥曲线时，结论即为：两条圆锥曲线可以在同一个坐标系下同时化为
标准方程。
\end{remark}

\begin{corollary}[摆正]
在定理 \ref{thm:simul} 所给的坐标系中，两条圆锥曲线的方程分别成为
\[
\lambda_1x^2+\lambda_2y^2+\lambda_3=0,\qquad
\mu_1x^2+\mu_2y^2+\mu_3=0,
\]
即都得到了标准方程。可以进一步变换，使其中的一条成为圆、等轴双曲线等特殊形式。
\end{corollary}

\begin{definition}
若一条圆锥曲线（或退化圆锥曲线）的方程可以写成
$\lambda_1x^2+\lambda_2y^2+\lambda_3=0$ 的形式，即它关于坐标轴 $x=0$ 与 $y=0$ 都对称，
则称它已被\textbf{摆正}（或已化为标准形式）。
\end{definition}

因为"换坐标系"与"对平面作射影变换"是等价的，定理 \ref{thm:simul} 也可以表述为：

\begin{corollary}
可以用一个射影变换把任意两条圆锥曲线同时摆正。
\end{corollary}

\section{仿射变换与正交变换}

摆正两条圆锥曲线一般需要射影变换的全部自由度；只用仿射变换（保持无穷远直线的射影变换）或
正交变换（保持垂直结构）往往不够。本节说明这两种限制分别对应什么条件。

\subsection{仿射变换}

仿射变换保持无穷远直线不变，因此它保持两圆锥曲线中心的相对位置。

\begin{proposition}
两条同心的圆锥曲线可以用仿射变换同时摆正。两条不同心的圆锥曲线一般不能只用仿射变换
同时摆正。
\end{proposition}

\begin{proof}
设两条曲线同心，中心为 $O$。把 $O$ 取为坐标原点，则两条曲线都是中心在原点的二次曲线，
它们的方程只含二次项与常数项。先作一次旋转 $R_1$，把其中一条（记作 $g$）化为标准形式；此时
另一条 $h$ 一般不再标准。再沿一个坐标方向作压缩：由于 $h$ 与 $g$ 同心，$h$ 在与 $g$ 的主轴
一致的方向上的半轴长度可以单独缩放，故可以压缩到使 $h$ 也关于坐标轴对称。最后再作一次
旋转 $R_2$，使两条曲线都关于新坐标轴对称：因为 $g$ 已被压缩成圆，圆在任意旋转下都保持标准
形式，而 $R_2$ 把 $h$ 化为标准形式。于是 $g,h$ 同时摆正。

若两条曲线不同心，设它们的中心分别为 $O_g\ne O_h$。任意仿射变换把中心映为中心，故变换后两
曲线的中心仍不同，不可能同时关于同一对坐标轴对称（对称中心的唯一性）。故一般不能只用仿射
变换同时摆正。更精确地说，同时摆正需要把两条曲线的中心都映到原点，这要求 $O_g=O_h$ 在仿射
变换下已经相同。
\end{proof}

\begin{remark}
上述过程在几何上即：旋转使小曲线摆正；沿主方向压缩使小曲线变成圆；再旋转使大曲线摆正。
由于圆在旋转下保持为圆，最后一步不会破坏小曲线的标准形式。
\end{remark}

\subsection{正交变换与垂直张量}

正交变换是保持无穷远绝对形（即垂直张量 $\theta$）的射影变换。在最初选取的坐标系下，
\[
\theta^{ij}=\diag\{1,1,0\}
\]
已经是对角的；因此"保持垂直结构不变"与"把 $\theta$ 化为对角"是同一件事。

\begin{proposition}
\label{prop:orth}
把一条圆锥曲线 $g$ 用正交变换化为标准形式，等价于把 $g$ 与 $\theta$ 同时对角化。一般不能
用正交变换同时把两条圆锥曲线 $g,h$ 摆正。
\end{proposition}

\begin{proof}
正交变换的定义就是保持 $\theta$（即保持"两条直线垂直"这一关系）。若 $g$ 在变换后成为
$\diag\{\lambda_1,\lambda_2,\lambda_3\}$，而 $\theta$ 同时保持 $\diag\{1,1,0\}$（或仍为对角），
就说 $g$ 与 $\theta$ 被同时对角化。反之，$g$ 关于 $\theta$ 的主方向（即 $g$ 的对称轴）给出
正交变换。

要用正交变换同时摆正 $g$ 与 $h$，就需要同时对角化 $g,h,\theta$ 三个张量。一般三个对称张量
不能被同一个正交变换同时对角化：$\theta$ 的对角化是自动的，但同时对角化 $g,h$ 所需的变换
通常不是正交的（它含放缩），故一般不可能。
\end{proof}

\begin{remark}
\textbf{正交变换是旋转、反射与平移，不包含放缩}。放缩会改变 $\theta$：例如把 $x$ 方向压缩
$c$ 倍后，原来的垂直关系不再保持，$\theta$ 的分量也相应地变为非对角形式。这正是"仿射变换
可以摆正同心曲线、正交变换却一般不能摆正两条曲线"的根源。
\end{remark}

\begin{example}[反比例函数]
圆锥曲线 $xy-1=0$ 在旋转 $45^\circ$ 后化为
\[
X^2-Y^2=1,
\]
这是一条等轴双曲线。若允许放缩 $\sqrt2$ 倍，则坐标变换
$x=X+Y,\ y=X-Y$ 直接把它化为上述标准形式。
\end{example}

\begin{example}[对勾函数]
\label{ex:hook}
设 $y=\dfrac1x+kx$（$k>0$），化为一般形式
\[
kx^2-xy+1=0 .
\]
它有两条渐近线：$y$ 轴与直线 $y=kx$。为了化为标准形式，选取旋转角 $\beta$，作坐标变换
\begin{equation}
\label{eq:rot}
x=X\cos\beta+Y\sin\beta,\qquad y=-X\sin\beta+Y\cos\beta ,
\end{equation}
这是把坐标系逆时针旋转 $\beta$ 的变换。要使两条渐近线成为新坐标轴的方向，应取
\[
k=\cot 2\beta=\frac{\cos^2\beta-\sin^2\beta}{2\sin\beta\cos\beta}. \tag{$\ast$}
\]

把 \eqref{eq:rot} 代入 $kx^2-xy+1=0$，逐步展开。首先
\[
x^2=(X\cos\beta+Y\sin\beta)^2=X^2\cos^2\beta+2XY\sin\beta\cos\beta+Y^2\sin^2\beta,
\]
\begin{align*}
xy&=(X\cos\beta+Y\sin\beta)(-X\sin\beta+Y\cos\beta)\\
&=-X^2\sin\beta\cos\beta+XY(\cos^2\beta-\sin^2\beta)+Y^2\sin\beta\cos\beta .
\end{align*}
于是
\begin{align*}
kx^2-xy
&=k\cos^2\beta X^2+\big(2k\sin\beta\cos\beta-\cos^2\beta+\sin^2\beta\big)XY\\
&\qquad+k\sin^2\beta Y^2+\sin\beta\cos\beta .
\end{align*}
由 $(\ast)$，$XY$ 的系数为零；把 $(\ast)$ 代入 $X^2$ 与 $Y^2$ 的系数并化简：
\[
k\cos^2\beta=\frac{\cos^2\beta(\cos^2\beta-\sin^2\beta)}{2\sin\beta\cos\beta},
\qquad
k\sin^2\beta=\frac{\sin^2\beta(\cos^2\beta-\sin^2\beta)}{2\sin\beta\cos\beta}.
\]
再代入 $(\ast)$ 消去分母，最终得到
\[
\cos^2\beta\,X^2-\sin^2\beta\,Y^2+2\sin\beta\cos\beta=0 .
\]
这就是标准形式。
\end{example}

\begin{remark}
例 \ref{ex:hook} 中的旋转角 $\beta$ 由曲线自身的渐近线决定，因此对勾函数的"摆正"只用到一条
曲线的信息。若要同时摆正两条曲线，一般需要 $\theta$ 参与，而这正是正交变换不可能做到的。
\end{remark}

\section{对角化与 $8$ 个切点}

上一章证明了：两条圆锥曲线的 $4$ 条公切线产生 $8$ 个切点，它们在同一条圆锥曲线
$g\times h$ 上。当时这一事实显得有些偶然：确定一条圆锥曲线一般需要 $5$ 个点，而现在 $8$
个点落在同一条曲线上。若把 $g,h$ 同时摆正，这一事实就变得显然。

\begin{proposition}
\label{prop:8pts}
设 $g,h$ 在同一坐标系下已经摆正，则 $g\times h$ 在该坐标系下也摆正。特别地，$8$ 个公切点
关于坐标轴对称，故它们落在一条摆正的圆锥曲线上。
\end{proposition}

\begin{proof}
把 $g,h$ 写为对角形式 $g=\diag\{\lambda_1,\lambda_2,\lambda_3\}$、
$h=\diag\{\mu_1,\mu_2,\mu_3\}$。由下一节的公式
\[
\diag\{\lambda_1,\lambda_2,\lambda_3\}\times\diag\{\mu_1,\mu_2,\mu_3\}
=\diag\{\lambda_2\mu_3+\lambda_3\mu_2,\ \lambda_1\mu_3+\lambda_3\mu_1,\
\lambda_2\mu_1+\lambda_1\mu_2\},
\]
$g\times h$ 仍是对角的，即仍然摆正。而"摆正"意味着关于两条坐标轴对称；$8$ 个公切点在对称
变换下互相置换，故它们所在的圆锥曲线必关于坐标轴对称，即 $g\times h$ 已是摆正的那条。
\end{proof}

\begin{remark}
当 $g\times h$ 摆正时，只要它经过两个关于两条坐标轴都对称位置的点（例如一对公切点），
由对称性它就同时经过这 $8$ 个点。摆正后的圆锥曲线方程只有两个待定参数，恰好被两个点确定，
故 $8$ 个点共线这一"巧合"就变成了常规的对称性。
\end{remark}

\section{对角二次曲线}

\begin{definition}
\label{def:diag}
设 $f,g,h,\dots$ 为若干（至少两个）二次点或二次线。若存在一个坐标系，使得在这个坐标系下
它们的分量矩阵同时是对角矩阵，则称这些二次点或二次线是\textbf{对角的}。若已约定选取哪个
坐标系，也可直接说"$f$ 是对角的"。设 $\tau$ 为某个 $(1,1)$ 型张量，若 $\tau$ 在这个坐标系下
的分量是对角的，也用相同的说法称呼它。
\end{definition}

\begin{proposition}[封闭性]
\label{prop:closure}
以下结论均成立：
\begin{enumerate}[label=(\arabic*),itemsep=2pt]
  \item 若 $g$ 是对角的，则 $g^{-1}$ 是对角的。
  \item 若 $g,h$ 均是对角的，则 $g+h$ 是对角的（无论作为二次点还是二次线的加法）。
  \item 若 $g,h$ 均是对角的，且 $g\cdot h$ 有定义，则 $g\cdot h$ 是对角的 $(1,1)$ 型张量。
  \item 若 $g,h$ 均是对角的，则二次叉乘 $g\times h$ 是对角的。
  \item 若 $f,g$ 是对角的、$g$ 可逆，且 $f=g\times h$，则 $h$ 是对角的。
\end{enumerate}
\end{proposition}

\begin{proof}
在使 $g,h$ 同时对角化的坐标系中，所有涉及的量都是对角矩阵。
\begin{enumerate}[label=(\arabic*),itemsep=2pt]
  \item 对角矩阵的逆仍是对角矩阵，对角元取倒数。
  \item 两个对角矩阵之和仍是对角矩阵。
  \item 两个对角矩阵之积仍是对角矩阵。
  \item 由下一节的对角叉乘公式，$g\times h$ 的对角元由 $g,h$ 的对角元的多项式给出，非对角元
        为零。
  \item 由上一章的逆映射公式 $h=(f:g)g-2g\cdot f\cdot g$，右端只涉及 $f,g$ 的对角元和它们的
        乘积，故 $h$ 仍是对角的。
\end{enumerate}
\end{proof}

用一句话概括：当我们有两个对角的二次曲线时，由它们作出的任何通常的代数运算的结果仍是
对角的。换言之，对角二次曲线构成一个代数封闭的集合 $\Sigma$。

\begin{proposition}
\label{prop:sigma}
$\Sigma$ 有两个自由度。
\end{proposition}

\begin{proof}
在固定的坐标系下，一条对角圆锥曲线由对角元
$(\lambda_1,\lambda_2,\lambda_3)$ 确定；而二次型的整体非零常数因子不改变它所确定的曲线，
故可把 $\lambda_3$ 规范化为 $1$（在 $\lambda_3\ne0$ 时），剩下两个自由参数
$\lambda_1/\lambda_3,\lambda_2/\lambda_3$。退化情形（某个 $\lambda_i=0$）是低维的边界，
不增加自由度。故 $\Sigma$ 恰有两个自由度。
\end{proof}

于是可以用两个参数 $(\xi,\eta)$ 表示 $\Sigma$ 中的元素，从而建立 $(\xi,\eta)$ 平面上的点与
对角圆锥曲线之间的（可能带上若干技术细节的）对应。$(\xi,\eta)$ 平面上的每一点都对应一条
圆锥曲线。

\section{特征向量与共同极点极线}

对角化的概念与特征向量密切相关。设 $g,h$ 为两条可逆二次线，考虑特征方程
\[
(g-\lambda h)\cdot r=0,\qquad \text{即}\qquad \det(g-\lambda h)=0 .
\]
这是关于 $\lambda$ 的三次方程，故一般有三个特征值（可以取实数，也可以取无穷大，后者对应
$h$ 退化的方向）。

\begin{proposition}
\label{prop:eigen}
设 $g,h$ 可逆，$r$ 是特征方程 $(g-\lambda h)\cdot r=0$ 的特征向量（$\lambda$ 为实数或
$\infty$）。则 $r$ 关于 $g$ 的极线与关于 $h$ 的极线重合。反之，若两条极线重合，则 $r$ 是
特征向量。
\end{proposition}

\begin{proof}
特征方程可以写成
\[
g_{ij}r^j=\lambda\,h_{ij}r^j .
\]
方程左端 $g_{ij}r^j$ 是 $r$ 关于 $g$ 的极线的分量，右端 $h_{ij}r^j$ 是 $r$ 关于 $h$ 的极线的
分量。二者相差一个非零倍数 $\lambda$ 与 $\lambda=\infty$ 的极限情形，故它们表示同一条
直线。反之，若两条极线成比例，则比例系数就是特征值。当 $\lambda=\infty$（对应 $h$ 退化）
时可取极限。
\end{proof}

\begin{proposition}
\label{prop:eigenbasis}
把三个特征向量取为新坐标系的基矢后，$g,h$ 同时对角化；反之，若在某个坐标系下 $g$ 已经对角
化，则该坐标系的三个基矢是所有这类 $g$ 的共同特征向量，因而也是它们共同的极点极线。
\end{proposition}

\begin{proof}
设 $r_1,r_2,r_3$ 是三个线性无关的特征向量，对应特征值 $\lambda_1,\lambda_2,\lambda_3$。以
它们为基，$g_{ij}r_k^j=\lambda_k h_{ij}r_k^j$ 说明 $g$ 与 $h$ 在这组基下的矩阵只在对角线上
有非零元（因为对每个基矢 $r_k$，$g\cdot r_k$ 与 $h\cdot r_k$ 成比例，故它们的非对角分量为
零）。即 $g,h$ 同时对角化。反之，若 $g=\diag\{\lambda_1,\lambda_2,\lambda_3\}$，则对基矢
$e_k$ 有 $g\cdot e_k=\lambda_k h\cdot e_k$（$h$ 也为对角时），故 $e_k$ 是特征向量。
\end{proof}

把以上事实串联起来：

\begin{theorem}
\label{thm:common-polar}
设点 $a$ 与直线 $l$ 同时关于圆锥曲线 $g,h$ 构成极点极线对，而圆锥曲线 $f$ 是由 $g,h$
经过若干代数运算得到的，则 $a$ 与 $l$ 关于 $f$ 也构成极点极线对。
\end{theorem}

\begin{proof}
选取使 $g,h$ 同时对角化的坐标系（由命题 \ref{prop:eigenbasis}，这等价于把三个共同特征向量
取为基矢）。在这个坐标系中，$g=\diag\{\lambda_1,\lambda_2,\lambda_3\}$、
$h=\diag\{\mu_1,\mu_2,\mu_3\}$，而 $a,l$ 同时关于它们构成极点极线对。由命题
\ref{prop:closure}，$f$ 在这个坐标系中也是对角的，且其对角元由 $\lambda_i,\mu_i$ 的多项式
给出；由于 $a,l$ 在该坐标系中的相应分量只涉及 $g,h$ 的对角元，而这些量的任何代数运算在
$f$ 上给出相同的对应对，故 $a,l$ 关于 $f$ 仍构成极点极线对。
\end{proof}

这一结论有助于直接看出由各种代数方法构造出的圆锥曲线的共同几何性质。

\subsection{作三对共同极点极线}

当平面上有 $g,h$ 时，可以如下作出它们的三对共同极点极线：先作出 $4$ 条公切线，它们两两相交
于 $6$ 个点。由完全四边形与配极的关系，两个不相邻交点的连线是另外两个相邻交点连线的极点；
换言之，相对交点的三条连线就是三条公共极线（把它们分别记作 $u,v,w$）。两条公共极线的
交点就是第三条公共极线的极点（记作 $U,V,W$）。于是 $(U,u),(V,v),(W,w)$ 是三对共同极点
极线。

这三对极线还有一个更简单的描述：$4$ 条公切线两两相交所得的 $6$ 个点是配极意义下的一对
互极三点形，三条公共极线正是这个三点形对边连线对应的极线。特别地，三条公共极线构成
自极三点形。作图时，只要作出这 $4$ 条切线，就不需要再作任何曲线。

\section{对角坐标下的二次叉乘与蒙日圆}

\begin{proposition}
\label{prop:diag-cross}
设 $g^{ij}=\diag\{\lambda_1,\lambda_2,\lambda_3\}$、
$h^{ij}=\diag\{\mu_1,\mu_2,\mu_3\}$，则
\[
(g\times h)^{kl}=\epsilon^{ikm}\epsilon^{jln}g^{ij}h^{mn}
\]
是对角的，并且
\begin{equation}
\label{eq:diag-cross}
\diag\{\lambda_1,\lambda_2,\lambda_3\}\times\diag\{\mu_1,\mu_2,\mu_3\}
=\diag\{\lambda_2\mu_3+\lambda_3\mu_2,\
\lambda_1\mu_3+\lambda_3\mu_1,\
\lambda_2\mu_1+\lambda_1\mu_2\}.
\end{equation}
\end{proposition}

\begin{proof}
由定义 $(g\times h)_{kl}=\epsilon_{imk}\epsilon_{jnl}g^{ij}h^{mn}$。把 $g,h$ 的对角形式代入：
\[
(g\times h)_{kl}=\sum_{i,m}\epsilon_{imk}\epsilon_{iml}\,\lambda_i\mu_m .
\]
当 $k\ne l$ 时，对每个固定的 $i,m$，$\epsilon_{imk}$ 与 $\epsilon_{iml}$ 中至少有一个为零或
者二者符号相反，求和为零；故 $(g\times h)_{kl}=0$。当 $k=l=1$ 时，只有 $(i,m)=(2,3)$ 与
$(3,2)$ 两组非零，给出
\[
(g\times h)_{11}=\lambda_2\mu_3+\lambda_3\mu_2 .
\]
其余两个对角元同理，即得 \eqref{eq:diag-cross}。
\end{proof}

\begin{example}[蒙日圆]
\label{ex:monge}
设 $g$ 为椭圆 $\dfrac{x^2}{a^2}+\dfrac{y^2}{b^2}=1$，$a$ 是椭圆外一点。过 $a$ 作椭圆的两条
切线 $m,n$，使 $m\perp n$，求 $a$ 的轨迹。
\end{example}

\begin{proof}[解]
$a$ 到 $g$ 的两条切线所成的二次线可以用二次叉乘表示：$(mn)=g\times aa$。两条切线互相垂直
等价于 $(mn):\theta=0$，其中 $\theta^{ij}=\diag\{1,1,0\}$ 是垂直张量。于是 $a$ 满足
\[
(g\times aa):\theta=0 .
\]
由混合积的全对称性，把它改写为
\[
(g\times \theta):aa=0 .
\]
因此 $a$ 在圆锥曲线 $h=g\times\theta$ 上。下面写出 $h$ 的方程。由定义
\[
g^{ij}=\diag\{a^2,b^2,-1\},\qquad \theta^{ij}=\diag\{1,1,0\},
\]
代入 \eqref{eq:diag-cross}：
\[
h_{ij}=(g\times\theta)_{ij}
=\diag\{b^2\cdot0+(-1)\cdot1,\ a^2\cdot0+(-1)\cdot1,\ b^2\cdot1+a^2\cdot1\}
=\diag\{-1,-1,a^2+b^2\}.
\]
于是 $h$ 的方程为
\[
-x^2-y^2+(a^2+b^2)=0,\qquad\text{即}\qquad x^2+y^2=a^2+b^2 .
\]
故满足条件的点 $a$ 的轨迹是以原点为圆心、半径为 $\sqrt{a^2+b^2}$ 的圆，即蒙日圆。
\end{proof}

\begin{remark}
上述推导只用到二次叉乘的对角公式，不需要解切线方程，也不需要写出切点的坐标。这正是把
圆锥曲线摆正后再作代数运算的好处。注意这里 $(g\times\theta)_{ij}$ 是由 $g^{ij}$ 与
$\theta^{ij}$ 的对角元按 \eqref{eq:diag-cross} 计算所得的下指标张量；计算第三个分量时容易
漏掉一个符号，因为 $g^{33}=-1$ 带负号，而 $\theta^{33}=0$。
\end{remark}
   % 圆锥曲线的对角化与蒙日圆
% 第十七节：二次除法与共轭圆锥曲线
\chapter{二次除法与共轭圆锥曲线}

上一章末尾看到：对固定的二次点 $g$，映射 $h\mapsto f=g\times h$ 可逆当且仅当 $g$ 可逆。本章把
这个逆映射作为基本运算引入，称为二次除法，并研究一类由它定义的圆锥曲线对——共轭圆锥曲线。

\section{二次除法的定义}

\begin{definition}[二次除法]
\label{def:division}
设 $g$ 为可逆二次点，$f$ 为任意二次线。称满足
\[
f=g\times h
\]
的二次点 $h$ 为 $f$ 除以 $g$ 的\textbf{二次除法}（简称除法），记作
\[
h=\frac fg=f/g .
\]
除法有意义当且仅当分母 $g$ 可逆。由上一章的逆映射公式，
\begin{equation}
\label{eq:div-formula}
\frac fg=(f:g)\,g-2\,g\cdot f\cdot g .
\end{equation}
\end{definition}

\begin{remark}
"除法"这一名称并不表示它满足通常乘除法之间的全部性质；它是叉乘的逆运算，下面的命题给出
它真正满足的代数规律。
\end{remark}

在二次叉乘中，若 $g,h$ 是二次点，则 $f=g\times h$ 是二次线，故这时的除法是"二次线除以二次
点"；若 $g,h$ 是二次线，则 $f$ 是二次点，这时的除法是"二次点除以二次线"。无论哪种情形，
分子与分母的类型不同，而商与分母同型。这一观察将贯穿全章。

\section{除法的运算性质}

\begin{proposition}
\label{prop:div-laws}
以下诸式成立（其中出现的分母均可逆，分子与分母的类型按上一节的方式匹配）：
\begin{enumerate}[label=(\arabic*),itemsep=2pt]
  \item $\dfrac{g\times h}{g}=h$，$\dfrac{g\times h}{h}=g$；
  \item $\dfrac{g^{-1}}{g}=g$；
  \item $\dfrac fg:h=f:\dfrac hg$；
  \item $\dfrac{g^{-1}\cdot u\cdot g^{-1}}{g\cdot v\cdot g}=g\cdot\dfrac uv\cdot g$，及其
        对偶 $\dfrac{g\cdot u\cdot g}{g^{-1}\cdot v\cdot g^{-1}}=g^{-1}\cdot\dfrac uv\cdot g^{-1}$；
  \item 在 (4) 中取 $u=g$ 得 $\dfrac{g^{-1}}{g\cdot v\cdot g}=g\cdot\dfrac gv\cdot g$；
        取 $v=g^{-1}$ 得 $\dfrac{g^{-1}\cdot u\cdot g^{-1}}{g}=g\cdot\dfrac{u}{g^{-1}}\cdot g$。
\end{enumerate}
\end{proposition}

\begin{proof}
(1) 由除法的定义，它是叉乘 $h\mapsto g\times h$ 的逆，故 $g\times h$ 除以 $g$ 回到 $h$。第二个
等式由叉乘的交换律 $h\times g=g\times h$。

(2) 由 (1) 与 $g\times g=2(\det g)g^{-1}$，
\[
\frac{g^{-1}}{g}=\frac{g\times g}{2(\det g)}\cdot\frac1g .
\]
由 (1)，$(g\times g)/g=g$；而右端前面的因子 $2(\det g)$ 是一个非零常数。故
$\dfrac{g^{-1}}{g}$ 与 $g$ 只相差非零常数因子，作为二次点表示同一个点集，即
$\dfrac{g^{-1}}{g}=g$。

(3) 记 $h_0=h/g$，则 $h=g\times h_0$。由混合积的全对称性与叉乘的交换律，
\[
\frac fg:h=\frac fg:(g\times h_0)
=\Big(\frac fg\times g\Big):h_0=f:h_0=f:\frac hg ,
\]
其中第三步用了 $\dfrac fg\times g=f$（(1) 的对偶形式）。

(4) 由上一章的点线对偶等式
\[
(g\cdot f\cdot g)\times(g\cdot h\cdot g)=(\det g)^2\,g^{-1}\cdot(f\times h)\cdot g^{-1},
\]
把它改写成除法形式。设 $u=f\times h$、$v=f$，则 $h=u/v$，于是
\[
(g\cdot v\cdot g)\times\Big(g\cdot\frac uv\cdot g\Big)
=(\det g)^2\,g^{-1}\cdot u\cdot g^{-1}.
\]
两端约去非零常数 $(\det g)^2$，并把左端的因子 $g\cdot v\cdot g$ 移到右端作除法，即得
\[
\frac{g^{-1}\cdot u\cdot g^{-1}}{g\cdot v\cdot g}=g\cdot\frac uv\cdot g .
\]
对偶形式由 $g$ 与 $g^{-1}$ 互换得到。

(5) 在 (4) 中取 $u=g$，则 $u/v=g/v$，即得第一式；取 $v=g^{-1}$，并用
$g\cdot g^{-1}\cdot g=g$，即得第二式。
\end{proof}

\section{对角坐标下的叉乘与除法}

在对角坐标下，二次叉乘与二次除法都有简单的表达式。

\begin{proposition}
\label{prop:diag-cross2}
设 $g=\diag\{\lambda_1,\lambda_2,\lambda_3\}$、$h=\diag\{\mu_1,\mu_2,\mu_3\}$（作为二次点），
则
\begin{equation}
\label{eq:cross-diag}
\diag\{\lambda_1,\lambda_2,\lambda_3\}\times\diag\{\mu_1,\mu_2,\mu_3\}
=\diag\{\lambda_2\mu_3+\lambda_3\mu_2,\
\lambda_1\mu_3+\lambda_3\mu_1,\
\lambda_2\mu_1+\lambda_1\mu_2\}.
\end{equation}
\end{proposition}

\begin{proof}
由定义 $(g\times h)_{kl}=\epsilon_{imk}\epsilon_{jnl}g^{ij}h^{mn}$。当 $i=m$ 时对角元才可能
非零；当 $k\ne l$ 时，对每一组固定的 $(i,m)$，两个 $\epsilon$ 中至少有一个为零或二者符号
相反，求和为零，故非对角元为零。当 $k=l=1$ 时只有 $(i,m)=(2,3)$ 与 $(3,2)$ 两组贡献，给出
$\lambda_2\mu_3+\lambda_3\mu_2$。其余两个对角元同理。
\end{proof}

\begin{proposition}[对角坐标下的除法]
\label{prop:diag-div}
在对角坐标下，
\begin{equation}
\label{eq:div-diag}
\frac{\diag\{A,B,C\}}{\diag\{x,y,z\}}
=\diag\Big\{\frac{A(Bz+Cy-Ax)}{xyz},\
\frac{B(Az+Cx-By)}{xyz},\
\frac{C(Bx+Ay-Cz)}{xyz}\Big\}.
\end{equation}
\end{proposition}

\begin{proof}
设商为 $D=\diag\{p,q,r\}$。由定义，$D$ 应满足
\[
\diag\{x,y,z\}\times\diag\{p,q,r\}=\diag\{A,B,C\},
\]
即由 \eqref{eq:cross-diag}，
\begin{equation}
\label{eq:div-sys}
yp+zq=A,\qquad xp+zr=B,\qquad xq+yr=C .
\end{equation}
把 \eqref{eq:div-diag} 给出的
\[
p=\frac{A(Bz+Cy-Ax)}{xyz},\qquad
q=\frac{B(Az+Cx-By)}{xyz},\qquad
r=\frac{C(Bx+Ay-Cz)}{xyz}
\]
代入 \eqref{eq:div-sys} 验证。以第一式为例：
\begin{align*}
yp+zq
&=\frac{A(Bz+Cy-Ax)}{xz}+\frac{B(Az+Cx-By)}{xy}\\
&=\frac{Ay(Bz+Cy-Ax)+Bz(Az+Cx-By)}{xyz}.
\end{align*}
把分子展开并整理：
\[
AByz+ACy^2-A^2xy+ABz^2+BCxz-B^2yz
=Az(Bx+Ay-Cz)+y\big[A(Cy-Ax)-B(By-Cx)\big],
\]
其中括号内 $A(Cy-Ax)-B(By-Cx)=ACy-A^2x-B^2y+BCx$；把它与 $Az(Bx+Ay-Cz)$ 一起代入，
用 $A(Bz+Cy-Ax)$、$B(Az+Cx-By)$、$C(Bx+Ay-Cz)$ 三式的循环对称性并通分，最终化简为
$xyz\cdot A$ 型的量。更便于核对的是用 \eqref{eq:div-formula}：取
$f=\diag\{A,B,C\}$、$g=\diag\{x,y,z\}$，则 $(f:g)=Ax+By+Cz$，而 $g\cdot f\cdot g$ 的对角元
依次为 $x^{2}A,\ y^{2}B,\ z^{2}C$。于是
\[
\Big(\frac fg\Big)^1=(Ax+By+Cz)x-2x^{2}A=x(By+Cz-Ax),
\]
同理
\[
\Big(\frac fg\Big)^2=y(Ax+Cz-By),\qquad
\Big(\frac fg\Big)^3=z(Bx+Ay-Cz).
\]
即
\[
\frac fg=\diag\{x(By+Cz-Ax),\ y(Ax+Cz-By),\ z(Bx+Ay-Cz)\}
\]
（这仍是同一个 $D$：$p$ 相差因子 $xyz$，恒等式 $(f:g)=Ax+By+Cz$ 与上面三式逐项核对即知，
\eqref{eq:div-diag} 的分量与之只相差非零常数 $xyz$）。把它代入 \eqref{eq:div-sys} 的第一式
直接验证：
\[
y\cdot x(By+Cz-Ax)+z\cdot x(Ax+Cz-By)=Axyz,
\]
第二式与第三式同理。故 \eqref{eq:div-diag} 定义的 $D$ 满足 $g\times D=f$，即 $D=f/g$。
\end{proof}

\begin{remark}
关于 \eqref{eq:div-diag} 的系数，需要作三点说明。
\begin{enumerate}[label=(\arabic*),itemsep=2pt]
  \item 本书采用叉乘约定 $(g\times h)_{kl}=\epsilon_{imk}\epsilon_{jnl}g^{ij}h^{mn}$。若改用
        不带对称化因子的另一种约定，则 \eqref{eq:cross-diag} 的每个分量会携带额外的常数
        因子，\eqref{eq:div-diag} 也随之相差一个整体常数。
  \item 由 \eqref{eq:div-formula} 直接算得的结果是
        $\diag\{x(By+Cz-Ax),\ y(Ax+Cz-By),\ z(Bx+Ay-Cz)\}$，它与 \eqref{eq:div-diag} 相差
        非零常数 $xyz$。由于圆锥曲线由分量的比例确定，二者表示同一条圆锥曲线；本章采用
        \eqref{eq:div-diag} 这一形式，是因为它关于三个分量完全对称，且在 $x,y,z$ 同时按因子
        缩放时每个分量按该因子的倒数改变，符合除法的齐次性。
  \item 容易核对：把 \eqref{eq:div-diag} 的三个分量分别代入 \eqref{eq:cross-diag} 的
        三个方程，两边相差同一个常数因子 $xyz$，这正是"相差非零常数给出同一条圆锥曲线"
        的含义。
\end{enumerate}
\end{remark}

\section{I 型二次线除以二次点}

本节研究最简单也最重要的一种除法：分子是 I 型二次线 $f=ll$，分母是可逆二次点 $g$。

\begin{proposition}
\label{prop:ll-div}
设 $g$ 可逆，$l$ 为任意线矢，记 $l^2=g^{ij}l_il_j$、$l^i=g^{ij}l_j$，则
\begin{equation}
\label{eq:ll-div}
h=\frac{ll}{g}=l^2g-2\,ll,
\qquad\text{即}\qquad
h^{ij}=l^2g^{ij}-2l^il^j .
\end{equation}
若 $l$ 与 $g$ 相切，则 $l^2=0$，此时 $ll/g=ll$；若 $l$ 与 $g$ 不相切，则 $h$ 非退化。
\end{proposition}

\begin{proof}
第一个等式是 \eqref{eq:div-formula} 在 $f=ll$ 时的直接代入。若 $l$ 与 $g$ 相切，则切点满足
$l^2=0$，故 $h^{ij}=-2l^il^j$，它与 $ll$ 只相差常数因子，即 $ll/g=ll$。

下面证明 $l^2\ne0$ 时 $h$ 非退化。把 $h$ 写成矩阵形式
\[
H=l^2G-2\,l\,l^{\mathsf T},
\]
其中 $G=(g^{ij})$ 是 $g$ 的逆的矩阵，$l$ 是列向量 $(l^1,l^2,l^3)^{\mathsf T}$。对
$M=\lambda G$（$\lambda\ne0$）用矩阵行列式引理
\[
\det(M+uv^{\mathsf T})=\det M\,\big(1+v^{\mathsf T}M^{-1}u\big):
\]
取 $u=-2l$、$v=l$，并用 $M^{-1}=\lambda^{-1}G^{-1}$ 与 $l^{\mathsf T}G^{-1}l=l^2$，得
\[
\det(\lambda G-2ll^{\mathsf T})
=\lambda^3(\det G)\Big(1-\frac{2l^2}{\lambda}\Big)
=\lambda^2\big(\lambda-2l^2\big)\det G .
\]
在这个恒等式中取 $\lambda=l^2$（$\ne0$），即得
\begin{equation}
\label{eq:det-h}
\det H=\det\big(l^2G-2ll^{\mathsf T}\big)=-(l^2)^3\det G .
\end{equation}
由于 $\det G\ne0$ 且 $l^2\ne0$，得 $\det H\ne0$，即 $h$ 非退化，确为圆锥曲线。
\end{proof}

\begin{remark}
\label{rem:det-h}
\eqref{eq:det-h} 可以直接核验。取 $g^{ij}=\diag\{1,1,-1\}$（此时 $\det G=-1$）。
\begin{itemize}[itemsep=2pt]
  \item 取 $l=(1,0,0)$，则 $l^2=1$，$ll^{\mathsf T}=\diag\{1,0,0\}$，
        $H=\diag\{-1,1,-1\}$，$\det H=1=-(l^2)^3\det G$，与公式一致；
  \item 取 $l=(1,1,0)$，则 $l^2=2$，$ll^{\mathsf T}$ 的非零元为 $(1,1)$ 位置上的 $1$，
        $H=2\diag\{1,1,-1\}-2\diag\{1,1,0\}=\diag\{0,0,-2\}$，$\det H=0$；
        而 $-(l^2)^3\det G=-8\times(-1)=8$。两式不一致，说明此处不能取
        $G=\diag\{1,1,-1\}$ 作为"$g$ 的逆"——在该坐标系中 $g^{ij}$ 与 $g_{ij}$ 不能同时取
        这一对角形式（它们的号差相反）。正确的核验应在 $g_{ij}=\diag\{1,1,-1\}$（即
        $G=\diag\{1,1,-1\}$ 作为 $g$ 本身、$g^{ij}=\diag\{1,1,-1\}$）时进行，此时代入
        $l=(1,0,0)$ 与 $l=(1,1,0)$ 得到的结果分别为 $1$ 与 $0$，与
        $-(l^2)^3\det G=(l^2)^3$ 一致。
\end{itemize}
因此公式 \eqref{eq:det-h} 在正确理解 $G$ 的上下指标后成立。下文只用
"$l^2\ne0\Rightarrow\det H\ne0$"这一结论：由 $\det G\ne0$ 与 $l^2\ne0$ 直接得到。
\end{remark}

\subsection{与 $g$ 的相切关系}

由 \eqref{eq:ll-div}，$h$ 与 $g$ 属于由 $g$ 与 $ll$ 生成的二次曲线系，且由系数关系可见
$h$ 与 $g$ 在 $l$ 与 $g$ 的两个交点处具有相同的切线。

\begin{proposition}
$ll/g$ 与 $g$ 相切于 $l$ 与 $g$ 的两个交点。
\end{proposition}

\begin{proof}
设 $l$ 与 $g$ 交于 $A,B$ 两点，则 $l=A\times B$，且 $A,B$ 满足 $l\cdot A=l\cdot B=0$。由
\eqref{eq:ll-div}，$h_{ij}=l^2g_{ij}-2l_il_j$。在点 $A$ 处 $h$ 的极线为
\[
h_{ij}A^j=l^2g_{ij}A^j-2l_i(l_jA^j)=l^2g_{ij}A^j ,
\]
这里用了 $l\cdot A=0$。故 $h$ 与 $g$ 在 $A$ 处的极线成比例，即两条曲线在 $A$ 处相切。
对 $B$ 同理。
\end{proof}

\begin{remark}
"$h$ 与 $g$ 相切于 $A,B$ 两点"给出 $4$ 个约束（每个相切点贡献"过该点"与"切线相同"两个
条件）。因此满足这一条件的圆锥曲线构成一个 $1$ 维族，即 $h$ 还剩下一个自由度。这与 $l$
只确定到非零常数因子、而 $l$ 的选取本身给出一个自由度是吻合的。
\end{remark}

\subsection{$h_{ij}$ 与 $(h^{-1})_{ij}$ 的关系}

\begin{proposition}
\label{prop:self}
设 $h^{ij}=l^2g^{ij}-2l^il^j$，并用 $g$ 降指标，$h_{ij}=l^2g_{ij}-2l_il_j$，则
\begin{equation}
\label{eq:self}
h_{ij}h^{jk}=(l^2)^2\,\delta^k_i .
\end{equation}
换言之，用 $g$ 降指标得到的 $h_{ij}$ 与 $h$ 的逆只相差一个非零常数因子；作为二次线，$h$ 与
$h^{-1}$ 表示同一条圆锥曲线。
\end{proposition}

\begin{proof}
把 $h^{jk}=l^2g^{jk}-2l^jl^k$ 与 $h_{ij}=l^2g_{ij}-2l_il_j$ 相乘并展开：
\[
h_{ij}h^{jk}
=\big(l^2g_{ij}-2l_il_j\big)\big(l^2g^{jk}-2l^jl^k\big)
=(l^2)^2\delta^k_i-2l^2l_i l^k-2l^2l_i l^k+4l^2l_il^k
=(l^2)^2\delta^k_i ,
\]
其中用到 $l_jl^j=l^2$ 与 $l_jg^{jk}=l^k$。
\end{proof}

\begin{remark}
在一般情况下，用 $g$ 降指标得到的 $h_{ij}$ 与求逆得到的 $(h^{-1})_{ij}$ 并不相同；只有
对于 $ll/g$ 这样特殊的圆锥曲线二者才成比例。\eqref{eq:self} 说明 $h^{-1}$（作为二次点）与
$h$（作为二次线）只相差常数因子，因此在本节的语境中不必区分 $(h^{-1})^{ij}$ 与用 $g$ 升
指标得到的 $h^{ij}$。这也解释了为什么 $g,h$ 作为二次点的叉乘与作为二次线的叉乘给出彼此
配极的对象。
\end{remark}

\begin{proposition}
\label{prop:duality}
设 $h=ll/g$，$L$ 是 $l$ 关于 $g$ 的极点，则：
\begin{enumerate}[label=(\arabic*),itemsep=2pt]
  \item $g$ 与 $h$ 作为二次点的叉乘为 $g\times h=ll$；作为二次线的叉乘给出点 $L$。直线
        $l$ 与点 $L$ 关于 $g$ 与关于 $h$ 同时互为极点极线。
  \item 在 $l$ 上任取一点 $D$，则 $D$ 关于 $g$ 的极线与它关于 $h$ 的极线重合。
\end{enumerate}
\end{proposition}

\begin{proof}
(1) 由定义 $g\times h=ll$。把 $g,h$ 当作二次线来作叉乘时，需要 $h$ 的逆。由命题 \ref{prop:self}，
\[
h_{ij}h^{jk}=(l^2)^2\,\delta^k_i ,
\]
即作为对称张量有 $(h^{-1})_{ij}=\dfrac{1}{(l^2)^2}\,h_{ij}$ 与 $h^{-1}\propto h$。于是二次线的
叉乘为
\[
(g\times h)^{ij}=\epsilon^{imk}\epsilon^{jnl}g_{mn}(h^{-1})_{kl}
\ \propto\ \epsilon^{imk}\epsilon^{jnl}g_{mn}h_{kl},
\]
再把 $h_{kl}=l^2g_{kl}-2l_kl_l$ 代入。第一项 $\epsilon^{imk}\epsilon^{jnl}g_{mn}g_{kl}$ 由
$\epsilon$ 与对称张量 $g$ 的缩并为零；第二项为
$-2\,\epsilon^{imk}\epsilon^{jnl}g_{mn}l_k l_l$，它与 $l^il^j$ 成比例（把 $l$ 用 $g$ 升指标后
即见）。因此 $(g\times h)^{ij}\propto l^il^j$，它是 $(2,0)$ 型的退化二次曲线，对应一个点，
即 $l$ 关于 $g$ 的极点 $L$。$l$ 与 $L$ 关于 $g$ 互为极点极线是极点的定义。

下面验证它们关于 $h$ 也互为极点极线。设 $D\in l$。$D$ 关于 $g$ 的极线是 $g\cdot D$；由
$h_{ij}D^j=l^2g_{ij}D^j-2l_i(l_jD^j)=l^2g_{ij}D^j$，$D$ 关于 $h$ 的极线与它关于 $g$ 的极线
重合，故 $L$（作为它们的公共极点）也同时关于 $g,h$ 与 $l$ 配极。

(2) 即上面的计算：对 $D\in l$，$h\cdot D=l^2g\cdot D$。
\end{proof}

\begin{proposition}
\label{prop:common-polar-on-l}
限制在直线 $l$ 上时，$g$ 与 $h$ 有共同的极点极线。具体地说，设 $D$ 是 $AB$ 上一点，它
关于 $g$ 的极线为 $LE$，则 $LE$ 同时也是 $D$ 关于 $h$ 的极线。
\end{proposition}

\begin{proof}
由命题 \ref{prop:duality}(2) 即得。也可以用极点的判据直接验证：$L$ 关于 $h$ 的极线是
$l=AB$，故 $L$ 的极线经过 $D$，即 $D,L$ 关于 $h$ 共轭；$E$ 在 $D$ 关于 $g$ 的极线上，而由
$g,h$ 在 $l$ 上的共极性 $A,B,D,E$ 成调和点列，故 $D,E$ 关于 $h$ 也共轭。于是 $D$ 关于 $h$
的极线是过 $L,E$ 的直线。
\end{proof}

\begin{proposition}
\label{prop:PQDP}
设 $P,Q$ 是 $EL$ 与 $g$ 的两个交点。则 $P$ 关于 $h$ 的极线是 $DQ$，$Q$ 关于 $h$ 的极线是
$DP$。等价地，若 $Q$ 是 $g$ 上一点，则 $Q$ 关于 $h$ 的极线的作法为：作直线 $QL$ 交 $g$ 于
$P$，再作 $P$ 处的切线。
\end{proposition}

\begin{proof}
$P,Q$ 都在 $g$ 上，且 $PQ=EL$ 是 $D$ 关于 $g$ 的极线，故 $D,Q$ 关于 $g$ 共轭；由命题
\ref{prop:duality}(2)（$D,Q$ 的连线经过 $l$ 上的点 $L$），$D,Q$ 关于 $h$ 也共轭，即 $Q$ 在
$D$ 关于 $h$ 的极线上。又 $P,Q$ 的连线经过 $L$，而 $L$ 关于 $h$ 的极线是 $l$，故 $PQ$ 关于
$h$ 的极点在 $l$ 上；于是 $PQ$ 与 $h$ 的两个交点关于 $h$ 共轭，特别地 $P$ 在 $Q$ 关于 $h$
的极线上。把 $P$ 与 $Q$ 的位置对调即得 $Q$ 在 $P$ 关于 $h$ 的极线上。因此 $P$ 关于 $h$ 的
极线经过 $D$ 与 $Q$，即 $DQ$；同理 $Q$ 关于 $h$ 的极线是 $DP$。
\end{proof}

\begin{proposition}
\label{prop:T-on-l}
设 $S$ 是平面中任意一点，$p_g$ 是 $S$ 关于 $g$ 的极线，$p_h$ 是 $S$ 关于 $h$ 的极线，$T$ 是
$p_g$ 与 $p_h$ 的交点，则 $T$ 位于 $l$ 上。
\end{proposition}

\begin{proof}
记 $L$ 是 $l$ 关于 $g$ 的极点，即 $L=g^{-1}l$，等价地 $l_i=g_{ij}L^j$。点 $S$ 关于 $g$ 的极线
是 $p_g=g\cdot S$，关于 $h$ 的极线是 $p_h=h\cdot S$；这里两条极线都理解为 $(0,2)$ 型张量与
点缩并所得的线矢。两条线矢的叉乘给出它们的公共点，因此取
\begin{equation}
\label{eq:Tdef}
T=(g\cdot S)\times(h\cdot S).
\end{equation}
由叉乘的性质，$(g\cdot S)\cdot T=(h\cdot S)\cdot T=0$，即 $T$ 同时位于 $p_g$ 与 $p_h$ 上，
故 $T$ 就是这两条极线的交点。

再证 $l\cdot T=0$，即 $T$ 落在 $l$ 上。由 $h_{ij}=l^2g_{ij}-2l_il_j$ 与 $l_j=g_{jk}L^k$，
对任意点 $S$ 有
\begin{equation}
\label{eq:hS}
h\cdot S=l^2\,(g\cdot S)-2(l\cdot S)\,L .
\end{equation}
把它代入 \eqref{eq:Tdef}，并用 $(g\cdot S)\times(g\cdot S)=0$，得
\begin{equation}
\label{eq:Tform}
T=-2(l\cdot S)\,(g\cdot S)\times L .
\end{equation}
于是只需证明 $l\cdot\bigl((g\cdot S)\times L\bigr)=0$。写成分量：
\[
l\cdot\bigl((g\cdot S)\times L\bigr)
=l_i\epsilon^{ijk}(g\cdot S)_jL_k .
\]
把 $l_i=g_{im}L^m$ 代入，得到 $\epsilon^{ijk}g_{im}L^mL_k(g\cdot S)_j$。由于 $g_{im}$ 关于
$i,m$ 对称而 $L^mL_k$ 关于 $m,k$ 对称，上式中的 $g_{im}L^mL_k$ 关于 $m,k$ 对称，而
$\epsilon^{ijk}$ 关于 $m,k$（即第 $2,3$ 个指标）反对称，故缩并为零。因此 $l\cdot T=0$。

综合两步，$T$ 既是 $p_g$ 与 $p_h$ 的交点，又落在 $l$ 上。
\end{proof}

\begin{remark}
上面的论证说明交点应取为 $(g\cdot S)\times(h\cdot S)$，而不能把 $g,h$ 直接依次作用在 $S$ 上：
后者的结果一般并不落在两条极线的交点上。关键在于这里的 $g\cdot S$ 与 $h\cdot S$ 都应当理解为
直线（线矢），它们的叉乘才给出所求的点。当 $S$ 取遍平面时，$T$ 取遍 $l$ 上的点列。
\end{remark}

\section{共轭圆锥曲线}

\begin{definition}[共轭圆锥曲线]
\label{def:conjugate}
设 $g,h$ 是两条圆锥曲线（可逆二次点）。若存在一条直线 $l$，使得
\[
g\times h=ll,
\]
则称 $g,h$ 是\textbf{共轭}的。
\end{definition}

\begin{proposition}[自由度计数]
与给定圆锥曲线 $g$ 共轭的圆锥曲线有 $2$ 个自由度。
\end{proposition}

\begin{proof}
平面上的直线由 $2$ 个参数确定，故 $l$ 有 $2$ 个自由度。另一方面，圆锥曲线共有 $5$ 个自由度
（一般二次型的 $6$ 个分量模去整体常数因子）；而 $g\times h=ll$ 这一条件给出 $3$ 个约束
（$g\times h$ 的分量张成一个 $6$ 维空间，$ll$ 形式只占其中由 $3$ 个参数描述的 $3$ 维子集），
$5-3=2$。两者一致。
\end{proof}

\subsection{共轭双曲线}

\begin{example}[共轭双曲线]
\label{ex:conj-hyp}
取两条双曲线
\[
g:\ \frac{x^2}{a^2}-\frac{y^2}{b^2}=1,\qquad
h:\ -\frac{x^2}{a^2}+\frac{y^2}{b^2}=1 .
\]
它们有共同的渐近线 $y=\pm\dfrac ba x$。计算它们的二次叉乘：由
\[
g^{ij}=\diag\{a^2,-b^2,-1\},\qquad h^{ij}=\diag\{-a^2,b^2,-1\},
\]
用对角叉乘公式 \eqref{eq:cross-diag}，
\[
(g\times h)_{11}=(-b^2)(-1)+(-1)(b^2)=b^2-b^2=0,
\]
\[
(g\times h)_{22}=a^2(-1)+(-1)(-a^2)=-a^2+a^2=0,
\]
\[
(g\times h)_{33}=(-b^2)(-a^2)+a^2(b^2)=a^2b^2+a^2b^2=2a^2b^2 .
\]
于是
\[
(g\times h)_{ij}=\diag\{0,0,2a^2b^2\},
\]
对角线上有两个零，故它正是 $ll$ 的形式（$l$ 为无穷远直线），符合共轭圆锥曲线的定义。这里
第三个分量的符号容易算错：它由 $-b^2$ 与 $-a^2$ 的乘积加上 $a^2$ 与 $b^2$ 的乘积得到，
两项同号，故为 $2a^2b^2$。
\end{example}

\begin{remark}
在例 \ref{ex:conj-hyp} 中，有限平面上的原点对应点 $L$，直线 $AB$ 是无穷远直线，而
$LA,LB$ 是两条双曲线共同的渐近线。这提示我们可以用射影几何的方法研究共轭双曲线，而不必
借助坐标系。解析几何中的"共轭双曲线"正是本节"共轭圆锥曲线"概念的最初等的特例。
\end{remark}

\subsection{$(1,1)$ 型张量 $\tau$}

考虑 $(1,1)$ 型张量
\begin{equation}
\label{eq:tau}
\tau^i{}_k=h^{ij}g_{jk}=l^2\delta^i_k-2l^il_k ,
\end{equation}
其中第二个等号用了 $h^{ij}=l^2g^{ij}-2l^il^j$ 与 $g^{ij}g_{jk}=\delta^i_k$。

\begin{proposition}
\label{prop:tau}
$\tau$ 可以看作把点映为点的线性映射，并且
\[
\tau^2=l^4\,I .
\]
因此 $\tau$ 是（三维中）对合变换的推广。
\end{proposition}

\begin{proof}
直接计算：
\[
(\tau^2)^i{}_l=\tau^i{}_k\tau^k{}_l
=(l^2\delta^i_k-2l^il_k)(l^2\delta^k_l-2l^kl_l)
=l^4\delta^i_l-2l^2l^il_l-2l^2l^il_l+4l^2l^il_l=l^4\delta^i_l . \qedhere
\]
\end{proof}

\begin{remark}
在射影变换的意义下，乘以非零常数不影响所定义的变换，故 $\tau^2=l^4I$ 意味着 $\tau$ 的平方
是恒等变换，即 $\tau$ 是（三维中的）对合变换的推广。在二维（例如限制在过 $L$ 的直线上）时，
$\tau$ 给出真正的对合。
\end{remark}

\begin{proposition}[$\tau$ 的作图]
\label{prop:tau-construction}
任取一点 $C$，记 $D=\tau C$。则 $D$ 可以如下作出：过 $C$ 作 $g$ 的两条切线，它们与 $l$ 交
于 $I,J$；再作 $I,J$ 到 $g$ 的切线，两条切线的交点就是 $D$。
\end{proposition}

\begin{proof}
由 $\tau=h\cdot g$，$D=h\cdot(g\cdot C)$：先取 $C$ 关于 $g$ 的极线，再取这条极线关于 $h$
的极点。故 $D$ 是 $C$ 关于 $g$ 的极线关于 $h$ 的极点。

设 $C$ 关于 $g$ 的极线为 $c^g$。过 $C$ 作 $g$ 的两条切线，切点为 $P,Q$，则
$c^g=PQ$，且 $c^g$ 与 $l$ 交于两点 $I,J$。因为在直线 $c^g$ 上，$g$ 与 $h$ 有共同的极点
极线（命题 \ref{prop:common-polar-on-l}），故 $I,J$ 到 $g$ 的两条切线之交 $E$ 也是 $I,J$ 到
$h$ 的切线之交，即 $E$ 是 $c^g$ 关于 $h$ 的极点。而 $c^g$ 关于 $h$ 的极点正是
$D=h\cdot(g\cdot C)$。故 $E=D$，即所作的点就是 $\tau C$。
\end{proof}

\begin{proposition}
\label{prop:tau-line}
无论 $C$ 位于何处，$C,D,L$ 总在同一条直线上。
\end{proposition}

\begin{proof}
由 \eqref{eq:tau}，$\tau C=l^2C-2(l\cdot C)l$。因此 $\tau C-C$ 与 $l$ 成比例。于是
\[
D=\tau C=C-\lambda\,l
\]
对某个标量 $\lambda$ 成立。在射影意义下，$C$ 与 $l$ 所张成的直线（即过 $C$ 与 $L$ 的直线，
因为 $L$ 在 $l$ 上）包含 $D$，即 $C,D,L$ 共线。
\end{proof}

\begin{proposition}
\label{prop:tau-involution}
设 $m$ 是一条过 $L$ 的直线，并设 $C_0=m\times l$（$m$ 与 $l$ 的交点）。把 $C$ 限制在 $m$ 上，
则 $C\mapsto D=\tau C$ 给出 $m$ 上的一个对合关系；这个对合的两个不动点分别是 $L$ 与 $C_0$，
且 $m$ 与 $g$ 的两个交点关于 $\tau$ 对合。
\end{proposition}

\begin{proof}
由命题 \ref{prop:tau-line}，$C\in m$ 时 $D\in m$，故 $\tau$ 在 $m$ 上给出一个射影变换。由
$\tau^2=l^4I$，它是（射影意义下的）对合。

不动点：设 $\tau C$ 与 $C$ 成比例。由 $\tau C=C-2(l\cdot C)l$，这等价于 $(l\cdot C)l$ 与
$C$ 成比例，即 $(l\cdot C)=0$ 或 $C$ 与 $l$ 成比例。前者的解恰是 $m$ 与 $l$ 的交点
$C_0$；后者的解是 $C=L$。故不动点正是 $L$ 与 $C_0$。

最后设 $A,B$ 是 $m$ 与 $g$ 的两个交点。在 $m$ 上任取一点 $N$（$N\notin l$，这样的 $N$ 总
存在），并把 $m$ 上的点参数化为 $C(t)=L+tN$；于是 $L$ 对应 $t=0$。把 $l$ 规范化使
$l\cdot l=1$。记
\[
H=L\cdot L,\qquad q=N\cdot N,\qquad (L\cdot l)^2=H-q .
\]
（最后一式是 $N$ 的这一取法所满足的规范化关系；若取 $N$ 使 $N\cdot l=0$，则它与前三式一起
确定 $H,q$ 的取值。）由 $l$ 与 $L$ 一般不正交，记 $L\cdot l$ 为第六个量。

计算 $\tau C=C-2(l\cdot C)l$。对 $C(t)=L+tN$，
\[
\tau C=L+tN-2\,(L\cdot l+t\,N\cdot l)\,l
=\alpha L+(\beta t)N+\text{（一个与 }t\text{ 无关的 $l$ 方向项）},
\]
把右端重新按 $L,N$（用 $l=(L-N)/(L\cdot l)$ 型的分解）整理成 $s'L+N'$ 的形式，可得
\[
\tau:\ t\longmapsto \frac{1-\alpha}{t},\qquad
\alpha=1-(N\cdot l)(L\cdot l)^{-1} ,
\]
即 $\tau$ 在坐标 $t$ 下是 $t\mapsto-\beta/t$ 型的对合（由 $M^2=\mathrm{id}$，
$\det M=1$）。它的不动点是 $t=0$（$L$）与 $t=\infty$（$m$ 上落在 $l$ 中的那个点，
即 $m\times l$），符合对合的两个不动点。

现在计算 $g$ 在 $m$ 上的限制。把 $C(t)=L+tN$ 代入并使用 $L\cdot N=H-q$ 与 $q=N\cdot N$：
\[
g(C(t),C(t))=H+2t\,(H-q)+t^2q=(1+t)^2(H-q)+q\,(t^2-2t-1)\ \text{ 型},
\]
它的两个零点 $t_1,t_2$ 就对应 $A,B$。直接计算
\[
g(\tau C,\tau C)=\frac{\beta^2}{t^2}\,g(C,C)
\]
（把 $\tau C=C-2(l\cdot C)l$ 代入，用 $l\cdot l=1$ 与 $g(C,l)=l\cdot C$ 整理），故 $\tau$ 保持
$g$ 在 $m$ 上的零点集，即 $\tau(t_1)=t_2$。换句话说 $A,B$ 关于 $\tau$ 对合。
\end{proof}

\section{坐标例题}

\subsection{例 1：共轭双曲线的切点共线}

\begin{example}
\label{ex:conj-tangents}
设两条共轭双曲线
\[
g:\ \frac{x^2}{a^2}-\frac{y^2}{b^2}=1,\qquad h:\ -\frac{x^2}{a^2}+\frac{y^2}{b^2}=1 ,
\]
点 $A=(x,y)$ 在 $g$ 上，点 $B=(-x,-y)$。过 $A$ 作 $g$ 的两条切线，切于 $C,D$；过 $B$ 作
$h$ 的两条切线，切于 $E,F$。求证 $C,D,E,F$ 共线。
\end{example}

\begin{proof}
首先注意：由 $A\in g$ 得 $\frac{x^2}{a^2}-\frac{y^2}{b^2}=1$，故 $B=(-x,-y)$ 满足
\[
-\frac{(-x)^2}{a^2}+\frac{(-y)^2}{b^2}=1,
\]
即 $B\in h$，两条切线的确存在（当 $A,B$ 不在无穷远处时）。

在射影坐标下把两条曲线的方程写成
\[
g:\ \frac{X^2}{a^2}-\frac{Y^2}{b^2}=Z^2,\qquad
h:\ -\frac{X^2}{a^2}+\frac{Y^2}{b^2}=Z^2 .
\]
点 $A=(x,y,1)$ 关于 $g$ 的极线由双线性形式给出：
\[
\frac{xX}{a^2}-\frac{yY}{b^2}=Z,
\]
它与 $g$ 的两个交点就是 $C,D$。点 $B=(-x,-y,1)$ 关于 $h$ 的极线由 $h$ 的双线性形式给出：
\[
\frac{(-x)X}{a^2}-\frac{(-y)Y}{b^2}=Z,
\]
即
\[
-\frac{xX}{a^2}+\frac{yY}{b^2}=Z .
\]
这条极线与 $h$ 的两个交点就是 $E,F$。

下面证明这两条极线重合。取 $h$ 的方程 $-\frac{X^2}{a^2}+\frac{Y^2}{b^2}-Z^2=0$，其双线性
形式为 $-\frac{Xx'}{a^2}+\frac{Yy'}{b^2}-Zz'$。点 $B=(-x,-y,1)$ 的极线是
\[
-\frac{X(-x)}{a^2}+\frac{Y(-y)}{b^2}-Z=0,
\qquad\text{即}\qquad
\frac{xX}{a^2}-\frac{yY}{b^2}-Z=0 .
\]
这与 $A$ 关于 $g$ 的极线完全相同。因此 $C,D$ 与 $E,F$ 都在同一条直线（$A$ 关于 $g$ 的
极线）上，四点共线。
\end{proof}

\begin{remark}
题设中若同时要求 $0<x<a$ 与 $0<y<b$，则 $A$ 不可能在双曲线 $g$ 上（因为双曲线上点的
横坐标满足 $|x|\ge a$）。正确的条件是 $A$ 在 $g$ 上的某一支上，$B=(-x,-y)$ 自动落在 $h$
的对应支上。上面的计算只用到 $A\in g$，结论仍然成立。
\end{remark}

\begin{remark}
上述证明只用到"共轭双曲线的极线重合"这一事实。也可以直接用解析式验证：把 $c^g$ 与 $g$、
$h$ 分别联立，所得四个切点的连线就是 $c^g$ 本身。
\end{remark}

\subsection{例 2：$\tau$ 的具体计算}

\begin{example}
\label{ex:tau-compute}
设椭圆 $g:\dfrac{x^2}4+y^2=1$，直线 $l:y=2x-1$，点 $C(-3,-4)$。过 $C$ 作椭圆的两条切线，
它们与 $l$ 分别交于 $E,F$；过 $E,F$ 再作椭圆的切线，两条切线交于 $D$。求 $D$ 的坐标。
\end{example}

\begin{proof}[解]
由本节的理论，$D$ 可以由公式
\[
d=\tau c=\frac{ll}{g}\cdot g^{-1}\cdot c ,
\qquad\text{即}\qquad
d^i=l^2c^i-2(c\cdot l)\,l^i
\]
计算。这里把 $g$ 作为度规。写出各量：
\[
g^{ij}=\diag\{4,1,-1\},\qquad c^i=(-3,-4,1),\qquad l_i=(2,-1,-1).
\]
于是
\[
c\cdot l=(-3)\times2+(-4)\times(-1)+1\times(-1)=-6+4-1=-3 .
\]
用 $g^{ij}$ 升指标：
\[
l^i=g^{ij}l_j=(4\times2,\ 1\times(-1),\ -1\times(-1))=(8,-1,1),
\]
\[
l^2=g^{ij}l_il_j=4\times2^2+1\times(-1)^2+(-1)\times(-1)^2
=16+1-1=16 .
\]
代入 $\tau$ 的公式：
\begin{align*}
d&=l^2c-2(c\cdot l)l=16(-3,-4,1)-2(-3)(8,-1,1)\\
&=(-48,-64,16)+(48,-6,6)=(0,-70,22).
\end{align*}
故
\[
D=\Big(0,-\frac{70}{22}\Big)=\Big(0,-\frac{35}{11}\Big)\approx(0,-3.18).
\]
\end{proof}

\begin{remark}
在本题的数据下，直线 $l:y=2x-1$ 与椭圆交于两点而不相切：把 $y=2x-1$ 代入
$\frac{x^2}4+y^2=1$，得
\[
\frac{x^2}{4}+(2x-1)^2=1
\ \Longleftrightarrow\
\frac{17x^2}{4}-4x=0
\ \Longleftrightarrow\
x(17x-16)=0,
\]
即交点分别为 $(0,-1)$ 与 $\big(\frac{16}{17},\frac{15}{17}\big)$。作图时 $E,F$ 是这两点附近
$l$ 上的点，$D$ 由张量公式给出。计算 $l^2$ 时要注意第三个分量的符号：$g^{33}=-1$ 给出
$-1$，故 $l^2=16+1-1=16$。
\end{remark}
   % 二次除法与共轭圆锥曲线

\part{圆锥曲线上的线性空间}
% 第十八节：指标分裂与二次嵌入
\chapter{指标分裂与二次嵌入}

\section{二次点列与一次点列}

在二次点列与二次对合一节中已经看到，圆锥曲线 $g$ 上的点集（二次点列）与直线上的点集
（一次点列）有许多相似之处：

\begin{enumerate}[label=(\arabic*),leftmargin=2.4em]
\item 两者具有相同的拓扑。一次点列补入无穷远点后是紧致的，二次点列同样是紧致的。
\item 在两者上都可以定义交比、对合与调和点列。
\item 在 $g$ 上任取一点 $p$，把 $g$ 上的点 $a$ 映为直线 $pa$，就得到二次点列与以 $p$ 为
      中心的一次线束之间的一一对应。
\item 射影变换的复合仍是射影变换。设有分别以 $p,q$ 为中心的线束 $\Sigma_p,\Sigma_q$。
      取 $\Sigma_p$ 中的直线 $pA$，令 $f_1$ 把它映为它与 $g$ 的交点 $A$，再令 $f_2$ 把 $A$
      映为 $\Sigma_q$ 中的直线 $qA$，于是得到复合映射 $f=f_2\circ f_1:\Sigma_p\to\Sigma_q$。
      由于 $g$ 不是线性子空间，$f_1$ 与 $f_2$ 都不是线性映射，$f$ 看上去不必是射影变换；
      然而由交比的不变性可以证明，$f$ 确实是射影变换。
\end{enumerate}

这四条性质的共同根据是：二次点列与一次点列之间存在保持交比的一一对应，因而前者可以
自然地参与经典射影几何中关于射影变换的讨论。这一结论在前面已经证明，但那个证明并不自然，
像是某种巧合。另一方面，二次点列与一次点列仍有明显的差别：

\begin{enumerate}[label=(\arabic*),leftmargin=2.4em]
\item 一次点列是平面的线性子空间，而二次点列不是平面的线性子空间。
\item 上面第 (4) 条中的 $f_1,f_2$ 虽然与射影变换相似，毕竟不是射影变换。其根源是：
      一次点列与二次点列之间的一一映射是非线性的。
\end{enumerate}

非线性的后果是二次点列的代数意义不明确，它像是一次点列间射影变换的中间产物。
一次点列的代数意义很清楚：它就是三维线性空间中一个二维子空间的元素，而"二维线性空间"
这一背景正提供了交比、对合与调和点列这些结构。二次点列却是
\[
\{a\mid g_{ij}a^ia^j=0\},
\]
即由一个方程确定的点集，它根本不是二维线性空间，却同样带有这些结构。对此可以给出如下
说法：

\begin{quote}
就本质而言，二次点列\emph{是}一个二维线性空间，这使得它具有交比、对合与调和点列；
但在把这个二维线性空间嵌入三维线性空间时，它被"弯曲"了，因而丢失了一部分线性性质。
\end{quote}

也就是说，二次点列是"部分线性、部分非线性"的：其中的线性部分使它与一次点列相似，
非线性部分使它与一次点列不同。为了明确讨论，需要先对"部分线性、部分非线性"作出定义。

\section{嵌入与线性嵌入}

\begin{definition}[嵌入]
设 $M,N$ 是流形，$\phi:M\to N$ 是连续的单射，则称 $\phi$ 为嵌入。
（这里只取数学定义中我们会用到的部分。）
\end{definition}

矢量空间是流形的特例。设 $U$ 是二维矢量空间，$V$ 是三维矢量空间，则 $U,V$ 都是流形，
并且各自带有加法。为区别二者，暂把 $U$ 中的加法记为 $\oplus$，其元素记为
$\mathfrak a,\mathfrak b,\mathfrak c,\mathfrak d$；$V$ 中的加法仍记为 $+$，其元素记为
$a,b,c,d$。

\begin{definition}[线性嵌入与非线性嵌入]
设 $\phi:U\to V$ 是嵌入。若对任意 $\mathfrak a,\mathfrak b\in U$ 恒有
\[
\phi(\mathfrak a\oplus\mathfrak b)=\phi(\mathfrak a)+\phi(\mathfrak b),
\]
则称 $\phi$ 为线性嵌入；否则称为非线性嵌入。
\end{definition}

线性嵌入保持矢量空间的线性运算，非线性嵌入则不然。一次点列作为 $U$ 到 $V$ 的嵌入是线性
嵌入；后文将看到，二次点列是非线性嵌入。

\section{指标的分解与旋量空间}

先用一个例子说明指标的分解是怎样产生的。在四维时空中，洛伦兹度规
$t^2-x^2-y^2-z^2$ 下的四维矢量 $a^i=(t,x,y,z)$ 可以写成 $2\times2$ 矩阵
\[
A=\begin{pmatrix} t+x & y+\ii z\\ y-\ii z & t-x\end{pmatrix},
\]
其行列式恰好等于 $t^2-x^2-y^2-z^2$。现在处理的是平面几何，用的是三维矢量空间，
只需去掉坐标 $z$：此时洛伦兹度规成为 $t^2-x^2-y^2$，点矢 $a^i=(t,x,y)$ 对应矩阵
\[
A=\begin{pmatrix} t+x & y\\ y & t-x\end{pmatrix},
\qquad \det A=t^2-x^2-y^2 .
\]
这里的"洛伦兹度规"是一个不定二次型，也就是一条圆锥曲线。因此，当平面上给定一条圆锥曲线
$g$ 时，只要选取坐标系使 $g$ 的方程成为 $t^2-x^2-y^2=0$，就可以用上述方法把平面上的
每一点写成 $2\times2$ 的对称矩阵。

矩阵的本质是张量，而张量可以用指标描述。矩阵 $A$ 的两个指标记为 $\alpha,\dot\alpha$，
取值范围都是 $1,2$，于是 $A$ 写成 $a^{\alpha\dot\alpha}$，与原指标 $i$（取值范围 $1,2,3$）
不同。对偶矢量 $a_i=g_{ij}a^j=(t,-x,-y)$ 同样可写成对称张量
\[
a_{\alpha\dot\alpha}=\begin{pmatrix} t-x & -y\\ -y & t+x\end{pmatrix}.
\]
进一步，任意张量，例如 ${f^{ij}}_k$，都可以改写成
${f^{\alpha\dot\alpha\beta\dot\beta}}_{\gamma\dot\gamma}$：原张量的每一个取值为 $1,2,3$
的指标，都被两个取值为 $1,2$ 的指标代替。

\begin{definition}[指标分裂与旋量空间]
把张量的每个指标替换为一对旋量指标 $\alpha\dot\alpha$（$\alpha,\dot\alpha\in\{1,2\}$）
的写法，称为\emph{指标分裂}。分裂后指标所在的二维矢量空间称为 $g$ 的\emph{旋量空间}，
记为 $S$。
\end{definition}

二维矢量空间中的二阶对称张量表示一个对合变换，因此可以说：平面上的每一个点矢都对应 $S$
中的一个对合变换。这里出现两个旋量指标 $\alpha$ 与 $\dot\alpha$，它们分别属于两个互为
共轭的二维旋量空间；在实的三维空间中这两个空间可以等同，故下文对不带点与带点的指标用同一
套规则处理。

$a$ 与 $a^{\alpha\dot\alpha}$ 之间的关系是线性的，因此可以用一个线性映射在两者之间转换：
\[
a^{\alpha\dot\alpha}=\sigma^{\alpha\dot\alpha}_i a^i .
\]
$\sigma^{\alpha\dot\alpha}_i$ 称为\emph{转换符}。按 $i=1,2,3$（即 $t,x,y$）的次序，它的
三个非零分量可以写成矩阵
\[
(\sigma^{\alpha\dot\alpha}_i)=
\Bigl(\,
\begin{pmatrix}1&0\\0&1\end{pmatrix},\
\begin{pmatrix}1&0\\0&-1\end{pmatrix},\
\begin{pmatrix}0&1\\1&0\end{pmatrix}
\,\Bigr),
\]
即
\[
\sigma^{11}_1=\sigma^{22}_1=1,\qquad
\sigma^{11}_2=1,\quad \sigma^{22}_2=-1,\qquad
\sigma^{12}_3=\sigma^{21}_3=1,
\]
其余分量为零。反过来的转换为
\[
a^i=\tfrac12\sigma^i_{\alpha\dot\alpha}a^{\alpha\dot\alpha},
\]
其中非零分量为
\[
\sigma^1_{11}=\sigma^1_{22}=1,\qquad
\sigma^2_{11}=1,\quad \sigma^2_{22}=-1,\qquad
\sigma^3_{12}=\sigma^3_{21}=1 .
\]
这里出现的系数 $\frac12$ 来自对称对 $(\alpha\dot\alpha)$ 在求和时的双重计数，下一节将把
这一约定讲清楚。由于 $a$ 与 $a^{\alpha\dot\alpha}$ 只相差一个转换符，为节省字母，下文把
矩阵 $a^{\alpha\dot\alpha}$ 也记为 $a$，只靠指标区分两种情形。

\section{对称旋量张量与二次嵌入}

$a^{\alpha\dot\alpha}$ 是对称的，而二维矢量空间中的对称张量总可以写成对称积的形式：
\[
a^{\alpha\dot\alpha}=\mathfrak a^\alpha\mathfrak b^{\dot\alpha}
+\mathfrak b^\alpha\mathfrak a^{\dot\alpha}=(\mathfrak a\mathfrak b)^{\alpha\dot\alpha}.
\]
旋量 $\mathfrak a,\mathfrak b$ 称为 $a$ 的\emph{旋量坐标}，它们正是这个对合变换的不动点。

点 $a$ 在 $g$ 上的条件是 $g_{ij}a^ia^j=0$。当 $a^i=(t,x,y)$ 时，该条件写为
$t^2-x^2-y^2=0$，即 $\det a^{\alpha\dot\alpha}=0$，也就是 $a^{\alpha\dot\alpha}$ 退化。
此时它不能写成两个不同旋量的对称积，只能写成一个旋量与自身的乘积：
\[
a=\mathfrak a\mathfrak a,\qquad\text{即}\qquad
a^{\alpha\dot\alpha}=\mathfrak a^\alpha\mathfrak a^{\dot\alpha}.
\]
此时称 $\mathfrak a$ 为 $a$ 的旋量坐标。

于是，圆锥曲线 $g$ 上的点 $a$ 与旋量空间 $S$ 中的点 $\mathfrak a$ 有一一对应；
更准确地说，$a$ 与 $\mathfrak a$ 的等价类 $[\mathfrak a]=\{\lambda\mathfrak a\}$
一一对应，即射影旋量空间与二次点列一一对应。映射
\[
\phi:S\to V,\qquad \phi(\mathfrak a)=\mathfrak a\mathfrak a
\]
把 $\mathfrak a$ 映为 $g$ 上的点，它是嵌入；又因为 $\phi$ 关于 $\mathfrak a$ 是二次的，
它还是非线性嵌入。

\begin{definition}[二次嵌入]
称嵌入 $\phi(\mathfrak a)=\mathfrak a\mathfrak a$ 为\emph{二次嵌入}，
即旋量空间 $S$ 可以二次嵌入到圆锥曲线 $g$ 上。
\end{definition}

用 $\mathfrak a$ 表达 $g$ 上的点有两点好处。其一，$\mathfrak a$ 作为二维矢量所构造出的
交比、对合与调和点列，正是前面在二次点列上定义的那些结构；其二，若用三维矢量 $a$ 表示
$g$ 上的点，就必须附加约束 $a^2=0$，当同时涉及多个圆锥曲线上的点时，这些约束会带来
麻烦，而用旋量则不必引入约束。

\section{旋量的几何表示}

旋量 $\mathfrak a$ 属于旋量空间 $S$，并不属于平面 $V$，因此严格说来不能在平面上画出旋量；
但它在平面上的像 $\phi(\mathfrak a)$ 是 $g$ 上的点，属于 $V$。因此可以用 $\phi(\mathfrak a)$
代表 $\mathfrak a$。

设 $g$ 外有一点 $p$，其旋量坐标为 $\mathfrak a,\mathfrak b$，即 $p=(\mathfrak a\mathfrak b)$，
则 $\phi(\mathfrak a)$ 与 $\phi(\mathfrak b)$ 是 $p$ 到 $g$ 的两条切线的切点
（证明见下一节命题 1）。若 $p$ 在 $g$ 内部，则 $\mathfrak a,\mathfrak b$ 的坐标为复数，
不能在实平面上画出。

设 $p=(\mathfrak a\mathfrak b)$，$q=(\mathfrak c\mathfrak d)$。$g$ 上四点 $a,b,c,d$
成调和点列等价于 $(\mathfrak a\mathfrak b):(\mathfrak c\mathfrak d)=0$；而由代数上的直觉
（下一节证明）应有
\[
(\mathfrak a\mathfrak b):(\mathfrak c\mathfrak d)=p\cdot q .
\]
于是 $a,b,c,d$ 成调和点列也就等价于 $p,q$ 共轭。这一结论前面证明时并不自然，
而在这里是自然得到的。

\section{旋量坐标}

通过乘以系数，总可以把旋量化为标准规范，例如取
$\mathfrak a=(u,1)$，$\mathfrak b=(v,1)$。于是平面上的任意一点 $p=(\mathfrak a\mathfrak b)$
都可用 $u,v$ 表出，这给出平面的另一种坐标系，称为\emph{旋量坐标系}。

具体地，由 $\mathfrak a=(u,1),\mathfrak b=(v,1)$ 得
\[
a^{\alpha\dot\alpha}
=\mathfrak a^\alpha\mathfrak b^{\dot\alpha}+\mathfrak b^\alpha\mathfrak a^{\dot\alpha}
=\begin{pmatrix} 2uv & u+v\\ u+v & 2\end{pmatrix}.
\]
与 $\begin{pmatrix} t+x & y\\ y & t-x\end{pmatrix}$ 对比，得到方程组
\[
t+x=2uv,\qquad t-x=2,\qquad y=u+v,
\]
解得
\[
t=uv+1,\qquad x=uv-1,\qquad y=u+v .
\]
取仿射规范把 $t$ 化为 $1$（此时 $g:t^2-x^2-y^2=0$ 对应单位圆），就得到直角坐标与
旋量坐标之间的转换关系
\[
x=\frac{uv-1}{uv+1},\qquad y=\frac{u+v}{uv+1},
\]
其中 $uv+1\neq0$。这是圆的旋量坐标，稍加修改即可得到椭圆的旋量坐标。

旋量坐标特别适合处理切线问题。当 $u$ 固定时 $\mathfrak a$ 随之固定，点 $p$ 只能在过
$\phi(\mathfrak a)$ 的切线上变动；对 $v$ 同理。因此\emph{固定任何一个旋量坐标就得到圆锥
曲线的一条切线}。此外，令 $u=v$，则 $\mathfrak a=\mathfrak b$，此时 $p$ 在 $g$ 上，
转换关系成为
\[
x=\frac{u^2-1}{u^2+1},\qquad y=\frac{2u}{u^2+1},
\]
它可以看作圆的参数方程。这一参数方程可由圆的经典参数方程与万能公式得到：
\[
u=\tan\frac\theta2,
\]
其中 $\theta$ 是 $OP$ 与 $x$ 轴的夹角。
   % 指标分裂与二次嵌入
% 第十九节：再谈二次点列
\chapter{再谈二次点列}

引入指标分裂后，若要用旋量运算解决问题，就必须把原有的代数工具——点乘、叉乘、缩并，
以及 $g_{ij}$、$\epsilon_{ijk}$ 等——全部翻译成旋量语言。为此先要熟悉二维矢量空间中的
交比、对合与调和点列。

\section{矢量运算的旋量翻译}

联系矢量与旋量的工具是转换符 $\sigma^{\alpha\dot\alpha}_i$ 与 $\sigma^i_{\alpha\dot\alpha}$，
其非零分量为
\[
\sigma^{11}_1=\sigma^{22}_1=1,\quad \sigma^{11}_2=1,\quad \sigma^{22}_2=-1,\quad
\sigma^{12}_3=\sigma^{21}_3=1,
\]
以及
\[
\sigma^1_{11}=\sigma^1_{22}=1,\quad \sigma^2_{11}=1,\quad \sigma^2_{22}=-1,\quad
\sigma^3_{12}=\sigma^3_{21}=1 .
\]
直接由分量计算可得
\begin{equation}\label{eq:sigma-contract}
\sigma^i_{\alpha\dot\alpha}\sigma^{\alpha\dot\alpha}_j=2\delta^i_j .
\end{equation}

\begin{remark}[归一化]
\eqref{eq:sigma-contract} 右端的系数 $2$ 来自对称对 $(\alpha\dot\alpha)$ 在求和时的双重计数。
若把两个转换符同时乘以归一化因子 $1/\sqrt2$，即令
$\hat\sigma^{\alpha\dot\alpha}_i=\frac1{\sqrt2}\sigma^{\alpha\dot\alpha}_i$、
$\hat\sigma^i_{\alpha\dot\alpha}=\frac1{\sqrt2}\sigma^i_{\alpha\dot\alpha}$，则
$\hat\sigma^i_{\alpha\dot\alpha}\hat\sigma^{\alpha\dot\alpha}_j=\delta^i_j$。由于射影几何中
点只由坐标的比确定，这一常数因子不影响任何几何结论。下面保留未归一化的转换符，并在每个
等式中把相应的因子写清楚。
\end{remark}

\eqref{eq:sigma-contract} 的意义在于它给出了缩并运算的旋量表达。以 $(1,1)$ 型张量的缩并
为例：
\[
\tau^i_i=\tau^j_i\delta^i_j=\tfrac12\,\tau^j_i\sigma^i_{\alpha\dot\alpha}\sigma^{\alpha\dot\alpha}_j
=\tau^{\alpha\dot\alpha}_{\alpha\dot\alpha},
\]
其中转换符的作用是把矢量指标换为旋量指标。推广之：\emph{矢量的缩并总可以替换为旋量的
二重缩并}。反之则未必：旋量的二重缩并不一定能替换为矢量的缩并。原因在于，一个矢量指标
分裂出的两个旋量指标是对称的，反过来两个对称的旋量指标可以合成为一个矢量指标；若两个旋量
指标不对称，就不能合成为矢量指标。概括地说，矢量代数总可以表示为旋量代数，旋量代数不一定
能表示为矢量代数，因而旋量对几何的刻画比矢量更细致。

旋量空间中的 Kronecker 记号记为 $\delta^\alpha_\gamma$。矢量空间的 $\delta^i_j$ 用旋量指标
改写后记为 $\delta^{\alpha\dot\alpha}_{\gamma\dot\delta}$，两者只靠指标个数区分。由定义
$\delta^{\alpha\dot\alpha}_{\gamma\dot\delta}
=\sigma_i^{\alpha\dot\alpha}\sigma^i_{\gamma\dot\delta}
=\sigma_1^{\alpha\dot\alpha}\sigma^1_{\gamma\dot\delta}
+\sigma_2^{\alpha\dot\alpha}\sigma^2_{\gamma\dot\delta}
+\sigma_3^{\alpha\dot\alpha}\sigma^3_{\gamma\dot\delta}$，直接计算得到
\[
\delta^{\alpha\dot\alpha}_{\gamma\dot\delta}
=\delta^\alpha_\gamma\delta^{\dot\alpha}_{\dot\delta}
+\delta^{\dot\alpha}_\gamma\delta^\alpha_{\dot\delta},
\]
也就是说，它是在两个 Kronecker 记号的基础上对指标做了对称化。在只涉及一整套不带点的指标
时，上式也常写成 $\delta^{\alpha\beta}_{\gamma\delta}
=\delta^\alpha_\gamma\delta^\beta_\delta+\delta^\beta_\gamma\delta^\alpha_\delta$。

\section{度规的旋量表示}

矢量空间的度规 $g_{ij}$（旋量正是在 $g$ 的基础上构造出来的）用旋量指标改写后记为
$g_{\alpha\dot\alpha\beta\dot\beta}$。它应当能用旋量空间中的基本张量 $\delta^\alpha_\gamma$
与 $e_{\alpha\beta}$ 表示。由于 $g$ 只涉及下指标，用到的旋量基本张量只有 $e_{\alpha\beta}$。
为寻找表达式，先列出它的三条对称性：

\begin{enumerate}[label=(\arabic*),leftmargin=2.4em]
\item 交换第一对内的两个指标，值不变：
      $g_{\alpha\dot\alpha\beta\dot\beta}=g_{\dot\alpha\alpha\beta\dot\beta}$；
\item 交换第二对内的两个指标，值不变：
      $g_{\alpha\dot\alpha\beta\dot\beta}=g_{\alpha\dot\alpha\dot\beta\beta}$；
\item 同时交换前一对与后一对指标，值不变：
      $g_{\alpha\dot\alpha\beta\dot\beta}=g_{\beta\dot\beta\alpha\dot\alpha}$。
\end{enumerate}

其中前两条来源于指标分裂，第三条来源于度规是对称张量。在这三条对称性的限制下，可以设
\begin{equation}\label{eq:g-spinor}
g_{\alpha\dot\alpha\beta\dot\beta}
=e_{\alpha\beta}e_{\dot\alpha\dot\beta}+e_{\alpha\dot\beta}e_{\dot\alpha\beta},
\end{equation}
直接验证可知它满足上述三条对称性，而且确实是正确的构造。

\begin{definition}[旋量形式下的度规]
满足 \eqref{eq:g-spinor} 的张量 $g_{\alpha\dot\alpha\beta\dot\beta}$ 称为度规 $g_{ij}$ 的
旋量表示。
\end{definition}

\eqref{eq:g-spinor} 有多个用途。第一个用途是给旋量降指标。设 $s^{\beta\dot\beta}$ 是一个
对称旋量张量，则
\[
g_{\alpha\dot\alpha\beta\dot\beta}s^{\beta\dot\beta}
=e_{\alpha\beta}e_{\dot\alpha\dot\beta}s^{\beta\dot\beta}
+e_{\alpha\dot\beta}e_{\dot\alpha\beta}s^{\beta\dot\beta}
=s_{\alpha\dot\alpha}+s_{\dot\alpha\alpha},
\]
其中第二个等号用 $e$ 降指标。若 $s_{\alpha\dot\alpha}$ 本身是对称的，则结果就是
$s_{\alpha\dot\alpha}$；否则还附加了一次对称化。

第二个用途是定义点乘。设 $a^i,b^i$ 是点矢，则
\[
a\cdot b=g_{ij}a^ib^j
=g_{\alpha\dot\alpha\beta\dot\beta}a^{\alpha\dot\alpha}b^{\beta\dot\beta}
=a^{\beta\dot\beta}b_{\beta\dot\beta}.
\]

\section{二次点列的三个命题}

借助新工具，可以更简单地证明二次点列与二次对合中的一些命题。

\begin{proposition}\label{prop:tangent-points}
设 $p$ 是 $g$ 外一点，其旋量坐标为 $\mathfrak a,\mathfrak b$，即 $p=(\mathfrak a\mathfrak b)$。
则 $\phi(\mathfrak a)$ 与 $\phi(\mathfrak b)$ 是 $p$ 到 $g$ 的两条切线的切点。
\end{proposition}

\begin{proof}
按定义
\[
p_{\alpha\dot\alpha}=\mathfrak a_\alpha\mathfrak b_{\dot\alpha}
+\mathfrak b_\alpha\mathfrak a_{\dot\alpha},
\]
而 $\phi(\mathfrak a)$ 的分量为
$\phi(\mathfrak a)^{\alpha\dot\alpha}=\mathfrak a^\alpha\mathfrak a^{\dot\alpha}$。
计算缩并
\[
p\cdot\phi(\mathfrak a)
=p_{\alpha\dot\alpha}\phi(\mathfrak a)^{\alpha\dot\alpha}
=\bigl(\mathfrak a_\alpha\mathfrak b_{\dot\alpha}
+\mathfrak b_\alpha\mathfrak a_{\dot\alpha}\bigr)\mathfrak a^\alpha\mathfrak a^{\dot\alpha}
=(\mathfrak a_\alpha\mathfrak a^\alpha)(\mathfrak b_{\dot\alpha}\mathfrak a^{\dot\alpha})
+(\mathfrak b_\alpha\mathfrak a^\alpha)(\mathfrak a_{\dot\alpha}\mathfrak a^{\dot\alpha}).
\]
由于 $\mathfrak a_\alpha\mathfrak a^\alpha=e_{\alpha\beta}\mathfrak a^\beta\mathfrak a^\alpha=0$
且 $\mathfrak a_{\dot\alpha}\mathfrak a^{\dot\alpha}=0$（$e$ 反对称），两项都为零，故
\[
p\cdot\phi(\mathfrak a)=0 .
\]
这正是说 $\phi(\mathfrak a)$ 位于 $p$ 的极线上。又因为
$\phi(\mathfrak a)=\mathfrak a\mathfrak a$ 是 $g$ 上的点，所以 $\phi(\mathfrak a)$ 是 $p$ 的
极线与 $g$ 的交点，即从 $p$ 引向 $g$ 的切线的切点。对 $\mathfrak b$ 的论证完全相同。
由于 $\mathfrak a,\mathfrak b$ 线性无关，$\phi(\mathfrak a)\neq\phi(\mathfrak b)$，故它们是
两条不同切线的切点。
\end{proof}

\begin{remark}
把上述论证改写成对偶的形式，就得到本章开头提到的结论：若 $p=(\mathfrak a\mathfrak b)$、
$q=(\mathfrak c\mathfrak d)$，则
$(\mathfrak a\mathfrak b):(\mathfrak c\mathfrak d)=p\cdot q$。因此 $g$ 上四点 $a,b,c,d$ 成
调和点列，当且仅当 $p,q$ 共轭。命题 3 给出了它的另一种写法。
\end{remark}

\begin{proposition}\label{prop:crossratio-spinor}
设 $a,b,c,d$ 是 $g$ 上的四个定点。则对 $g$ 上任一点 $p$，四条直线
$a\times p,\ b\times p,\ c\times p,\ d\times p$ 的交比与 $p$ 无关，且等于旋量
$\mathfrak a,\mathfrak b,\mathfrak c,\mathfrak d$ 的交比。
\end{proposition}

\begin{proof}
直线 $a\times p$ 的旋量式为 $(\mathfrak a\mathfrak p)_{\alpha\dot\alpha}$，换成矢量式为
\[
l_i=(\mathfrak a\mathfrak p)_{\alpha\dot\alpha}\sigma^{\alpha\dot\alpha}_i .
\]
右端关于 $\mathfrak a$ 是线性的，因此
\[
f:\mathfrak a\longmapsto(\mathfrak a\mathfrak p)_{\alpha\dot\alpha}\sigma^{\alpha\dot\alpha}_i
\]
是从旋量空间到以 $p$ 为中心的线束的线性映射。线性映射保持交比不变，所以无论 $p$ 如何选取，
四条直线的交比都等于 $\mathfrak a,\mathfrak b,\mathfrak c,\mathfrak d$ 的交比。于是可以把
这一交比定义为二次点列的交比。
\end{proof}

值得注意的是，$p$ 必须在 $g$ 上：只有此时 $p=\mathfrak p\mathfrak p$，直线 $a\times p$ 才能
写成关于 $\mathfrak a$ 的线性表达式。若 $p$ 不在 $g$ 上，$p$ 需要两个旋量坐标，映射就不再
是线性的。

\begin{proposition}\label{prop:harmonic}
$g$ 上四点 $a,b,c,d$ 成调和点列，当且仅当直线 $a\times b$ 与 $c\times d$ 共轭。
\end{proposition}

\begin{proof}
直线 $a\times b$ 与 $c\times d$ 共轭，等价于它们的极点 $p,q$ 共轭。由命题
\ref{prop:tangent-points}，$p=(\mathfrak a\mathfrak b)$，$q=(\mathfrak c\mathfrak d)$，
故共轭条件 $p\cdot q=0$ 化为
\[
g_{\alpha\dot\alpha\beta\dot\beta}
(\mathfrak a\mathfrak b)^{\alpha\dot\alpha}(\mathfrak c\mathfrak d)^{\beta\dot\beta}=0 .
\]
由于 $g$ 本身对称，$(\mathfrak a\mathfrak b)$ 与 $(\mathfrak c\mathfrak d)$ 中的对称符号
可以去掉，得到
\[
g_{\alpha\dot\alpha\beta\dot\beta}\,
\mathfrak a^\alpha\mathfrak b^{\dot\alpha}\mathfrak c^\beta\mathfrak d^{\dot\beta}=0 .
\]
代入 \eqref{eq:g-spinor}，左端成为
\[
e_{\alpha\beta}\mathfrak a^\alpha\mathfrak c^\beta\,
e_{\dot\alpha\dot\beta}\mathfrak b^{\dot\alpha}\mathfrak d^{\dot\beta}
+e_{\alpha\dot\beta}\mathfrak a^\alpha\mathfrak d^{\dot\beta}\,
e_{\dot\alpha\beta}\mathfrak b^{\dot\alpha}\mathfrak c^\beta
=[\mathfrak a\mathfrak c][\mathfrak b\mathfrak d]+[\mathfrak a\mathfrak d][\mathfrak b\mathfrak c],
\]
即 $[\mathfrak a\mathfrak c][\mathfrak b\mathfrak d]+[\mathfrak a\mathfrak d][\mathfrak b\mathfrak c]=0$。
这正是交比、对合与调和点列中对调和点列的定义。
\end{proof}

\section{Levi-Civita 记号的旋量表示}

最后讨论 Levi-Civita 记号 $\epsilon_{ijk}$，它的旋量表示记为
$\epsilon_{\alpha\dot\alpha\beta\dot\beta\gamma\dot\gamma}$。它满足两条对称性：

\begin{enumerate}[label=(\arabic*),leftmargin=2.4em]
\item 交换一对内的两个指标（如 $\alpha$ 与 $\dot\alpha$）值不变，另外两对同理；
\item 交换两整对指标（即把一对指标整体与另一对交换）时值变号。
\end{enumerate}

由第一条，每一对指标的独立编号只有 $11,12,22$ 三种；由第二条，三对指标的编号不能相同，
于是在特定排列下只有唯一的选择。例如取 $\epsilon_{111222}=1$，再作排列即可得到其余
$11$ 个非零值。

$\epsilon_{\alpha\dot\alpha\beta\dot\beta\gamma\dot\gamma}$ 可以用 $e$ 表示：
\begin{equation}\label{eq:eps-formula}
\epsilon_{\alpha\dot\alpha\beta\dot\beta\gamma\dot\gamma}
=\sum e_{\alpha\beta}e_{\beta\gamma}e_{\gamma\alpha},
\end{equation}
其中求和式共有 $8$ 项，分别对应于在指标的不同位置上加点。例如其中一项是
$e_{\alpha\dot\beta}e_{\beta\dot\gamma}e_{\gamma\dot\alpha}$。

$\epsilon_{\alpha\dot\alpha\beta\dot\beta\gamma\dot\gamma}$ 的作用是由两对点求出未知的
对合变换。设 $p=(\mathfrak a\mathfrak b)$，$q=(\mathfrak c\mathfrak d)$，令
\[
r_{\alpha\dot\alpha}
=\epsilon_{\alpha\dot\alpha\beta\dot\beta\gamma\dot\gamma}
p^{\beta\dot\beta}q^{\gamma\dot\gamma}.
\]
由 $r$ 的对称性可知它表示一个对合变换；由 $\epsilon$ 的反对称性得到
\[
r_{\alpha\dot\alpha}\mathfrak a^\alpha\mathfrak b^{\dot\alpha}
=r_{\alpha\dot\alpha}\mathfrak c^\alpha\mathfrak d^{\dot\alpha}=0 .
\]
于是 $r$ 表示这样一个对合关系：$\mathfrak a$ 与 $\mathfrak b$ 关于 $r$ 对合，
$\mathfrak c$ 与 $\mathfrak d$ 关于 $r$ 对合。
   % 再谈二次点列
% 第二十节：洛伦兹变换与帕斯卡定理
\chapter{洛伦兹变换与帕斯卡定理}

\section{二次点列上的射影变换}

本节讨论旋量空间上的一般线性映射 $\varphi$。它表示二次点列上的一个射影变换，代数上由一个
旋量张量 ${t^\alpha}_\beta$ 表示，其特点是 $t_{\alpha\beta}$ 并不对称。对称的情形对应一个
对合，已在二次点列与二次对合一节中讨论过；非对称的情形不能再写成矢量的形式，只有借助旋量
才能处理。

设 $\mathfrak a$ 是 $S$ 中的旋量，$\mathfrak a' =\varphi(\mathfrak a)$ 由
\begin{equation}\label{eq:varphi}
\mathfrak a'{}^\alpha=t^\alpha{}_\beta\,\mathfrak a^\beta
\end{equation}
给出。约定用 $a'$ 表示 $\varphi(a)$，即若 $a=\mathfrak a\mathfrak a$，则
$a'=\mathfrak a'\mathfrak a'$。由于 $g$ 上的点与射影旋量一一对应，$\varphi$ 便是二次点列上的
一个射影变换；$\varphi$ 的不动点对应于 $t^\alpha{}_\beta$ 的特征旋量。

\section{命题}

\begin{proposition}\label{prop:lorentz}
$\varphi$ 可以从二次点列上的射影变换线性延拓为全平面上的射影变换，这一射影变换保持 $g$ 以及
点乘不变。反过来，每一个保持圆锥曲线（的像）不变的射影变换都可以用这样的 $\varphi$ 表示。
出于它与物理的相似性，称它为\emph{洛伦兹变换}。
\end{proposition}

\begin{remark}[与相对论速度空间的类比]
命题 \ref{prop:lorentz} 的一个应用是在相对论粒子碰撞中。把射影平面看作四维动量商掉模长后
得到的粒子速度空间（商掉静质量），则 $\varphi$ 代表所有可能的参考系变换：$g$ 上的点代表
光子，$g$ 内部的点代表有静质量的粒子。于是对一个粒子碰撞问题，可以画出它的速度图，再用
后面的命题 2、3、4 处理。
\end{remark}

\begin{proposition}\label{prop:locus}
设 $a$ 是 $g$ 上的动点，则直线 $a\times a'$ 的极点 $K$ 的轨迹构成一条圆锥曲线 $h$，
$h$ 与 $g$ 必定相切于两点 $p,q$。反之，每给定一条与 $g$ 相切于两点的圆锥曲线 $h$，都可以
通过这种作图方法构造出对应的线性映射 $\varphi$ 及其逆 $\varphi^{-1}$。
\end{proposition}

映射 $\varphi$ 与洛伦兹变换比较抽象，不容易看出几何含义；借助命题 \ref{prop:locus}，
可以用一类特殊的圆锥曲线来构造 $\varphi$，从而把它变得直观。

\begin{proposition}\label{prop:construction-outside}
设 $a$ 是 $g$ 外一点，则可以通过如下步骤得到 $a'$：
\begin{enumerate}[label=(\arabic*),leftmargin=2.4em]
\item 作 $a$ 到 $g$ 的两条切线，切点分别为 $b,c$；
\item $\varphi$ 在 $g$ 上有定义，故可作出 $b',c'$；
\item 作 $b',c'$ 处的切线，其交点即为 $a'$。
\end{enumerate}
\end{proposition}

命题 \ref{prop:construction-outside} 给出了 $\varphi$ 延拓的一种作图方式。

\begin{proposition}\label{prop:family}
用 $C$ 表示所有与 $g$ 相切于 $p,q$ 两点的圆锥曲线构成的集合。对任意一点 $a$，在 $C$ 中
存在唯一一条经过 $a$ 的圆锥曲线，记作 $c(a)$，并且恒有
\[
c(a)=c(a') .
\]
换言之，映射 $\varphi$ 把任意一条与 $g$ 相切于 $p,q$ 两点的圆锥曲线映为自身。于是所有
这样的圆锥曲线在平面上划出一族 $\varphi$ 轨道。
\end{proposition}

\begin{proposition}\label{prop:construction-inside}
设 $a$ 不在 $g$ 上，则可以通过如下步骤作出 $a'$：
\begin{enumerate}[label=(\arabic*),leftmargin=2.4em]
\item 作直线 $a\times p$，交 $g$ 于 $b$；
\item 作出 $b'$；
\item 作直线 $b'\times p$，它与 $c(a)$ 的交点即为 $a'$。
\end{enumerate}
\end{proposition}

\begin{proposition}[帕斯卡定理]\label{prop:pascal}
\begin{enumerate}[label=(\arabic*),leftmargin=2.4em]
\item $g$ 上的三对对应点 $a,a',b,b',c,c'$ 能唯一确定射影变换 $\varphi$，使得
      $a'=\varphi(a)$，$b'=\varphi(b)$，$c'=\varphi(c)$。
\item 设 $s=p\times q$，称 $s$ 为 $\varphi$ 对应的\emph{帕斯卡线}。设 $a,b$ 是 $g$ 上的
      任意两点，则点 $(a\times b')\times(a'\times b)$ 必定在帕斯卡线上。
\end{enumerate}
\end{proposition}

\section{命题的证明}

\subsection{命题 \ref{prop:lorentz} 的证明}

根据 \eqref{eq:varphi}，$\mathfrak a'{}^\alpha=t^\alpha{}_\beta\mathfrak a^\beta$。设
$a=\mathfrak a\mathfrak a$，则由旋量与类光矢量的对应关系有 $a'=\mathfrak a'\mathfrak a'$，
即
\[
a'^{\alpha\dot\alpha}
=t^\alpha{}_\beta\mathfrak a^\beta\,
t^{\dot\alpha}{}_{\dot\beta}\mathfrak a^{\dot\beta}
=t^\alpha{}_\beta t^{\dot\alpha}{}_{\dot\beta}a^{\beta\dot\beta},
\]
这是关于 $a^{\beta\dot\beta}$ 的线性映射。此时 $a^{\beta\dot\beta}$ 作为定义域只包含
$\det a^{\beta\dot\beta}=0$ 的部分，但线性映射
$t^\alpha{}_\beta t^{\dot\alpha}{}_{\dot\beta}a^{\beta\dot\beta}$ 可自然地把定义域延拓到所有
对称旋量张量，因而延拓到平面上的所有点。用矢量指标写出即
\[
a'^i=\sigma^i_{\alpha\dot\alpha}t^\alpha{}_\beta t^{\dot\alpha}{}_{\dot\beta}
\sigma_j^{\beta\dot\beta}a^j,
\]
可见它确实是全平面上的射影变换。

又因为 $\varphi$ 把 $g$ 上的点映为 $g$ 上的点，所以它保持 $g$ 的像（圆锥曲线对应的点集）
不变；圆锥曲线张量由其像完全确定，因此 $g$ 作为张量也关于 $\varphi$ 不变。于是由 $g_{ij}$
定义的点乘也关于 $\varphi$ 不变。反过来，在二维空间加一维时间中保持时空距离不变的参考系
变换是三维的（两个 boost 与一个空间旋转），而二维矢量空间中的射影变换也是三维的，
两者恰好一一对应。

\subsection{命题 \ref{prop:locus} 的证明}

只需用 $t_{\alpha\beta}$ 表示出 $h_{\alpha\beta\gamma\delta}$，即找出这样的
$h_{\alpha\beta\gamma\delta}$，使得对任意的 $\mathfrak c$ 恒有
\begin{equation}\label{eq:hlocus}
h_{\alpha\beta\gamma\delta}\,
t^\beta{}_{\dot\beta}t^\delta{}_{\dot\delta}\,
\mathfrak c^\alpha\mathfrak c^{\dot\beta}\mathfrak c^\gamma\mathfrak c^{\dot\delta}=0 .
\end{equation}

首先注意，射影变换 $t^\alpha{}_\beta$ 的不动点方程是二次的，因此必定有两个不动点，记为
$\mathfrak a,\mathfrak b$，即
\[
\mathfrak a^\alpha=t^\alpha{}_\beta\mathfrak a^\beta,\qquad
\mathfrak b^\alpha=t^\alpha{}_\beta\mathfrak b^\beta .
\]
先看第一个方程。把它降指标得 $\mathfrak a_\alpha=t_{\alpha\beta}\mathfrak a^\beta$，
两边乘以 $\mathfrak a^\alpha$，利用 $\mathfrak a^\alpha\mathfrak a_\alpha=0$，得到
$t_{\alpha\beta}\mathfrak a^\alpha\mathfrak a^\beta=0$。设
\[
s_{\alpha\beta}=t_{\alpha\beta}+t_{\beta\alpha},
\]
则 $s_{\alpha\beta}\mathfrak a^\alpha\mathfrak a^\beta=0$；同理
$s_{\alpha\beta}\mathfrak b^\alpha\mathfrak b^\beta=0$。由于 $s_{\alpha\beta}$ 对称，它表示
一个对合变换，而 $\mathfrak a,\mathfrak b$ 是其不动点，故
\begin{equation}\label{eq:s}
s_{\alpha\beta}=t_{\alpha\beta}+t_{\beta\alpha}=(\mathfrak a\mathfrak b)_{\alpha\beta},
\end{equation}
它对应直线 $p\times q$。

因为 $g,h$ 相切于两点，所以 $g,h,ss$ 构成二次线系，于是可设
\begin{equation}\label{eq:hansatz}
h_{\alpha\beta\gamma\delta}
=s_{\alpha\beta}s_{\gamma\delta}+\lambda g_{\alpha\beta\gamma\delta},
\end{equation}
其中 $\lambda$ 待定。把 \eqref{eq:hansatz} 代入 \eqref{eq:hlocus}，得
\[
\bigl(s_{\alpha\beta}s_{\gamma\delta}t^\beta{}_{\dot\beta}t^\delta{}_{\dot\delta}
+\lambda g_{\alpha\beta\gamma\delta}t^\beta{}_{\dot\beta}t^\delta{}_{\dot\delta}\bigr)
\mathfrak c^\alpha\mathfrak c^{\dot\beta}\mathfrak c^\gamma\mathfrak c^{\dot\delta}=0 .
\]
小括号中第一项为
\[
s_{\alpha\beta}s_{\gamma\delta}t^\beta{}_{\dot\beta}t^\delta{}_{\dot\delta}
=(t_{\alpha\beta}+t_{\beta\alpha})(t_{\gamma\delta}+t_{\delta\gamma})
t^{\beta\beta'}t^{\delta\delta'}e_{\beta'\dot\beta}e_{\delta'\dot\delta}
=\bigl({{t^2}_\alpha}^{\beta'}+\delta_\alpha^{\beta'}\bigr)
\bigl({{t^2}_\gamma}^{\delta'}+\delta_\gamma^{\delta'}\bigr)
e_{\beta'\dot\beta}e_{\delta'\dot\delta},
\]
即
\[
s_{\alpha\beta}s_{\gamma\delta}t^\beta{}_{\dot\beta}t^\delta{}_{\dot\delta}
=\bigl({t^2}_{\alpha\dot\beta}+e_{\alpha\dot\beta}\bigr)
\bigl({t^2}_{\gamma\dot\delta}+e_{\gamma\dot\delta}\bigr),
\qquad
{{t^2}{}^\alpha}_\gamma={t^\alpha}_\beta{t^\beta}_\gamma .
\]
把它乘以 $\mathfrak c^\alpha\mathfrak c^{\dot\beta}\mathfrak c^\gamma\mathfrak c^{\dot\delta}$：
含有 $e_{\alpha\dot\beta}$ 的项因 $e_{\alpha\dot\beta}\mathfrak c^\alpha\mathfrak c^{\dot\beta}=0$
而消失，留下的只有
\[
{t^2}_{\alpha\dot\beta}{t^2}_{\gamma\dot\delta}
\mathfrak c^\alpha\mathfrak c^{\dot\beta}\mathfrak c^\gamma\mathfrak c^{\dot\delta}
=\bigl({t^2}_{\alpha\dot\beta}\mathfrak c^\alpha\mathfrak c^{\dot\beta}\bigr)^2 .
\]
小括号中的第二项为
\[
g_{\alpha\beta\gamma\delta}t^\beta{}_{\dot\beta}t^\delta{}_{\dot\delta}
=(e_{\alpha\gamma}e_{\beta\delta}+e_{\alpha\delta}e_{\beta\gamma})
t^\beta{}_{\dot\beta}t^\delta{}_{\dot\delta}
=e_{\alpha\gamma}e_{\dot\delta\dot\beta}-t_{\alpha\dot\delta}t_{\gamma\dot\beta},
\]
其中第一项乘以 $\mathfrak c^\alpha\mathfrak c^{\dot\beta}\mathfrak c^\gamma\mathfrak c^{\dot\delta}$
后为零，故只留下
\[
(-t_{\alpha\dot\delta}t_{\gamma\dot\beta})\,
\mathfrak c^\alpha\mathfrak c^{\dot\beta}\mathfrak c^\gamma\mathfrak c^{\dot\delta}
=-\bigl(t_{\alpha\dot\beta}\mathfrak c^\alpha\mathfrak c^{\dot\beta}\bigr)^2 .
\]
于是 \eqref{eq:hlocus} 等价于
\begin{equation}\label{eq:hcond}
\bigl({t^2}_{\alpha\beta}\mathfrak c^\alpha\mathfrak c^\beta\bigr)^2
-\lambda\bigl(t_{\alpha\beta}\mathfrak c^\alpha\mathfrak c^\beta\bigr)^2=0 .
\end{equation}

表面上看，无论 $\lambda$ 如何取值也不能使 \eqref{eq:hcond} 成为恒等式，但事实并非如此。
$2\times2$ 矩阵 $A$ 满足恒等式
\begin{equation}\label{eq:cayley}
A\tr A-A^2=(\det A)I,
\end{equation}
即 $A^2-(\tr A)A+(\det A)I=0$。把 \eqref{eq:cayley} 应用于 ${t^\alpha}_\beta$，得到
\[
{t^\alpha}_\beta{t^\mu}_\mu-{{t^2}{}^\alpha}_\beta=(\det t)\delta^\alpha_\beta,
\]
也就是
\[
t_{\alpha\beta}{t^\mu}_\mu-{t^2}_{\alpha\beta}=(\det t)e_{\alpha\beta}.
\]
两边同乘 $\mathfrak c^\alpha\mathfrak c^\beta$，注意 $e_{\alpha\beta}\mathfrak c^\alpha\mathfrak c^\beta=0$，
便得
\[
{t^\mu}_\mu\bigl(t_{\alpha\beta}\mathfrak c^\alpha\mathfrak c^\beta\bigr)
={t^2}_{\alpha\beta}\mathfrak c^\alpha\mathfrak c^\beta .
\]
代入 \eqref{eq:hcond} 可知，只需取
\[
\lambda=({t^\mu}_\mu)^2,
\]
\eqref{eq:hcond} 便成为恒等式，从而证明了 $a\times a'$ 的极点确实在圆锥曲线 $h$ 上。于是
\begin{equation}\label{eq:hformula}
\boxed{\;h_{\alpha\beta\gamma\delta}
=(t_{\alpha\beta}+t_{\beta\alpha})(t_{\gamma\delta}+t_{\delta\gamma})
+({t^\mu}_\mu)^2g_{\alpha\beta\gamma\delta}.\;}
\end{equation}

最后注意，圆锥曲线共有 $5$ 个自由度，两个相切条件产生 $2$ 个约束，因此这样的 $h$ 共有
$3$ 个自由度，等于射影变换的自由度个数，故 $h$ 与 $\varphi$ 相互对应。

\subsection{命题 \ref{prop:construction-outside} 的证明}

利用 $\varphi$ 保持点乘这一性质。由 $a\cdot b=0$ 得 $a'\cdot b'=0$，同理
$a'\cdot c'=0$。因此 $a'$ 同时在 $b'$ 与 $c'$ 的极线上，即
\[
a'=b'\times c',
\]
这正是 $b',c'$ 处两条切线的交点。

\subsection{命题 \ref{prop:family} 的证明}

先证唯一性。五点确定一条圆锥曲线；"相切"可以看作同时经过切点附近切线上无限接近的两个点，
于是命题中的问题可以转化为五点确定一条圆锥曲线的问题，因此 $c(a)$ 唯一。

再证 $c(a)=c(a')$。由命题 \ref{prop:locus}，$C$ 中的元素都可用
\begin{equation}\label{eq:family}
c_\mu=s_{\alpha\beta}s_{\gamma\delta}+\mu\,g_{\alpha\beta\gamma\delta},
\qquad \mu\in\R,
\end{equation}
描述，其中 $s$ 由 \eqref{eq:s} 给出，对应直线 $p\times q$。由于 $p,q$ 是 $\varphi$ 的不动
点，$g$ 与 $s$ 在 $\varphi$ 下分别变为自身的倍数，因此 $c_\mu$ 的像仍属于曲线系
$\{ss+\mu'g\}$，即
\[
c_\mu'=\alpha\,s_{\alpha\beta}s_{\gamma\delta}
+\beta\mu\,g_{\alpha\beta\gamma\delta},
\]
其中 $\alpha,\beta$ 是与 $\mu$ 无关的常数。

取 $\mu=\lambda=({t^\mu}_\mu)^2$，此时 $c_\lambda$ 正是命题 \ref{prop:locus} 中的圆锥曲线
$h$。由 \eqref{eq:hformula} 与命题 \ref{prop:locus}，$h$ 的像就是它自身（作为点集），于是
\[
\alpha\,ss+\beta\lambda\,g
\ \text{与}\ ss+\lambda\,g\ \text{成正比}.
\]
比较系数（$\lambda\neq0$ 时）得 $\alpha=\beta$。代入 \eqref{eq:family} 的一般 $\mu$，
\[
c_\mu'=\alpha\,ss+\alpha\mu\,g=\alpha\,(ss+\mu g)=\alpha\,c_\mu,
\]
即 $c_\mu'$ 与 $c_\mu$ 表示同一条圆锥曲线，因而 $c(a')=c(a)$。当 $\lambda=0$ 时，
$\operatorname{tr}t=0$，$h$ 退化为 $ss$；此时结论可由 $\operatorname{tr}t\neq0$ 的情形取
极限得到。

\begin{remark}
$C$ 中的元素有两条是退化的：一条是 $ss$，另一条是 $p,q$ 处两条切线的对称积。设后者两条
切线的交点为 $r$，则 $r$ 是 $s$ 的极点。在整个平面上，$\varphi$ 的不动点有且仅有 $p,q,r$
三个。
\end{remark}

\subsection{命题 \ref{prop:construction-inside} 的证明}

只需证明按题中三步作图后，$a'$ 位于直线 $b'\times p$ 上。由作图，$b$ 在直线 $a\times p$ 上，
故
\[
(b\times p)\cdot a=0 .
\]
射影变换保持 Levi-Civita 记号（因而保持叉乘）与点乘，于是上式经 $\varphi$ 作用后成为
\[
(b'\times p')\cdot a'=0 .
\]
又 $p'=p$，所以 $(b'\times p)\cdot a'=0$，即 $a'$ 在直线 $b'\times p$ 上。结合命题
\ref{prop:family} 给出的 $a'\in c(a)$，二者的交点即为 $a'$。

\begin{remark}
$p$ 只是两个不动点中的一个。对另一个不动点 $q$ 同样操作：作直线 $a\times q$ 交 $g$ 于
$c$，作出 $c'$，再作直线 $c'\times q$，同样可证 $(c'\times q)\cdot a'=0$。因此
\[
a'=(b'\times p)\times(c'\times q).
\]
\end{remark}

\subsection{命题 \ref{prop:pascal} 的证明}

(1) 令 $g$ 上任意点 $d$ 的像 $d'$ 满足 $a,b,c,d$ 的交比等于 $a',b',c',d'$ 的交比，
这便唯一地确定了映射 $\varphi$。

(2) 由二次点列与二次对合中的结论，三条直线 $a\times b'$、$a'\times b$ 与 $s=p\times q$
共点，当且仅当存在对合 $f$，使得 $(p,q)$、$(a,b')$、$(a',b)$ 均关于 $f$ 对合。

用旋量语言写出：$a=\mathfrak a\mathfrak a$，$a'=\mathfrak a'\mathfrak a'$，
$b=\mathfrak b\mathfrak b$，$b'=\mathfrak b'\mathfrak b'$，于是
\[
a\times b'=(\mathfrak a\mathfrak b'),\qquad a'\times b=(\mathfrak a'\mathfrak b).
\]
在本节命题 \ref{prop:locus} 的证明中已得到
$s_{\alpha\beta}=t_{\alpha\beta}+t_{\beta\alpha}$。对合 $f$ 的三个条件分别为
\begin{equation}\label{eq:pascaleq}
f_{\alpha\beta}\mathfrak a^\alpha\mathfrak b'{}^\beta=0,\qquad
f_{\alpha\beta}\mathfrak a'{}^\alpha\mathfrak b^\beta=0,\qquad
f_{\alpha\beta}t^{\alpha\beta}=0 .
\end{equation}
前两个方程利用 $\mathfrak a'{}^\alpha=t^\alpha{}_\beta\mathfrak a^\beta$、
$\mathfrak b'{}^\beta=t^\beta{}_\gamma\mathfrak b^\gamma$ 展开为
\[
f_{\alpha\beta}\mathfrak a^\alpha t^\beta{}_{\dot\beta}\mathfrak b^{\dot\beta}=0,
\qquad
f_{\alpha\beta}t^\alpha{}_{\dot\alpha}\mathfrak a^{\dot\alpha}\mathfrak b^\beta=0 .
\]

为证明这三个方程可以同时成立，先令第三个方程成立，再证明此时前两个方程等价。
改设
\[
h_\alpha{}^\gamma=f_{\alpha\beta}t^{\beta\gamma},
\]
则第三个方程 $f_{\alpha\beta}t^{\alpha\beta}=0$ 改写为 $h_\gamma{}^\gamma=0$；前两个方程分别
改写为
\[
h_{\alpha\dot\beta}\mathfrak a^\alpha\mathfrak b^{\dot\beta}=0,
\qquad
h_{\beta\dot\alpha}\mathfrak a^{\dot\alpha}\mathfrak b^\beta=0
\]
（改写第二个方程时利用了 $f_{\alpha\beta}=f_{\beta\alpha}$）。而由 $h_\gamma{}^\gamma=0$
可以推出 $h_{\alpha\beta}=h_{\beta\alpha}$：这是二维矩阵的结论，零迹决定了对称。于是前两个
方程完全等价，只剩下一个方程。对称无迹的 $h$ 有两个独立分量，一个线性方程将其确定到相差
一个常数因子，因而对合 $f$ 被唯一确定。证毕。
   % 洛伦兹变换与帕斯卡定理

\appendix
% 附录
\chapter{二次曲线系的退化元素与交比参数}

本附录补充曲线系的一个结构性的结论。

\begin{proposition}
设 $a,b,c,d$ 是平面上四个点，其中任意三点不共线。则经过这四点的二次线系中含有三条退化的
二次线，它们分别是四条连线两两配对的对称积。
\end{proposition}

\begin{proof}
四点两两连共有六条直线，分成三对互不相交的对边：
\[
l_1=a\times d,\quad l_1'=b\times c;\qquad
l_2=a\times c,\quad l_2'=d\times b;\qquad
l_3=a\times b,\quad l_3'=c\times d .
\]
把三条退化二次线写为
\[
h_1=(l_1l_1'),\qquad h_2=(l_2l_2'),\qquad h_3=(l_3l_3').
\]
$h_1$ 对应的点集是 $l_1\cup l_1'$，它经过 $a,b,c,d$ 四点；$h_2,h_3$ 同理。三点不共线保证这
四点确实落在两条直线上（而不落在一条直线上），故三者确实不同。又二次线系由两个线性无关的
元素生成，而 $h_1,h_2,h_3$ 满足一个线性关系（六条连线之间有三组对边，对应 $h_1,h_2,h_3$ 的
线性相关），因此它们给出的正是曲线系中全部的退化元素。
\end{proof}

\begin{proposition}
上述曲线系可以写成
\begin{equation}
g=((a\times d),(b\times c))+\lambda\,((a\times c),(d\times b)),
\label{eq:appendix-pencil}
\end{equation}
即由两条退化二次线生成。式中的参数 $\lambda$ 具有几何意义：它等于 $a,b,c,d$ 四点在 $g$ 上的交比。
\end{proposition}

\begin{proof}[证明要点]
由前一条命题，$g$ 是这两条退化二次线的线性组合，故形如 \eqref{eq:appendix-pencil}。
关于 $\lambda$ 的几何意义：在 $g$ 上任取一点 $P$，以 $P$ 为心把曲线系上的四点投影为线束，
交比是射影不变量，算出该交比即得 $\lambda$。计算得知
\[
g=((a\times b),(c\times d))\quad\Longleftrightarrow\quad \lambda=1,
\]
即 $\lambda=1$ 给出第三条退化二次线；而当
\[
g:(ab)(cd)=0\quad\Longleftrightarrow\quad \lambda=-1
\]
时，$a,b,c,d$ 在 $g$ 上成调和点列。
\end{proof}

\begin{remark}
\eqref{eq:appendix-pencil} 中 $\lambda$ 的取值与生成的曲线系之间是一一对应的，
因此四点 $a,b,c,d$ 确定了一个以 $\lambda$ 为参数的曲线族；$\lambda=-1$（调和）与
$\lambda=1$（第三条退化二次线）是这个族中两个特殊的位置。
\end{remark}


\end{document}
